% Generated by the blueprint tool (blueprint-pdf/1) from a blueprint snapshot; edit the
% Markdown sources, not this file.  Compile with `latexmk -pdf` or by running
% `pdflatex` twice (the second run settles the contents and the page references).
% Images, if any, are read from assets/ next to this file.
\documentclass[11pt]{article}
\usepackage[T1]{fontenc}
\usepackage[utf8]{inputenc}
\usepackage{lmodern}
\usepackage{textcomp}
\usepackage[a4paper,margin=2.5cm,headheight=14pt]{geometry}
\usepackage{amsmath,amssymb,amsthm}
\usepackage{amscd}  % KaTeX's CD diagrams; part of amsmath, so always there
\usepackage{mathtools}
\usepackage{bm}
\usepackage{graphicx}
\usepackage{xcolor}
\usepackage{array,booktabs,tabularx}
\usepackage{fancyvrb}
\usepackage[normalem]{ulem}
\usepackage{microtype}
% Optional, and only with their fonts: a small TeX installation has the
% packages without the fonts, which fails at the first formula.
\IfFileExists{mathrsfs.sty}{\IfFileExists{rsfs10.tfm}{\usepackage{mathrsfs}}{}}{}
\IfFileExists{stmaryrd.sty}{\IfFileExists{stmary10.tfm}{\usepackage{stmaryrd}}{}}{}
\IfFileExists{cancel.sty}{\usepackage{cancel}}{}
\usepackage[hyphens]{url}
\usepackage[pdfusetitle,bookmarksnumbered=false]{hyperref}
\usepackage{bookmark}
\definecolor{bplink}{HTML}{1F4E9C}
\definecolor{bpmuted}{HTML}{5F676E}
\hypersetup{colorlinks=true,linkcolor=bplink,urlcolor=bplink,citecolor=bplink}
\setlength{\parindent}{0pt}
\setlength{\parskip}{0.55em plus 0.2em minus 0.1em}
% Lists as tight as the page's.
\makeatletter
\def\@listi{\leftmargin\leftmargini\parsep 0pt\topsep 0.3em plus 0.1em\itemsep 0.1em plus 0.1em}
\let\@listI\@listi
\def\@listii{\leftmargin\leftmarginii\labelwidth\leftmarginii\advance\labelwidth-\labelsep\parsep 0pt\topsep 0.1em\itemsep 0.1em}
% An ordered list that does not start at 1.
\newcommand{\bpstartat}[1]{\setcounter{\@enumctr}{\numexpr#1-1\relax}}
\makeatother
\setcounter{tocdepth}{3}
\emergencystretch=3em
\allowdisplaybreaks

% A numbered statement, set as a paper sets one: "Theorem 2.1 (Title)."
% #1 is the lead ("Theorem~2.1"), #2 what follows it in roman.
\newenvironment{bpstatement}[2]{\par\addvspace{0.6\baselineskip}\noindent\textbf{#1}#2\textbf{.}\enspace\ignorespaces}
  {\par\addvspace{0.4\baselineskip}}
\newenvironment{bpstatementit}[2]{\begin{bpstatement}{#1}{#2}\itshape\ignorespaces}{\end{bpstatement}}
% A step: a relation with prose of its own, read under the object it starts from.
\newenvironment{bpstep}[1]{\par\addvspace{0.3\baselineskip}\begin{list}{}{\leftmargin=1.6em\rightmargin=0pt\topsep=0pt}\item\relax
  {\small\color{bpmuted}#1}\par\nopagebreak}{\end{list}\par}
% Notes about where an object is written out.
\newenvironment{bpnote}{\par\small\color{bpmuted}\itshape}{\par}
\newcommand{\bpomitted}[1]{\par{\small\color{bpmuted}\itshape #1\par}}
\newcommand{\bpbodyheading}[1]{\par\addvspace{0.4\baselineskip}\noindent\textbf{#1}\par\nopagebreak}
\newcommand{\bprule}{\par\medskip\noindent\hfil\rule{0.3\linewidth}{0.4pt}\hfil\par\medskip}
\newcommand{\bpboundary}[1]{\par\nopagebreak{\small\color{bpmuted}#1\par}\nopagebreak}
% A reference to something this export leaves out, and one that resolves nowhere.
\newcommand{\bpout}[1]{#1\textsuperscript{\hyperlink{bp-outside}{\ensuremath{\ast}}}}
\newcommand{\bpbroken}[1]{\textcolor{red!70!black}{[#1?]}}
\newcommand{\bpslug}[1]{\texttt{#1}}
% A bibliography entry: label, anchor, entry; the label hangs in the margin.
\newcommand{\bpbibitem}[3]{\par\setbox0\hbox{[#1]\enspace}\ifdim\wd0<5em\setbox0\hbox to 5em{[#1]\hfil}\fi
  \hangindent=\wd0 \hangafter=1 \noindent\box0\phantomsection\label{#2}#3\par\smallskip}
% A character the export knows no way to set.
\newcommand{\bpglyph}[1]{\fbox{\scriptsize #1}}
% Images, at their own size unless they are wider than the line.
\newsavebox{\bpimagebox}
\newcommand{\bpimage}[1]{\sbox{\bpimagebox}{\includegraphics{#1}}%
  \ifdim\wd\bpimagebox>\linewidth\resizebox{\linewidth}{!}{\usebox{\bpimagebox}}\else\usebox{\bpimagebox}\fi}
\newcommand{\bpmissingimage}[1]{\fbox{\parbox{0.8\linewidth}{\centering\small\color{bpmuted}Image not included: #1}}}
\DefineVerbatimEnvironment{bpverbatim}{Verbatim}{fontsize=\small,xleftmargin=1em}
\renewcommand{\qedsymbol}{\ensuremath{\blacksquare}}
\hypersetup{pdftitle={Weil Conjectures — Appendices},pdfsubject={Part Appendices; all levels},pdfcreator={blueprint}}
% KaTeX commands the website knows
\providecommand{\R}{\mathbb{R}}\providecommand{\N}{\mathbb{N}}\providecommand{\Z}{\mathbb{Z}}
\providecommand{\Q}{\mathbb{Q}}\providecommand{\C}{\mathbb{C}}
\providecommand{\Reals}{\mathbb{R}}\providecommand{\reals}{\mathbb{R}}\providecommand{\Complex}{\mathbb{C}}
\providecommand{\cnums}{\mathbb{C}}\providecommand{\natnums}{\mathbb{N}}\providecommand{\Bbb}[1]{\mathbb{#1}}
\providecommand{\bold}[1]{\mathbf{#1}}\providecommand{\frak}[1]{\mathfrak{#1}}\providecommand{\mathbbm}[1]{\mathbb{#1}}
\providecommand{\lang}{\langle}\providecommand{\rang}{\rangle}\providecommand{\bra}[1]{\langle #1\rvert}
\providecommand{\ket}[1]{\lvert #1\rangle}\providecommand{\braket}[1]{\langle #1\rangle}
\providecommand{\set}[1]{\{#1\}}\providecommand{\Set}[1]{\left\{#1\right\}}
\providecommand{\infin}{\infty}\providecommand{\isin}{\in}\providecommand{\plusmn}{\pm}\providecommand{\sub}{\subset}
\providecommand{\sube}{\subseteq}\providecommand{\supe}{\supseteq}\providecommand{\exist}{\exists}
\providecommand{\bull}{\bullet}\providecommand{\larr}{\leftarrow}\providecommand{\rarr}{\rightarrow}
\providecommand{\lrarr}{\leftrightarrow}\providecommand{\Larr}{\Leftarrow}\providecommand{\Rarr}{\Rightarrow}
\providecommand{\Lrarr}{\Leftrightarrow}\providecommand{\uarr}{\uparrow}\providecommand{\darr}{\downarrow}
\providecommand{\harr}{\leftrightarrow}\providecommand{\lArr}{\Leftarrow}\providecommand{\rArr}{\Rightarrow}
\providecommand{\hArr}{\Leftrightarrow}\providecommand{\alef}{\aleph}\providecommand{\alefsym}{\aleph}
\providecommand{\real}{\Re}\providecommand{\image}{\Im}\providecommand{\weierp}{\wp}
\providecommand{\Alpha}{\mathrm{A}}\providecommand{\Beta}{\mathrm{B}}\providecommand{\Epsilon}{\mathrm{E}}
\providecommand{\Zeta}{\mathrm{Z}}\providecommand{\Eta}{\mathrm{H}}\providecommand{\Iota}{\mathrm{I}}
\providecommand{\Kappa}{\mathrm{K}}\providecommand{\Mu}{\mathrm{M}}\providecommand{\Nu}{\mathrm{N}}
\providecommand{\Omicron}{\mathrm{O}}\providecommand{\Rho}{\mathrm{P}}\providecommand{\Tau}{\mathrm{T}}
\providecommand{\Chi}{\mathrm{X}}\providecommand{\omicron}{o}\providecommand{\gt}{>}\providecommand{\lt}{<}
\providecommand{\And}{\&}\providecommand{\hdashline}{\hline}\providecommand{\operatornamewithlimits}[1]{\operatorname*{#1}}
\providecommand{\llbracket}{[\![}\providecommand{\rrbracket}{]\!]}\providecommand{\cancel}[1]{#1}
\providecommand{\mathscr}{\mathcal}
% the characters beyond ASCII this document uses
\DeclareUnicodeCharacter{2160}{\textup{I}}
\DeclareUnicodeCharacter{2161}{\textup{II}}
\DeclareUnicodeCharacter{2162}{\textup{III}}
% the project's macros ([katex.macros] in blueprint.toml)
\providecommand{\A}{}\renewcommand{\A}{\mathbf A}
\providecommand{\AEn}{}\renewcommand{\AEn}{A_{E,n}}
\providecommand{\Aaf}{}\renewcommand{\Aaf}{\mathbf A}
\providecommand{\Ab}{}\renewcommand{\Ab}{\mathsf{Ab}}
\providecommand{\AffSch}{}\renewcommand{\AffSch}{\operatorname{AffSch}}
\providecommand{\Aone}{}\renewcommand{\Aone}{\mathbf A^1}
\providecommand{\Bl}{}\renewcommand{\Bl}{\operatorname{Bl}}
\providecommand{\Br}{}\renewcommand{\Br}{\mathrm{Br}}
\providecommand{\C}{}\renewcommand{\C}{\mathbb C}
\providecommand{\Cech}{}\renewcommand{\Cech}{\check{H}}
\providecommand{\Coext}{}\renewcommand{\Coext}{\operatorname{Coext}}
\providecommand{\Cone}{}\renewcommand{\Cone}{\operatorname{Cone}}
\providecommand{\Cons}{}\renewcommand{\Cons}{\operatorname{Cons}}
\providecommand{\Cref}{}\renewcommand{\Cref}[1]{\text{#1}}
\providecommand{\D}{}\renewcommand{\D}{D}
\providecommand{\DLAFT}{}\renewcommand{\DLAFT}{\operatorname{FT}}
\providecommand{\DLAFr}{}\renewcommand{\DLAFr}{\operatorname{Fr}}
\providecommand{\DLAQ}{}\renewcommand{\DLAQ}{{\overline{\mathbf Q}_\ell}}
\providecommand{\DLRQ}{}\renewcommand{\DLRQ}{{\overline{\mathbf Q}_\ell}}
\providecommand{\Dbc}{}\renewcommand{\Dbc}{D^b_c}
\providecommand{\Dc}{}\renewcommand{\Dc}{D_{\mathrm{cons}}}
\providecommand{\Dctf}{}\renewcommand{\Dctf}{D^b_{ctf}}
\providecommand{\Dic}{}\renewcommand{\Dic}{D_{\IndCons}}
\providecommand{\Dplus}{}\renewcommand{\Dplus}{D^{+}}
\providecommand{\End}{}\renewcommand{\End}{\operatorname{End}}
\providecommand{\Et}{}\renewcommand{\Et}{\mathrm{\acute{e}t}}
\providecommand{\Ex}{}\renewcommand{\Ex}{\mathrm{Ex}}
\providecommand{\Ext}{}\renewcommand{\Ext}{\operatorname{Ext}}
\providecommand{\F}{}\renewcommand{\F}{\mathbb F}
\providecommand{\FT}{}\renewcommand{\FT}{\operatorname{FT}}
\providecommand{\Fbar}{}\renewcommand{\Fbar}{{\overline{\mathbf F}_q}}
\providecommand{\Fell}{}\renewcommand{\Fell}{\mathbf F_\ell}
\providecommand{\Fib}{}\renewcommand{\Fib}{\operatorname{Fib}}
\providecommand{\Fix}{}\renewcommand{\Fix}{\operatorname{Fix}}
\providecommand{\Fp}{}\renewcommand{\Fp}{{\mathbf F_{p}}}
\providecommand{\Fq}{}\renewcommand{\Fq}{{\mathbf F_q}}
\providecommand{\Fqn}{}\renewcommand{\Fqn}{{\mathbf F_{q^{n}}}}
\providecommand{\Frob}{}\renewcommand{\Frob}{\operatorname{Frob}}
\providecommand{\GL}{}\renewcommand{\GL}{\operatorname{GL}}
\providecommand{\Gal}{}\renewcommand{\Gal}{\operatorname{Gal}}
\providecommand{\Hens}{}\renewcommand{\Hens}{\operatorname{Hens}}
\providecommand{\Hilb}{}\renewcommand{\Hilb}{\operatorname{Hilb}}
\providecommand{\Hom}{}\renewcommand{\Hom}{\operatorname{Hom}}
\providecommand{\IndCons}{}\renewcommand{\IndCons}{\operatorname{IndCons}}
\providecommand{\Kar}{}\renewcommand{\Kar}{\operatorname{Kar}}
\providecommand{\LT}{}\renewcommand{\LT}{\operatorname{LT}}
\providecommand{\Lie}{}\renewcommand{\Lie}{\operatorname{Lie}}
\providecommand{\Loc}{}\renewcommand{\Loc}{\operatorname{Loc}}
\providecommand{\Mod}{}\renewcommand{\Mod}{\mathsf{Mod}}
\providecommand{\N}{}\renewcommand{\N}{\mathbb{N}}
\providecommand{\Needspace}{}\renewcommand{\Needspace}{\par}
\providecommand{\O}{}\renewcommand{\O}{{\mathcal{O}}}
\providecommand{\Ob}{}\renewcommand{\Ob}{\mathrm{Ob}}
\providecommand{\Oe}{}\renewcommand{\Oe}{\mathcal O_E}
\providecommand{\PP}{}\renewcommand{\PP}{\mathbf{P}}
\providecommand{\Perf}{}\renewcommand{\Perf}{\operatorname{Perf}}
\providecommand{\Pic}{}\renewcommand{\Pic}{\mathrm{Pic}}
\providecommand{\Pm}{}\renewcommand{\Pm}{\mathbf P^{m}}
\providecommand{\Pone}{}\renewcommand{\Pone}{\mathbf P^{1}}
\providecommand{\Pqn}{}\renewcommand{\Pqn}{\mathbf P^{n}}
\providecommand{\ProFin}{}\renewcommand{\ProFin}{\operatorname{ProFin}}
\providecommand{\Proet}{}\renewcommand{\Proet}{\mathrm{pro\acute{e}t}}
\providecommand{\Q}{}\renewcommand{\Q}{\mathbb Q}
\providecommand{\Qb}{}\renewcommand{\Qb}{\overline{\mathbf Q}_{\ell}}
\providecommand{\Qbar}{}\renewcommand{\Qbar}{\overline{\mathbf Q}}
\providecommand{\QbarX}{}\renewcommand{\QbarX}{\overline{\mathbf Q}_{\ell,X}}
\providecommand{\Qbarl}{}\renewcommand{\Qbarl}{{\overline{\mathbf Q}_{\ell}}}
\providecommand{\Qell}{}\renewcommand{\Qell}{{\mathbf{Q}_{\ell}}}
\providecommand{\Ql}{}\renewcommand{\Ql}{\mathbf Q_{\ell}}
\providecommand{\Qlb}{}\renewcommand{\Qlb}{{\overline{\mathbf Q}_\ell}}
\providecommand{\Qlbar}{}\renewcommand{\Qlbar}{\overline{\mathbf Q}_\ell}
\providecommand{\R}{}\renewcommand{\R}{\mathbb{R}}
\providecommand{\RG}{}\renewcommand{\RG}{\operatorname{\mathrm{R}\Gamma}}
\providecommand{\RGamma}{}\renewcommand{\RGamma}{R\Gamma}
\providecommand{\RGammac}{}\renewcommand{\RGammac}{R\Gamma_c}
\providecommand{\RHom}{}\renewcommand{\RHom}{R\mathcal{H}om}
\providecommand{\RHomX}{}\renewcommand{\RHomX}{R\mathcal H om_X}
\providecommand{\Rat}{}\renewcommand{\Rat}{\operatorname{IsRat}}
\providecommand{\Res}{}\renewcommand{\Res}{\operatorname{Res}}
\providecommand{\Rlim}{}\renewcommand{\Rlim}{R\!\varprojlim}
\providecommand{\SIm}{}\renewcommand{\SIm}{\operatorname{SIm}}
\providecommand{\Sch}{}\renewcommand{\Sch}{\mathsf{Sch}}
\providecommand{\Sh}{}\renewcommand{\Sh}{\operatorname{Sh}}
\providecommand{\Span}{}\renewcommand{\Span}{\operatorname{span}}
\providecommand{\Spec}{}\renewcommand{\Spec}{\operatorname{Spec}}
\providecommand{\Tr}{}\renewcommand{\Tr}{\operatorname{Tr}}
\providecommand{\Trsm}{}\renewcommand{\Trsm}{\operatorname{Tr}^{\mathrm{sm}}}
\providecommand{\Wt}{}\renewcommand{\Wt}{\operatorname{Wt}}
\providecommand{\Z}{}\renewcommand{\Z}{\mathbb{Z}}
\providecommand{\Zar}{}\renewcommand{\Zar}{\mathrm{zar}}
\providecommand{\Zell}{}\renewcommand{\Zell}{{\mathbf{Z}_{\ell}}}
\providecommand{\Zl}{}\renewcommand{\Zl}{{\mathbf Z_\ell}}
\providecommand{\affet}{}\renewcommand{\affet}[1]{{#1}^{\mathrm{aff}}_{\mathrm{ét}}}
\providecommand{\affproet}{}\renewcommand{\affproet}[1]{{#1}^{\mathrm{aff}}_{\mathrm{proét}}}
\providecommand{\boxt}{}\renewcommand{\boxt}{\boxtimes_A^{\mathrm L}}
\providecommand{\cF}{}\renewcommand{\cF}{\mathcal{F}}
\providecommand{\cHom}{}\renewcommand{\cHom}{R\mathcal{H}\!om}
\providecommand{\cI}{}\renewcommand{\cI}{\mathcal{I}}
\providecommand{\cL}{}\renewcommand{\cL}{\mathcal L}
\providecommand{\cM}{}\renewcommand{\cM}{\mathcal M}
\providecommand{\cN}{}\renewcommand{\cN}{\mathcal N}
\providecommand{\cO}{}\renewcommand{\cO}{\mathcal O}
\providecommand{\cT}{}\renewcommand{\cT}{\mathcal T}
\providecommand{\checked}{}\renewcommand{\checked}{\color{red}\aftergroup\restorecolor}
\providecommand{\coeff}{}\renewcommand{\coeff}{\operatorname{coeff}}
\providecommand{\coker}{}\renewcommand{\coker}{\operatorname{coker}}
\providecommand{\colim}{}\renewcommand{\colim}{\operatorname{colim}}
\providecommand{\constructedby}{}\renewcommand{\constructedby}[1]{\par\smallskip\noindent\textcolor{purple!60!black}{\textsf{Constructed by:} }
  \footnotesize\sloppy\nolinkurl{#1}\par}
\providecommand{\deepresult}{}\renewcommand{\deepresult}[1]{\par\smallskip\noindent\textcolor{orange!45!black}{
  \textsf{Blueprint role: named deep theorem node; see the cluster proof architecture.}}
  \par\noindent\textcolor{orange!45!black}{\textsf{Mathematical source:} }#1\par}
\providecommand{\dt}{}\renewcommand{\dt}{\mathbin{\otimes^{\mathbf L}}}
\providecommand{\dw}{}\renewcommand{\dw}{\operatorname{dw}}
\providecommand{\et}{}\renewcommand{\et}[1]{{#1}_{\mathrm{ét}}}
\providecommand{\ev}{}\renewcommand{\ev}{\operatorname{ev}}
\providecommand{\fib}{}\renewcommand{\fib}{\operatorname{fib}}
\providecommand{\formalizationnote}{}\renewcommand{\formalizationnote}[1]{\begin{quote}\small\raggedright\textbf{Formalization note.} #1\end{quote}}
\providecommand{\fpqc}{}\renewcommand{\fpqc}{X_{\mathrm{fpqc}}}
\providecommand{\gammacont}{}\renewcommand{\gammacont}{\Gamma_{\mathrm{cont}}}
\providecommand{\gm}{}\renewcommand{\gm}{\mathbf G_m}
\providecommand{\hocolim}{}\renewcommand{\hocolim}{\operatorname{hocolim}}
\providecommand{\holim}{}\renewcommand{\holim}{\operatorname{holim}}
\providecommand{\hrB}{}\renewcommand{\hrB}{\@ifnextchar]{\hrEnd}{\hrStep}}
\providecommand{\hrEnd}{}\renewcommand{\hrEnd}{\@firstoftwo{}}
\providecommand{\hrStep}{}\renewcommand{\hrStep}{\@firstoftwo{\hrB}}
\providecommand{\id}{}\renewcommand{\id}{\operatorname{id}}
\providecommand{\kunFl}{}\renewcommand{\kunFl}{\mathbf F_\ell}
\providecommand{\kunZn}{}\renewcommand{\kunZn}{\mathbf Z/n\mathbf Z}
\providecommand{\kuncd}{}\renewcommand{\kuncd}{\operatorname{cd}}
\providecommand{\kundim}{}\renewcommand{\kundim}{\dim}
\providecommand{\kunotimes}{}\renewcommand{\kunotimes}[1]{\otimes^{\mathbf L}_{#1}}
\providecommand{\kunsupp}{}\renewcommand{\kunsupp}{\operatorname{Supp}}
\providecommand{\lotimes}{}\renewcommand{\lotimes}{\otimes^{\mathrm L}}
\providecommand{\m}{}\renewcommand{\m}{{\mathfrak{m}}}
\providecommand{\n}{}\renewcommand{\n}{{\mathfrak{n}}}
\providecommand{\nagataGap}{}\renewcommand{\nagataGap}[1]{\par\noindent\textbf{Formalization gap.} #1\par}
\providecommand{\openobligation}{}\renewcommand{\openobligation}[1]{\par\smallskip\noindent\textcolor{red!65!black}{\textsf{Open obligation:} }
  \footnotesize\sloppy\nolinkurl{#1}\par}
\providecommand{\p}{}\renewcommand{\p}{{\mathfrak{p}}}
\providecommand{\pr}{}\renewcommand{\pr}{\operatorname{pr}}
\providecommand{\preshv}{}\renewcommand{\preshv}[1]{\mathrm{PreShv}(#1)}
\providecommand{\pro}{}\renewcommand{\pro}[1]{\mathrm{Pro}(#1)}
\providecommand{\proet}{}\renewcommand{\proet}{X_{\mathrm{proét}}}
\providecommand{\q}{}\renewcommand{\q}{{\mathfrak{q}}}
\providecommand{\rad}{}\renewcommand{\rad}{\operatorname{rad}}
\providecommand{\rank}{}\renewcommand{\rank}{\operatorname{rank}}
\providecommand{\red}{}\renewcommand{\red}{\color{red}}
\providecommand{\restorecolor}{}\renewcommand{\restorecolor}{\color{black}}
\providecommand{\retiredobligation}{}\renewcommand{\retiredobligation}[1]{\par\smallskip\noindent\textcolor{red!75!black}{\textsf{Retired invalid route:} }
  \footnotesize\sloppy\nolinkurl{#1}\par}
\providecommand{\rk}{}\renewcommand{\rk}{\operatorname{rk}}
\providecommand{\sep}{}\renewcommand{\sep}{\textnormal{sep}}
\providecommand{\shv}{}\renewcommand{\shv}[1]{\mathrm{Shv}(#1)}
\providecommand{\spec}{}\renewcommand{\spec}[1]{\mathrm{Spec}(#1)}
\providecommand{\stacks}{}\renewcommand{\stacks}[1]{\href{https://stacks.math.columbia.edu/tag/#1}{\textnormal{Stacks, Tag~#1}}}
\providecommand{\stacksproject}{}\renewcommand{\stacksproject}[1]{{\cite[\href{https://stacks.math.columbia.edu/tag/#1}{Tag #1}]{stacks-project}}}
\providecommand{\statusnote}{}\renewcommand{\statusnote}[1]{\begin{quote}\small\raggedright\textbf{Status.} #1\end{quote}}
\providecommand{\tildeideal}{}\renewcommand{\tildeideal}[2]{\widetilde{(#1, #2)}}
\providecommand{\tildering}{}\renewcommand{\tildering}[3]{{#1}_{\widetilde{(#2, #3)}}}
\providecommand{\tr}{}\renewcommand{\tr}{\operatorname{tr}}
\providecommand{\wt}{}\renewcommand{\wt}{\operatorname{wt}}

\begin{document}
\pagestyle{plain}
\begin{center}
{\LARGE\bfseries Weil Conjectures\par}
\vspace{0.7em}
{\large Part: Appendices\par}
\vspace{0.5em}
{\small\color{bpmuted}Appendices included\par}
{\small\color{bpmuted}Snapshot 2026-09-26, revision 8e4cb88e2a\par}
\end{center}
\vspace{0.5em}
{\small\tableofcontents}
\bigskip\hrule\bigskip

% object: part-Appendices (blueprint/part-Appendices/_section.md)
\part*{\texorpdfstring{Appendices}{Appendices}}
\phantomsection\label{bp1}%
\addcontentsline{toc}{part}{\texorpdfstring{Appendices}{Appendices}}
% text: prose-PA-01
\phantomsection\label{bp2}%
To define \(Rf_!\) in Part I we must compactify every separated morphism of finite type, and to prove proper base change we must dominate every proper morphism by a projective one. These are Nagata's theorem and Chow's lemma (Appendix \hyperref[bp3]{A}). Nagata's theorem, followed by normalization, also gives the smooth compactification of a curve used in Part III (Lemma \bpout{14.1}).

Proper base change ends with a change of separably closed ground field. There we need the cohomology of curves for all torsion coefficients, including torsion divisible by the characteristic (Theorem \bpout{4.20}). Purity starts from sheaves of rank one. Their geometric monodromy is finite because Frobenius fixes no class in \(H^1\) of a curve (Proposition \bpout{14.3}). For coefficients prime to the characteristic, both facts come from the Jacobian. Indeed, \(H^1(C,\mu_n)\) is the \(n\)-torsion of the Jacobian, and a fixed class would give \(\ell\)-power torsion points of unbounded order among its finitely many rational points. Appendices \hyperref[bp75]{Appendix B}, \hyperref[bp119]{Appendix C} and \hyperref[bp172]{Appendix D} prove both facts.

Rankin's argument in Chapter \bpout{14} rests on the determinant weights of tensor and exterior powers (Proposition \bpout{14.19}). These come from the structure of the monodromy group (Lemma \bpout{14.15}), and Appendix \hyperref[bp318]{F} proves the facts about algebraic groups over \(\overline{\mathbf Q}_\ell\) behind it.

Appendix \hyperref[bp221]{E} is of a different kind. Following Varshavsky [\hyperref[bp347]{Var05}], it computes the local terms of contracting correspondences in a second way; Chapter \bpout{12} does not depend on it.

% object: chap-Nagata-compactification-and-Chows-lemma (blueprint/chap-Nagata-compactification-and-Chows-lemma/_section.md), content hidden
\section*{\texorpdfstring{A\quad Nagata compactification and Chow's lemma}{A Nagata compactification and Chow's lemma}}
\phantomsection\label{bp3}%
\addcontentsline{toc}{section}{\texorpdfstring{A\quad Nagata compactification and Chow's lemma}{A Nagata compactification and Chow's lemma}}
% text: prose-A-01
\phantomsection\label{bp4}%
The main text defines cohomology with compact support as in topology: to push forward along \(f:X\to Y\), embed \(X\) as an open subset of something proper over \(Y\), extend by zero and push forward, so that \(Rf_!=Rp_*j_!\) (Proposition \bpout{5.26}). This definition needs every morphism to have a compactification, and it needs any two compactifications to be dominated by a third, so that proper base change can show that \(Rf_!\) does not depend on the choice. The second point is elementary once we use compactifications in which \(X\) is scheme-theoretically dense: the closure of \(X\) in the fibre product of two of them dominates both (Proposition \hyperref[bp21]{A.7}). The first point is Nagata's theorem, Theorem \hyperref[bp9]{A.2}: every separated morphism of finite type \(X\to S\) with \(S\) quasi-compact and quasi-separated factors as a quasi-compact open immersion \(X\to\overline
X\) followed by a proper morphism \(\overline X\to S\).

Why is this hard? Locally a compactification is easy to find: if \(S\) and \(X\) are affine, we take the closure of \(X\) in projective space over \(S\). The difficulty is to glue. For example, cover \(X=\mathbf A^1_k\) by \(U_1=\mathbf A^1\setminus\{0\}\) and \(U_2=\mathbf A^1\setminus\{1\}\), compactify each by \(\mathbf P^1\), and glue the two copies of \(\mathbf
P^1\) along \(U_1\cap U_2\). The result is then \(\mathbf P^1\) with the three points \(0\), \(1\) and \(\infty\) doubled, which is not separated. For curves blow-ups are not needed, and they would not help either, since blowing up a point on a regular curve does nothing. In dimension at least two, however, the closures of \(U_1\setminus U_2\) and \(U_2\setminus U_1\) can meet in the boundary, and we pull them apart by blow-ups with centers in the boundary, which leave \(X\) unchanged (Lemma \hyperref[bp55]{A.20} and Proposition \hyperref[bp63]{A.23}). Induction over an affine cover then gives the theorem over a Noetherian base (Section \hyperref[bp65]{A.4}), and Noetherian approximation gives the general case, which we need because proper base change is proved over arbitrary qcqs bases.

Proper base change itself is computed on projective space, and Chow's lemma carries it to every proper morphism (Section \bpout{4.5}). We use it in the form of Theorem \hyperref[bp73]{A.25}: for a proper morphism \(X\to S\) there is a proper surjection \(\widetilde X\to X\) with \(\widetilde X\) projective over \(S\).

% object: sec-Compactifications (blueprint/chap-Nagata-compactification-and-Chows-lemma/sec-Compactifications/_section.md), content hidden
\subsection*{\texorpdfstring{A.1\quad Compactifications}{A.1 Compactifications}}
\phantomsection\label{bp5}%
\addcontentsline{toc}{subsection}{\texorpdfstring{A.1\quad Compactifications}{A.1 Compactifications}}
% text: prose-A-02
\phantomsection\label{bp6}%
To show that \(Rp_*j_!\) does not depend on the compactification, we must compare two of them, and in general neither maps to the other. The plane \(\mathbf A^2\) is open in \(\mathbf P^2\) and in \(\mathbf P^1\times\mathbf P^1\), but the identity of \(\mathbf A^2\) extends to a morphism in neither direction: the map \((x:y:z)\mapsto((x:z),(y:z))\) is undefined at \((1:0:0)\) and \((0:1:0)\), and its inverse is undefined at \((\infty,\infty)\). What exists is a third compactification mapping to both, here the closure of \(\mathbf A^2\) in \(\mathbf P^2\times\mathbf P^1\times\mathbf P^1\). So we compare two compactifications through a common refinement. A refinement \(h\) is required to satisfy \(h^{-1}(X)=X\). This is what lets us compute \(Rh_*j'_!\) by restriction to \(X\) in the proof of Proposition \bpout{5.26}.

Without density, refinements are not unique. For example, the point \(X=\operatorname{Spec}k\) is open in the compactification \(\operatorname{Spec}k\sqcup\mathbf P^1_k\), which is not dense. Every automorphism of \(\mathbf P^1_k\) gives a refinement of it to itself. Dense compactifications do not have this problem: morphisms between them are unique, and with their refinements they form a cofiltered category (Proposition \hyperref[bp21]{A.7}).

% object: definition-compactification-refinement (blueprint/chap-Nagata-compactification-and-Chows-lemma/sec-Compactifications/definition-compactification-refinement/_section.md)
\phantomsection\label{bp7}%
\begin{bpstatement}{Definition~A.1}{ (Compactification and refinement)}
Let \(f:X\to S\) be a morphism. A compactification is a factorization

\[
X\xrightarrow{j}\overline X\xrightarrow{p}S
\]

with \(j\) a quasi-compact open immersion and \(p\) proper. We say that it is \emph{dense} if \(j\) is scheme-theoretically dense, i.e., if the map \(\mathcal O_{\overline X}\to j_*\mathcal O_X\) is injective. A refinement of compactifications is an \(S\)-morphism \(h:\overline X'\to\overline X\) with \(hj'=j\) and \(h^{-1}(X)=X\), under the specified identifications with \(X\).
\end{bpstatement}

% text: prose-A-03
\phantomsection\label{bp8}%
For example, \(\mathbf A^1_k\subset\mathbf P^1_k\) is a dense compactification of \(\mathbf A^1_k\to\operatorname{Spec}k\), since a regular function on an open of \(\mathbf P^1_k\) that vanishes on \(\mathbf A^1_k\) is zero. The theorem says that such a factorization always exists.

% object: theorem-nagata-s-compactification-theorem (blueprint/chap-Nagata-compactification-and-Chows-lemma/sec-Compactifications/theorem-nagata-s-compactification-theorem/_section.md)
\phantomsection\label{bp9}%
\begin{bpstatementit}{Theorem~A.2}{ (Nagata's compactification theorem)}
Let \(S\) be quasi-compact and quasi-separated. Every separated morphism \(f:X\to S\) of finite type has a dense compactification.
\end{bpstatementit}

% text: prose-A-04
\phantomsection\label{bp10}%
This is [Stacks, \href{https://stacks.math.columbia.edu/tag/0F41}{0F41}], and we give the proof in Section \hyperref[bp65]{A.4}. We need it in this generality because proper base change is proved over arbitrary qcqs bases; for schemes of finite type over \(k\) the Noetherian case of the proof suffices.

Density is a condition on structure sheaves. For instance, a nilpotent section supported on the boundary is invisible on \(X\), but it prevents \(X\) from being schematically dense. The condition does two things for us. First, morphisms between dense compactifications are unique and common refinements exist (Proposition \hyperref[bp21]{A.7}). Second, passing to a dense compactification does not change the direct image with proper support (Lemma \hyperref[bp17]{A.5}).

% object: lemma-schematic-closure-uniqueness (blueprint/chap-Nagata-compactification-and-Chows-lemma/sec-Compactifications/lemma-schematic-closure-uniqueness/_section.md)
\phantomsection\label{bp11}%
\begin{bpstatementit}{Lemma~A.3}{ (Schematic closure)}
Let \(i:U\to T\) be a quasi-compact immersion, with \(T\) quasi-separated. Its scheme-theoretic image \(\overline U\subseteq T\) is a closed subscheme, \(U\to\overline U\) is an open immersion, and \(U\) is scheme-theoretically dense in \(\overline U\). The formation of \(\overline U\) commutes with flat base change \(T'\to T\), in particular with restriction to opens of \(T\).
\end{bpstatementit}

% proof: proof/lemma-schematic-closure-uniqueness
\begin{proof}[{Proof}]\phantomsection\label{bp12}%
The scheme-theoretic image is the smallest closed subscheme through which \(i\) factors. Its ideal is \(\ker(\mathcal O_T\to i_*\mathcal O_U)\), which is quasi-coherent because \(i\) is quasi-compact and quasi-separated. Its formation commutes with flat base change \(T'\to T\): pushforward along the quasi-compact quasi-separated \(i\) commutes with flat base change, and pullback along a flat morphism is exact, so the pullback of the ideal is the ideal of the base change [Stacks, \href{https://stacks.math.columbia.edu/tag/081I}{081I}]. Restriction to an open of \(T\) is a special case. On an open where the immersion is closed the image is therefore \(U\) itself, and this shows that \(U\) is open in its image. By construction, a section of \(\mathcal O_{\overline U}\) vanishing on \(U\) is zero, which is the definition of schematic density. For these statements see [Stacks, \href{https://stacks.math.columbia.edu/tag/01R8}{01R8}, \href{https://stacks.math.columbia.edu/tag/01RG}{01RG}].
\end{proof}

% text: prose-A-05
\phantomsection\label{bp13}%
Density is what makes extensions from \(U\) unique.

% object: lemma-uniqueness-maps-dense-open (blueprint/chap-Nagata-compactification-and-Chows-lemma/sec-Compactifications/lemma-uniqueness-maps-dense-open/_section.md)
\phantomsection\label{bp14}%
\begin{bpstatementit}{Lemma~A.4}{ (Uniqueness of extensions from a dense open)}
Let \(T\) be an \(S\)-scheme, let \(U\subseteq T\) be a scheme-theoretically dense open, and let \(Y\to S\) be separated. Two \(S\)-maps \(T\to Y\) that agree on \(U\) are equal.
\end{bpstatementit}

% proof: proof/lemma-uniqueness-maps-dense-open
\begin{proof}[{Proof}]\phantomsection\label{bp15}%
The equalizer of the two maps is a closed subscheme of \(T\), since it is the inverse image of the diagonal of \(Y/S\). Its ideal restricts to zero on \(U\), because the two maps agree there. As a section of \(\mathcal O_T\) that vanishes on a scheme-theoretically dense open is zero, the ideal is zero and the equalizer is all of \(T\).
\end{proof}

% text: prose-A-06
\phantomsection\label{bp16}%
By Lemma \hyperref[bp11]{A.3}, every compactification contains a dense one, the schematic closure of \(X\). For this replacement to be harmless we need that it does not change \(Rp_*j_!\).

% object: lemma-passing-schematic-closure (blueprint/chap-Nagata-compactification-and-Chows-lemma/sec-Compactifications/lemma-passing-schematic-closure/_section.md)
\phantomsection\label{bp17}%
\begin{bpstatementit}{Lemma~A.5}{ (Passing to the schematic closure)}
Let \(X\xrightarrow{j}\overline X\xrightarrow{p}S\) be any compactification with \(\overline X\) Noetherian, and let \(i:Z\hookrightarrow\overline X\) be the schematic closure of \(X\), with \(j_Z:X\hookrightarrow Z\) the induced open immersion. For a condensed coefficient ring \(\Lambda\), we use the pullbacks of its coefficient sheaf to these schemes as coefficients. For every \(K\in D(X_{\mathrm{pro}\acute{e}t},\Lambda_X)\) there are canonical identifications

\[
j_!K\simeq i_*j_{Z,!}K,\qquad
Rp_*j_!K\simeq R(p\circ i)_*j_{Z,!}K.
\]
\end{bpstatementit}

% proof: proof/lemma-passing-schematic-closure
\begin{proof}[{Proof}]\phantomsection\label{bp18}%
The first identification is a statement about supports, and the second follows from it by applying \(Rp_*\). The only point that needs care on the pro-étale site is the exactness of \(i_*\).

Since \(j_!K\) vanishes outside \(X\subseteq Z\), it should be pushed forward from \(Z\). To say this in the derived category we need \(i_*\) to be exact, so that \(Ri_*=i_*\). By Proposition \bpout{2.58}, the functors attached to open and closed immersions on the pro-étale site are exact when the complementary open is quasi-compact. Since \(\overline X\) is Noetherian, the complement of \(Z\) is quasi-compact, so all functors below are exact and may be applied degreewise to unbounded complexes. The inverse image of \(j_!K\) on \(Z\) is \(j_{Z,!}K\), so we have the natural map \(j_!K\to i_*i^*j_!K=i_*j_{Z,!}K\). It is the identity on the open \(X\), and both sides restrict to zero on the closed complement of \(X\). By the localization sequence for an open and its closed complement, this map is therefore an isomorphism. The second identification then follows by applying \(Rp_*\) and using \(Rp_*Ri_*\simeq R(p\circ i)_*\).
\end{proof}

% object: remark-non-noetherian-compactifications (blueprint/chap-Nagata-compactification-and-Chows-lemma/sec-Compactifications/remark-non-noetherian-compactifications/_section.md)
\phantomsection\label{bp19}%
\begin{bpstatement}{Remark~A.6}{ (Non-Noetherian compactifications)}
The same proof works for qcqs \(\overline X\) whenever \(\overline X\setminus Z\) is quasi-compact. For a non-Noetherian compactification this can fail, and then Proposition \bpout{2.58} does not apply to pushforward along the closed immersion on the pro-étale site. The schematic closure and dense refinements exist in the generality of Theorem \hyperref[bp9]{A.2}.
\end{bpstatement}

% text: prose-A-07
\phantomsection\label{bp20}%
For dense compactifications the comparison is formal. We obtain common refinements by Lemma \hyperref[bp11]{A.3}, and morphisms are unique by Lemma \hyperref[bp14]{A.4}.

% object: proposition-dense-compactifications-form-cofiltered-category (blueprint/chap-Nagata-compactification-and-Chows-lemma/sec-Compactifications/proposition-dense-compactifications-form-cofiltered-category/_section.md)
\phantomsection\label{bp21}%
\begin{bpstatementit}{Proposition~A.7}{ (Dense compactifications form a cofiltered category)}
Let \(S\) be quasi-compact and quasi-separated, and let \(f:X\to S\) be separated and of finite type. Any finite family of dense compactifications of \(f\) has a common dense refinement. Morphisms between dense compactifications that restrict to the identity on \(X\) are unique, proper, and automatically satisfy \(h^{-1}(X)=X\). The category of dense compactifications of \(f\) is therefore cofiltered.
\end{bpstatementit}

% proof: proof/proposition-dense-compactifications-form-cofiltered-category
\begin{proof}[{Proof}]\phantomsection\label{bp22}%
A common refinement of two compactifications is the closure of \(X\) in their fibre product, as in the example of \(\mathbf A^2\) above. The properties of morphisms then follow from density.

Given two compactifications, let \(Z\) be the schematic closure of the image of \(X\) in \(\overline X_1\times_S\overline X_2\) under the diagonal map. The map from \(X\) is a quasi-compact immersion, because over the open \(X\subseteq\overline X_1\) it is the graph of a morphism to \(\overline X_2\), and \(\overline X_2\) is separated over \(S\). By Lemma \hyperref[bp11]{A.3}, \(X\) is an open subscheme of \(Z\) and is scheme-theoretically dense in \(Z\). Since the product is proper over \(S\), the scheme \(Z\) and its projections are proper.

The projections are refinements, that is, \(X\) is its own inverse image in \(Z\). By Lemma \hyperref[bp11]{A.3} the formation of the schematic closure commutes with restriction to opens. Over \(X\subseteq\overline X_1\) we may therefore compute it on the open \(X\times_S\overline X_2\). There the immersion is the graph of \(X\to\overline X_2\), which is closed, so the closure is this graph. This means that the inverse image of \(X\) under the first projection is \(X\), on which the first projection is the identity. The same holds for the second projection. The assertion for finite families follows by induction.

Let \(h:\overline X'\to\overline X\) be a morphism of dense compactifications restricting to the identity on \(X\). To see that \(h\) is proper, note that it factors through its graph. This graph is closed in \(\overline X'\times_S\overline X\) because \(\overline X\) is separated over \(S\), and the projection of this product to \(\overline X\) is proper because \(\overline X'\) is proper over \(S\).

Next, the inclusion \(X\to h^{-1}(X)\) is a section of the separated map \(h^{-1}(X)\to X\). It is therefore closed, as well as open. Since \(X\) is schematically dense in \(\overline X'\), it is schematically dense in the open \(h^{-1}(X)\). A closed subscheme containing it must be all of \(h^{-1}(X)\), and we conclude that \(h^{-1}(X)=X\). Two parallel morphisms are equal by Lemma \hyperref[bp14]{A.4}. Finally, the category is nonempty by Nagata's theorem. Since any two objects have a common refinement and parallel morphisms are equal, it is cofiltered.
\end{proof}

% text: prose-A-08
\phantomsection\label{bp23}%
Proper base change is applied after a change of base \(S'\to S\), so we also need to know what happens to a compactification under base change, where density need not survive.

% object: remark-base-change-compactification (blueprint/chap-Nagata-compactification-and-Chows-lemma/sec-Compactifications/remark-base-change-compactification/_section.md)
\phantomsection\label{bp24}%
\begin{bpstatement}{Remark~A.8}{ (Base change of a compactification)}
For every \(S'\to S\), the base change of a compactification is an open immersion followed by a proper morphism, which suffices to define the direct image with compact support of \(X\times_SS'\to S'\). When the base-changed ambient scheme is Noetherian, as for schemes of finite type over \(k\), Lemma \hyperref[bp17]{A.5} shows that this direct image agrees with the one defined by the schematic closure.

Schematic closure commutes with flat base change (Lemma \hyperref[bp11]{A.3}), and so does schematic density, but density can be lost under an arbitrary base change. In that case the schematic closure of \(X\times_SS'\) in \(\overline X\times_SS'\) is a dense compactification, and a base-changed common refinement can be replaced by its closure in the same way. The comparisons in the main text are made on a further common refinement, where they are canonical by uniqueness.
\end{bpstatement}


% object: sec-Elimination-of-indeterminacy (blueprint/chap-Nagata-compactification-and-Chows-lemma/sec-Elimination-of-indeterminacy/_section.md), content hidden
\subsection*{\texorpdfstring{A.2\quad Elimination of indeterminacy}{A.2 Elimination of indeterminacy}}
\phantomsection\label{bp25}%
\addcontentsline{toc}{subsection}{\texorpdfstring{A.2\quad Elimination of indeterminacy}{A.2 Elimination of indeterminacy}}
% text: prose-A-09
\phantomsection\label{bp26}%
Every modification in the proof of Nagata's theorem is a blow-up whose center avoids a fixed open, so that this open is not changed. Here the open \(U\) is the scheme being compactified, or a piece of it, and only the boundary may be modified. For instance, the coordinate axes minus the origin are disjoint closed subsets of \(U=\mathbf A^2\setminus\{0\}\), but their closures in \(\mathbf A^2\) meet at the origin. The blow-up of the origin is an isomorphism over \(U\), and in it the strict transforms of the axes meet the exceptional line in the two distinct points of slope \(0\) and \(\infty\), so they become disjoint. Lemma \hyperref[bp55]{A.20} proves the same in general.

We use blow-ups rather than arbitrary proper modifications that are isomorphisms over \(U\), because a blow-up is given by an ideal. This makes them easy to combine. Two of them are dominated by the blow-up of the product ideal, and a composite is again a blow-up (Lemma \hyperref[bp35]{A.13}). The ideal is required to be of finite type, so that the blow-up is proper.

% object: definition-admissible-blow-up (blueprint/chap-Nagata-compactification-and-Chows-lemma/sec-Elimination-of-indeterminacy/definition-admissible-blow-up/_section.md)
\phantomsection\label{bp27}%
\begin{bpstatement}{Definition~A.9}{ (Admissible blow-up)}
Let \(T\) be quasi-compact and quasi-separated, and let \(U\subseteq T\) be a quasi-compact open. For a quasi-coherent ideal of finite type \(\mathcal I\subseteq\mathcal O_T\) with \(\mathcal I|_U=\mathcal O_U\), the blow-up

\[
b:\operatorname{Bl}_{\mathcal I}(T)=
\operatorname{Proj}_T\Bigl(\bigoplus_{m\geq0}\mathcal I^m\Bigr)
\longrightarrow T
\]

is called \(U\)-admissible. Its center is disjoint from \(U\).
\end{bpstatement}

% text: prose-A-10
\phantomsection\label{bp28}%
On an affine open \(\operatorname{Spec}(R)\), write \(I=(a_1,\ldots,a_r)\). The charts of the blow-up have rings

\[
R[I/a_i]=\bigl((\bigoplus_{m\geq0}I^mt^m)_{a_it}\bigr)_0.
\]

This is the degree-zero part of the localization at \(a_it\) of the Rees algebra, where \(t\) is a variable of degree one. It is the quotient of \(R[x_1,\ldots,x_r]/(a_ix_j-a_j)_j\) by its \(a_i\)-power torsion, that is, by the elements killed by some power of \(a_i\), and this torsion can be nonzero. In the chart \(R[I/a_i]\), the element \(a_i\) is a nonzerodivisor and generates the extended ideal \(IR[I/a_i]\). In other words, the center becomes an effective Cartier divisor. The blow-up is an isomorphism away from the center. Since \(I\) is finitely generated, the blow-up of an affine base is a closed subscheme of a projective space over it, and therefore \(b\) is proper. These constructions glue under restriction to opens.

We first use admissible blow-ups to restore density. A compactification may contain components or embedded components supported in the boundary. Blowing up the boundary removes them, because after the blow-up the boundary is cut out locally by a nonzerodivisor.

% object: lemma-making-open-schematically-dense (blueprint/chap-Nagata-compactification-and-Chows-lemma/sec-Elimination-of-indeterminacy/lemma-making-open-schematically-dense/_section.md)
\phantomsection\label{bp29}%
\begin{bpstatementit}{Lemma~A.10}{ (Making an open schematically dense)}
Let \(T\) be quasi-compact and quasi-separated, and let \(U\subseteq T\) be a quasi-compact open. There is a \(U\)-admissible blow-up \(T'\to T\) such that the inverse image of \(U\) is schematically dense in \(T'\) and its complement is the support of an effective Cartier divisor.
\end{bpstatementit}

% proof: proof/lemma-making-open-schematically-dense
\begin{proof}[{Proof}]\phantomsection\label{bp30}%
We blow up an ideal that cuts out the complement of \(U\). Its pullback is then locally generated by a nonzerodivisor \(a\), and an open whose complement is cut out by a nonzerodivisor is schematically dense, because a section killed by a power of \(a\) is zero.

Choose a quasi-coherent ideal \(\mathcal I\) of finite type with \(V(\mathcal I)=T\setminus U\). Such an ideal exists because \(U\) is a quasi-compact open in a qcqs scheme [Stacks, \href{https://stacks.math.columbia.edu/tag/01PH}{01PH}]. We blow up this ideal, and the extended ideal defines the exceptional Cartier divisor. On the affine charts its local equation \(a\) is a nonzerodivisor. Hence each chart ring injects into its localization at \(a\). The complement of the divisor is the inverse image of \(U\), and it is therefore schematically dense.
\end{proof}

% object: example-blow-up-embedded-point (blueprint/chap-Nagata-compactification-and-Chows-lemma/sec-Elimination-of-indeterminacy/example-blow-up-embedded-point/_section.md)
\phantomsection\label{bp31}%
\begin{bpstatement}{Example~A.11}{ (Blowing up removes an embedded point)}
Let \(T=\operatorname{Spec}R\) with \(R=k[x,y]/(xy,y^2)\), the line \(y=0\) with an embedded point at the origin, and let \(U=D(x)\). On \(U\) the element \(y\) is zero, since \(xy=0\) and \(x\) is invertible, so \(U\) is not schematically dense in \(T\). On the other hand, the ideal \(I=(x)\) is the unit ideal on \(U\), so its blow-up is \(U\)-admissible. Since the ideal is principal, the chart formula above gives \(R[I/x]=R/\{b:x^nb=0\text{ for some }n\}\). Now every element of \(R\) is uniquely \(f(x)+cy\) with \(c\in k\), and \(x^n(f+cy)=x^nf\), so the \(x\)-power torsion is \(ky\) and the blow-up is \(\operatorname{Spec}k[x]\). Thus the embedded point is gone, \(U\) is dense, and the boundary is the Cartier divisor \(x=0\), as in Lemma \hyperref[bp29]{A.10}.
\end{bpstatement}

% text: prose-A-11
\phantomsection\label{bp32}%
After a blow-up, we replace a scheme over \(T\) by the part of its base change that is not concentrated on the exceptional divisor.

% object: definition-strict-transform (blueprint/chap-Nagata-compactification-and-Chows-lemma/sec-Elimination-of-indeterminacy/definition-strict-transform/_section.md)
\phantomsection\label{bp33}%
\begin{bpstatement}{Definition~A.12}{ (Strict transform)}
Let \(b:T'\to T\) be the blow-up of \(T\) in a quasi-coherent ideal of finite type, with exceptional divisor \(E\subseteq T'\), and let \(W\to T\) be a morphism. The \emph{strict transform} of \(W\) is the closed subscheme of \(W\times_TT'\) whose ideal consists of the local sections of \(\mathcal O_{W\times_TT'}\) supported on the inverse image of \(E\).
\end{bpstatement}

% text: prose-A-12
\phantomsection\label{bp34}%
Over a chart of \(T'\) where \(E\) has equation \(a\) and \(W\times_TT'\) is affine with ring \(B\), the ring of the strict transform is

\[
B/\{b\in B:\ a^nb=0\text{ for some }n\geq0\}.
\]

If \(U\) is the complement of the center, this is exactly the schematic closure of \(W_U\) in \(W\times_TT'\), the object that occurs in the elimination of indeterminacy. By [Stacks, \href{https://stacks.math.columbia.edu/tag/080E}{080E}], the strict transform is the blow-up of \(W\) in the pulled-back ideal, so it is proper over \(W\).

In the proofs below we construct blow-ups locally and one after another. By the next lemma, finitely many such constructions combine into a single \(U\)-admissible blow-up of \(T\).

% object: lemma-products-composites-u-admissible-blow (blueprint/chap-Nagata-compactification-and-Chows-lemma/sec-Elimination-of-indeterminacy/lemma-products-composites-u-admissible-blow/_section.md)
\phantomsection\label{bp35}%
\begin{bpstatementit}{Lemma~A.13}{ (Products and composites of $U$-admissible blow-ups)}
Finitely many \(U\)-admissible blow-ups of \(T\) admit a common \(U\)-admissible blow-up. A composite of two \(U\)-admissible blow-ups is again such a blow-up. A \(U\)-admissible blow-up on a quasi-compact open \(V\subseteq T\) containing \(U\) extends to a \(U\)-admissible blow-up of \(T\).
\end{bpstatementit}

% proof: proof/lemma-products-composites-u-admissible-blow
\begin{proof}[{Proof}]\phantomsection\label{bp36}%
All three assertions are statements about ideals, because products of ideals dominate blow-ups, composites are computed by products of ideals, and extension from an open amounts to extending a quasi-coherent ideal of finite type. Let \(\mathcal I,\mathcal J\) be ideals, both equal to \(\mathcal O\) on \(U\), so that their product is too. Then the blow-up of their product factors through both blow-ups. Indeed, on it the product \(\mathcal I\mathcal J\) is invertible, hence so is each factor, and the universal property of each blow-up then gives the factorization. The first assertion follows by induction. For the second, blowing up first \(\mathcal I\) and then the pullback of \(\mathcal J\) yields the blow-up of \(\mathcal I\mathcal J\) [Stacks, \href{https://stacks.math.columbia.edu/tag/080A}{080A}], which proves the composition assertion when the second center is pulled back from \(T\). In general, however, the second center lives only on the first blow-up, and the composite is a blow-up of \(T\) by the calculation with ideals of finite type in [Stacks, \href{https://stacks.math.columbia.edu/tag/080B}{080B}]. For the extension assertion we extend a quasi-coherent ideal of finite type from a quasi-compact open, keeping the center outside \(U\); such an extension exists because on a qcqs scheme every quasi-coherent sheaf is the union of its quasi-coherent subsheaves of finite type, and the precise result on admissible extensions is [Stacks, \href{https://stacks.math.columbia.edu/tag/080M}{080M}].
\end{proof}

% text: prose-A-13
\phantomsection\label{bp37}%
Blow-ups commute with flat base change [Stacks, \href{https://stacks.math.columbia.edu/tag/0805}{0805}], because flat pullback preserves the inclusions \(\mathcal I^m\subseteq\mathcal O_T\). Under a base change that is not flat, the pullbacks of these inclusions can have kernels. In that case the blow-up of the pulled-back ideal is the strict transform of the base change, which can be strictly smaller than the base change itself.

The gluing argument has to extend a morphism defined on a dense open \(U\subseteq X\) as far as possible. When the target \(Y\) is not proper, the extension can stop before reaching all of \(X\). We need a way to say that a partial map is already extended as far as it can be.

Take \(X=Y=\mathbf A^1_k\) and \(U=\mathbf A^1\setminus\{0\}\). The inclusion \(t\mapsto t\) extends over \(0\). Its graph, the diagonal minus the origin, is not closed in \(X\times Y\). The map \(t\mapsto1/t\) does not extend, even after a blow-up, since blowing up a point on a regular curve does nothing. Its graph \(V(ty-1)\) is closed. In general, a closed graph means that no limit point lies over \(X\setminus U\). A map with closed graph is therefore not the restriction of a map on a larger open. Both the hypothesis and the conclusion of Theorem \hyperref[bp43]{A.16} are stated in this form. For proper \(Y\) the condition forces the domain to be all of \(X\), which gives Corollary \hyperref[bp45]{A.17}.

% object: definition-quasi-domination (blueprint/chap-Nagata-compactification-and-Chows-lemma/sec-Elimination-of-indeterminacy/definition-quasi-domination/_section.md)
\phantomsection\label{bp38}%
\begin{bpstatement}{Definition~A.14}{ (Quasi-domination)}
Let \(X\) and \(Y\) be \(S\)-schemes, with \(Y\to S\) separated. A \emph{quasi-domination of \(X\) over \(Y\)} consists of a dense open \(U\subseteq X\) and an \(S\)-map \(f:U\to Y\) whose graph is closed in \(X\times_SY\).
\end{bpstatement}

% text: prose-A-14
\phantomsection\label{bp39}%
Since \(Y\to S\) is separated, the graph is always closed in \(U\times_SY\); the condition is about its behaviour over \(X\setminus U\). Every map \(X\to Y\) is a quasi-domination with domain \(X\). Conversely, if \(Y\to S\) is proper, every quasi-domination has domain \(X\), because the projection of the closed graph to \(X\) is closed with image the dense open \(U\).

The extension theorem below is proved first for targets that embed in affine space over the base. We reduce a general target to this form by Zariski's Main Theorem.

% object: proposition-relative-form-zariski-s-main (blueprint/chap-Nagata-compactification-and-Chows-lemma/sec-Elimination-of-indeterminacy/proposition-relative-form-zariski-s-main/_section.md)
\phantomsection\label{bp40}%
\begin{bpstatementit}{Proposition~A.15}{ (The relative form of Zariski's Main Theorem)}
A separated, quasi-compact, locally quasi-finite morphism \(Z\to T\) locally of finite type factors as an open immersion \(Z\hookrightarrow T_0\) followed by an affine morphism \(T_0\to T\). Consequently \(Z\to T\) is quasi-affine.
\end{bpstatementit}

% proof: proof/proposition-relative-form-zariski-s-main
\begin{proof}[{Proof}]\phantomsection\label{bp41}%
The scheme \(T_0\) is the normalization of \(T\) in \(Z\), that is, the relative spectrum of the integral closure of \(\mathcal O_T\) in the direct image of \(\mathcal O_Z\); it is affine over \(T\) by construction. By Zariski's Main Theorem in the form [Stacks, \href{https://stacks.math.columbia.edu/tag/02LR}{02LR}], locally on \(T\) the canonical map from \(Z\) to the normalization \(T_0\) of \(T\) in \(Z\) is an open immersion, and \(T_0\to T\) is affine. These factorizations agree on overlaps and glue to the stated factorization. Over an affine open of \(T\) the scheme \(T_0\) is affine, so \(Z\) is a quasi-compact open of an affine scheme there, which means that \(Z\to T\) is quasi-affine.
\end{proof}

% text: prose-A-15
\phantomsection\label{bp42}%
Quasi-affine targets are the ones where a partial section can be extended by hand, by clearing denominators of finitely many coordinates. Using the proposition, we reduce the general extension theorem to this case, once we know that the relevant closure is quasi-finite over the base.

% object: theorem-extension-partial-section (blueprint/chap-Nagata-compactification-and-Chows-lemma/sec-Elimination-of-indeterminacy/theorem-extension-partial-section/_section.md)
\phantomsection\label{bp43}%
\begin{bpstatementit}{Theorem~A.16}{ (Extension of a partial section)}
Let \(T\) be qcqs, let \(U\subseteq V\subseteq T\) be quasi-compact opens with \(U\) dense in \(T\), and let \(p:Y\to T\) be separated and of finite type. Suppose that \(f:U\to Y\) is a section of \(p\) over \(U\) and that its image is closed in \(p^{-1}(V)\). There are a \(V\)-admissible blow-up \(\pi:T'\to T\), an open \(W\subseteq T'\) containing \(\pi^{-1}(U)\), and a map \(g:W\to Y\) such that

\[
p\circ g=\pi|_W,\qquad
g|_{\pi^{-1}(U)}=f\circ\pi|_{\pi^{-1}(U)},
\]

and the graph of \(g\) is closed in \(Y\times_TT'\).
\end{bpstatementit}

% text: prose-A-16
\phantomsection\label{bp44}%
We quote this theorem from [\hyperref[bp348]{Con07}, Theorem 2.4], with the auxiliary base equal to \(T\), and explain below the structure of its proof. Note that \(Y\) is qcqs, since \(p\) is separated and of finite type and \(T\) is qcqs. The hypothesis on \(f\) says that \(f\) is a quasi-domination of \(V\) over \(Y\). We apply the theorem only when \(V=U\), the case of [\hyperref[bp348]{Con07}, Remark 2.5]. Then the hypothesis holds automatically, because a section of a separated morphism is a closed immersion.

The proof first treats a quasi-affine \(Y\to T\). If \(Y\to T\) is quasi-affine and \(T\) is affine, then \(Y\) is a quasi-affine scheme of finite type over \(T\), and it embeds in some \(\mathbf A^n_T\) [Stacks, \href{https://stacks.math.columbia.edu/tag/04II}{04II}]. The coordinates of \(f\) are then finitely many fractions, and a blow-up clearing their denominators extends \(f\). One first makes \(U\) schematically dense, so that the extended coordinates still satisfy the equations of \(Y\). Using a finite affine cover of \(T\), we then obtain the quasi-affine case over a qcqs base. In fact the proof needs this case for opens \(U\subseteq V\) with \(U\) not necessarily dense, because the reduction below applies it to the opens \(f^{-1}(Y_i)\subseteq V\), which need not be dense. For this, a further construction in Conrad's proof keeps the centers outside \(V\).

The general case reduces to the quasi-affine one, and this reduction is where Zariski's Main Theorem enters. The plan is to extend the section on each affine open of \(Y\) separately. The closure of the section is then quasi-finite over the base, hence quasi-affine, and the quasi-affine case applies once more.

In detail, choose a finite affine open cover \(Y=\bigcup_iY_i\). Since \(T\) is quasi-separated, each \(Y_i\to T\) is quasi-compact, and as \(Y_i\) is affine it is quasi-affine. Hence the quasi-affine case applies to each \(f|_{f^{-1}(Y_i)}\), and we take a common \(V\)-admissible blow-up \(T_1\to T\). Let \(Z\) be the schematic closure of the image of \(f\) in \(Y\times_TT_1\). It meets each \(Y_i\times_TT_1\) inside the closed graph of an extension, which has at most one point over each point of \(T_1\). Since there are finitely many \(Y_i\), the fibres of \(Z\to T_1\) are finite. As \(Z\to T_1\) is separated and of finite type, it is quasi-finite, and it is quasi-affine by Proposition \hyperref[bp40]{A.15}. The quasi-affine case then extends the section into \(Z\), and the composite of the two blow-ups is \(V\)-admissible by Lemma \hyperref[bp35]{A.13}.

% object: corollary-extension-proper-target (blueprint/chap-Nagata-compactification-and-Chows-lemma/sec-Elimination-of-indeterminacy/corollary-extension-proper-target/_section.md)
\phantomsection\label{bp45}%
\begin{bpstatementit}{Corollary~A.17}{ (Extension to a proper target)}
Let \(X\) be qcqs over \(S\), let \(U\subseteq X\) be a dense quasi-compact open and let \(Y\to S\) be proper. Every \(S\)-map \(f:U\to Y\) extends to an \(S\)-map \(X'\to Y\) for some \(U\)-admissible blow-up \(X'\to X\).
\end{bpstatementit}

% proof: proof/corollary-extension-proper-target
\begin{proof}[{Proof}]\phantomsection\label{bp46}%
A map \(U\to Y\) is the same as a section over \(U\) of the projection \(Y\times_SX\to X\). The theorem extends this section to an open of a blow-up, and properness of \(Y\) forces this open to be everything.

Apply Theorem \hyperref[bp43]{A.16}, with \(T=X\) and \(V=U\), to the section \((f,\mathrm{id}):U\to Y\times_SX\) of the proper projection to \(X\). This gives an extension on an open \(W\subseteq X'\) whose graph is closed in \(Y\times_SX'\). The open \(W\) is dense, because it contains the inverse image of \(U\), which is dense in \(X'\). Since \(Y\to S\) is proper, the observation after Definition \hyperref[bp38]{A.14} shows that \(W=X'\).
\end{proof}

% text: prose-A-17
\phantomsection\label{bp47}%
When the target is not proper, the extension is defined only on an open of the blow-up, and we can still say what it looks like. Suppose the target is a separated map \(W\to T\) of finite type that is an isomorphism over \(U\), and the section is its inverse over \(U\). Then the extension identifies a blow-up of \(W\) with an open of a blow-up of \(T\).

% object: theorem-elimination-indeterminacy (blueprint/chap-Nagata-compactification-and-Chows-lemma/sec-Elimination-of-indeterminacy/theorem-elimination-indeterminacy/_section.md)
\phantomsection\label{bp48}%
\begin{bpstatementit}{Theorem~A.18}{ (Elimination of indeterminacy)}
Let \(T\) be qcqs, let \(U\subseteq T\) be a quasi-compact open, and let \(W\to T\) be separated and of finite type, with \(W_U\to U\) an isomorphism. There is a \(U\)-admissible blow-up \(T'\to T\) such that the inverse image of \(U\) is schematically dense in \(T'\) and the strict transform \(W'\to T'\) of \(W\to T\) is an open immersion.
\end{bpstatementit}

% proof: proof/theorem-elimination-indeterminacy
\begin{proof}[{Proof}]\phantomsection\label{bp49}%
This is [Stacks, \href{https://stacks.math.columbia.edu/tag/081S}{081S}]. The idea is to view the inverse of \(W_U\to U\) as a partial section of \(W\to T\) and to extend it by Theorem \hyperref[bp43]{A.16}; the strict transform is then the graph of the extension.

Choose a finite type ideal \(\mathcal I_0\) cutting out \(T\setminus U\). Its blow-up \(T_1\to T\) makes \(U\) schematically dense by Lemma \hyperref[bp29]{A.10}. Over \(T_1\), the base change of \(W\to T\) to \(T_1\) is separated and of finite type, and the inverse of \(W_U\to U\) is a section of it over \(U\). We may therefore apply Theorem \hyperref[bp43]{A.16} over \(T_1\), with \(V=U\), which gives a further blow-up and an extension of this section, with closed graph, over an open containing the inverse image of \(U\). By Lemma \hyperref[bp35]{A.13} the composite is a \(U\)-admissible blow-up \(\operatorname{Bl}_{\mathcal I}(T)\). Let \(b:T'\to T\) be the blow-up of \(\mathcal I\mathcal I_0\), which dominates both. Pulling back to \(T'\) gives an open \(W_0\subseteq T'\) containing \(b^{-1}(U)\) and a map \(W_0\to W\) over \(T\) whose graph \(\Gamma\) is closed in \(W\times_TT'\). As the zero set of \(\mathcal I\mathcal I_0\) is exactly \(T\setminus U\), \(b^{-1}(U)\) is schematically dense in \(T'\), hence in \(W_0\).

We included the factor \(\mathcal I_0\) so that the center of \(b\) is exactly \(T\setminus U\). Then the strict transform \(W'\) is the schematic closure of \(W_U\) in \(W\times_TT'\). Over \(U\), the subscheme \(W_U\) is the graph of the original section, so it lies in \(\Gamma\) and \(W'\subseteq\Gamma\). Conversely, \(\Gamma\simeq W_0\) contains the inverse image of \(U\) as a schematically dense open. Hence the schematic closure of \(W_U\) in \(\Gamma\) is \(\Gamma\). We conclude that \(W'=\Gamma\simeq W_0\), and \(W'\to T'\) is an open immersion.
\end{proof}

% text: prose-A-18
\phantomsection\label{bp50}%
The blow-up is not uniquely determined. However, on any such modification with \(U\) schematically dense, the graph of the extension is forced to be the closure of the original section. By Lemma \hyperref[bp14]{A.4}, the identifications obtained on common admissible refinements therefore agree.

In the gluing argument below, a compactification of an open \(V\) of \(M\) is related to \(M\) itself by a proper map defined on a larger open of the compactification.

% object: lemma-extending-graph-from-compactification (blueprint/chap-Nagata-compactification-and-Chows-lemma/sec-Elimination-of-indeterminacy/lemma-extending-graph-from-compactification/_section.md)
\phantomsection\label{bp51}%
\begin{bpstatementit}{Lemma~A.19}{ (Extending a graph from a compactification)}
Let \(S\) be qcqs, let \(M\to S\) be separated of finite type, and let \(V\subseteq M\) be a quasi-compact open with a compactification \(V\subseteq Y\) over \(S\). There are a \(V\)-admissible blow-up \(Y'\to Y\), an open \(V\subseteq V'\subseteq Y'\) and a proper map \(\psi:V'\to M\) with \(\psi^{-1}(V)=V\) and \(\psi|_V=\mathrm{id}_V\).
\end{bpstatementit}

% proof: proof/lemma-extending-graph-from-compactification
\begin{proof}[{Proof}]\phantomsection\label{bp52}%
The map \(\psi\) is the projection to \(M\) of the closure of the graph of \(V\to M\). By elimination of indeterminacy this closure becomes an open of a blow-up of \(Y\), and \(\psi\) is proper because the graph is closed.

Let \(G\) be the schematic closure of the graph of \(V\to M\) in \(Y\times_SM\). Its projection to \(Y\) is separated of finite type, and it is an isomorphism over \(V\). By Theorem \hyperref[bp48]{A.18}, after a \(V\)-admissible blow-up \(Y'\to Y\) its strict transform \(G'\) is an open \(V'\subseteq Y'\). This strict transform is the closure of the graph of \(V\to M\) in \(Y'\times_SM\), since both are the closure of the same graph outside the exceptional divisor [Stacks, \href{https://stacks.math.columbia.edu/tag/0F3Y}{0F3Y}]. So \(G'\) is closed in \(Y'\times_SM\), and since the projection of \(Y'\times_SM\) to \(M\) is proper, the map \(V'\simeq G'\to M\) is proper. It remains to identify the inverse image of \(V\). For this, note that over \(V\subseteq M\) the graph of the inclusion \(V\to Y'\) is already closed in \(Y'\times_SV\), because \(Y'\) is separated over \(S\). The inverse image of \(V\) under \(\psi\) is therefore this graph, that is, \(V\) itself. See also [Stacks, \href{https://stacks.math.columbia.edu/tag/0F3Z}{0F3Z}].
\end{proof}


% object: sec-Separation-by-blow-ups (blueprint/chap-Nagata-compactification-and-Chows-lemma/sec-Separation-by-blow-ups/_section.md), content hidden
\subsection*{\texorpdfstring{A.3\quad Separation by blow-ups}{A.3 Separation by blow-ups}}
\phantomsection\label{bp53}%
\addcontentsline{toc}{subsection}{\texorpdfstring{A.3\quad Separation by blow-ups}{A.3 Separation by blow-ups}}
% text: prose-A-19
\phantomsection\label{bp54}%
The inductive step of Nagata's theorem is Proposition \hyperref[bp63]{A.23}. Suppose that \(M=M_1\cup M_2\) with \(M_1\cap M_2\) dense, and that each \(M_i\) has a compactification \(X_i\); we want a compactification of \(M\). The natural candidate glues \(M\) to each \(X_i\) along \(M_i\), after we remove from \(X_i\) the closure of the part lying over \(M\setminus M_i\). For the result to be separated, however, and for every valuation to extend into one of the two pieces, the closures in \(X_i\) of the two disjoint closed subsets over \(M\setminus M_1\) and \(M\setminus M_2\) must be disjoint. They need not be, and this is where Lemma \hyperref[bp55]{A.20} enters: it makes them disjoint by an admissible blow-up, as with the coordinate axes in Section \hyperref[bp25]{A.2}.

We need two other lemmas. Lemma \hyperref[bp57]{A.21} keeps an isomorphism near such a closure, and Lemma \hyperref[bp60]{A.22} glues two modifications into one separated scheme. Both use elimination of indeterminacy from the previous section. Recall that it turns a map that is an isomorphism over a quasi-compact open \(U\) into an open immersion after a \(U\)-admissible blow-up.

% object: lemma-separation-closures (blueprint/chap-Nagata-compactification-and-Chows-lemma/sec-Separation-by-blow-ups/lemma-separation-closures/_section.md)
\phantomsection\label{bp55}%
\begin{bpstatementit}{Lemma~A.20}{ (Separation of closures)}
Let \(T\) be qcqs and let \(V\subseteq T\) be a quasi-compact open. If \(F_1,F_2\subseteq V\) are disjoint constructible closed subsets, there is a \(V\)-admissible blow-up of \(T\) on which their closures are disjoint.
\end{bpstatementit}

% text: prose-A-20
\phantomsection\label{bp56}%
We quote this from [Stacks, \href{https://stacks.math.columbia.edu/tag/0F3V}{0F3V}] and sketch the argument, which works with closed subschemes of \(T\) rather than closed subsets of \(V\). Let \(\mathcal I_1,\mathcal I_2\) be ideals of finite type on \(T\) cutting out closed subschemes \(Z_1,Z_2\) that are disjoint over \(V\) and restrict to \(F_1,F_2\) there. Producing such ideals on all of \(T\) is where constructibility and finite presentation are needed. Since \(\mathcal I_1+\mathcal I_2\) is the unit ideal over \(V\), its blow-up is \(V\)-admissible. Moreover, on each affine chart the extended ideal is generated by a single element of \(\mathcal I_1\) or of \(\mathcal I_2\). Now, on a chart where it is generated by \(f\in\mathcal I_1\), the strict transform of \(Z_1\) is obtained from its inverse image by dividing out the \(f\)-power torsion. Since \(f\) vanishes on \(Z_1\), every element of the ring of the inverse image is killed by \(f\). Hence the strict transform of \(Z_1\) does not meet that chart. The same holds for \(Z_2\) on the other charts, and therefore the strict transforms of \(Z_1\) and \(Z_2\) are disjoint [Stacks, \href{https://stacks.math.columbia.edu/tag/080P}{080P}]. The constructibility needed to produce such ideals causes no trouble for us, because in the Noetherian applications below every closed subset is constructible and its reduced ideal is of finite type.

For the gluing we also need to know that a proper map that is an isomorphism near a closed subset of \(V\) remains one near its closure in \(T\), after a blow-up.

% object: lemma-extending-isomorphism-over-closure-closed (blueprint/chap-Nagata-compactification-and-Chows-lemma/sec-Separation-by-blow-ups/lemma-extending-isomorphism-over-closure-closed/_section.md)
\phantomsection\label{bp57}%
\begin{bpstatementit}{Lemma~A.21}{ (Extending an isomorphism over the closure of a closed subset)}
Let \(q:P\to T\) be a proper morphism of qcqs schemes, let \(V\subseteq T\) be a quasi-compact open, and let \(F\subseteq V\) be closed. Suppose that \(q\) is an isomorphism over a neighbourhood of \(F\) in \(V\). After a \(V\)-admissible blow-up, the strict transform of \(q\) is an isomorphism over a neighbourhood of the closure of \(F\).
\end{bpstatementit}

% proof: proof/lemma-extending-isomorphism-over-closure-closed
\begin{proof}[{Proof}]\phantomsection\label{bp58}%
We separate the closure of \(F\) from the closure of the locus where \(q\) fails to be an isomorphism, and then apply elimination of indeterminacy on the complement of the second closure.

Inside \(V\), the set \(F\) is disjoint from the closed complement of the largest open over which \(q\) is an isomorphism. Since \(V\) is spectral, both closed subsets can be enlarged to disjoint constructible closed subsets. To these we apply Lemma \hyperref[bp55]{A.20} so that their closures in \(T\) become disjoint. Then, on the complement \(T_0\) of the second closure, the restriction of \(q\) is an isomorphism over an open \(V_0\) containing \(F\). After making this open schematically dense, we apply Theorem \hyperref[bp48]{A.18} on \(T_0\). The resulting strict transform is both open and proper over the modification, and hence it is an isomorphism there. By the extension part of Lemma \hyperref[bp35]{A.13}, this blow-up of \(T_0\) extends to a \(V\)-admissible blow-up of \(T\), which is an isomorphism over an open containing the closure of \(F\). We have only sketched the argument; the precise version, including the choice of constructible enlargements, is [Stacks, \href{https://stacks.math.columbia.edu/tag/0F3W}{0F3W}].
\end{proof}

% text: prose-A-21
\phantomsection\label{bp59}%
The last ingredient replaces two separated schemes containing a common open by modifications that glue to a separated scheme.

% object: lemma-gluing-after-admissible-blow-ups (blueprint/chap-Nagata-compactification-and-Chows-lemma/sec-Separation-by-blow-ups/lemma-gluing-after-admissible-blow-ups/_section.md)
\phantomsection\label{bp60}%
\begin{bpstatementit}{Lemma~A.22}{ (Gluing after admissible blow-ups)}
Let \(S\) be qcqs, and let \(U\subseteq T_i\) (\(i=1,2\)) be open immersions of separated \(S\)-schemes of finite type, with \(U\) quasi-compact. There are \(U\)-admissible blow-ups \(T_i'\to T_i\) and open immersions \(T_i'\to T\) agreeing on \(U\), with \(T\) separated and of finite type over \(S\). If \(U\) is dense in both \(T_i\), then it is dense in \(T\).
\end{bpstatementit}

% proof: proof/lemma-gluing-after-admissible-blow-ups
\begin{proof}[{Proof}]\phantomsection\label{bp61}%
We glue the \(T_i\) along the closure of \(U\) in \(T_1\times_ST_2\), after blow-ups that make this closure open in both. The glued scheme is separated because the closure is closed in the product.

Let \(G\) be the schematic closure of \(U\) in \(T_1\times_ST_2\). Then Theorem \hyperref[bp48]{A.18}, applied to the two projections, makes the strict transforms of \(G\) opens in suitable blow-ups of the two sides. To obtain a single modification, we blow up \(G\) in the product of the two pulled-back centers, which gives a common refinement. By Lemma \hyperref[bp35]{A.13}, these further centers extend from the opens to each side, and the resulting scheme \(G'\) is an open in both \(T_i'\).

The map \(G'\to T_1'\times_ST_2'\) is a closed immersion. It is an immersion, and it is proper because \(G'\to G\) is proper, \(G\subseteq T_1\times_ST_2\) is closed, and the modified product is separated over the original product. We glue the two \(T_i'\) along \(G'\), and the glued scheme \(T\) is covered by \(T_1'\) and \(T_2'\). To see that \(T\) is separated, note that the diagonal of \(T\) restricted to \(T_i'\times_ST_j'\) is the diagonal of \(T_i'\) for \(i=j\) and the graph \(G'\) for \(i\neq j\). Both are closed, so \(T\) is separated. It is of finite type because it has a finite cover by opens of finite type.

Suppose that \(U\) is dense in both \(T_i\). The complement of the exceptional divisor of a blow-up is schematically dense, as in the proof of Lemma \hyperref[bp29]{A.10}, and it contains the inverse image of \(U\) as a dense open. Hence \(U\) is dense in each \(T_i'\), and so in \(T\). This construction is [Stacks, \href{https://stacks.math.columbia.edu/tag/0F3X}{0F3X}].
\end{proof}

% text: prose-A-22
\phantomsection\label{bp62}%
The gluing lemma produces a separated scheme but says nothing about properness. In the inductive step we prove properness by the valuative criterion, using the separation of closures.

% object: proposition-compactifying-union-two-opens (blueprint/chap-Nagata-compactification-and-Chows-lemma/sec-Separation-by-blow-ups/proposition-compactifying-union-two-opens/_section.md)
\phantomsection\label{bp63}%
\begin{bpstatementit}{Proposition~A.23}{ (Compactifying a union of two opens)}
Let \(S\) be Noetherian and let \(M\to S\) be separated of finite type. Suppose that \(M=M_1\cup M_2\), that each \(M_i\) has a compactification over \(S\), and that \(D=M_1\cap M_2\) is dense in \(M\). Then \(M\) has a compactification over \(S\).
\end{bpstatementit}

% proof: proof/proposition-compactifying-union-two-opens
\begin{proof}[{Proof}]\phantomsection\label{bp64}%
The idea is to glue \(M\) to each compactification \(X_i\) of \(M_i\) after removing from \(X_i\) the closure of the part lying over \(M\setminus M_i\). This gives separated schemes \(W_i\) containing \(M\), and a common modification combines them into one separated \(W\). It remains to check that \(W\) is proper, and for this we use the valuative criterion. Separating the closures of the boundaries, we obtain (A.1), and from it we deduce that every valuation of \(D\) extends into one of the \(W_i\).

\emph{Separating the boundaries.} Since the subsets to be removed from \(X_i\) are defined by points of \(M\) outside \(M_i\), we must relate each \(X_i\) to \(M\) itself. Choose dense compactifications \(M_i\subseteq X_i\). By Lemma \hyperref[bp51]{A.19}, after admissible blow-ups there are opens \(M_i\subseteq V_i\subseteq X_i\) and proper maps \(\psi_i:V_i\to M\) that are isomorphisms above \(M_i\). Put, for \(\{i,j\}=\{1,2\}\),

\[
Z_i=M\setminus M_j,\qquad
Z_{i,r}=\psi_i^{-1}(Z_r),\qquad
B_{i,r}=\overline{Z_{i,r}}\subseteq X_i.
\]

The subsets \(Z_{i,1},Z_{i,2}\) are disjoint and closed in \(V_i\). By Lemma \hyperref[bp55]{A.20}, after a further \(V_i\)-admissible blow-up we may assume that \(B_{i,1}\cap B_{i,2}=\varnothing\), and subsequent \(V_i\)-admissible blow-ups preserve this disjointness.

Let \(X_{12}\) be the schematic closure of \(V_1\times_MV_2\) in \(X_1\times_SX_2\), with proper projections \(p_i\). Above \(V_1\), the first projection is the base change of \(\psi_2\), so it is an isomorphism near \(Z_{1,2}\). By Lemma \hyperref[bp57]{A.21} it becomes an isomorphism near \(B_{1,2}\) after a blow-up. We replace \(X_{12}\) by its strict transform, or equivalently by the closure of the new graph. Since \(p_1\) is now an isomorphism over a neighbourhood of \(B_{1,2}\) and \(Z_{1,2}\) is dense in \(B_{1,2}\), the inverse image of \(Z_{1,2}\) is dense in \(p_1^{-1}(B_{1,2})\). The second projection maps that inverse image into \(Z_{2,2}\). By density it maps all of \(p_1^{-1}(B_{1,2})\) into \(B_{2,2}\), which is disjoint from \(B_{2,1}\). This gives the disjointness used in the valuative argument below:

\[
X_{12}\cap(B_{1,2}\times_SB_{2,1})=\varnothing.
\tag{A.1}
\]

\emph{The separated schemes \(W_i\).} We glue

\[
W_i=M\coprod_{M_i}(X_i\setminus B_{i,j}).
\]

We claim that these schemes are separated over \(S\). The scheme \(W_i\) is covered by \(M\) and \(X_i\setminus B_{i,j}\), both separated, and they are glued along \(M_i\). It therefore suffices to show that the graph of the inclusion \(M_i\to X_i\setminus B_{i,j}\), which is the mixed part of the diagonal, is closed in \(M\times_S(X_i\setminus B_{i,j})\).

Suppose that a point of the graph of \(M_i\) specializes in \(M\times_S(X_i\setminus B_{i,j})\). Since \(\psi_i\) is proper, its graph is closed in \(M\times_SX_i\). Hence the specialized point lies in the graph of \(\psi_i\), over a point of \(V_i\). If its image in \(M\) lay in \(M\setminus M_i=Z_j\), the lifted point would lie in \(Z_{i,j}\subseteq B_{i,j}\), contradicting the second coordinate. So the specialization remains in the graph. The graph is an immersion with constructible image, and in the Noetherian situation a constructible subset that is closed under specialization is closed. Therefore the graph is a closed immersion, which proves separatedness.

\emph{Valuations of \(D\).} For properness we use the valuative criterion, and we show first that every valuation of \(D\) extends into one of the \(W_i\). Let \(R\) be a valuation ring with fraction field \(K\), and let \(\operatorname{Spec}(K)\to D\) be an \(S\)-map over a given \(\operatorname{Spec}(R)\to S\). It extends to both \(X_i\), since they are proper, and the pair of extensions factors through \(X_{12}\), because the closure of the graph contains its generic point and \(\operatorname{Spec}(K)\) is schematically dense in \(\operatorname{Spec}(R)\). By (A.1), the closed point cannot map into both of the removed closed subsets \(B_{1,2}\) and \(B_{2,1}\). We may therefore choose \(i\) such that the closed point does not map into \(B_{i,j}\). Since an open containing the closed point of a local scheme contains the whole scheme, the extension factors through \(X_i\setminus B_{i,j}\subseteq W_i\).

\emph{Gluing and properness.} Lemma \hyperref[bp60]{A.22}, applied to \(M\subseteq W_i\), gives a separated \(W\) of finite type, covered by admissible modifications \(W_i'\). Since the modification is proper, the lift of the valuation into \(W_i\) lifts further into \(W_i'\), hence into \(W\). Such lifts are unique because \(W\) is separated. The valuative criterion [Stacks, \href{https://stacks.math.columbia.edu/tag/0894}{0894}], in the form that only requires lifts of maps from \(\operatorname{Spec}(K)\) into the dense open \(D\), then shows that \(W\to S\) is universally closed. As \(W\) is also of finite type and separated, it is proper over \(S\). Its open \(M\) yields the desired compactification, which becomes dense after we take the schematic closure of \(M\). This is [Stacks, \href{https://stacks.math.columbia.edu/tag/0F40}{0F40}].
\end{proof}


% object: sec-Proof-of-Nagatas-theorem (blueprint/chap-Nagata-compactification-and-Chows-lemma/sec-Proof-of-Nagatas-theorem/_section.md), content hidden
\subsection*{\texorpdfstring{A.4\quad Proof of Nagata's theorem}{A.4 Proof of Nagata's theorem}}
\phantomsection\label{bp65}%
\addcontentsline{toc}{subsection}{\texorpdfstring{A.4\quad Proof of Nagata's theorem}{A.4 Proof of Nagata's theorem}}
% text: prose-A-23
\phantomsection\label{bp66}%
Over a Noetherian base, the theorem follows from Proposition \hyperref[bp63]{A.23} by induction over an affine cover whose members overlap densely. Each affine piece is compactified by its closure in a projective space. The proposition then compactifies the union of the first \(r\) pieces with the next one. It requires the overlap to be dense, because its valuative criterion only tests maps from spectra of valuation rings whose generic point maps into the overlap. Lemma \hyperref[bp67]{A.24} provides a cover whose total intersection contains all generic points of \(X\).

Over a general qcqs base we use Noetherian approximation. This works because a morphism of finite presentation is the base change of one over a Noetherian scheme, and compactifications are stable under base change. A morphism that is only of finite type, however, need not descend in this way. We therefore first embed \(X\) as a closed subscheme of a scheme of finite presentation over \(S\), and take the closure of \(X\) in a compactification of the latter.

% object: lemma-affine-covers-dense-intersection (blueprint/chap-Nagata-compactification-and-Chows-lemma/sec-Proof-of-Nagatas-theorem/lemma-affine-covers-dense-intersection/_section.md)
\phantomsection\label{bp67}%
\begin{bpstatementit}{Lemma~A.24}{ (Affine covers with dense intersection)}
For a scheme \(X\) of finite type over a Noetherian scheme \(S\), there is a finite affine cover \(X=\bigcup_{r=1}^m U_r\) whose total intersection is dense in \(X\). If \(X\to S\) is separated, then each \(U_r\to S\) has a projective compactification.
\end{bpstatementit}

% proof: proof/lemma-affine-covers-dense-intersection
\begin{proof}[{Proof}]\phantomsection\label{bp68}%
The scheme \(X\) has finitely many irreducible components. By [Stacks, \href{https://stacks.math.columbia.edu/tag/01ZX}{01ZX}], every point of \(X\) has an affine open neighbourhood that also contains all the generic points of those components, and since \(X\) is quasi-compact, finitely many such opens cover \(X\). Their total intersection still contains all generic points, so it is dense. For the second assertion, note that each \(U_r\) is affine of finite type over \(S\), so by [Stacks, \href{https://stacks.math.columbia.edu/tag/04II}{04II}] we may embed it as a locally closed subscheme of some \(\mathbf A^{n_r}_S\). Its schematic closure in \(\mathbf P^{n_r}_S\) is proper, and by Lemma \hyperref[bp11]{A.3} the scheme \(U_r\) is open in that closure.
\end{proof}

% text: prose-A-24
\phantomsection\label{bp69}%
We prove the theorem in three steps of increasing generality: a Noetherian base, a finitely presented morphism over any qcqs base, and a morphism of finite type.

% proof: proof/theorem-nagata-s-compactification-theorem
\begin{proof}[{Proof of Theorem A.2}]\phantomsection\label{bp70}%
If \(S\) is Noetherian, we choose a cover as in Lemma \hyperref[bp67]{A.24} and argue by induction on the unions \(X_r=U_1\cup\cdots\cup U_r\). For \(r=1\), the lemma gives a compactification of \(U_1\). If \(X_r\) has a compactification, we apply Proposition \hyperref[bp63]{A.23} to \(M=X_{r+1}\), with \(M_1=X_r\) and \(M_2=U_{r+1}\). Each piece has a compactification, by induction and by the lemma. The overlap \(X_r\cap U_{r+1}\) is dense in \(X_{r+1}\), since it contains the total intersection, which contains all generic points of \(X\). Hence \(X_{r+1}\) has a compactification. The last union is \(X\) itself, and taking the schematic closure of \(X\) makes the compactification dense.

Now let \(S\) be a general qcqs scheme, and suppose first that \(X\to S\) is separated and of finite presentation. By absolute Noetherian approximation we can write

\[
S=\varprojlim_i S_i
\]

with \(S_i\) Noetherian and affine transition maps [Stacks, \href{https://stacks.math.columbia.edu/tag/01ZA}{01ZA}]. Since \(X\to S\) is of finite presentation, it descends to a model \(X_i\to S_i\) for some \(i\), and after increasing \(i\) the model is separated [Stacks, \href{https://stacks.math.columbia.edu/tag/01ZM}{01ZM}, \href{https://stacks.math.columbia.edu/tag/01ZQ}{01ZQ}]. We then compactify this model by the Noetherian case and base change the compactification to \(S\). The result is a compactification, because open immersions and proper morphisms are stable under base change, and taking the schematic closure of \(X\) makes the result dense.

Finally, let \(X\to S\) be separated of finite type. By [Stacks, \href{https://stacks.math.columbia.edu/tag/01ZJ}{01ZJ}], there is a closed immersion \(X\hookrightarrow X_0\) with \(X_0\to S\) separated of finite presentation. We compactify \(X_0\) by the case already proved, and take the schematic closure of \(X\) in that compactification. This closure is proper over \(S\). Its restriction to the open \(X_0\) is the closed subscheme \(X\) itself, so \(X\) is open in the closure, and it is schematically dense there by Lemma \hyperref[bp11]{A.3}.
\end{proof}


% object: sec-Chows-lemma (blueprint/chap-Nagata-compactification-and-Chows-lemma/sec-Chows-lemma/_section.md), content hidden
\subsection*{\texorpdfstring{A.5\quad Chow's lemma}{A.5 Chow's lemma}}
\phantomsection\label{bp71}%
\addcontentsline{toc}{subsection}{\texorpdfstring{A.5\quad Chow's lemma}{A.5 Chow's lemma}}
% text: prose-A-25
\phantomsection\label{bp72}%
Let \(f:X\to S\) be proper. In Section \bpout{4.5} proper base change is first proved for projective morphisms. By Lemma \bpout{4.29} it then descends along a proper surjection \(\pi:\widetilde X\to X\) as soon as it holds for \(\pi\) and for the composite \(f\pi\). Chow's lemma says precisely that such a \(\pi\) exists with both \(\pi\) and \(f\pi\) projective. Some modification of \(X\) is unavoidable here, since there are proper varieties that are not projective.

% object: theorem-chow-s-lemma-projective-form (blueprint/chap-Nagata-compactification-and-Chows-lemma/sec-Chows-lemma/theorem-chow-s-lemma-projective-form/_section.md)
\phantomsection\label{bp73}%
\begin{bpstatementit}{Theorem~A.25}{ (Chow's lemma, projective form)}
Let \(f:X\to S\) be proper with \(S\) qcqs. There are a proper surjection \(\pi:\widetilde X\to X\) and a closed immersion \(i_S:\widetilde X\hookrightarrow\mathbf P^n_S\) such that \(f\pi=q_Si_S\), where \(q_S:\mathbf P^n_S\to S\) is the projection. Moreover,

\[
i_X=(\pi,i_S):\widetilde X\hookrightarrow
X\times_S\mathbf P^n_S=\mathbf P^n_X
\]

is a closed immersion. Both \(\pi\) and \(f\pi\) are projective.
\end{bpstatementit}

% proof: proof/theorem-chow-s-lemma-projective-form
\begin{proof}[{Proof}]\phantomsection\label{bp74}%
We take the immersion into projective space given by the Stacks Project and show that properness of \(X\) makes it closed. The second assertion then follows by writing \(i_X\) as a graph followed by a base change of \(i_S\).

The version of Chow's lemma in [Stacks, \href{https://stacks.math.columbia.edu/tag/0202}{0202}] gives a proper surjection \(\pi:\widetilde X\to X\) and an immersion \(i_S:\widetilde X\to\mathbf P^n_S\) for every separated \(X\to S\) of finite type with \(S\) qcqs. It is deduced from the Noetherian case by the approximation argument in the proof of Theorem \hyperref[bp9]{A.2}, using the embedding of \(X\) into a scheme of finite presentation [Stacks, \href{https://stacks.math.columbia.edu/tag/01ZJ}{01ZJ}]. In our situation, moreover, \(f\) and \(\pi\) are proper, so \(\widetilde X\) is proper over \(S\). Now a map from a proper \(S\)-scheme to a separated \(S\)-scheme is proper, because its graph is closed and the projection of the product to the target is proper. Hence \(i_S\) is a proper immersion. Its image is therefore closed, and \(i_S\) is a closed immersion.

For the second assertion, we factor \(i_X\) as the graph of \(\pi\) followed by the base change of \(i_S\):

\[
\widetilde X\longrightarrow\widetilde X\times_SX
\longrightarrow\mathbf P^n_S\times_SX.
\]

The first arrow is closed because \(X\to S\) is separated, and the second is closed because \(i_S\) is. Their composite is \(i_X\). Since \(\pi\) is \(i_X\) followed by the projection \(\mathbf P^n_X\to X\), it is projective, and so is \(f\pi=q_Si_S\).
\end{proof}



% object: chap-Abelian-varieties (blueprint/chap-Abelian-varieties/_section.md), content hidden
\section*{\texorpdfstring{B\quad Abelian varieties}{B Abelian varieties}}
\phantomsection\label{bp75}%
\addcontentsline{toc}{section}{\texorpdfstring{B\quad Abelian varieties}{B Abelian varieties}}
% text: prose-B-01
\phantomsection\label{bp76}%
Two steps of the main proof come down to the torsion points of a Jacobian. By the Kummer sequence, \(H^1(C,\mu_n)\) of a curve is the \(n\)-torsion of its Jacobian. The first step, the finiteness of \(H^1\) needed for proper base change, is therefore a count of torsion points. The second step is the finiteness of geometric monodromy in rank one, where purity begins. It needs Frobenius to fix no class in \(H^1\) of a curve over \(\mathbf F_q\) (Proposition \bpout{14.3}).

For a complex torus \(\mathbf C^g/\Lambda\) the count is trivial: the \(n\)-torsion is \(\tfrac1n\Lambda/\Lambda\simeq(\mathbf Z/n)^{2g}\). For an abelian variety \(A\) of dimension \(g\) over an arbitrary field \(k\), however, there is no lattice, and the torsion points must be counted algebraically. This is what we do: we prove that for \(n\) invertible in \(k\) the group scheme \(A[n]\) is finite étale of rank \(n^{2g}\) (Proposition \hyperref[bp108]{B.12}). It follows that the Tate module \(T_\ell A=\varprojlim_mA[\ell^m](\bar k)\) is free of rank \(2g\) over \(\mathbf Z_\ell\) (Proposition \hyperref[bp114]{B.14}). Moreover, over \(\mathbf F_q\), no power of Frobenius has a nonzero fixed vector on \(T_\ell A\) (Proposition \hyperref[bp117]{B.15}). The reason is short: a fixed vector would be a compatible system of \(\ell\)-power torsion points of unbounded order in the finite group \(A_0(\mathbf F_{q^m})\).

The count rests on one formula. If \(\mathcal N\) is a symmetric ample invertible sheaf on \(A\), then

\[
[n]^*\mathcal N\simeq\mathcal N^{\otimes n^2}
\]

(Proposition \hyperref[bp105]{B.11}). On an elliptic curve with \(\mathcal N=\mathcal O(e)\) this says that the divisor \([n]^*(e)\) is linearly equivalent to \(n^2(e)\). When \(n\) is prime to the characteristic, \([n]\) is étale and \([n]^*(e)\) is the reduced divisor \(E[n]\). Comparing degrees then gives \(\#E[n]=n^2\).

In general the formula is used twice. First, it forces \([n]\) to have finite fibres. Indeed, \(\mathcal N^{\otimes n^2}\) has positive degree on every complete curve in \(A\), while \([n]^*\mathcal N\) is trivial on a curve contracted by \([n]\). Second, we compute the Euler characteristics \(\chi(A,[n]^*\mathcal N^{\otimes m})\) in two ways, and comparing their leading coefficients gives \(\deg[n]=n^{2g}\).

The formula itself comes from the theorem of the cube (Theorem \hyperref[bp91]{B.6}). Following [Stacks, \href{https://stacks.math.columbia.edu/tag/0BEZ}{0BEZ}], we prove the theorem of the cube together with the rigidity lemma. As a by-product, from the rigidity lemma we deduce that \(A\) is commutative; see also [\hyperref[bp349]{Mil86a}].

In this appendix all group schemes are of finite type over a field \(k\) of arbitrary characteristic.

% object: sec-The-theorem-of-the-cube (blueprint/chap-Abelian-varieties/sec-The-theorem-of-the-cube/_section.md), content hidden
\subsection*{\texorpdfstring{B.1\quad The theorem of the cube}{B.1 The theorem of the cube}}
\phantomsection\label{bp77}%
\addcontentsline{toc}{subsection}{\texorpdfstring{B.1\quad The theorem of the cube}{B.1 The theorem of the cube}}
% text: prose-B-02
\phantomsection\label{bp78}%
The theorem of the cube says that an invertible sheaf \(\mathcal L\) on \(X\times Y\times Z\), with \(X,Y\) complete varieties and \(Z\) connected, that is trivial on the three slices through a point \((x,y,z)\) is trivial. To prove it, we view \(\mathcal L\) as a family of invertible sheaves on \(X\times Y\) parametrized by \(Z\). Following [Stacks, \href{https://stacks.math.columbia.edu/tag/0BEZ}{0BEZ}], we then consider the set of points of \(Z\) over which the sheaf is trivial. This set is closed, by semicontinuity. It is also open, by a computation of \(H^0\) and \(H^1\) near such a point. Since it contains \(z\) and \(Z\) is connected, it is all of \(Z\). The openness uses the triviality on the two slices \(x\times Y\times Z\) and \(X\times y\times Z\), and with only two factors it fails.

% object: example-cube-two-factors (blueprint/chap-Abelian-varieties/sec-The-theorem-of-the-cube/example-cube-two-factors/_section.md)
\phantomsection\label{bp79}%
\begin{bpstatement}{Example~B.1}{ (Two factors do not suffice)}
Let \(E\) be an elliptic curve over an algebraically closed field, with origin \(O\), and let \(\Delta\subset E\times E\) be the diagonal. On \(E\times E\) put

\[
\mathcal L=\mathcal O_{E\times E}\bigl(\Delta-E\times\{O\}-\{O\}\times E\bigr).
\]

Consider the slice \(E\times\{t\}\). It meets the diagonal in the point \(t\), and the divisor \(\{O\}\times E\) meets it in the point \(O\). The divisor \(E\times\{O\}\), on the other hand, is pulled back from the second factor, so it contributes the trivial sheaf for every \(t\). Hence the restriction is \(\mathcal O_E(t-O)\). By symmetry the restriction to \(\{t\}\times E\) is also \(\mathcal O_E(t-O)\). In particular \(\mathcal L\) is trivial on both slices through \((O,O)\).

However, \(\mathcal L\) is not trivial, because for \(t\neq O\) the sheaf \(\mathcal O_E(t-O)\) is nontrivial. Indeed, a rational function with divisor \(t-O\) would be a map \(E\to\mathbf P^1\) of degree one, which is impossible in genus one. The set of \(t\) over which the fibre is trivial is the point \(O\). It is closed, as Lemma \hyperref[bp80]{B.2} below predicts, but not open.
\end{bpstatement}

% object: lemma-triviality-fibres (blueprint/chap-Abelian-varieties/sec-The-theorem-of-the-cube/lemma-triviality-fibres/_section.md)
\phantomsection\label{bp80}%
\begin{bpstatementit}{Lemma~B.2}{ (Triviality on fibres)}
Let \(P\) be a proper geometrically integral \(k\)-scheme, let \(T\) be a \(k\)-scheme of finite type, let \(f:P\times T\to T\) be the projection, and let \(\mathcal L\) be an invertible sheaf on \(P\times T\). For \(t\in T\) write \(P_t=P\times_k\kappa(t)\) and \(\mathcal L_t=\mathcal L|_{P_t}\).

\begin{enumerate}
\item \(\mathcal L_t\) is trivial if and only if \(H^0(P_t,\mathcal L_t)\ne0\) and \(H^0(P_t,\mathcal L_t^{-1})\ne0\).
\item The set of \(t\in T\) such that \(\mathcal L_t\) is trivial is closed.
\end{enumerate}
\end{bpstatementit}

% proof: proof/lemma-triviality-fibres
\begin{proof}[{Proof}]\phantomsection\label{bp81}%
The proof uses only that \(P_t\) is proper and integral with \(H^0(P_t,\mathcal O_{P_t})=\kappa(t)\), which holds by [Stacks, \href{https://stacks.math.columbia.edu/tag/0BUG}{0BUG}] because \(P\) is proper and geometrically integral.

For (1), a trivial sheaf and its inverse both have the nonzero section \(1\), so only the converse needs an argument. On \(\mathbf P^1\), for instance, \(\mathcal O(d)\) and \(\mathcal O(-d)\) both have nonzero sections only when \(d=0\), and the general argument is the same. Let \(s\) and \(s'\) be nonzero sections of \(\mathcal L_t\) and \(\mathcal L_t^{-1}\). Since \(P_t\) is integral, the product \(ss'\) is a nonzero element of \(H^0(P_t,\mathcal O)=\kappa(t)\), hence a unit. Since \(ss'\) would vanish wherever \(s\) does, the section \(s\) vanishes nowhere, and it trivializes \(\mathcal L_t\).

For (2), note that by (1) the set in question is the intersection of the loci where \(\dim H^0(P_t,\mathcal L_t)\) and \(\dim H^0(P_t,\mathcal L_t^{-1})\) are positive. These loci are closed because the functions \(t\mapsto\dim H^0(P_t,\mathcal L_t^{\pm1})\) are upper semicontinuous [Stacks, \href{https://stacks.math.columbia.edu/tag/0BDN}{0BDN}], as \(\mathcal L\) is flat over \(T\) and \(f\) is proper.
\end{proof}

% text: prose-B-03
\phantomsection\label{bp82}%
Part (2) is the closed half of the theorem of the cube. If the sheaf is trivial over the members of an open cover of \(T\), then the evaluation map identifies it with a pullback from \(T\), since sections along the proper factor \(P\) are constant.

% object: corollary-triviality-over-open-cover-base (blueprint/chap-Abelian-varieties/sec-The-theorem-of-the-cube/corollary-triviality-over-open-cover-base/_section.md)
\phantomsection\label{bp83}%
\begin{bpstatementit}{Corollary~B.3}{ (Triviality over an open cover of the base)}
With the notation of Lemma \hyperref[bp80]{B.2}, suppose that \(\mathcal L\) is trivial on \(P\times U_\alpha\) for an open cover \((U_\alpha)\) of \(T\). Then \(f_*\mathcal L\) is invertible and the evaluation map \(f^*f_*\mathcal L\to\mathcal L\) is an isomorphism.
\end{bpstatementit}

% proof: proof/corollary-triviality-over-open-cover-base
\begin{proof}[{Proof}]\phantomsection\label{bp84}%
Both assertions are local on \(T\), so we may replace \(T\) by one of the \(U_\alpha\) and assume that \(\mathcal L\simeq\mathcal O_{P\times T}\). Since \(H^0(P,\mathcal O_P)=k\) [Stacks, \href{https://stacks.math.columbia.edu/tag/0BUG}{0BUG}], and \(T\) is flat over \(k\), flat base change gives \(f_*\mathcal O_{P\times T}=\mathcal O_T\). Hence the chosen trivialization identifies \(f_*\mathcal L\) with the invertible sheaf \(\mathcal O_T\). Under the same identification the evaluation map \(f^*f_*\mathcal L\to\mathcal L\) becomes the identity of \(\mathcal O_{P\times T}\).
\end{proof}

% text: prose-B-04
\phantomsection\label{bp85}%
The open half of the theorem is the substantial step. It uses the fact that \(H^1\) of the structure sheaf of a product is the sum of the \(H^1\) of the factors.

% object: lemma-h-1-structure-sheaf-product (blueprint/chap-Abelian-varieties/sec-The-theorem-of-the-cube/lemma-h-1-structure-sheaf-product/_section.md)
\phantomsection\label{bp86}%
\begin{bpstatementit}{Lemma~B.4}{ ($H^1$ of the structure sheaf of a product)}
Let \(K\) be a field, let \(P\) and \(Q\) be proper \(K\)-schemes with \(H^0(P,\mathcal O_P)=H^0(Q,\mathcal O_Q)=K\), and let \(p_0\in P(K)\) and \(q_0\in Q(K)\). Restriction to the slices \(P\times q_0\) and \(p_0\times Q\) is an isomorphism

\[
H^1(P\times Q,\mathcal O)\longrightarrow
H^1(P,\mathcal O)\oplus H^1(Q,\mathcal O).
\]
\end{bpstatementit}

% proof: proof/lemma-h-1-structure-sheaf-product
\begin{proof}[{Proof}]\phantomsection\label{bp87}%
By the Künneth formula [Stacks, \href{https://stacks.math.columbia.edu/tag/0BED}{0BED}], and since \(H^0\) of \(\mathcal O\) is \(K\) on both factors, the source is \(\operatorname{pr}_1^*H^1(P,\mathcal O)\oplus
\operatorname{pr}_2^*H^1(Q,\mathcal O)\). Restriction to \(P\times q_0\) is the identity on the first summand. It kills the second, because on the second summand it factors through \(H^1\) of the point \(q_0\). Restriction to \(p_0\times Q\) does the opposite.
\end{proof}

% object: lemma-triviality-near-fibre (blueprint/chap-Abelian-varieties/sec-The-theorem-of-the-cube/lemma-triviality-near-fibre/_section.md)
\phantomsection\label{bp88}%
\begin{bpstatementit}{Lemma~B.5}{ (Triviality near a fibre)}
Let \(P\) and \(Q\) be proper geometrically integral \(k\)-schemes with points \(p_0\in P(k)\) and \(q_0\in Q(k)\), let \(T\) be a \(k\)-scheme of finite type, and let \(\mathcal L\) be an invertible sheaf on \(P\times Q\times T\) whose restrictions to \(p_0\times Q\times T\) and to \(P\times q_0\times T\) are trivial. If \(\mathcal L\) is trivial on the fibre over a point \(t\in T\), then \(\mathcal L\) is trivial on \(P\times Q\times U\) for an open neighbourhood \(U\) of \(t\).
\end{bpstatementit}

% proof: proof/lemma-triviality-near-fibre
\begin{proof}[{Proof}]\phantomsection\label{bp89}%
We want a section of \(\mathcal L\) over \(X\times U\), where \(X=P\times Q\), that does not vanish on the fibre over \(t\). Its zero locus then misses a neighbourhood of that fibre. Such a section exists as soon as \(f_*\mathcal L\) is free of rank one near \(t\), with formation compatible with passage to the fibre.

To get this, we represent \(Rf_*\mathcal L\), and the pushforwards of the restrictions of \(\mathcal L\) to the two slices, by complexes of free modules whose differentials vanish at \(t\). By the Künneth formula, restriction from \(H^1\) of the fibre to \(H^1\) of the two slices is an isomorphism. This forces the first differential of \(Rf_*\mathcal L\) to vanish after shrinking \(T\). Then \(H^0\) of \(Rf_*\mathcal L\) is free of rank one, as we need.

\emph{Setup.} Since \(X_{\bar k}=P_{\bar k}\times Q_{\bar k}\) is a product of integral varieties, \(X\) is geometrically integral [Stacks, \href{https://stacks.math.columbia.edu/tag/05P3}{05P3}]. We may assume that \(T=\operatorname{Spec}R\) is affine, and we let \(\mathfrak p\) be the prime of \(t\). Below, “shrinking \(T\)” means replacing \(R\) by \(R_g\) for some \(g\notin\mathfrak p\); everything already constructed survives this. Consider the projections

\[
f:X\times T\to T,\qquad f_1:P\times T\to T,\qquad f_2:Q\times T\to T,
\]

and the restrictions \(\mathcal L_1\) and \(\mathcal L_2\) of \(\mathcal L\) to \(P\times q_0\times T\) and \(p_0\times Q\times T\), which are trivial. The complexes \(E=Rf_*\mathcal L\) and \(E_i=Rf_{i*}\mathcal L_i\) are perfect, and their formation commutes with base change [Stacks, \href{https://stacks.math.columbia.edu/tag/0B91}{0B91}]. Restriction of sections gives a map \(E\to E_1\oplus E_2\), also compatible with base change. Since \(T\) is affine, \(H^0(X\times T,\mathcal L)=H^0(E)\) and likewise for \(E_i\).

\emph{Minimal complexes.} A perfect complex over a local ring can be represented by a complex of free modules whose differentials vanish modulo the maximal ideal; the ranks are then the dimensions of the cohomology of the fibre. We use this to read off the ranks from the fibre over \(t\), where everything is trivial. To do so, put \(h^j=\dim H^j(X_t,\mathcal O)\) and \(h_i^j=\dim H^j\) of \(\mathcal O\) on the fibres \(P_t\) and \(Q_t\). Since \(\mathcal L\) is trivial on \(X_t\) and the \(\mathcal L_i\) are trivial, \(E\otimes^{\mathbf L}_R\kappa(t)\) and \(E_i\otimes^{\mathbf L}_R\kappa(t)\) have cohomology of dimensions \(h^j\) and \(h_i^j\), with \(h^0=h_i^0=1\). After shrinking \(T\) we may therefore represent \(E\) and \(E_i\) by complexes of finite free modules

\[
R\xrightarrow{d}R^{h^1}\to R^{h^2}\to\cdots,\qquad
R\xrightarrow{d_i}R^{h_i^1}\to R^{h_i^2}\to\cdots
\]

whose differentials vanish modulo \(\mathfrak p\) [Stacks, \href{https://stacks.math.columbia.edu/tag/0BCD}{0BCD}]. We represent the map \(E\to E_i\) by a map of complexes [Stacks, \href{https://stacks.math.columbia.edu/tag/064B}{064B}], with components \(a_i:R\to R\) in degree \(0\) and \(b_i:R^{h^1}\to R^{h_i^1}\) in degree \(1\). Since the differentials vanish modulo \(\mathfrak p\), the map \(b_i\otimes\kappa(t)\) is the restriction map on \(H^1\) of the fibre over \(t\).

We want to show that \(d=0\) after shrinking \(T\), since then \(H^0(E)=\ker d=R\) is free of rank one. The maps \(b_i\) will let us transfer this question to the \(d_i\), which are easier because \(\mathcal L_i\) is trivial.

\emph{The differentials \(d_i\) vanish.} Fix a trivialization \(\mathcal L_1\simeq\mathcal O\), and let \(u\in H^0(E_1)=\ker d_1\subset R\) be the section corresponding to \(1\). Its image in \(H^0(E_1\otimes^{\mathbf L}\kappa(t))=H^0(P_t,\mathcal O)\) is \(1\ne0\), so \(u\notin\mathfrak p\). Since \(d_1(u)=u\,d_1(1)=0\), after inverting \(u\) we get \(d_1=0\). The same argument gives \(d_2=0\).

\emph{The differential \(d\) vanishes.} This step uses the two slices together. If the map \((b_1,b_2)\) is an isomorphism, then \(d\) is determined by \(d_1\) and \(d_2\), and this condition can be checked on the fibre. The map \((b_1,b_2)\otimes\kappa(t)\) is restriction from \(H^1(P_t\times Q_t,\mathcal O)\) to the slices \(P_t\times q_0\) and \(p_0\times Q_t\). It is an isomorphism by Lemma \hyperref[bp86]{B.4}, and in particular \(h^1=h_1^1+h_2^1\). Hence the square matrix \((b_1,b_2)\) has determinant outside \(\mathfrak p\). After inverting this determinant, \((b_1,b_2)\) is an isomorphism. Since the maps form a map of complexes, we have \(b_i\circ d=d_i\circ a_i=0\), and therefore \(d=0\).

\emph{Conclusion.} A generator of \(H^0\) gives the trivialization. We have \(H^0(X\times T,\mathcal L)=H^0(E)=\ker d=R\), and its generator \(\sigma\) maps to a generator of \(H^0(E\otimes^{\mathbf L}\kappa(t))=H^0(X_t,\mathcal L_t)\). Since \(\mathcal L_t\simeq\mathcal O_{X_t}\) and \(H^0(X_t,\mathcal O)=\kappa(t)\), the section \(\sigma\) vanishes nowhere on \(X_t\). Its zero locus is closed in \(X\times T\) and disjoint from \(X_t\). Since \(f\) is proper, its image in \(T\) is a closed set not containing \(t\). On the complement \(U\) of this image, \(\sigma\) trivializes \(\mathcal L\).
\end{proof}

% text: prose-B-05
\phantomsection\label{bp90}%
With the closed and the open half in hand, it remains to glue the local trivializations.

% object: theorem-theorem-cube (blueprint/chap-Abelian-varieties/sec-The-theorem-of-the-cube/theorem-theorem-cube/_section.md)
\phantomsection\label{bp91}%
\begin{bpstatementit}{Theorem~B.6}{ (Theorem of the cube)}
Let \(X\) and \(Y\) be proper geometrically integral \(k\)-schemes with points \(x\in X(k)\) and \(y\in Y(k)\), let \(Z\) be a connected \(k\)-scheme of finite type with a point \(z\in Z(k)\), and let \(\mathcal L\) be an invertible sheaf on \(X\times Y\times Z\). If the restrictions of \(\mathcal L\) to \(x\times Y\times Z\), \(X\times y\times Z\) and \(X\times Y\times z\) are trivial, then \(\mathcal L\) is trivial.
\end{bpstatementit}

% proof: proof/theorem-theorem-cube
\begin{proof}[{Proof}]\phantomsection\label{bp92}%
The local trivializations are glued by Corollary \hyperref[bp83]{B.3}, which shows that \(\mathcal L\) is pulled back from \(Y\times Z\); the slice \(x\times Y\times Z\) then identifies the sheaf on \(Y\times Z\).

Let \(Z_0\subset Z\) be the set of points over which \(\mathcal L\) is trivial on the fibre of \(p:X\times Y\times Z\to Z\). It contains \(z\). It is closed by Lemma lemma-triviality-fibres applied to \(X\times Y\), which is proper and geometrically integral [Stacks, \href{https://stacks.math.columbia.edu/tag/05P3}{05P3}]. It is open by Lemma \hyperref[bp88]{B.5}. Since \(Z\) is connected, \(Z_0=Z\). Now Lemma \hyperref[bp88]{B.5} gives an open cover of \(Z\) over whose members \(\mathcal L\) is trivial.

Let \(\pi:X\times Y\times Z\to Y\times Z\) be the projection and \(\sigma\) its section through \(x\). We apply Corollary \hyperref[bp83]{B.3} with \(P=X\), \(T=Y\times Z\) and the open cover by the products of \(Y\) with these opens. It shows that \(\mathcal M=\pi_*\mathcal L\) is invertible and that \(\pi^*\mathcal M\to\mathcal L\) is an isomorphism. Hence \(\mathcal M\simeq\sigma^*\pi^*\mathcal M\simeq\sigma^*\mathcal L\). This is the restriction of \(\mathcal L\) to \(x\times Y\times Z\), which is trivial by hypothesis. We conclude that \(\mathcal L\simeq\pi^*\mathcal M\) is trivial.
\end{proof}

% text: prose-B-06
\phantomsection\label{bp93}%
The second tool, which we use to show that an abelian variety is commutative (Lemma \hyperref[bp100]{B.9}), is a rigidity statement. It says that a family of maps from a proper variety that is constant on one member is constant on all nearby members.

% object: lemma-rigidity (blueprint/chap-Abelian-varieties/sec-The-theorem-of-the-cube/lemma-rigidity/_section.md)
\phantomsection\label{bp94}%
\begin{bpstatementit}{Lemma~B.7}{ (Rigidity)}
Let \(X\) be a proper geometrically integral \(k\)-scheme, let \(Y\) be a \(k\)-scheme of finite type, let \(Z\) be a \(k\)-scheme, and let \(h:X\times Y\to Z\) be a morphism. If \(h\) maps \(X\times y_0\) to a single point for some \(y_0\in Y\), then there are an open neighbourhood \(V\) of \(y_0\) and a morphism \(g:V\to Z\) with \(h=g\circ\operatorname{pr}_2\) on \(X\times V\).
\end{bpstatementit}

% proof: proof/lemma-rigidity
\begin{proof}[{Proof}]\phantomsection\label{bp95}%
Near \(X\times y_0\) the map \(h\) lands in an affine open of \(Z\), and a map to an affine scheme is given by functions. Since \(X\) is proper and geometrically integral, functions on \(X\times V\) are functions on \(V\). Hence \(h\) factors through \(V\).

Let \(W\subset Z\) be an affine open containing \(h(X\times y_0)\). Since \(X\) is proper, the image in \(Y\) of the closed set \(h^{-1}(Z\setminus W)\) is closed, and its complement contains \(y_0\). Let \(V\) be an affine open neighbourhood of \(y_0\) in this complement. Then \(h\) maps \(X\times V\) into the affine scheme \(W\). It is therefore determined by the ring map \(\mathcal O(W)\to H^0(X\times V,\mathcal O)\) [Stacks, \href{https://stacks.math.columbia.edu/tag/01I1}{01I1}]. By flat base change [Stacks, \href{https://stacks.math.columbia.edu/tag/02KH}{02KH}] and \(H^0(X,\mathcal O_X)=k\) [Stacks, \href{https://stacks.math.columbia.edu/tag/0BUG}{0BUG}], we have \(H^0(X\times V,\mathcal O)=\mathcal O(V)\). The ring map is thus \(\operatorname{pr}_2^*\) composed with a ring map \(\mathcal O(W)\to\mathcal O(V)\), and this ring map defines \(g\).
\end{proof}


% object: sec-Multiplication-by-n (blueprint/chap-Abelian-varieties/sec-Multiplication-by-n/_section.md), content hidden
\subsection*{\texorpdfstring{B.2\quad Multiplication by $n$}{B.2 Multiplication by n}}
\phantomsection\label{bp96}%
\addcontentsline{toc}{subsection}{\texorpdfstring{B.2\quad Multiplication by $n$}{B.2 Multiplication by n}}
% text: prose-B-07
\phantomsection\label{bp97}%
Let \(\mathcal L\) be an invertible sheaf on an abelian variety \(A\). We pull \(\mathcal L\) back to \(A\times A\times A\) along the partial sums of the coordinates. There the theorem of the cube gives a second order recursion for \(n\mapsto[n]^*\mathcal L\), whose solution is

\[
[n]^*\mathcal L\simeq\mathcal L^{\otimes n(n+1)/2}\otimes([-1]^*\mathcal L)^{\otimes n(n-1)/2}
\]

(Proposition \hyperref[bp105]{B.11}). For symmetric \(\mathcal L\) the formula gives \(\mathcal L^{\otimes n^2}\). If \([-1]^*\mathcal L\simeq\mathcal L^{-1}\), as for an invertible sheaf of degree zero on an elliptic curve, it gives \(\mathcal L^{\otimes n}\).

We apply the formula to a symmetric ample sheaf (Lemma \hyperref[bp103]{B.10}). It shows that \([n]\) is finite and flat of degree \(n^{2g}\), and that \([n]\) is étale when \(n\) is invertible in \(k\) (Proposition \hyperref[bp108]{B.12}). The hypothesis on \(n\) cannot be dropped. For example, in characteristic \(p\) the differential of \([p]\) is zero, and as a consequence an elliptic curve has at most \(p\) points of order dividing \(p\) over \(\bar k\), although \(E[p]\) has rank \(p^2\).

% object: definition-abelian-variety (blueprint/chap-Abelian-varieties/sec-Multiplication-by-n/definition-abelian-variety/_section.md)
\phantomsection\label{bp98}%
\begin{bpstatement}{Definition~B.8}{ (Abelian variety)}
An \emph{abelian variety} over \(k\) is a smooth proper geometrically connected \(k\)-group scheme. If \(A\) is an abelian variety, \([n]\) denotes multiplication by the integer \(n\), and \(A[n]=\ker[n]\) for \(n>0\). An invertible sheaf \(\mathcal L\) on \(A\) is \emph{symmetric} if \([-1]^*\mathcal L\simeq\mathcal L\).
\end{bpstatement}

% text: prose-B-08
\phantomsection\label{bp99}%
We use three facts: commutativity, which follows from rigidity; a symmetric ample sheaf, which follows from projectivity; and the formula for \([n]^*\mathcal L\), which follows from the cube. Rigidity and the cube require \(A\) to be geometrically integral, which we check together with commutativity.

% object: lemma-abelian-variety-commutative (blueprint/chap-Abelian-varieties/sec-Multiplication-by-n/lemma-abelian-variety-commutative/_section.md)
\phantomsection\label{bp100}%
\begin{bpstatementit}{Lemma~B.9}{ (Commutativity)}
An abelian variety \(A\) over \(k\) is geometrically integral and commutative.
\end{bpstatementit}

% proof: proof/lemma-abelian-variety-commutative
\begin{proof}[{Proof}]\phantomsection\label{bp101}%
An abelian variety is geometrically reduced because it is smooth, and geometrically irreducible because the identity component of a group scheme over a field is geometrically irreducible [Stacks, \href{https://stacks.math.columbia.edu/tag/0B7R}{0B7R}] and \(A\) is connected.

For commutativity, consider \(h:A\times A\to A\), \(h(x,y)=xyx^{-1}y^{-1}\). It maps \(A\times\{e\}\) to \(e\). By Lemma \hyperref[bp94]{B.7}, applied with the first factor as the proper one, there are an open \(V\ni e\) and \(g:V\to A\) with \(h|_{A\times V}=g\circ\operatorname{pr}_2\). Composing with \((e,\operatorname{id}):V\to A\times V\) gives \(g=h(e,-)=e\), so \(h=e\) on \(A\times V\). To pass from \(A\times V\) to \(A\times A\), we compare \(h\) with the constant map \(e\). Since \(A\times A\) is smooth, hence reduced, and irreducible, the nonempty open \(A\times V\) is scheme-theoretically dense. Two maps to a separated scheme that agree on a scheme-theoretically dense open agree everywhere (Lemma \hyperref[bp14]{A.4}). Hence \(h=e\) [Stacks, \href{https://stacks.math.columbia.edu/tag/056D}{056D}, \href{https://stacks.math.columbia.edu/tag/01RH}{01RH}].
\end{proof}

% text: prose-B-09
\phantomsection\label{bp102}%
From now on we write the group law additively. Projectivity is the one structural fact that we quote.

% object: lemma-symmetric-ample-sheaf (blueprint/chap-Abelian-varieties/sec-Multiplication-by-n/lemma-symmetric-ample-sheaf/_section.md)
\phantomsection\label{bp103}%
\begin{bpstatementit}{Lemma~B.10}{ (A symmetric ample sheaf)}
An abelian variety \(A\) over \(k\) is projective, and it has a symmetric ample invertible sheaf.
\end{bpstatementit}

% proof: proof/lemma-symmetric-ample-sheaf
\begin{proof}[{Proof}]\phantomsection\label{bp104}%
Every group scheme of finite type over a field is quasi-projective [Stacks, \href{https://stacks.math.columbia.edu/tag/0BF7}{0BF7}]. A quasi-projective proper scheme is projective [Stacks, \href{https://stacks.math.columbia.edu/tag/0B45}{0B45}]. If \(\mathcal L\) is ample, then so is \([-1]^*\mathcal L\), since \([-1]\) is an automorphism, and the tensor product \(\mathcal N=\mathcal L\otimes[-1]^*\mathcal L\) of two ample sheaves is ample. It is symmetric because \([-1]\) is an involution, which interchanges the two factors.
\end{proof}

% object: proposition-pullback-multiplication-n (blueprint/chap-Abelian-varieties/sec-Multiplication-by-n/proposition-pullback-multiplication-n/_section.md)
\phantomsection\label{bp105}%
\begin{bpstatementit}{Proposition~B.11}{ (Pullback of invertible sheaves by \{$[n]$\})}
Let \(A\) be an abelian variety over \(k\) and let \(\mathcal L\) be an invertible sheaf on \(A\). For every integer \(n\geq0\),

\[
[n]^*\mathcal L\simeq
\mathcal L^{\otimes n(n+1)/2}\otimes
([-1]^*\mathcal L)^{\otimes n(n-1)/2}.
\]

In particular \([n]^*\mathcal N\simeq\mathcal N^{\otimes n^2}\) if \(\mathcal N\) is symmetric.
\end{bpstatementit}

% proof: proof/proposition-pullback-multiplication-n
\begin{proof}[{Proof}]\phantomsection\label{bp106}%
The idea is to produce one trivial sheaf \(\mathcal M\) on \(A\times A\times A\) by the theorem of the cube, and to pull it back along maps \(A\to A\times A\times A\) chosen so that the result is a relation between \([n+1]^*\mathcal L\), \([n]^*\mathcal L\) and \([n-1]^*\mathcal L\). The theorem of the cube applies because \(A\) is geometrically integral (Lemma \hyperref[bp100]{B.9}).

Consider on \(A\times A\times A\) the invertible sheaf

\[
\mathcal M=m_{123}^*\mathcal L\otimes m_{12}^*\mathcal L^{-1}\otimes
m_{13}^*\mathcal L^{-1}\otimes m_{23}^*\mathcal L^{-1}\otimes
p_1^*\mathcal L\otimes p_2^*\mathcal L\otimes p_3^*\mathcal L,
\]

where the \(p_i\) are the projections and \(m_{12}\), \(m_{13}\), \(m_{23}\), \(m_{123}\) are the sums of the indicated coordinates. To apply the theorem of the cube, we restrict \(\mathcal M\) to the three slices. On \(e\times A\times A\), with coordinates \((x_2,x_3)\), the seven factors restrict to \(\mathcal L\) pulled back along \(x_2+x_3\), \(x_2\), \(x_3\), \(x_2+x_3\), the constant map \(e\), \(x_2\) and \(x_3\), with exponents \(1,-1,-1,-1,1,1,1\). Of these, the pullback along the constant map is trivial, and the others cancel. The same holds on \(A\times e\times A\) and \(A\times A\times e\). Hence \(\mathcal M\) is trivial by Theorem \hyperref[bp91]{B.6}. Finally, pulling the trivial sheaf \(\mathcal M\) back along \(x\mapsto(x,x,[n-1]x)\) for \(n\ge0\) yields

\[
[n+1]^*\mathcal L\otimes\mathcal L^{\otimes2}
\otimes[n-1]^*\mathcal L
\simeq [2]^*\mathcal L\otimes([n]^*\mathcal L)^{\otimes2}.
\]

To see this, note that the composites of \(x\mapsto(x,x,[n-1]x)\) with \(m_{123}\), \(m_{12}\), \(m_{13}\), \(m_{23}\), \(p_1\), \(p_2\), \(p_3\) are \([n+1]\), \([2]\), \([n]\), \([n]\), \(\mathrm{id}\), \(\mathrm{id}\), \([n-1]\). The cases \(n=0\) and \(n=1\) of the formula are clear. For \(n=0\), since \([0]^*\mathcal L\) is trivial, the relation reads \(\mathcal L^{\otimes3}\otimes[-1]^*\mathcal L\simeq[2]^*\mathcal L\). This is the case \(n=2\) of the formula. In general, the relation expresses \([n+1]^*\mathcal L\) in terms of \([n]^*\mathcal L\) and \([n-1]^*\mathcal L\), and the exponents in the formula satisfy the same second order recursion. Induction on \(n\) therefore gives the formula [Stacks, \href{https://stacks.math.columbia.edu/tag/0BFF}{0BFF}]. In particular, for symmetric \(\mathcal N\) the two exponents add up to \(n^2\).
\end{proof}

% text: prose-B-10
\phantomsection\label{bp107}%
For a symmetric ample \(\mathcal N\), the degree of \([n]\) can be read off from the growth of \(\chi(A,\mathcal N^{\otimes m})\), since pulling back by \([n]\) multiplies \(m\) by \(n^2\) and \(\chi\) grows like \(m^g\).

% object: proposition-multiplication-n (blueprint/chap-Abelian-varieties/sec-Multiplication-by-n/proposition-multiplication-n/_section.md)
\phantomsection\label{bp108}%
\begin{bpstatementit}{Proposition~B.12}{ (Multiplication by $n$)}
Let \(A\) be an abelian variety of dimension \(g\) over \(k\). For every \(n>0\), the morphism \([n]\) is surjective and finite locally free of degree \(n^{2g}\), and its kernel \(A[n]\) has rank \(n^{2g}\) over \(k\). If \(n\) is invertible in \(k\), then \([n]\) and \(A[n]\to\operatorname{Spec}k\) are finite étale.
\end{bpstatementit}

% proof: proof/proposition-multiplication-n
\begin{proof}[{Proof}]\phantomsection\label{bp109}%
We use the isomorphism \([n]^*\mathcal N\simeq\mathcal N^{\otimes n^2}\) twice: on curves in the fibres of \([n]\) it shows that \([n]\) is finite, and in Euler characteristics it computes the degree. In between, translation invariance shows that all fibres have the same length, which makes \([n]\) finite locally free. Finally, whether \([n]\) is étale is decided by its differential.

Since finiteness, flatness and the degree may be checked after base change to an algebraic closure, we may assume that \(k\) is algebraically closed. Choose a symmetric ample \(\mathcal N\) (Lemma \hyperref[bp103]{B.10}); then \([n]^*\mathcal N\simeq\mathcal N^{\otimes n^2}\) by Proposition \hyperref[bp105]{B.11}. If a fibre of \([n]\) had positive dimension, it would contain a proper integral curve. Since \([n]\) maps that curve to a point, the restriction of \([n]^*\mathcal N\) to it is trivial, whereas the restriction of \(\mathcal N^{\otimes n^2}\) has positive degree. Hence every fibre has dimension zero, and since \([n]\) is proper, it is finite. Its image is a closed subset of dimension \(g\), hence equal to the integral variety \(A\) (Lemma \hyperref[bp100]{B.9}). After base change to an algebraic closure of its residue field, every fibre of \([n]\) becomes a translate of \(A[n]\). Indeed, if \(y=nx_0\), then translation by \(x_0\) carries \(A[n]=[n]^{-1}(0)\) isomorphically onto \([n]^{-1}(y)\), because \([n]\) is a homomorphism. Hence all fibres have the same degree \(r=\dim_k\Gamma(A[n],\mathcal O_{A[n]})\).

To get flatness from the constancy of the fibre degree, we use the following general fact. A finite morphism \(f\) to a reduced Noetherian scheme whose fibres all have degree \(r\) is finite locally free of rank \(r\). Indeed, near any point Nakayama's lemma gives a surjection \(\mathcal O^{\oplus r}\to f_*\mathcal O\), since the fibre at that point has dimension \(r\). This surjection is an isomorphism at each generic point, where the local ring is a field and both sides are vector spaces of dimension \(r\). Hence on an affine open every coordinate of a section of its kernel lies in all minimal primes, and is therefore zero because the ring is reduced. Since \(A\) is smooth, hence reduced, \([n]\) is finite locally free.

Recall that for a locally free sheaf \(\mathcal E\) of rank \(r\) on \(A\), the function \(m\mapsto\chi(A,\mathcal E\otimes\mathcal N^{\otimes m})\) is a polynomial of degree \(g\) with leading coefficient \(r\deg_{\mathcal N}(A)/g!\). Here \(\deg_{\mathcal N}(A)=(\mathcal N^g)\) is the degree of \(A\) with respect to \(\mathcal N\), a positive integer because \(\mathcal N\) is ample. We compare the coefficients of \(m^g\) in two expressions for \(\chi(A,[n]^*\mathcal N^{\otimes m})\). Since \([n]\) is finite, this Euler characteristic equals \(\chi(A,[n]_*[n]^*\mathcal N^{\otimes m})\), which by the projection formula is \(\chi(A,\mathcal N^{\otimes m}\otimes[n]_*\mathcal O_A)\). As \([n]_*\mathcal O_A\) is locally free of rank \(\deg[n]\), the coefficient is \(\deg[n]\deg_{\mathcal N}(A)/g!\). Since \([n]^*\mathcal N^{\otimes m}\simeq\mathcal N^{\otimes n^2m}\), it is also \(n^{2g}\deg_{\mathcal N}(A)/g!\). Hence

\[
\deg[n]\,\deg_{\mathcal N}(A)
=n^{2g}\deg_{\mathcal N}(A).
\]

Since \(\deg_{\mathcal N}(A)>0\), we conclude that \(\deg[n]=n^{2g}\); compare [Stacks, \href{https://stacks.math.columbia.edu/tag/0BFG}{0BFG}]. It remains to treat étaleness. The differential of the group law at \((0,0)\) is addition of tangent vectors. By induction, the differential of \([n]\) at the identity is \(n\operatorname{id}\). Since \([n]\) commutes with translations, the same holds at every point. If \(n\) is invertible, \([n]\) is therefore étale. Base change along the identity section gives the assertion for \(A[n]\).
\end{proof}

% text: prose-B-11
\phantomsection\label{bp110}%
It follows that the group \(A(\bar k)\) is divisible. Indeed, since \([n]\) is finite and surjective, its fibre over a point of \(A(\bar k)\) is a nonempty finite \(\bar k\)-scheme, and so it has a \(\bar k\)-point.


% object: sec-Tate-modules (blueprint/chap-Abelian-varieties/sec-Tate-modules/_section.md), content hidden
\subsection*{\texorpdfstring{B.3\quad Tate modules}{B.3 Tate modules}}
\phantomsection\label{bp111}%
\addcontentsline{toc}{subsection}{\texorpdfstring{B.3\quad Tate modules}{B.3 Tate modules}}
% text: prose-B-12
\phantomsection\label{bp112}%
Fix a prime \(\ell\) invertible in \(k\). With this prime, compatible systems of \(\ell\)-power torsion points of \(A\) form the Tate module \(T_\ell A\), a \(\mathbf Z_\ell\)-module with an action of \(\operatorname{Gal}(\bar k/k)\). For example, for a complex torus \(\mathbf C^g/\Lambda\) it is \(\Lambda\otimes\mathbf Z_\ell\), the first homology with \(\ell\)-adic coefficients. In general, Proposition \hyperref[bp114]{B.14} shows that \(T_\ell A\) is free of rank \(2g\) over \(\mathbf Z_\ell\). The input is the orders of the groups \(A[\ell^n](\bar k)\), which are known from Proposition \hyperref[bp108]{B.12}, and these orders force each of these groups to be free over \(\mathbf Z/\ell^n\) of rank \(2g\). Moreover, the transition maps are surjective because \(A(\bar k)\) is divisible.

Although Frobenius can have many fixed points on the finite levels, no power of it has a nonzero fixed vector in the Tate module (Proposition \hyperref[bp117]{B.15}). For example, for the elliptic curve \(E_0:y^2=x^3-x\) over \(\mathbf F_3\), the three roots of \(x^3-x\) lie in \(\mathbf F_3\). So all of \(E[2]\) is rational, and Frobenius acts trivially on it. On \(T_2E\), however, a fixed vector would be an unbounded compatible system of \(2\)-power torsion points in the finite group \(E_0(\mathbf F_{3^m})\). For the Jacobian of a curve, Proposition \bpout{14.3} turns this into the statement that Frobenius has no fixed class in \(H^1\) of the curve.

% object: definition-tate-module (blueprint/chap-Abelian-varieties/sec-Tate-modules/definition-tate-module/_section.md)
\phantomsection\label{bp113}%
\begin{bpstatement}{Definition~B.13}{ (Tate module)}
Let \(A\) be an abelian variety over \(k\) and let \(\ell\) be a prime invertible in \(k\). The \emph{Tate module} of \(A\) is

\[
T_\ell A=\varprojlim_n A[\ell^n](\bar k),
\]

the limit taken along the maps \([\ell]:A[\ell^{n+1}](\bar k)\to
A[\ell^n](\bar k)\), and we put \(V_\ell A=T_\ell A\otimes_{\mathbf Z_\ell}\mathbf Q_\ell\).
\end{bpstatement}

% object: proposition-integral-tate-module (blueprint/chap-Abelian-varieties/sec-Tate-modules/proposition-integral-tate-module/_section.md)
\phantomsection\label{bp114}%
\begin{bpstatementit}{Proposition~B.14}{ (The integral Tate module)}
Let \(A\) be an abelian variety of dimension \(g\) over \(k\), and let \(\ell\) be a prime invertible in \(k\). There are noncanonical isomorphisms

\[
A[\ell^n](\bar k)\simeq(\mathbf Z/\ell^n)^{2g},\qquad
T_\ell A\simeq\mathbf Z_\ell^{2g}.
\]

The canonical projection induces an isomorphism \(T_\ell A/\ell^nT_\ell A\simeq A[\ell^n](\bar k)\). The groups \(A[\ell^n](\bar k)\), the transition maps and the projections from \(T_\ell A\) are \(\operatorname{Gal}(\bar k/k)\)-equivariant, and they do not change when \(\bar k\) is replaced by a larger algebraically closed field.
\end{bpstatementit}

% proof: proof/proposition-integral-tate-module
\begin{proof}[{Proof}]\phantomsection\label{bp115}%
A finite abelian group of order \(\ell^{2gn}\) killed by \(\ell^n\) need not be free over \(\mathbf Z/\ell^n\); for \(g\geq1\) the group \((\mathbf Z/\ell^2)^{2g-2}\oplus(\mathbf Z/\ell)^{4}\) has order \(\ell^{4g}\) and is killed by \(\ell^2\). For \(A[\ell^n]\), however, this cannot happen because we know the orders at every level together with the exact sequences given by \([\ell]\). Indeed, we use them to lift a basis of the \(\ell\)-torsion step by step and check by counting that the lifts form a basis at each level.

\emph{The \(\ell\)-torsion and the exact sequences.} The group \(A[\ell](\bar k)\) is a vector space of dimension \(2g\) over \(\Fell\), since it is killed by \(\ell\) and has order \(\ell^{2g}\) by Proposition \hyperref[bp108]{B.12}. Moreover there are exact sequences

\[
0\longrightarrow A[\ell](\bar k)
\longrightarrow A[\ell^{n+1}](\bar k)
\xrightarrow{[\ell]}A[\ell^n](\bar k)\longrightarrow0,
\]

whose surjectivity holds because \(A(\bar k)\) is divisible: every point of \(A[\ell^n](\bar k)\) is \(\ell\) times a point, which then lies in \(A[\ell^{n+1}](\bar k)\).

\emph{Lifting bases.} Choose a basis of \(A[\ell](\bar k)\) and lift it successively along \([\ell]\); the lifts form a basis at each level. Suppose that \(x_1,\ldots,x_{2g}\in A[\ell^{n+1}](\bar k)\) lift a basis of \(A[\ell^n](\bar k)\) obtained in this way, so that \(\ell^nx_1,\ldots,\ell^nx_{2g}\) is the chosen basis of \(A[\ell](\bar k)\). To see that the \(x_i\) are independent, note first that all coefficients of a relation between them are divisible by \(\ell^n\), as we see by applying \([\ell]\) to it. The relation then becomes a relation between the elements \(\ell^nx_i\) of the chosen basis of \(A[\ell](\bar k)\), so the coefficients are even divisible by \(\ell^{n+1}\). As \(A[\ell^{n+1}](\bar k)\) has order \(\ell^{2g(n+1)}\), the \(x_i\) therefore form a basis of it over \(\mathbf Z/\ell^{n+1}\). Passing to the inverse limit of these compatible bases, we obtain both isomorphisms and the description of \(T_\ell A/\ell^nT_\ell A\).

\emph{Equivariance and independence of \(\bar k\).} The groups and maps in the statement commute with the action of \(\operatorname{Gal}(\bar k/k)\), because the group schemes \(A[\ell^n]\) and the maps \([\ell]\) are defined over \(k\). In particular the projection \(T_\ell A/\ell^nT_\ell A\to A[\ell^n](\bar k)\) is equivariant. The chosen bases need not be equivariant, which is why the isomorphisms with \((\mathbf Z/\ell^n)^{2g}\) and \(\mathbf Z_\ell^{2g}\) are noncanonical. Finally, the points of \(A[\ell^n]\) and the multiplication maps do not change under a further algebraically closed extension, because \(A[\ell^n]\) is finite étale and hence a finite disjoint union of points over \(\bar k\).
\end{proof}

% text: prose-B-13
\phantomsection\label{bp116}%
Over a finite field, Frobenius acts on \(T_\ell A\), and by the last proposition this action is compatible with its action on each \(A[\ell^n](\bar k)\). The next proposition shows that no power of Frobenius has a nonzero fixed vector. For the Jacobian of a curve, this is the input to Proposition \bpout{14.3}.

% object: proposition-no-frobenius-fixed-vector (blueprint/chap-Abelian-varieties/sec-Tate-modules/proposition-no-frobenius-fixed-vector/_section.md)
\phantomsection\label{bp117}%
\begin{bpstatementit}{Proposition~B.15}{ (No Frobenius-fixed vector)}
Let \(A_0/\mathbf F_q\) be an abelian variety, let \(A=A_0\times_{\mathbf F_q}\overline{\mathbf F}_q\), and let \(F\) be geometric Frobenius on \(T_\ell A\). For every \(m>0\),

\[
(T_\ell A)^{F^m}=0,
\qquad (V_\ell A)^{F^m}=0.
\]

Consequently \(F^m-1\) is an automorphism of the finite-dimensional \(\mathbf Q_\ell\)-vector space \(V_\ell A\), and hence of \(V_\ell A\otimes_{\mathbf Q_\ell}E\) for every field extension \(E\) of \(\mathbf Q_\ell\).
\end{bpstatementit}

% proof: proof/proposition-no-frobenius-fixed-vector
\begin{proof}[{Proof}]\phantomsection\label{bp118}%
A fixed vector is a compatible system of torsion points, each rational over \(\mathbf F_{q^m}\). But the rational points form a finite group, so the orders of the \(\ell\)-power torsion points in it are bounded, and compatibility under \([\ell]\) then forces the system to vanish. The statement for \(V_\ell A\) then follows by clearing denominators.

Since an automorphism and its inverse have the same fixed vectors, we may use arithmetic Frobenius instead. On \(A(\overline{\mathbf F}_q)\) its \(m\)th power fixes exactly \(A_0(\mathbf F_{q^m})\). This group is finite, because on each affine open of a finite affine cover a rational point is determined by finitely many coordinates in the finite field. Hence the \(\ell\)-primary part of \(A_0(\mathbf F_{q^m})\) is killed by some \(\ell^e\). Let \((t_n)\) be a fixed compatible sequence. Each \(t_{n+e}\) lies in this \(\ell\)-primary part, and since the transition maps are \([\ell]\), we find

\[
t_n=\ell^e t_{n+e}=0
\]

for every \(n\). This proves that \((T_\ell A)^{F^m}=0\).

Since \(T_\ell A\) is free by Proposition \hyperref[bp114]{B.14}, it injects into \(V_\ell A\). A suitable power of \(\ell\) times a fixed vector of \(V_\ell A\) is a fixed vector of \(T_\ell A\), and so it vanishes. Hence \((V_\ell A)^{F^m}=0\). In other words, \(F^m-1\) is injective on the finite-dimensional space \(V_\ell A\), and therefore bijective.
\end{proof}



% object: chap-The-Picard-scheme-of-a-curve (blueprint/chap-The-Picard-scheme-of-a-curve/_section.md), content hidden
\section*{\texorpdfstring{C\quad The Picard scheme of a curve}{C The Picard scheme of a curve}}
\phantomsection\label{bp119}%
\addcontentsline{toc}{section}{\texorpdfstring{C\quad The Picard scheme of a curve}{C The Picard scheme of a curve}}
% text: prose-C-01
\phantomsection\label{bp120}%
For a curve \(C\) proper and connected over a separably closed field and \(n\) invertible in that field, the Kummer sequence shows that \(H^1(C,\mu_n)\) is the \(n\)-torsion of the Picard group of \(C\). But the Picard group is an abstract group. In particular, nothing about it says that its \(n\)-torsion is finite, let alone isomorphic to \((\mathbf Z/n)^{2g}\). Nor does it say why Frobenius should have no fixed vector on the \(\ell\)-adic limit, as Proposition \bpout{14.3} requires. Both questions become statements of Appendix \hyperref[bp75]{B} once we know that \(\operatorname{Pic}^0(C)\), the group of invertible sheaves of degree zero, is the group of points of an abelian variety.

Over \(\mathbf C\) this is the Abel–Jacobi theorem, \(\operatorname{Pic}^0(C)\simeq
H^0(C,\Omega^1)^\vee/H_1(C,\mathbf Z)\). Over an arbitrary field we construct the Jacobian as a Picard scheme, which represents invertible sheaves in families.

Let \(C\) be a smooth proper connected curve of genus \(g\) over a separably closed field \(k\). The idea is that most invertible sheaves of degree \(g\) are given by divisors. Suppose \(\mathcal L\) has degree \(g\) and \(h^0(\mathcal L)=1\). Then the zeros of its unique section up to scalars form an effective divisor of degree \(g\), and this divisor determines \(\mathcal L\). Effective divisors of degree \(g\) are the points of the Hilbert scheme \(\operatorname{Hilb}^{g}_{C/k}\), a smooth proper variety of dimension \(g\) (Theorem \hyperref[bp139]{C.6}).

By Riemann–Roch, every invertible sheaf of degree \(g\) has \(h^0\geq1\), with equality exactly when \(h^1=0\). This condition is open in families, by semicontinuity, and it defines an open chart \(W\subset\operatorname{Hilb}^{g}_{C/k}\) of the Picard functor (Proposition \hyperref[bp142]{C.7}). Riemann–Roch also moves every invertible sheaf into this chart by a twist defined over \(k\) (Lemma \hyperref[bp145]{C.8}). Hence the translates of \(W\) cover the Picard functor, and they glue to the Picard scheme. Its components are proper because the Abel map \(\operatorname{Hilb}^{g}_{C/k}\to\operatorname{Pic}^g\) is surjective and its source is proper.

The Picard functor itself is a sheafified quotient of \(T\mapsto\operatorname{Pic}(C\times T)\). We have to divide out the invertible sheaves pulled back from \(T\), since these are trivial locally on \(T\). When \(C\) has a rational point, we describe the Picard functor by invertible sheaves trivialized along that point (Lemma \hyperref[bp125]{C.2}).

The outcome is Theorem \hyperref[bp154]{C.11}: the component of degree zero, the Jacobian \(J(C)\), is an abelian variety of dimension \(g\) with \(J(C)(k)=\operatorname{Pic}^0(C)\) and tangent space \(H^1(C,\mathcal O_C)\) at the origin. Over a finite field we descend the Jacobian to an abelian variety \(J_0\) over \(\mathbf F_q\). Then Propositions \hyperref[bp167]{C.16} and \hyperref[bp170]{Proposition C.17} give an isomorphism \(\varprojlim_nH^1(C,\mathbf Z/\ell^n(1))\simeq T_\ell J_0\) compatible with Frobenius, which is the form in which Proposition \bpout{14.3} uses the Jacobian; see also [\hyperref[bp350]{Mil86b}].

% object: sec-The-Picard-functor (blueprint/chap-The-Picard-scheme-of-a-curve/sec-The-Picard-functor/_section.md), content hidden
\subsection*{\texorpdfstring{C.1\quad The Picard functor}{C.1 The Picard functor}}
\phantomsection\label{bp121}%
\addcontentsline{toc}{subsection}{\texorpdfstring{C.1\quad The Picard functor}{C.1 The Picard functor}}
% text: prose-C-02
\phantomsection\label{bp122}%
We want a scheme whose points over \(T\) are families of invertible sheaves on \(C\) parametrized by \(T\). The obvious candidate \(T\mapsto\operatorname{Pic}(C_T)\), however, contains too much. For example, for \(C=\mathbf P^1_k\) and connected \(T\), every invertible sheaf on \(\mathbf P^1_T\) is \(\mathcal O(d)\otimes f_T^*\mathcal M\) for a unique \(d\in\mathbf Z\) and a unique \(\mathcal M\in\operatorname{Pic}(T)\). In other words \(\operatorname{Pic}(\mathbf P^1_T)=\mathbf Z\oplus\operatorname{Pic}(T)\), whereas the Picard scheme should be the constant scheme \(\mathbf Z\). In fact the functor \(T\mapsto\operatorname{Pic}(C_T)\) is not even a Zariski sheaf, because the pullback \(f_T^*\mathcal M\) of a nontrivial invertible sheaf \(\mathcal M\) on \(T\) is trivial locally on \(T\) but not globally, so no scheme represents this functor.

The Picard functor is therefore defined as the sheafification of \(T\mapsto\operatorname{Pic}(C_T)/f_T^*\operatorname{Pic}(T)\) (Definition \hyperref[bp123]{C.1}). When \(C\) has a section, Lemma \hyperref[bp125]{C.2} removes the pullbacks from \(T\) in another way, by fixing a trivialization along the section. We use this rigidified description to construct the chart of Proposition \hyperref[bp142]{C.7} and the Picard scheme of Theorem \hyperref[bp154]{C.11}.

% object: definition-relative-picard-functor (blueprint/chap-The-Picard-scheme-of-a-curve/sec-The-Picard-functor/definition-relative-picard-functor/_section.md)
\phantomsection\label{bp123}%
\begin{bpstatement}{Definition~C.1}{ (Relative Picard functor)}
For a morphism of schemes \(f:C\to S\), let

\[
\operatorname{Pic}_{C/S}
=\bigl(T\longmapsto
\operatorname{Pic}(C_T)/f_T^*\operatorname{Pic}(T)\bigr)^{\mathrm{fppf}}.
\]

Here \(C_T=C\times_ST\), and the superscript denotes sheafification on the big fppf site of \(S\). When this functor is represented by a scheme, we call the representing scheme the \emph{Picard scheme}.
\end{bpstatement}

% text: prose-C-03
\phantomsection\label{bp124}%
Equivalently, one may sheafify \(T\mapsto\operatorname{Pic}(C_T)\), since an invertible sheaf on \(T\) becomes trivial on a Zariski cover. The sheafification is needed: invertible sheaves satisfy effective descent, but their automorphisms can prevent their isomorphism classes from forming a sheaf. For this reason neither the displayed quotient nor \(T\mapsto\operatorname{Pic}(C_T)\) need be an fppf sheaf [Stacks, \href{https://stacks.math.columbia.edu/tag/0B9L}{0B9L}]. If \(f\) has a section, the next lemma describes the sheafification directly. We also write \(\operatorname{Pic}_{C/S}\) for the Picard scheme when it exists.

% object: lemma-rigidification-along-section (blueprint/chap-The-Picard-scheme-of-a-curve/sec-The-Picard-functor/lemma-rigidification-along-section/_section.md)
\phantomsection\label{bp125}%
\begin{bpstatementit}{Lemma~C.2}{ (Rigidification along a section)}
Suppose that \(f\) has a section \(\sigma:S\to C\) and that

\[
\mathcal O_T\longrightarrow f_{T,*}\mathcal O_{C_T}
\quad\text{is an isomorphism for every }T\to S.
\tag{$\ast$}
\]

Then \(\operatorname{Pic}_{C/S}(T)\) is identified, functorially in \(T\), with the set of isomorphism classes of pairs \((\mathcal L,\alpha)\), with \(\mathcal L\) invertible on \(C_T\) and \(\alpha:\sigma_T^*\mathcal L\simeq\mathcal O_T\). Every such pair has only the identity automorphism. There is a split exact sequence

\[
0\longrightarrow\operatorname{Pic}(T)
\xrightarrow{f_T^*}\operatorname{Pic}(C_T)
\longrightarrow\operatorname{Pic}_{C/S}(T)
\longrightarrow0,
\]

whose retraction is \(\sigma_T^*\). The normalization map is

\[
\mathcal L\longmapsto
\mathcal L\otimes f_T^*(\sigma_T^*\mathcal L)^{-1}.
\]
\end{bpstatementit}

% proof: proof/lemma-rigidification-along-section
\begin{proof}[{Proof}]\phantomsection\label{bp126}%
The point is that rigidified pairs have no automorphisms other than the identity. Because of this they form a sheaf, and the quotient presheaf maps isomorphically onto this sheaf.

An automorphism of an invertible sheaf is multiplication by a global unit. By condition \((\ast)\) this unit comes from a unit on \(T\), and compatibility with a rigidification forces it to be \(1\). To see that the rigidified functor is a sheaf, suppose we are given rigidified classes on an fppf cover that agree on overlaps. We then choose local isomorphisms and rescale them so that they preserve the rigidifications. Such isomorphisms are unique, since rigidified pairs have no nontrivial automorphisms. In particular they satisfy the cocycle condition, and the pairs glue by effective descent for invertible sheaves.

The normalization map identifies the quotient presheaf \(T\mapsto\operatorname{Pic}(C_T)/f_T^*\operatorname{Pic}(T)\) with this rigidified functor. Since the latter is already a sheaf, the quotient presheaf needs no sheafification and equals \(\operatorname{Pic}_{C/S}\). The sequence splits since \(\sigma_T^*f_T^*=\mathrm{id}\); for this argument see [Stacks, \href{https://stacks.math.columbia.edu/tag/0B9N}{0B9N}].
\end{proof}

% text: prose-C-04
\phantomsection\label{bp127}%
For a smooth proper geometrically connected curve \(C\) over a field \(k\), we have \(H^0(C,\mathcal O_C)=k\). For every field extension \(k'/k\) the curve \(C_{k'}\) is again smooth, proper and integral, and flat base change for coherent cohomology gives \(H^0(C_{k'},\mathcal O_{C_{k'}})=k'\) and \(\dim_{k'}H^1(C_{k'},\mathcal O_{C_{k'}})=\dim_kH^1(C,\mathcal O_C)\). Every \(k\)-scheme is flat over \(k\), so the same flat base change shows that \((\ast)\) holds for \(C\to\operatorname{Spec}k\). If \(C\) has a \(k\)-point \(P\), the preceding lemma applies with \(\sigma=P\) (see [Stacks, \href{https://stacks.math.columbia.edu/tag/0B9U}{0B9U}]).


% object: sec-Riemann-Roch-and-Serre-duality (blueprint/chap-The-Picard-scheme-of-a-curve/sec-Riemann-Roch-and-Serre-duality/_section.md), content hidden
\subsection*{\texorpdfstring{C.2\quad Riemann–Roch and Serre duality}{C.2 Riemann–Roch and Serre duality}}
\phantomsection\label{bp128}%
\addcontentsline{toc}{subsection}{\texorpdfstring{C.2\quad Riemann–Roch and Serre duality}{C.2 Riemann–Roch and Serre duality}}
% text: prose-C-05
\phantomsection\label{bp129}%
The chart of the Picard scheme in degree \(g\) is cut out by the condition \(h^1=0\). Let \(\mathcal L\) be an invertible sheaf of degree \(g\) on a curve of genus \(g\). Then Riemann–Roch gives \(h^0(\mathcal L)-h^1(\mathcal L)=1\), so \(h^1(\mathcal L)=0\) if and only if \(\mathcal L\) has exactly one section up to scalars. In a family over \(T\), coherent base change shows that the condition “\(Rf_{T,*}\mathcal L\) is an invertible sheaf in degree zero” is open on \(T\) and compatible with base change. It therefore defines an open subfunctor of the Picard functor.

The condition is not closed, because \(h^0\) can jump up under specialization. For example, on a curve of genus two a divisor \(D\) of degree two has \(h^0(\mathcal O(D))=1\) unless \(\mathcal O(D)\simeq\omega_{C/k}\), in which case \(h^0=2\). The chart therefore misses the line of canonical divisors in \(\operatorname{Hilb}^{2}_{C/k}\), and we need its translates to cover the Picard scheme in degree two.

Fix a smooth proper geometrically connected curve \(C\) over \(k\), with structure morphism \(f:C\to\operatorname{Spec}k\), and set \(g=\dim_kH^1(C,\mathcal O_C)\). Such a curve is projective [Stacks, \href{https://stacks.math.columbia.edu/tag/0A26}{0A26}].

% object: theorem-riemann-roch-coherent-serre-duality (blueprint/chap-The-Picard-scheme-of-a-curve/sec-Riemann-Roch-and-Serre-duality/theorem-riemann-roch-coherent-serre-duality/_section.md)
\phantomsection\label{bp130}%
\begin{bpstatementit}{Theorem~C.3}{ (Riemann–Roch and coherent Serre duality)}
For an invertible sheaf \(\mathcal L\) on \(C\) we have

\[
\chi(C,\mathcal L)=\deg\mathcal L+1-g,
\qquad
H^1(C,\mathcal L)^\vee
\simeq H^0(C,\omega_{C/k}\otimes\mathcal L^{-1}),
\qquad \deg\omega_{C/k}=2g-2,
\qquad \dim_kH^0(C,\omega_{C/k})=g.
\]

In particular, if \(\deg\mathcal L>2g-2\), then

\[
H^1(C,\mathcal L)=0,
\qquad h^0(C,\mathcal L)=\deg\mathcal L+1-g.
\]
\end{bpstatementit}

% text: prose-C-06
\phantomsection\label{bp131}%
For \(C=\mathbf P^1_k\) and \(\mathcal L=\mathcal O(d)\) with \(d\geq0\), the sections are the binary forms of degree \(d\), a space of dimension \(d+1=\deg\mathcal L+1-0\), and \(\omega=\mathcal O(-2)\) has degree \(-2=2g-2\).

We quote the theorem from [Stacks, \href{https://stacks.math.columbia.edu/tag/0B5B}{0B5B}], where it is proved for smooth curves and the duality is given by the residue trace on differentials. We prove its compatibility with field extension in Corollary \hyperref[bp132]{C.4} below. For the displayed vanishing, recall that a nonzero section of an invertible sheaf on an integral curve has an effective divisor of zeros. Hence an invertible sheaf of negative degree has no nonzero section. If \(\deg\mathcal L>2g-2\), then \(\omega_{C/k}\otimes\mathcal L^{-1}\) has negative degree, and \(H^1(C,\mathcal L)=0\) by duality. The dimension formula then follows from Riemann–Roch. For \(\mathcal L=\mathcal O_C\) the duality gives \(\dim_kH^0(C,\omega_{C/k})=\dim_kH^1(C,\mathcal O_C)=g\).

The Picard scheme is constructed over a separably closed field, but its points are tested over arbitrary extensions. Hence Riemann–Roch and coherent base change must be compatible with extension of the ground field, and for duality we need this for the pairing itself, not only for the dimensions.

% object: corollary-serre-duality-field-extension (blueprint/chap-The-Picard-scheme-of-a-curve/sec-Riemann-Roch-and-Serre-duality/corollary-serre-duality-field-extension/_section.md)
\phantomsection\label{bp132}%
\begin{bpstatementit}{Corollary~C.4}{ (Serre duality and field extension)}
Let \(k'/k\) be a field extension. Under flat base change for coherent cohomology and the identification \(\omega_{C_{k'}/k'}\simeq\omega_{C/k}|_{C_{k'}}\), the Serre duality pairing for \(\mathcal L|_{C_{k'}}\) is obtained from that for \(\mathcal L\) by extension of scalars. In particular, the genus is unchanged by field extension.
\end{bpstatementit}

% proof: proof/corollary-serre-duality-field-extension
\begin{proof}[{Proof}]\phantomsection\label{bp133}%
We factor the duality on \(C\) through a finite map to \(\mathbf P^1_k\), where the trace is an explicit residue, and check that each step is compatible with extension of scalars.

A nonconstant rational function on \(C\) extends to a finite map \(\phi:C\to\mathbf P^1_k\), which is flat because \(\phi_*\mathcal O_C\) is torsion free over the discrete valuation rings of \(\mathbf P^1_k\). We may construct the duality on \(C\) as duality for \(\phi\), given by the dual of the finite locally free module \(\phi_*\mathcal L\), followed by duality on \(\mathbf P^1_k\), whose trace \(H^1(\mathbf P^1_k,\omega)\to k\) takes the coefficient of \(t^{-1}dt\) in a Čech cocycle for the standard cover. Both steps commute with extension of scalars, hence so does the resulting pairing. Flat base change shows that \(g=\dim_kH^1(C,\mathcal O_C)\) is unchanged by field extension.
\end{proof}

% text: prose-C-07
\phantomsection\label{bp134}%
Recall that the chart of the Picard scheme consists of the invertible sheaves \(\mathcal L\) for which \(Rf_{T,*}\mathcal L\) is an invertible sheaf in degree \(0\). We use coherent base change in the following form.

% object: lemma-coherent-base-change-invertible-sheaves (blueprint/chap-The-Picard-scheme-of-a-curve/sec-Riemann-Roch-and-Serre-duality/lemma-coherent-base-change-invertible-sheaves/_section.md)
\phantomsection\label{bp135}%
\begin{bpstatementit}{Lemma~C.5}{ (Coherent base change for invertible sheaves)}
Let \(T\) be a \(k\)-scheme and let \(\mathcal L\) be invertible on \(C_T\). Then \(Rf_{T,*}\mathcal L\) is locally represented by a bounded complex of finite locally free \(\mathcal O_T\)-modules. For any \(a:T'\to T\), the base change map \(La^*Rf_{T,*}\mathcal L\to Rf_{T',*}\mathcal L_{T'}\) is an isomorphism in the derived category of \(\mathcal O_{T'}\)-modules. The locus where this complex is an invertible module in degree \(0\) is open. On this locus, \(f_{T,*}\mathcal L\) is invertible, its formation commutes with arbitrary base change, and \(R^1f_{T,*}\mathcal L=0\).
\end{bpstatementit}

% proof: proof/lemma-coherent-base-change-invertible-sheaves
\begin{proof}[{Proof}]\phantomsection\label{bp136}%
This is base change for coherent cohomology along a proper morphism, applied to a flat finitely presented module, together with the openness criterion for a perfect complex (see [Stacks, \href{https://stacks.math.columbia.edu/tag/0B91}{0B91}, \href{https://stacks.math.columbia.edu/tag/0B9S}{0B9S}]). To construct the complex, one may use a projective embedding, sufficiently large twists and a finite affine Čech complex; by Serre vanishing we then obtain models that are finite locally free. Given such a model, openness follows from a calculation with matrices. Indeed, near a point \(t\) where the derived fibre of the complex has cohomology only in degree \(0\), of dimension \(1\), we may cancel the differentials that are invertible at \(t\), and what remains is an invertible module in degree \(0\).

This proves the lemma for locally Noetherian \(T\). If \(T\) is an arbitrary \(k\)-scheme, then locally on \(T\) the sheaf \(\mathcal L\) is pulled back from \(C_{T_0}\) for some \(k\)-scheme \(T_0\) of finite type, and the base change isomorphism for \(T\to T_0\) reduces every assertion to \(T_0\). The open locus is the preimage of the locus for \(T_0\), and by base change it does not depend on \(T_0\).
\end{proof}


% object: sec-Hilbert-schemes-of-points (blueprint/chap-The-Picard-scheme-of-a-curve/sec-Hilbert-schemes-of-points/_section.md), content hidden
\subsection*{\texorpdfstring{C.3\quad Hilbert schemes of points}{C.3 Hilbert schemes of points}}
\phantomsection\label{bp137}%
\addcontentsline{toc}{subsection}{\texorpdfstring{C.3\quad Hilbert schemes of points}{C.3 Hilbert schemes of points}}
% text: prose-C-08
\phantomsection\label{bp138}%
On a smooth curve a finite subscheme of length \(d\) is the same as an effective divisor of degree \(d\). These divisors are parametrized by the Hilbert scheme of points \(\operatorname{Hilb}^{d}_{C/k}\), which is the symmetric power \(C^{(d)}\). Iterated addition of points gives a finite surjection \(C^d\to \operatorname{Hilb}^{d}_{C/k}\), from which \(\operatorname{Hilb}^{d}_{C/k}\) inherits properness, irreducibility and dimension \(d\) (Theorem \hyperref[bp139]{C.6}). For example, for \(C=\mathbf P^1\), an effective divisor of degree \(d\) is the zero set of a binary form of degree \(d\) up to scalars, and \(\operatorname{Hilb}^{d}_{C/k}=\mathbf P^d\).

A divisor \(D\) gives the invertible sheaf \(\mathcal O(D)\) with its canonical section. Conversely, an invertible sheaf with a section up to scalars gives a divisor. When \(h^0=1\) the section is unique up to scalars, and then the divisor is determined by the sheaf. In this way the open part of \(\operatorname{Hilb}^{g}_{C/k}\) where \(h^0(\mathcal O(D))=1\) and \(h^1=0\) becomes an open subfunctor of the Picard functor (Proposition \hyperref[bp142]{C.7}).

For \(d\geq0\), let \(\operatorname{Hilb}^d_{C/k}\) be the functor on \(k\)-schemes sending \(T\) to the set of closed subschemes \(D\subset C_T\) that are finite locally free of rank \(d\) over \(T\).

% object: theorem-hilbert-schemes-points-smooth-curve (blueprint/chap-The-Picard-scheme-of-a-curve/sec-Hilbert-schemes-of-points/theorem-hilbert-schemes-points-smooth-curve/_section.md)
\phantomsection\label{bp139}%
\begin{bpstatementit}{Theorem~C.6}{ (Hilbert schemes of points on a smooth curve)}
The functor \(\operatorname{Hilb}^{d}_{C/k}\) is represented by a smooth proper geometrically integral \(k\)-scheme of dimension \(d\), with \(\operatorname{Hilb}^{0}_{C/k}=\operatorname{Spec}k\). Its \(T\)-points are precisely the relative effective Cartier divisors that are finite locally free of degree \(d\). It has a universal divisor \(D_d\subset C\times_k\operatorname{Hilb}^{d}_{C/k}\) and an associated invertible sheaf \(\mathcal O(D_d)\).
\end{bpstatementit}

% proof: proof/theorem-hilbert-schemes-points-smooth-curve
\begin{proof}[{Proof}]\phantomsection\label{bp140}%
Representability is a general theorem, once we know that finite sets of points lie in affine opens. We then obtain the geometric properties from the addition map \(\operatorname{Hilb}^{d}_{C/k}\times C\to \operatorname{Hilb}^{d+1}_{C/k}\), which we show to be finite locally free. Starting from \(\operatorname{Hilb}^{1}_{C/k}=C\), we get smoothness and the dimension by induction on \(d\), and properness and irreducibility from the finite surjection \(C^d\to \operatorname{Hilb}^{d}_{C/k}\).

Every finite set of points of \(C\) lies in an affine open. Indeed, if \(Q\) is a closed point outside the set, then \(\mathcal O_C(Q)\) has positive degree, hence is ample by Riemann–Roch, and so \(C\setminus\{Q\}\) is affine. Hence \(\operatorname{Hilb}^{d}_{C/k}\) is represented by a scheme, separated and locally of finite type over \(k\) [Stacks, \href{https://stacks.math.columbia.edu/tag/0B9A}{0B9A}]. The universal closed subscheme is the divisor \(D_d\).

Let \(T\) be a \(k\)-scheme and \(t\in T\) a point. Since the local rings of \(C_{\kappa(t)}\) at closed points are discrete valuation rings, every finite closed subscheme of a fibre is an effective Cartier divisor. By the fibrewise criterion for relative effective Cartier divisors, every \(D\in \operatorname{Hilb}^{d}_{C/k}(T)\) is a relative effective Cartier divisor on \(C_T/T\). This gives the description of \(T\)-points and the invertible sheaf \(\mathcal O(D_d)\). Multiplying ideals adds such divisors, and we obtain the addition map

\[
\operatorname{Hilb}^{d}_{C/k}\times C\longrightarrow \operatorname{Hilb}^{d+1}_{C/k},\qquad (D,x)\longmapsto D+\Gamma_x.
\]

It identifies \(\operatorname{Hilb}^{d}_{C/k}\times C\) with the universal divisor over \(\operatorname{Hilb}^{d+1}_{C/k}\): if a divisor \(E\) contains the graph \(\Gamma_x\), then dividing a local equation of \(E\) by one of \(\Gamma_x\) gives a relative effective Cartier divisor \(D\) with \(E=D+\Gamma_x\), of degree \(d\) by the exact sequence for a sum of divisors. In other words, the addition map is the projection to \(\operatorname{Hilb}^{d+1}_{C/k}\) of the universal divisor \(D_{d+1}\subset C\times \operatorname{Hilb}^{d+1}_{C/k}\). By definition of the Hilbert functor, \(D_{d+1}\) is finite locally free of degree \(d+1\) over \(\operatorname{Hilb}^{d+1}_{C/k}\). Hence the addition map is finite locally free of degree \(d+1\) [Stacks, \href{https://stacks.math.columbia.edu/tag/0B9G}{0B9G}], in particular faithfully flat.

We now argue by induction on \(d\), starting from \(\operatorname{Hilb}^{0}_{C/k}=\operatorname{Spec}k\) and \(\operatorname{Hilb}^{1}_{C/k}=C\). The scheme \(\operatorname{Hilb}^{d+1}_{C/k}\) is smooth because regularity descends along faithfully flat maps, also after extending \(k\). It has dimension \(d+1\) because the addition map is finite and surjective [Stacks, \href{https://stacks.math.columbia.edu/tag/0B9H}{0B9H}].

Properness comes from the finite surjective map \(C^d\to \operatorname{Hilb}^{d}_{C/k}\) given by iterated addition. This map makes \(\operatorname{Hilb}^{d}_{C/k}\) quasi-compact. It also makes \(\operatorname{Hilb}^{d}_{C/k}\) universally closed, because \(C^d\) is proper and surjectivity is stable under base change. Since \(C^d\) is geometrically integral, \(\operatorname{Hilb}^{d}_{C/k}\) is geometrically irreducible. Being smooth, it is then geometrically integral [Stacks, \href{https://stacks.math.columbia.edu/tag/0B9I}{0B9I}].

For \(C=\mathbf P^1\), the map \(C^d\to \operatorname{Hilb}^{d}_{C/k}=\mathbf P^d\) sends \(d\) points to the binary form with these zeros, that is, the product of \(d\) linear forms; it is the quotient by the symmetric group, of degree \(d!\).
\end{proof}

% text: prose-C-09
\phantomsection\label{bp141}%
Assume now that \(C(k)\) contains a point \(P\), and let us describe \(\operatorname{Pic}_{C/k}\) by rigidified invertible sheaves as in Lemma \hyperref[bp125]{C.2}. We define a subfunctor \(F\subset\operatorname{Pic}_{C/k}\) as follows. A rigidified class \(\mathcal L\) on \(C_T\) lies in \(F(T)\) if, in the derived category of \(\mathcal O_T\)-modules,

\[
Rf_{T,*}\mathcal L\simeq\mathcal N[0]
\quad\text{for some invertible }\mathcal N\text{ on }T.
\]

By Lemma \hyperref[bp135]{C.5}, this condition defines an open locus of \(T\) compatible with base change, so that the inclusion \(F\subset\operatorname{Pic}_{C/k}\) is representable by open immersions.

% object: proposition-divisors-degree-g-open-chart (blueprint/chap-The-Picard-scheme-of-a-curve/sec-Hilbert-schemes-of-points/proposition-divisors-degree-g-open-chart/_section.md)
\phantomsection\label{bp142}%
\begin{bpstatementit}{Proposition~C.7}{ (Divisors of degree $g$ as an open chart)}
The functor \(F\) is represented by an open subscheme \(W\subset \operatorname{Hilb}^{g}_{C/k}\). The morphism \(W\to\operatorname{Pic}_{C/k}\) associates to a divisor its rigidified invertible sheaf. Consequently, \(F\) is smooth of dimension \(g\).
\end{bpstatementit}

% proof: proof/proposition-divisors-degree-g-open-chart
\begin{proof}[{Proof}]\phantomsection\label{bp143}%
We construct mutually inverse maps: a divisor gives its invertible sheaf, normalized along \(P\), and a family in \(F\) has, up to units on \(T\), a unique section in each fibre, whose zero scheme is a family of divisors of degree \(g\).

\emph{From divisors to sheaves.} On \(C\times \operatorname{Hilb}^{g}_{C/k}\), normalize the universal invertible sheaf as

\[
\mathcal L_g=\mathcal O(D_g)\otimes
f_{\operatorname{Hilb}^g}^*(P_{\operatorname{Hilb}^g}^*\mathcal O(D_g))^{-1}.
\]

Let \(W\) be the open locus where \(Rf_{\operatorname{Hilb}^g,*}\mathcal L_g\) is a line bundle in degree \(0\); then \(\mathcal L_g\) defines a morphism \(W\to F\).

\emph{From sheaves to divisors.} Let \(\mathcal L\) be a rigidified class in \(F(T)\) and set \(\mathcal N=f_{T,*}\mathcal L\). By assumption \(Rf_{T,*}\mathcal L\simeq\mathcal N'[0]\) for an invertible \(\mathcal N'\), so \(\mathcal N\simeq\mathcal N'\) is invertible. The evaluation map \(f_T^*\mathcal N\to\mathcal L\) is a section \(s\) of \(\mathcal L\otimes f_T^*\mathcal N^{-1}\). It is determined by \(\mathcal L\) up to a unit on \(T\), and such a unit does not change its zero scheme \(D\).

This \(D\) is a relative effective Cartier divisor. Indeed, the formation of \(\mathcal N\) commutes with base change (Lemma \hyperref[bp135]{C.5}). Hence on each fibre \(C_{\kappa(t)}\) the section \(s\) is a nonzero element of the one-dimensional space of sections. The fibre is integral because \(C\) is geometrically integral, so the zero scheme of \(s\) on the fibre is an effective Cartier divisor. Moreover, the fibres of \(D\) have degree \(g\): for a point \(t\) of \(T\) we have \(h^0(\mathcal L_t)=1\) and \(H^1(C_t,\mathcal L_t)=0\), so Riemann–Roch gives

\[
\deg\mathcal L_t
=\chi(\mathcal L_t)-(1-g)=1-(1-g)=g.
\]

The morphism \(D\to T\) is proper, since \(D\) is closed in the proper scheme \(C_T\). It has finite fibres, and it is flat, as is every relative effective Cartier divisor. A proper, flat, finitely presented morphism with finite fibres is finite locally free. Its rank at \(t\) is the degree \(g\) of the fibre \(D_t\), the divisor of zeros of a section of \(\mathcal L_t\). Hence \(D\) defines a morphism \(T\to \operatorname{Hilb}^{g}_{C/k}\).

\emph{The constructions are inverse.} Normalizing \(\mathcal O(D)\) along \(P_T\) recovers \(\mathcal L\), because \(\mathcal O(D)\simeq\mathcal L\otimes f_T^*\mathcal N^{-1}\). The two constructions are therefore inverse to each other, and they commute with base change, which proves that \(W\simeq F\) (see [Stacks, \href{https://stacks.math.columbia.edu/tag/0B9X}{0B9X}]).
\end{proof}

% text: prose-C-10
\phantomsection\label{bp144}%
To cover the Picard functor by translates of \(F\), we have to move every invertible sheaf into \(F\) by a twist defined over \(k\).

% object: lemma-twisting-h-0-1-h (blueprint/chap-The-Picard-scheme-of-a-curve/sec-Hilbert-schemes-of-points/lemma-twisting-h-0-1-h/_section.md)
\phantomsection\label{bp145}%
\begin{bpstatementit}{Lemma~C.8}{ (Twisting to $h^0=1$ and $h^1=0$)}
Suppose that \(k\) is separably closed. For every field extension \(k'/k\) and every invertible sheaf \(\mathcal L\) on \(C_{k'}\), there is an invertible sheaf \(\mathcal M_0\) on \(C\) such that

\[
h^0(C_{k'},\mathcal L\otimes\mathcal M_{0,k'})=1,
\qquad H^1(C_{k'},\mathcal L\otimes\mathcal M_{0,k'})=0.
\]
\end{bpstatementit}

% proof: proof/lemma-twisting-h-0-1-h
\begin{proof}[{Proof}]\phantomsection\label{bp146}%
We first twist by a large power of an ample sheaf on \(C\), which kills \(H^1\) and leaves many sections. Then we remove \(k\)-rational points one at a time, each time lowering \(h^0\) by one and keeping \(H^1=0\), until \(h^0=1\). The twist is defined over \(k\) because the points are.

Choose an ample invertible sheaf \(\mathcal M\) on \(C\). For \(r\gg0\), Serre vanishing and Riemann–Roch show that \(\mathcal L\otimes\mathcal M_{k'}^{\otimes r}\) has vanishing \(H^1\) and a positive-dimensional space of sections. As long as \(h^0\) is larger than \(1\), we choose a nonzero section and a point \(x\in C(k)\) where that section is nonzero. Such a point exists: the zero divisor of the section is finite, so its image in \(C\) is a finite set, and the complement of this set is a nonempty open of \(C\). This open contains a point of \(C(k)\), because the \(k\)-rational points of a smooth \(k\)-scheme are dense when \(k\) is separably closed [Stacks, \href{https://stacks.math.columbia.edu/tag/056U}{056U}]. For the current sheaf \(\mathcal F\), consider the sequence

\[
0\longrightarrow\mathcal F(-x_{k'})
\longrightarrow\mathcal F
\longrightarrow\mathcal F|_{x_{k'}}\longrightarrow0
\]

The chosen section does not vanish at \(x_{k'}\), so the map \(H^0(C_{k'},\mathcal F)\to\mathcal F|_{x_{k'}}\) is onto the one-dimensional fibre. By the long exact sequence, subtracting \(x_{k'}\) lowers \(h^0\) by exactly \(1\) and preserves \(H^1=0\). We continue in this way until \(h^0=1\). Since all subtracted points come from \(k\), the resulting twist has the form \(\mathcal M_0=\mathcal M^{\otimes r}(-x_1-\cdots-x_s)\) on \(C\). Finally, normalizing along \(P\) does not change the dimensions of fibre cohomology, so the rigidified twist satisfies the same conditions (see [Stacks, \href{https://stacks.math.columbia.edu/tag/0B9Y}{0B9Y}]).
\end{proof}


% object: sec-The-Jacobian (blueprint/chap-The-Picard-scheme-of-a-curve/sec-The-Jacobian/_section.md), content hidden
\subsection*{\texorpdfstring{C.4\quad The Jacobian}{C.4 The Jacobian}}
\phantomsection\label{bp147}%
\addcontentsline{toc}{subsection}{\texorpdfstring{C.4\quad The Jacobian}{C.4 The Jacobian}}
% text: prose-C-11
\phantomsection\label{bp148}%
We glue the translates of the chart of Proposition \hyperref[bp142]{C.7} to the Picard scheme, whose component of degree zero is the Jacobian. Gluing uses twice that \(k\) is separably closed: to find a rational point for the rigidification, and to obtain the twists of Lemma \hyperref[bp145]{C.8} over \(k\). For an elliptic curve, the Abel map \(C=\operatorname{Hilb}^{1}_{C/k}\to\operatorname{Pic}^1_{C/k}\) is an isomorphism, and translation by \(\mathcal O_C(-P)\) identifies \(C\) with its own Jacobian. In general, the tangent space of the Jacobian at the origin is computed by deforming the trivial sheaf over the dual numbers.

% object: lemma-representability-picard-functor (blueprint/chap-The-Picard-scheme-of-a-curve/sec-The-Jacobian/lemma-representability-picard-functor/_section.md)
\phantomsection\label{bp149}%
\begin{bpstatementit}{Lemma~C.9}{ (Representability of the Picard functor)}
Let \(k\) be separably closed and let \(C/k\) be a smooth proper connected curve. The fppf Picard functor \(\operatorname{Pic}_{C/k}\) is represented by a scheme locally of finite type over \(k\), which is covered by the open subschemes representing the translates \(M^{-1}F\), for \(M\in\operatorname{Pic}_{C/k}(k)\), of the chart \(F\) of Proposition \hyperref[bp142]{C.7}, defined using a point \(P\in C(k)\).
\end{bpstatementit}

% proof: proof/lemma-representability-picard-functor
\begin{proof}[{Proof}]\phantomsection\label{bp150}%
The preceding lemmas assume \(C\) geometrically integral. To see that this holds, note that \(\bar k/k\) is purely inseparable. Hence \(C_{\bar k}\to C\) is a homeomorphism, and \(C_{\bar k}\) is connected and smooth, hence integral. We rigidify using a point \(P\in C(k)\). Such a point exists because \(k\) is separably closed and the \(k\)-points of the smooth curve \(C\) are dense [Stacks, \href{https://stacks.math.columbia.edu/tag/056U}{056U}].

By Proposition \hyperref[bp142]{C.7} we have a representable open subfunctor \(F\subset\operatorname{Pic}_{C/k}\). For a fixed class \(M\in\operatorname{Pic}_{C/k}(k)\), the translate \(M^{-1}F\), consisting of the classes \(\mathcal L\) with \(\mathcal L\otimes M\) in \(F\), is another representable open subfunctor.

We show that these translates cover \(\operatorname{Pic}_{C/k}\). Let a class on a \(k\)-scheme \(T\) and a point \(t\in T\) be given. We apply Lemma \hyperref[bp145]{C.8} to its fibre over the residue field \(k'=\kappa(t)\). This gives an invertible sheaf \(\mathcal M_0\) on \(C\), defined over \(k\), such that the twisted fibre has \(h^0=1\) and \(H^1=0\). By Lemma \hyperref[bp135]{C.5}, the twisted class then lies in \(F\) over an open neighbourhood of \(t\). Letting \(t\) vary, we see that the open loci where the corresponding translated class lies in \(F\) cover \(T\). Note that the translates are indexed by classes over \(k\), while the points of \(T\) have arbitrary residue fields. This is why Lemma \hyperref[bp145]{C.8} asks for a twist defined over \(k\) for every extension \(k'\).

The intersections of the charts are open in either chart, and their transition isomorphisms satisfy the cocycle condition by functoriality. An fppf sheaf covered by representable open subfunctors is representable [Stacks, \href{https://stacks.math.columbia.edu/tag/0B9Q}{0B9Q}, \href{https://stacks.math.columbia.edu/tag/0B9Z}{0B9Z}]. Hence gluing the charts gives a scheme representing \(\operatorname{Pic}_{C/k}\). It is locally of finite type because each chart is an open of \(\operatorname{Hilb}^{g}_{C/k}\).
\end{proof}

% text: prose-C-12
\phantomsection\label{bp151}%
We compute the tangent space with the dual numbers \(k[\varepsilon]\), \(\varepsilon^2=0\); the computation works for any scheme over \(k\).

% object: lemma-tangent-space-picard-functor (blueprint/chap-The-Picard-scheme-of-a-curve/sec-The-Jacobian/lemma-tangent-space-picard-functor/_section.md)
\phantomsection\label{bp152}%
\begin{bpstatementit}{Lemma~C.10}{ (First-order deformations of invertible sheaves)}
Let \(X\) be a scheme over a field \(k\) and \(X_\varepsilon=X\times_k\operatorname{Spec}k[\varepsilon]\). The kernel of the restriction map \(\operatorname{Pic}(X_\varepsilon)\to\operatorname{Pic}(X)\) is \(H^1(X,\mathcal O_X)\).
\end{bpstatementit}

% proof: proof/lemma-tangent-space-picard-functor
\begin{proof}[{Proof}]\phantomsection\label{bp153}%
Because \(\varepsilon^2=0\), the map \(a\mapsto1+\varepsilon a\) is a homomorphism, and it gives an exact sequence of sheaves of groups on \(X\)

\[
0\longrightarrow\mathcal O_X\longrightarrow
\mathcal O_{X_\varepsilon}^\times\longrightarrow\mathcal O_X^\times
\longrightarrow1,
\]

split by the inclusion \(\mathcal O_X^\times\subset
\mathcal O_{X_\varepsilon}^\times\). The schemes \(X\) and \(X_\varepsilon\) have the same underlying space, so \(\operatorname{Pic}(X_\varepsilon)=H^1(X,\mathcal O_{X_\varepsilon}^\times)\). Since the sequence is split, taking \(H^1\) gives \(H^1(X,\mathcal O_{X_\varepsilon}^\times)=H^1(X,\mathcal O_X)\oplus\operatorname{Pic}(X)\), in which the restriction map is the projection to the second summand. Hence the kernel is \(H^1(X,\mathcal O_X)\).
\end{proof}

% object: theorem-picard-scheme-jacobian-over-separably (blueprint/chap-The-Picard-scheme-of-a-curve/sec-The-Jacobian/theorem-picard-scheme-jacobian-over-separably/_section.md)
\phantomsection\label{bp154}%
\begin{bpstatementit}{Theorem~C.11}{ (The Picard scheme and the Jacobian over a separably closed field)}
Let \(k\) be separably closed and let \(C/k\) be a smooth proper connected curve of genus \(g\). The fppf Picard functor is represented by a smooth separated group scheme, locally of finite type over \(k\), with a decomposition

\[
\operatorname{Pic}_{C/k}
=\coprod_{d\in\mathbf Z}\operatorname{Pic}^d_{C/k}.
\]

Each component is a smooth proper geometrically connected variety of dimension \(g\), parametrizing invertible sheaves of degree \(d\). The identity component \(J(C)=\operatorname{Pic}^0_{C/k}\) is an abelian variety, and

\[
J(C)(k)=\operatorname{Pic}^0(C),
\qquad T_0J(C)=H^1(C,\mathcal O_C).
\]

The formation of this representing scheme, its degree components and \(J(C)\) commutes with extension of the ground field.
\end{bpstatementit}

% proof: proof/theorem-picard-scheme-jacobian-over-separably
\begin{proof}[{Proof}]\phantomsection\label{bp155}%
Representability is Lemma \hyperref[bp149]{C.9}. The tangent space is given by Lemma \hyperref[bp152]{C.10}. Fix the point \(P\in C(k)\) used to define \(F\).

For an invertible sheaf \(\mathcal L\) on \(C_T\), the Euler characteristic \(\chi(C_t,\mathcal L_t)\) is a locally constant function of \(t\in T\). This holds because \(C_T\to T\) is proper and flat and \(Rf_{T,*}\mathcal L\) is perfect (Lemma \hyperref[bp135]{C.5}). By Riemann–Roch the degree of \(\mathcal L_t\) is locally constant as well. This gives the displayed decomposition into open and closed pieces. The charts in \(\operatorname{Hilb}^{g}_{C/k}\) and their translates show that the scheme is smooth of dimension \(g\).

It remains to see that the scheme is separated. A \(k\)-rational point of a scheme locally of finite type over a field is closed. Hence the unit section of the group scheme is a closed immersion. The diagonal is the preimage of the unit section under \((x,y)\mapsto xy^{-1}\), so it is closed as well.

The universal divisor defines the Abel map

\[
\operatorname{Hilb}^{g}_{C/k}\longrightarrow\operatorname{Pic}^g_{C/k},
\qquad D\longmapsto\mathcal O_C(D).
\]

For a line bundle \(\mathcal L\) of degree \(g\) on \(C_{k'}\), with \(k'\) algebraically closed, Riemann–Roch gives \(h^0(\mathcal L)\geq\chi(\mathcal L)=g+1-g=1\). A nonzero section has an effective divisor \(D\) of degree \(g\) with \(\mathcal O(D)\simeq\mathcal L\). Hence the Abel map is surjective on points over algebraically closed fields, and this suffices for surjectivity. It follows that \(\operatorname{Pic}^g_{C/k}\) is quasi-compact. It is of finite type since the charts are locally of finite type. A separated scheme of finite type that is the surjective image of the proper scheme \(\operatorname{Hilb}^{g}_{C/k}\) is proper. It is geometrically irreducible because \(\operatorname{Hilb}^{g}_{C/k}\) is. Translation by \(\mathcal O_C((d-g)P)\) identifies \(\operatorname{Pic}^g_{C/k}\) with the component of each degree \(d\). In particular, the component of degree zero is smooth, proper, connected and commutative. Such a group scheme is an abelian variety (see [Stacks, \href{https://stacks.math.columbia.edu/tag/0BA0}{0BA0}]).

The tangent space \(T_0J(C)\) is the kernel of \(\operatorname{Pic}_{C/k}(k[\varepsilon])\to\operatorname{Pic}_{C/k}(k)\), where \(\varepsilon^2=0\). Since \(\operatorname{Pic}(k[\varepsilon])=0\), this is the kernel of \(\operatorname{Pic}(C_{k[\varepsilon]})\to\operatorname{Pic}(C)\), which is \(H^1(C,\mathcal O_C)\) by Lemma \hyperref[bp152]{C.10}.

For a field extension \(k'/k\), the Picard functor of \(C_{k'}\) is the restriction of \(\operatorname{Pic}_{C/k}\) to \(k'\)-schemes. Hence its representing scheme is the base change. The degree of \(\mathcal L_t\) is computed on the fibre \(C_t\), so the decomposition by degree, and with it \(J(C)\), is compatible with this base change.
\end{proof}

% text: prose-C-13
\phantomsection\label{bp156}%
We use the theorem in the following form: over an algebraically closed field, the Picard group of \(C\) is an extension of \(\mathbf Z\) by the group of points of an abelian variety.

% object: corollary-degree-jacobian (blueprint/chap-The-Picard-scheme-of-a-curve/sec-The-Jacobian/corollary-degree-jacobian/_section.md)
\phantomsection\label{bp157}%
\begin{bpstatementit}{Corollary~C.12}{ (Degree and the Jacobian)}
Let \(C\) be a smooth proper connected curve over an algebraically closed field \(k\). There is an exact sequence

\[
0\longrightarrow J(C)(k)
\longrightarrow\operatorname{Pic}(C)
\xrightarrow{\deg}\mathbf Z\longrightarrow0,
\]

and a point \(P\in C(k)\) splits it by \(d\mapsto\mathcal O_C(dP)\).
\end{bpstatementit}

% proof: proof/corollary-degree-jacobian
\begin{proof}[{Proof}]\phantomsection\label{bp158}%
The degree is additive under tensor product, so it is a homomorphism. By Theorem \hyperref[bp154]{C.11}, its kernel is \(\operatorname{Pic}^0(C)=J(C)(k)\). Since \(\deg\mathcal O_C(dP)=d\), the map \(d\mapsto\mathcal O_C(dP)\) is a section of the degree, which is therefore surjective.
\end{proof}

% object: example-degree-sequence-genus-zero-one (blueprint/chap-The-Picard-scheme-of-a-curve/sec-The-Jacobian/example-degree-sequence-genus-zero-one/_section.md)
\phantomsection\label{bp159}%
\begin{bpstatement}{Example~C.13}{ (The degree sequence in genus zero and one)}
On \(C=\mathbf P^1_k\) a divisor \(\sum_in_ia_i+n_\infty\infty\) of degree zero, with \(a_i\in k\), is the divisor of the rational function \(\prod_i(x-a_i)^{n_i}\), whose order at \(\infty\) is \(-\sum_in_i=n_\infty\). So \(\operatorname{Pic}^0(\mathbf P^1)=0\), and the degree is an isomorphism \(\operatorname{Pic}(\mathbf P^1)\simeq\mathbf Z\). This agrees with \(J(\mathbf P^1)=0\) for \(g=0\).

On an elliptic curve \(E\) with origin \(O\) the kernel is as large as the curve. Let \(\mathcal L\) have degree zero. Then \(\mathcal L(O)\) has degree one, and by Riemann–Roch with \(g=1\) it has a one-dimensional space of sections. Hence \(\mathcal L(O)\simeq\mathcal O_E(P)\) for a unique point \(P\), and \(P\mapsto\mathcal O_E(P-O)\) is a bijection \(E(k)\to\operatorname{Pic}^0(E)\). In particular \(\mathcal O_E(P-Q)\) is nontrivial for \(P\neq Q\). One can also see this directly: a rational function with divisor \(P-Q\) would be a map \(E\to\mathbf P^1\) of degree one, that is, an isomorphism.
\end{bpstatement}


% object: sec-Torsion-points (blueprint/chap-The-Picard-scheme-of-a-curve/sec-Torsion-points/_section.md), content hidden
\subsection*{\texorpdfstring{C.5\quad Torsion points}{C.5 Torsion points}}
\phantomsection\label{bp160}%
\addcontentsline{toc}{subsection}{\texorpdfstring{C.5\quad Torsion points}{C.5 Torsion points}}
% text: prose-C-14
\phantomsection\label{bp161}%
The Kummer sequence on a curve \(C\) over an algebraically closed field identifies \(H^1(C,\mu_n)\) with \(\operatorname{Pic}(C)[n]\) and \(H^2(C,\mu_n)\) with \(\operatorname{Pic}(C)/n\) (Proposition \hyperref[bp179]{D.2}). To compute these groups, we combine two facts. First, the exact sequence of Corollary \hyperref[bp157]{C.12} separates the discrete part \(\mathbf Z\) of the Picard group from the Jacobian. Second, Appendix \hyperref[bp75]{B} describes multiplication by \(n\) on an abelian variety over an algebraically closed field: it is surjective on points, and for \(n\) invertible in \(k\) its kernel is \((\mathbf Z/n)^{2g}\). Now apply the snake lemma to multiplication by \(n\) on the degree sequence. Since \(\mathbf Z\) is torsion free, the \(n\)-torsion of \(\operatorname{Pic}(C)\) is \(J(C)[n](k)\simeq(\mathbf Z/n)^{2g}\). Since \(J(C)(k)\) is divisible, the cokernel of \(n\) on \(\operatorname{Pic}(C)\) is \(\mathbf Z/n\), given by the degree.

For proper base change we also need these groups to be unchanged under extension of the algebraically closed field. This holds because \(J(C)[n]\) is finite étale (Corollary \hyperref[bp163]{C.15}).

% object: example-two-torsion-elliptic-curve (blueprint/chap-The-Picard-scheme-of-a-curve/sec-Torsion-points/example-two-torsion-elliptic-curve/_section.md)
\phantomsection\label{bp162}%
\begin{bpstatement}{Example~C.14}{ (Two-torsion of an elliptic curve)}
Let \(k\) be algebraically closed of characteristic different from two and let \(E\colon y^2=(x-e_1)(x-e_2)(x-e_3)\) with distinct \(e_i\in k\), with origin \(O\) at infinity and \(P_i=(e_i,0)\). The function \(x-e_i\) has a double zero at \(P_i\), because \(y\) is a local parameter there. It has a double pole at \(O\). Hence \(\operatorname{div}(x-e_i)=2P_i-2O\), and the classes \(\mathcal O_E(P_i-O)\) are killed by \(2\). They are nonzero and pairwise distinct, since \(\mathcal O_E(P_i-P_j)\) is trivial only for \(i=j\) (Example \hyperref[bp159]{C.13}). The divisor \(\operatorname{div}(y)=P_1+P_2+P_3-3O\) shows that the three classes add up to zero. So \(0\) and the classes of \(P_i-O\) form a subgroup \(\{0,P_1-O,P_2-O,P_3-O\}\simeq(\mathbf Z/2)^2\) of \(\operatorname{Pic}(E)[2]\). By Corollary \hyperref[bp163]{C.15} below, \(\operatorname{Pic}(E)[2]\) has order \(2^{2g}=4\), so this subgroup is all of it.

Under the Kummer sequence, the class \(P_i-O\), with the trivialization of its square given by \(x-e_i\), is the \(\mu_2\)-torsor obtained by adjoining a square root of \(x-e_i\). This torsor is unramified everywhere because the divisor of \(x-e_i\) is even. For the running example \(y^2=x^3-x\) over \(\overline{\mathbf F}_3\) the \(e_i\) are \(0,1,-1\in\mathbf F_3\). Hence the three torsors are defined over \(\mathbf F_3\), and Frobenius fixes all of \(E[2]\), as noted in Section \hyperref[bp111]{B.3}.
\end{bpstatement}

% object: corollary-torsion-jacobian-field-extension (blueprint/chap-The-Picard-scheme-of-a-curve/sec-Torsion-points/corollary-torsion-jacobian-field-extension/_section.md)
\phantomsection\label{bp163}%
\begin{bpstatementit}{Corollary~C.15}{ (Torsion of the Jacobian and field extension)}
Let \(k\) be algebraically closed, let \(C/k\) be a smooth proper connected curve of genus \(g\), and let \(n>0\) be invertible in \(k\). Then there is a noncanonical isomorphism

\[
J(C)[n](k)\simeq(\mathbf Z/n\mathbf Z)^{2g}.
\]

For every algebraically closed extension \(k'/k\), the natural map

\[
J(C)[n](k)\longrightarrow J(C_{k'})[n](k')
\]

is an isomorphism. Multiplication by \(n\) on \(J(C)(k)\) is surjective, and consequently

\[
\operatorname{Pic}(C)[n]=J(C)[n](k),
\qquad
\operatorname{Pic}(C)/n\operatorname{Pic}(C)
\xrightarrow[\deg]{\ \sim\ }\mathbf Z/n\mathbf Z.
\]
\end{bpstatementit}

% proof: proof/corollary-torsion-jacobian-field-extension
\begin{proof}[{Proof}]\phantomsection\label{bp164}%
The assertions about \(J(C)\) are those of Appendix \hyperref[bp75]{B} for this abelian variety, and the degree sequence carries them over to \(\operatorname{Pic}(C)\).

By Proposition \hyperref[bp108]{B.12}, the kernel \(J(C)[n]\) is finite étale of rank \(n^{2g}\). A finite étale scheme of rank \(r\) over the algebraically closed field \(k\) is a disjoint union of \(r\) copies of \(\operatorname{Spec}k\), a description stable under base extension. By Theorem \hyperref[bp154]{C.11} the formation of the Jacobian and of \([n]\) commutes with base change, so the natural map \(J(C)[n](k)\to J(C_{k'})[n](k')\) is a bijection.

If \(\ell^m\) is the exact power of a prime \(\ell\) dividing \(n\), then \(J(C)[\ell^m](k)\simeq(\mathbf Z/\ell^m)^{2g}\) by Proposition \hyperref[bp114]{B.14} applied to \(J(C)\). Combining the parts for the primes dividing \(n\), we get the stated \(\mathbf Z/n\mathbf Z\)-module structure.

By Proposition \hyperref[bp108]{B.12}, the group \(J(C)(k)\) is divisible. In the exact sequence of Corollary \hyperref[bp157]{C.12}, the group \(\mathbf Z\) has no \(n\)-torsion and \(J(C)(k)/n=0\). This proves the two assertions for \(\operatorname{Pic}(C)\).
\end{proof}


% object: sec-Descent-to-the-ground-field (blueprint/chap-The-Picard-scheme-of-a-curve/sec-Descent-to-the-ground-field/_section.md), content hidden
\subsection*{\texorpdfstring{C.6\quad Descent to the ground field}{C.6 Descent to the ground field}}
\phantomsection\label{bp165}%
\addcontentsline{toc}{subsection}{\texorpdfstring{C.6\quad Descent to the ground field}{C.6 Descent to the ground field}}
% text: prose-C-15
\phantomsection\label{bp166}%
To compare the Galois action on the cohomology of a curve over \(\bar k\) with a Tate module, we need the Jacobian over the ground field \(k\). For example, in Proposition \bpout{14.3} we work over a finite field, and Frobenius acts on \(\varprojlim_nH^1(C,\mathbf Z/\ell^n(1))\). It has no fixed vector there because this module is \(T_\ell J_0\) for an abelian variety \(J_0\) over \(\mathbf F_q\) (Proposition \hyperref[bp117]{B.15}). The construction so far works over a separably closed field and uses a rational point for the rigidification. However, a curve over \(k\) need not have a \(k\)-point; for instance, the conic \(x^2+y^2+z^2=0\) over \(\mathbf R\) has none. We therefore define \(J_0\) as the scheme representing the degree-zero part of the fppf Picard sheaf. Concretely, we obtain it from the Jacobian over \(k_s\) by Galois descent through a finite Galois extension \(k'/k\) (Proposition \hyperref[bp167]{C.16}). The Kummer sequence then identifies \(H^1(C,\mathbf Z/\ell^n(1))\) with \(J_0[\ell^n](\bar k)\) compatibly with the Galois action and the transition maps (Proposition \hyperref[bp170]{C.17}).

% object: proposition-jacobian-over-ground-field (blueprint/chap-The-Picard-scheme-of-a-curve/sec-Descent-to-the-ground-field/proposition-jacobian-over-ground-field/_section.md)
\phantomsection\label{bp167}%
\begin{bpstatementit}{Proposition~C.16}{ (The Jacobian over the ground field)}
Let \(C_0/k\) be a smooth proper geometrically connected curve of genus \(g\), and put \(C=C_0\times_k\bar k\). There is an abelian variety \(J_0/k\) of dimension \(g\) representing the degree-zero relative Picard fppf sheaf of \(C_0\). For every field extension \(k'/k\), the base change \(J_0\times_kk'\) represents the degree-zero Picard sheaf of \(C_0\times_kk'\); in particular \(J_0\times_k\bar k\) represents \(\operatorname{Pic}^0_{C/\bar k}\), and the canonical identification

\[
J_0(\bar k)=\operatorname{Pic}^0(C)
\]

is Galois-equivariant.
\end{bpstatementit}

% proof: proof/proposition-jacobian-over-ground-field
\begin{proof}[{Proof}]\phantomsection\label{bp168}%
We descend the Jacobian of \(C_{0,k_s}\), where \(k_s\subseteq\bar
k\) is a separable closure of \(k\). In general \(k_s/k\) is infinite, hence not an fppf cover. For this reason we move the descent datum on the Jacobian over \(k_s\) by limit arguments to a finite Galois extension, and descend along that extension. A second limit argument shows that the descended scheme represents the degree-zero Picard sheaf.

Let \(F=\operatorname{Pic}^0_{C_0/k}\) be the degree-zero part of the fppf Picard sheaf. Here we use that the fibrewise degree is locally constant, by coherent base change and Riemann–Roch. Since the Picard sheaf and the degree commute with base change, Theorem \hyperref[bp154]{C.11} gives an abelian variety \(J_s\) over \(k_s\) representing the restriction of \(F\) to \(k_s\)-schemes; \(J_s\) is projective. Over \(k_s\otimes_kk_s\) the two pullbacks of \(J_s\) represent the same functor, so by the Yoneda lemma we get an isomorphism between them. All identifications are induced by \(F\). Hence this isomorphism satisfies the cocycle condition over \(k_s\otimes_kk_s\otimes_kk_s\) and respects the group structures.

To move this descent datum to a finite extension, note that \(k_s\) is the filtered union of its finite separable subextensions, and \(k_s\otimes_kk_s\) and the triple tensor product are the corresponding unions. By the limit theorems for schemes of finite presentation [Stacks, \href{https://stacks.math.columbia.edu/tag/01ZM}{01ZM}], the defining equations of \(J_s\) in a projective space, the descent isomorphism, its inverse and the group operations are defined over some finite stage. The finitely many identities among them (inverse, cocycle, group axioms) hold after enlarging the stage. Finally, we enlarge it further to a Galois extension. In this way we obtain a finite Galois extension \(K/k\) in \(k_s\) and a projective \(K\)-scheme \(J_K\) with a group structure and a compatible descent datum for \(K/k\), whose base change to \(k_s\) is the datum on \(J_s\). Since \(\operatorname{Spec}K\to\operatorname{Spec}k\) is finite locally free and surjective, the descent datum on the projective scheme \(J_K\) is effective [Stacks, \href{https://stacks.math.columbia.edu/tag/0CCJ}{0CCJ}]. This gives a commutative group scheme \(J_0\) over \(k\) with \(J_0\times_kk_s\simeq J_s\). Properness and smoothness descend along \(K/k\), and connectedness and dimension may be checked over \(\bar
k\). Hence \(J_0\) is an abelian variety of dimension \(g\).

The functors \(J_0\) and \(F\) are fppf sheaves that agree after base change to \(k_s\). Since \(k_s/k\) is not finite, we need a limit argument to show that they agree. Both functors are locally of finite presentation. For \(J_0\) this is [Stacks, \href{https://stacks.math.columbia.edu/tag/01ZC}{01ZC}]. For \(F\) it follows from Lemma \hyperref[bp125]{C.2} once the finite stages contain a field of definition of a point of \(C_0\), because invertible sheaves, rigidifications and their isomorphisms are of finite presentation. Hence, for affine \(T\) over \(k\), a section of either functor over \(T_{k_s}\), or over \(T_{k_s\otimes_kk_s}\), comes from a finite stage. Likewise, two sections that agree over \(T_{k_s}\) agree over a finite stage.

The identification \(J_0(T_{k_s})\simeq F(T_{k_s})\) is compatible with the descent data. Let a section of one functor over \(T\) be given. We transport its image over \(T_{k_s}\) and descend the result to some \(T_L\) with \(L/k\) finite. After enlarging \(L\), its two pullbacks to \(T_{L\otimes_kL}\) agree. Since both functors are fppf sheaves, it then descends to \(T\). The same argument applied to equalities shows that the result is unique. We obtain inverse bijections \(J_0(T)\simeq F(T)\), natural in affine \(T\), and hence an isomorphism of sheaves.

The Picard sheaf and the degree commute with field extension. Hence \(J_0\times_kk'\) represents the degree-zero Picard sheaf of \(C_0\times_kk'\) for every \(k'/k\). Over \(\bar k\) the value of the Picard sheaf is the Picard group, by rigidification and \(\operatorname{Pic}(\bar k)=0\). This gives the displayed identification. For \(\gamma\in\operatorname{Gal}(\bar k/k)\), naturality with respect to \(\operatorname{Spec}(\gamma)\) says that this group isomorphism intertwines the Galois action on points with pullback of invertible sheaves. Here \(\operatorname{Gal}(\bar k/k)=\operatorname{Gal}(k_s/k)\) because \(\bar k/k_s\) is purely inseparable.
\end{proof}

% text: prose-C-16
\phantomsection\label{bp169}%
Since \(J_0\) is defined over \(k\), the Galois group acts on its torsion points. Using the Kummer sequence on \(C\), we compare this action with the Galois action on cohomology.

% object: proposition-h-1-curve-tate-module (blueprint/chap-The-Picard-scheme-of-a-curve/sec-Descent-to-the-ground-field/proposition-h-1-curve-tate-module/_section.md)
\phantomsection\label{bp170}%
\begin{bpstatementit}{Proposition~C.17}{ ($H^1$ of a curve and the Tate module of its Jacobian)}
In the preceding setting, suppose \(\ell\) is invertible in \(k\). For every \(n\geq1\) there is a canonical Galois-equivariant isomorphism, with cohomology computed on the pro-étale site,

\[
H^1(C,\mathbf Z/\ell^n(1))
\simeq J_0[\ell^n](\bar k).
\]

The coefficient map \(\mu_{\ell^{n+1}}\to\mu_{\ell^n}\), \(z\mapsto z^\ell\), corresponds to multiplication by \(\ell\) on the torsion points. Consequently

\[
\varprojlim_n H^1(C,\mathbf Z/\ell^n(1))\simeq T_\ell J_0
\]

as \(\mathbf Z_\ell\)-modules with Galois action.
\end{bpstatementit}

% proof: proof/proposition-h-1-curve-tate-module
\begin{proof}[{Proof}]\phantomsection\label{bp171}%
The isomorphism comes from the Kummer sequence and the identification of the preceding proposition. The point is that both are compatible with the Galois action and with the transition maps. Then they pass to the inverse limit.

We compute on the étale site and transport the result to the pro-étale site by Theorem \bpout{2.35}. Since that comparison is natural, it preserves all the maps and Galois actions below. The Kummer sequence gives

\[
0\longrightarrow
H^0(C,\mathbf G_m)/\ell^n
\longrightarrow H^1(C,\mu_{\ell^n})
\longrightarrow\operatorname{Pic}(C)[\ell^n]
\longrightarrow0.
\]

The first term vanishes. Indeed, \(C\) is proper and connected, so \(H^0(C,\mathbf G_m)=\bar k^{\times}\), and this group is divisible. For the last term, note that the degree map takes values in the torsion-free group \(\mathbf Z\), so \(\operatorname{Pic}(C)[\ell^n]=\operatorname{Pic}^0(C)[\ell^n]\), and Proposition \hyperref[bp167]{C.16} identifies this group with \(J_0[\ell^n](\bar k)\). It remains to check Galois equivariance. The map induced by \(\mu_{\ell^n}\to\mathbf G_m\) and the identification of Proposition \hyperref[bp167]{C.16} are natural in the curve and in the field, and hence Galois-equivariant.

To check the transition maps, we map the Kummer sequence for \(\ell^{n+1}\) to that for \(\ell^n\) by \([\ell]\) on the first \(\mathbf G_m\) and by the identity on the second. On the kernels this is \(z\mapsto z^\ell\). On the Picard group it sends an invertible sheaf to its \(\ell\)-th tensor power, which is \([\ell]\) on \(J_0\). Therefore the displayed isomorphisms identify the two inverse systems, including their transition maps.
\end{proof}



% object: chap-Cohomology-of-torsion-sheaves-on-curves (blueprint/chap-Cohomology-of-torsion-sheaves-on-curves/_section.md), content hidden
\section*{\texorpdfstring{D\quad Cohomology of torsion sheaves on curves}{D Cohomology of torsion sheaves on curves}}
\phantomsection\label{bp172}%
\addcontentsline{toc}{section}{\texorpdfstring{D\quad Cohomology of torsion sheaves on curves}{D Cohomology of torsion sheaves on curves}}
% text: prose-D-01
\phantomsection\label{bp173}%
The last step of proper base change for \(\mathbf P^1\) compares the cohomology of a curve over two separably closed fields, one inside the other. Smooth base change needs the same comparison for affine curves, with coefficients prime to the characteristic. Both rest on Theorem \bpout{4.20}, whose assertions we recall. Let \(k\) be algebraically closed, let \(C\) be a separated scheme of finite type over \(k\) with \(\dim C\leq1\), and let \(F\) be a classical torsion abelian sheaf on \(C_{\mathrm{pro}\acute{e}t}\). Then:

\begin{itemize}
\item \((C_{1})\) \(H^q(C,F)=0\) for \(q>2\);
\item \((C_{2})\) if \(C\) is affine, \(H^q(C,F)=0\) for \(q>1\);
\item \((C_{3})\) if \(\operatorname{char}k=p>0\) and \(F\) is \(p\)-primary, \(H^q(C,F)=0\) for \(q>1\);
\item \((C_{4})\) if \(F\) is constructible with finite stalks and prime to \(\operatorname{char}k\), all \(H^q(C,F)\) are finite;
\item \((C_{5})\) if \(C\) is proper and \(F\) is constructible with finite stalks, all \(H^q(C,F)\) are finite;
\item \((C_{6})\) if \(k'/k\) is an extension of algebraically closed fields and \(F\) is prime to \(\operatorname{char}k\), the pullback maps \(H^q(C,F)\to H^q(C_{k'},F|)\) are isomorphisms;
\item \((C_{7})\) if \(C\) is proper, the maps in \((C_6)\) are isomorphisms for every torsion \(F\);
\item \((C_{8})\) for any open \(U\subseteq C\), restriction \(H^2(C,F)\to H^2(U,F|_U)\) is surjective.
\end{itemize}

Over \(\mathbf C\) the vanishing and finiteness assertions are statements about a Riemann surface, which has the homotopy type of a CW complex of dimension two, or of one if it is not compact.

In algebraic geometry we compute the groups. Let \(C\) be a smooth proper connected curve of genus \(g\) and let \(n\) be invertible in \(k\). Three ingredients compute the cohomology of \(\mu_n\): the Kummer sequence \(0\to\mu_n\to\mathbf G_m\xrightarrow{n}\mathbf G_m\to0\), the vanishing of \(H^q(C,\mathbf G_m)\) for \(q\geq2\), which is a consequence of Tsen's theorem, and the Jacobian of Appendix \hyperref[bp119]{C}. They give

\[
H^0(C,\mu_n)=\mu_n(k),\qquad H^1(C,\mu_n)=J(C)[n](k)\simeq(\mathbf Z/n)^{2g},
\qquad H^2(C,\mu_n)=\mathbf Z/n
\]

(Proposition \hyperref[bp179]{D.2}). These groups do not change under extension of algebraically closed fields, because torsion points of an abelian variety do not. For \(\mathbf P^1\), which has \(g=0\), the computation fixes the normalization of the class of \(\mathcal O(1)\) that Remark \hyperref[bp181]{D.3} provides for the proof of proper base change. Removing a nonempty finite set of \(r\) points kills \(H^2\) and raises the rank of \(H^1\) to \(2g+r-1\) (Proposition \hyperref[bp184]{D.4}).

Proper base change, however, is needed for all torsion coefficients, including those divisible by the characteristic \(p\), and there the answer changes. On the affine line, \(H^1(\mathbf A^1_k,\mathbf F_p)=k[t]/\{h^p-h\}\) is infinite, and it grows when \(k\) is enlarged (Example \hyperref[bp195]{D.7}). Indeed, the Artin–Schreier torsors \(y^p-y=at\), for \(a\in k\), have pairwise distinct classes, and over a larger field \(k'\) there are new ones with \(a\in k'\setminus k\). So for \(p\)-torsion coefficients, finiteness and invariance under field extension hold only for proper curves. This is why \((C_5)\) and \((C_7)\) of Theorem \bpout{4.20} assume properness, and why Lemma \bpout{4.21} is stated for proper curves.

The Artin–Schreier sequence \(0\to\mathbf F_p\to\mathbf G_a\xrightarrow{\varphi-1}\mathbf G_a\to0\) reduces the cohomology of \(\mathbf F_p\) to the coherent cohomology of \(\mathcal O_C\). On the affine line it presents \(H^1\) as the cokernel of \(h\mapsto h^p-h\) on \(k[t]\), as above. On every curve it gives \(H^q(C,\mathbf F_p)=0\) for \(q\geq2\) (Theorem \bpout{4.20}, \((C_3)\)).

A general torsion sheaf is reduced to a constant one by two observations. First, a sheaf supported on finitely many points has no cohomology in positive degrees. Hence passing to the normalization of a curve, or extending a sheaf by zero from a dense open subset, changes cohomology only in degrees zero and one, and only by finite groups (Lemmas \hyperref[bp205]{D.10} and \hyperref[bp210]{Lemma D.11}). Second, a finite \(\mathbf F_\ell\)-local system becomes a successive extension of constant sheaves on a finite étale cover of degree \(d\) prime to \(\ell\). This cover is given by an \(\ell\)-Sylow subgroup of the monodromy. Dividing the trace of the cover by \(d\) shows that the local system is a direct summand of the pushforward of its pullback (Lemma \hyperref[bp215]{D.12}). We follow the proof of [Stacks, \href{https://stacks.math.columbia.edu/tag/03SC}{03SC}].

Throughout this appendix, cohomology means cohomology on \(C_{\mathrm{pro}\acute{e}t}\), and the sheaves are classical in the sense of Definition \bpout{2.30}: by Theorem \bpout{2.32} such a sheaf is the inverse image \(\nu_C^*G\) of a sheaf \(G\) on \(C_{\acute{e}t}\) along the comparison morphism. By Theorem \bpout{2.35}, its cohomology, together with the restriction and pullback maps, agrees with étale cohomology, so we carry out all calculations on \(C_{\acute{e}t}\) and transport them back by \(\nu_C^*\).

Let \(p\ge0\) be the characteristic of \(k\). In this appendix \(\ell\) is an arbitrary prime, possibly equal to \(p\). In characteristic zero the hypotheses in \((C_4)\) and \((C_6)\) of Theorem \bpout{4.20} that \(F\) be prime to the characteristic hold automatically.

% object: sec-Kummer-theory (blueprint/chap-Cohomology-of-torsion-sheaves-on-curves/sec-Kummer-theory/_section.md), content hidden
\subsection*{\texorpdfstring{D.1\quad Kummer theory}{D.1 Kummer theory}}
\phantomsection\label{bp174}%
\addcontentsline{toc}{subsection}{\texorpdfstring{D.1\quad Kummer theory}{D.1 Kummer theory}}
% text: prose-D-02
\phantomsection\label{bp175}%
Let \(C\) be a smooth curve over an algebraically closed field \(k\), and let \(\nu_C^*\mathbf G_m\) be the pullback of the multiplicative group \(\mathbf G_m\) of \(C_{\acute{e}t}\) to \(C_{\mathrm{pro}\acute{e}t}\). For \(n\) invertible in \(k\), the Kummer sequence \(0\to\mu_n\to\mathbf G_m\xrightarrow{n}\mathbf G_m\to0\) reduces the cohomology of \(\mu_n\) to that of \(\nu_C^*\mathbf G_m\), so we compute the latter first. In degree one, \(H^1(C,\mathbf G_m)\) is the Picard group, as in the Zariski topology, while the groups \(H^q(C,\mathbf G_m)\) vanish for \(q\ge2\) (Lemma \hyperref[bp176]{D.1}). In particular the Brauer group of \(C\) is zero. This matches, over \(\mathbf C\), the topology of a Riemann surface, whose integral cohomology vanishes above degree two, and whose \(H^2\) is accounted for by the degree of line bundles. In algebraic geometry the vanishing comes from Tsen's theorem, which says that forms of low degree over the function field \(k(C)\) have nontrivial zeros. Over a finite field Tsen's theorem fails: by class field theory the function field of a curve has nonzero Brauer group. This is one reason why we work over an algebraically closed field. With the lemma and the Jacobian, the Kummer sequence then computes the cohomology of \(\mu_n\) on a proper curve (Proposition \hyperref[bp179]{D.2}).

% object: lemma-cohomology-multiplicative-group (blueprint/chap-Cohomology-of-torsion-sheaves-on-curves/sec-Kummer-theory/lemma-cohomology-multiplicative-group/_section.md)
\phantomsection\label{bp176}%
\begin{bpstatementit}{Lemma~D.1}{ (Cohomology of the multiplicative group)}
For a smooth separated curve \(C\) of finite type over an algebraically closed field, one has

\[
H^1(C,\nu_C^*\mathbf G_m)=\operatorname{Pic}(C),
\qquad
H^q(C,\nu_C^*\mathbf G_m)=0\quad(q\geq2).
\]

The identifications are natural under restriction and extension of the algebraically closed ground field. In particular \(\operatorname{Br}(C)=H^2(C_{\acute{e}t},\mathbf G_m)=0\).
\end{bpstatementit}

% proof: proof/lemma-cohomology-multiplicative-group
\begin{proof}[{Proof}]\phantomsection\label{bp177}%
We compute on \(C_{\acute{e}t}\). In degree one, descent for invertible modules identifies \(\mathbf G_m\)-torsors with line bundles. In degrees at least two, by contrast, the vanishing comes from Tsen's theorem, which we prove first (see also [Stacks, \href{https://stacks.math.columbia.edu/tag/03RM}{03RM}]). There are three steps: Tsen's theorem, proved by counting dimensions with Riemann–Roch; the vanishing of the Brauer group and of the Galois cohomology of \(k(C)\); and the passage from the generic point to \(C\) by the divisor sequence.

Tsen's theorem says that \(k(C)\) is \(C_1\), that is, a homogeneous form of degree \(d\) in \(m>d\) variables with coefficients in \(k(C)\) has a nonzero zero.  To see this, let \(\overline C\) be the smooth proper model of \(k(C)\), of genus \(g\), choose an effective divisor \(D\) on \(\overline C\) bounding the poles of all coefficients, and let the unknowns range over \(H^0(\overline C,\mathcal O(NH))\) for an effective divisor \(H\) of positive degree.  The value of the form then lies in \(H^0(\overline C,\mathcal O(dNH+D))\).  Once \(N\deg H>2g-2\), we know both dimensions by Theorem \hyperref[bp130]{C.3}.  In terms of bases, the \(m\) unknowns amount to \(m(N\deg H+1-g)\) unknowns in \(k\). Each coordinate of the value, in a basis of the target space, is a homogeneous polynomial of degree \(d\) in these unknowns, so the vanishing of the value amounts to \(dN\deg H+\deg D+1-g\) homogeneous polynomial equations in them.  The number of unknowns minus the number of equations is \((m-d)N\deg H+(m-1)(1-g)-\deg D\), which is positive for large \(N\) because \(m>d\).  Since \(k\) is algebraically closed, homogeneous equations in more unknowns than equations have a nonzero solution [Stacks, \href{https://stacks.math.columbia.edu/tag/03RD}{03RD}].

The reduced norm form of a division algebra of degree \(d>1\) over \(k(C)\) has degree \(d\) in \(d^2>d\) variables and does not vanish on a nonzero element.  Therefore Tsen's theorem implies that \(\operatorname{Br}(k(C))=0\) [Stacks, \href{https://stacks.math.columbia.edu/tag/03RF}{03RF}].  The same holds for every finite extension of \(k(C)\), since it is again the function field of a curve over \(k\). This is enough for Galois cohomology, because a field all of whose finite extensions have trivial Brauer group has \(H^q(K,(K^{\mathrm{sep}})^\times)=0\) for all \(q\geq1\) [Stacks, \href{https://stacks.math.columbia.edu/tag/03R8}{03R8}]; hence this holds for \(K=k(C)\). To pass from \(k(C)\) to \(C\) we use the divisor sequence, which relates \(\mathbf G_m\) on \(C_{\acute{e}t}\) to the pushforward of the multiplicative group at the generic point and to skyscraper sheaves \(\mathbf Z\) at the closed points. Of these, the former has no higher cohomology by this Galois vanishing, and the latter have none because the residue fields are algebraically closed. Therefore, from the long exact cohomology sequence, we conclude that \(H^q(C_{\acute{e}t},\mathbf G_m)=0\) for \(q\geq2\).

By Theorem \bpout{2.35} these identifications carry over to \(C_{\mathrm{pro}\acute{e}t}\), and they are natural because the comparison maps, the identification of torsors and the restriction maps are functorial.
\end{proof}

% text: prose-D-03
\phantomsection\label{bp178}%
On a proper curve, the lemma and the Kummer sequence reduce the cohomology of \(\mu_n\) to the groups \(\operatorname{Pic}(C)[n]\) and \(\operatorname{Pic}(C)/n\), both computed in Appendix \hyperref[bp119]{C}.

% object: proposition-smooth-proper-curves-torsion-prime (blueprint/chap-Cohomology-of-torsion-sheaves-on-curves/sec-Kummer-theory/proposition-smooth-proper-curves-torsion-prime/_section.md)
\phantomsection\label{bp179}%
\begin{bpstatementit}{Proposition~D.2}{ (Smooth proper curves, torsion prime to the characteristic)}
Let \(C\) be smooth, proper and connected over an algebraically closed field \(k\), with genus \(g\), and let \(n\geq1\) be invertible in \(k\). There are canonical identifications

\[
\begin{aligned}
H^0(C,\mu_n)&=\mu_n(k),&
H^1(C,\mu_n)&=J(C)[n](k),\\
H^2(C,\mu_n)&=\mathbf Z/n,&
H^q(C,\mu_n)&=0\quad(q>2).
\end{aligned}
\]

The isomorphism in degree two sends the Kummer boundary of a line bundle \(L\) to \(\deg L\bmod n\).  These identifications commute with extension of algebraically closed fields.
\end{bpstatementit}

% proof: proof/proposition-smooth-proper-curves-torsion-prime
\begin{proof}[{Proof}]\phantomsection\label{bp180}%
Pulling back the Kummer sequence of \(C_{\acute{e}t}\) by \(\nu_C^*\), we get the exact sequence of sheaves on \(C_{\mathrm{pro}\acute{e}t}\)

\[
0\longrightarrow\mu_n\longrightarrow\nu_C^*\mathbf G_m
\xrightarrow{(\cdot)^n}\nu_C^*\mathbf G_m
\longrightarrow0.
\]

Here surjectivity of the \(n\)th power map is a local question. For a unit \(a\) the equation \(T^n=a\) defines an étale cover, since its derivative \(nT^{n-1}\) is a unit. Moreover, the global units on the proper connected curve form the group \(k^\times\), which is \(n\)-divisible because \(k\) is algebraically closed. Hence the long exact sequence starts with \(H^0(C,\mu_n)=\mu_n(k)\). Since the map \(H^0(C,\mathbf G_m)\to H^0(C,\mathbf G_m)\) is onto, \(H^1(C,\mu_n)\) injects into \(H^1(C,\mathbf G_m)\). Using Lemma \hyperref[bp176]{D.1}, we read off from the long exact sequence that

\[
H^1(C,\mu_n)=\operatorname{Pic}(C)[n],\qquad
H^2(C,\mu_n)=\operatorname{Pic}(C)/n,
\]

with all higher groups zero.

Corollary \hyperref[bp157]{C.12} gives the exact sequence

\[
0\longrightarrow J(C)(k)\longrightarrow\operatorname{Pic}(C)
\xrightarrow{\deg}\mathbf Z\longrightarrow0.
\]

Multiplication by \(n\) on \(J(C)(k)\) is surjective, with kernel \(J(C)[n](k)\), and \(J(C)[n]\) is finite étale of rank \(n^{2g}\) (Corollary \hyperref[bp163]{C.15}).  Applying the snake lemma to multiplication by \(n\) on this sequence, and using that \(\mathbf Z\) is torsion free, we obtain \(\operatorname{Pic}(C)[n]=J(C)[n](k)\) and \(\operatorname{Pic}(C)/n=\mathbf Z/n\) via the degree. In particular, \(H^1(C,\mu_n)\) is noncanonically isomorphic to \((\mathbf Z/n)^{2g}\) [Stacks, \href{https://stacks.math.columbia.edu/tag/03RQ}{03RQ}].

Under an algebraically closed extension \(k'/k\), the schemes \(J(C)\) and \(J(C)[n]\) base change to the corresponding objects for \(C_{k'}\). The map \(J(C)[n](k)\to J(C)[n](k')\) is bijective by Corollary \hyperref[bp163]{C.15}. The same argument applies to the finite étale scheme \(\mu_n\). Finally, pullback of line bundles preserves the degree. Hence the pullback maps on cohomology are isomorphisms.
\end{proof}

% object: remark-twists-projective-line (blueprint/chap-Cohomology-of-torsion-sheaves-on-curves/sec-Kummer-theory/remark-twists-projective-line/_section.md)
\phantomsection\label{bp181}%
\begin{bpstatement}{Remark~D.3}{ (Twists and the projective line)}
The degree identifies \(H^2(C,\mu_n)\) with \(\mathbf Z/n\) without any choice. Since \(\mu_n=(\mathbf Z/n)(1)\), twisting by \(-1\) gives

\[
H^2(C,\mathbf Z/n)=(\mathbf Z/n)(-1),
\]

and identifying this group with \(\mathbf Z/n\) requires the choice of a primitive \(n\)th root of unity.  For \(C=\mathbf P^1_k\) the Jacobian is zero, so the cohomology with constant coefficients is concentrated in degrees zero and two, and the Kummer boundary of \(\mathcal O(1)\) corresponds to \(1\in\mathbf Z/n\). Section \bpout{4.4} uses this normalization on the geometric fibre.
\end{bpstatement}


% object: sec-Affine-curves (blueprint/chap-Cohomology-of-torsion-sheaves-on-curves/sec-Affine-curves/_section.md), content hidden
\subsection*{\texorpdfstring{D.2\quad Affine curves}{D.2 Affine curves}}
\phantomsection\label{bp182}%
\addcontentsline{toc}{subsection}{\texorpdfstring{D.2\quad Affine curves}{D.2 Affine curves}}
% text: prose-D-04
\phantomsection\label{bp183}%
Every smooth connected curve over an algebraically closed field that is not proper is affine and is the complement of a finite nonempty set of closed points in a smooth proper connected curve.  To see this, take a projective model of its function field and normalize it, and recall that regular curves over a perfect field are smooth.

On the affine curve \(U=\overline C\setminus D\) the Picard group is a quotient of \(\operatorname{Pic}^0(\overline C)\). It is therefore divisible, and \(H^2(U,\mu_n)=0\); this matches the fact that, topologically, a punctured Riemann surface is homotopy equivalent to a graph. In degree one, on the other hand, new torsors appear, since a Kummer torsor on \(U\) need not extend over \(\overline C\). The obstruction to extending it at each puncture is the residue of the torsor, that is, the order modulo \(n\) of the corresponding rational section. For example, on \(\mathbf G_m=\mathbf P^1\setminus\{0,\infty\}\) the torsor \(y^n=x\) has residues \(1\) and \(-1\) at the two punctures. It generates \(H^1(\mathbf G_m,\mu_n)\simeq\mathbf Z/n\). In general the residues sum to zero. The residue sequence of Proposition \hyperref[bp184]{D.4} shows that \(H^1(U,\mu_n)\) is free of rank \(2g+\#D-1\) over \(\mathbf Z/n\). This is the first Betti number of a surface of genus \(g\) with \(\#D\) punctures.

% object: proposition-cohomology-smooth-affine-curve (blueprint/chap-Cohomology-of-torsion-sheaves-on-curves/sec-Affine-curves/proposition-cohomology-smooth-affine-curve/_section.md)
\phantomsection\label{bp184}%
\begin{bpstatementit}{Proposition~D.4}{ (Cohomology of a smooth affine curve)}
Let \(\overline C\) be a smooth proper connected curve of genus \(g\) over an algebraically closed field \(k\).  Let \(U=\overline C\setminus D\), where \(D\) is a finite nonempty set of closed points, and let \(n\) be invertible in \(k\).  There is a natural exact sequence

\[
0\longrightarrow H^1(\overline C,\mu_n)
\longrightarrow H^1(U,\mu_n)
\xrightarrow{\partial}(\mathbf Z/n)^D
\xrightarrow{\sum}\mathbf Z/n\longrightarrow0.
\tag{D.1}
\]

Moreover, \(H^0(U,\mu_n)=\mu_n(k)\) and \(H^q(U,\mu_n)=0\) for \(q>1\). The group \(H^1(U,\mu_n)\) is noncanonically \((\mathbf Z/n)^{2g+\#D-1}\), and the pullback maps on all these groups are isomorphisms under algebraically closed extension.
\end{bpstatementit}

% proof: proof/proposition-cohomology-smooth-affine-curve
\begin{proof}[{Proof}]\phantomsection\label{bp185}%
By the Kummer sequence, the vanishing above degree one follows from the divisibility of \(\operatorname{Pic}(U)\) and the vanishing of \(H^q(U,\mathbf G_m)\) for \(q\geq2\) (Lemma \hyperref[bp176]{D.1}). For \(H^1\) we describe \(\mu_n\)-torsors as line bundles \(L\) with a trivialization of \(L^{\otimes n}\). After extending \(L\) to \(\overline C\), the orders at the punctures measure the failure of the torsor to extend. These orders give the residue map \(\partial\). We prove exactness at each place, and then deduce the rank and the invariance under field extension.

Since \(\overline C\) is regular, we may use divisors to see that

\[
\operatorname{Pic}(U)=\operatorname{Pic}(\overline C)/
\langle\mathcal O_{\overline C}(x):x\in D\rangle.
\]

One of these line bundles has degree one. Every class in \(\operatorname{Pic}(\overline C)\) is therefore congruent, modulo the subgroup generated by it, to a class of degree zero. Hence \(J(\overline C)(k)\) surjects onto \(\operatorname{Pic}(U)\). As \(J(\overline C)(k)\) is \(n\)-divisible (Corollary \hyperref[bp163]{C.15}), so is \(\operatorname{Pic}(U)\). By the Kummer sequence and Lemma \hyperref[bp176]{D.1}, we have \(H^2(U,\mu_n)=\operatorname{Pic}(U)/n=0\) and \(H^q(U,\mu_n)=0\) for \(q\geq3\).

We construct the sequence (D.1) using line bundles.  Recall that a \(\mu_n\)-torsor on \(U\) is given by a pair \((L,\alpha:L^{\otimes n}\simeq\mathcal O_U)\) of a line bundle and a trivialization of its \(n\)th power.  To define its residues, we extend \(L\) to a line bundle \(\overline L\) on \(\overline C\) and let \(s=\alpha^{-1}(1)\) be the corresponding rational section of \(\overline L^{\otimes n}\).  The divisor of \(s\) is supported on \(D\), since \(s\) generates on \(U\). Its orders modulo \(n\) do not depend on the extension. Indeed, two extensions differ by \(\mathcal O(E')\) with \(E'\) supported on \(D\), and this changes the divisor of \(s\) by \(nE'\). These orders define \(\partial(L,\alpha)\). Their sum is zero modulo \(n\), since the divisor of \(s\) has degree \(\deg\overline L^{\otimes n}=n\deg\overline L\).

Suppose all orders are divisible by \(n\), and write \(\operatorname{div}(s)=nE\) with \(E\) supported on \(D\). Then \(s\) is a nowhere vanishing section of \((\overline L(-E))^{\otimes n}=\overline L^{\otimes n}(-nE)\). So the trivialization extends to a torsor with line bundle \(\overline L(-E)\) on \(\overline C\), and this torsor restricts to \((L,\alpha)\) on \(U\). This proves exactness at \(H^1(U,\mu_n)\).

To prove injectivity, note that a trivialization on \(U\) of a torsor on \(\overline C\) gives a rational section of its line bundle whose \(n\)th power has no zero or pole on \(\overline C\).  Therefore the divisor of the section is zero, and the section trivializes the torsor on \(\overline C\).

To see that \(\partial\) maps onto the kernel of \(\sum\), we lift a residue vector with sum zero modulo \(n\) to an integral divisor \(B\) supported on \(D\). Its degree is divisible by \(n\) because the residues sum to zero. We claim that, as a consequence, there is a line bundle \(\overline L\) with \(\overline L^{\otimes n}\simeq\mathcal O_{\overline C}(B)\). To find it, choose \(P\in\overline C(k)\) and put \(L_1=\mathcal O_{\overline C}((\deg B/n)P)\), so that the line bundle \(\mathcal O_{\overline C}(B)\otimes L_1^{\otimes-n}\) has degree zero. By divisibility of \(J(\overline C)(k)\) it is \(N^{\otimes n}\) for some \(N\) of degree zero, and \(\overline L=L_1\otimes N\) works. It remains to produce the torsor. The canonical section of \(\mathcal O_{\overline C}(B)\) has divisor \(B\) and is nowhere zero on \(U\). It therefore defines a trivialization of \(\overline L^{\otimes n}|_U\), that is, a Kummer torsor on \(U\) whose residues are the orders of \(B\). Finally, the summation map is surjective since \(D\) is nonempty. This proves that the sequence is exact (see also [Stacks, \href{https://stacks.math.columbia.edu/tag/03RR}{03RR}]).

The left term and the kernel of summation are free \(\mathbf Z/n\)-modules of ranks \(2g\) and \(\#D-1\). Hence the resulting short exact sequence splits noncanonically. For naturality, note that the formation of the compactification, the boundary points, the line bundles and the orders of vanishing commutes with base change. The leftmost pullback map is an isomorphism by Proposition \hyperref[bp179]{D.2}, and the residue groups and their summation maps are unchanged. By the five lemma applied to the two exact sequences, the middle pullback map is an isomorphism.
\end{proof}

% text: prose-D-05
\phantomsection\label{bp186}%
For the constant sheaf \(\mathbf F_\ell\), with \(\ell\) a prime different from \(\operatorname{char}k\), the two preceding propositions prove \((C_1)\)–\((C_8)\) on smooth curves. Indeed, a primitive \(\ell\)th root of unity in \(k\) identifies \(\mathbf F_\ell\) with \(\mu_\ell\), compatibly with the pullback maps to \(C_{k'}\), since the root of unity stays primitive in \(k'\). For \((C_8)\), note that an open in a smooth connected proper curve is either the whole curve or an affine curve, which has \(H^2=0\), and that an open in a smooth affine curve is affine.


% object: sec-The-Artin-Schreier-sequence (blueprint/chap-Cohomology-of-torsion-sheaves-on-curves/sec-The-Artin-Schreier-sequence/_section.md), content hidden
\subsection*{\texorpdfstring{D.3\quad The Artin–Schreier sequence}{D.3 The Artin–Schreier sequence}}
\phantomsection\label{bp187}%
\addcontentsline{toc}{subsection}{\texorpdfstring{D.3\quad The Artin–Schreier sequence}{D.3 The Artin–Schreier sequence}}
% text: prose-D-06
\phantomsection\label{bp188}%
Since the Kummer cover \(T^n=a\) is étale only when \(n\) is invertible, \(p\)-torsion in characteristic \(p\) needs a different sequence. In the Artin–Schreier sequence the role of \(\mathbf G_m\) is played by the additive group, whose étale cohomology is the coherent cohomology of the structure sheaf. Since coherent cohomology of a curve vanishes above degree one, this gives \(H^q(C,\mathbf F_p)=0\) for \(q\geq3\). For \(q=2\) we also need that \(\varphi-1\) is surjective on \(H^1(C,\mathcal O_C)\), which holds over an algebraically closed field. So \(H^2(C,\mathbf F_p)=0\) even on a proper curve, and there is no analogue of the class of a point in \(H^2\).

On a proper curve, \(H^1(C,\mathbf F_p)\) is the fixed space of the Frobenius-linear map \(\varphi\) on the \(g\)-dimensional space \(H^1(C,\mathcal O_C)\). Hence it has dimension at most \(g\) (Proposition \hyperref[bp189]{D.5}). The bound can be strict. For example, for an elliptic curve, \(\varphi\) on the line \(H^1(E,\mathcal O_E)\) is bijective when \(E\) is ordinary, and zero when \(E\) is supersingular, as is \(y^2=x^3-x\) over \(\overline{\mathbf F}_3\). Accordingly \(H^1(E,\mathbf F_p)\) has dimension one or zero. On an affine curve, by contrast, the cokernel of \(\varphi-1\) on the ring of functions contributes to \(H^1\), and this group is infinite (Example \hyperref[bp195]{D.7}).

% object: proposition-p-torsion-smooth-proper-curve (blueprint/chap-Cohomology-of-torsion-sheaves-on-curves/sec-The-Artin-Schreier-sequence/proposition-p-torsion-smooth-proper-curve/_section.md)
\phantomsection\label{bp189}%
\begin{bpstatementit}{Proposition~D.5}{ ($p$-torsion on a smooth proper curve in characteristic $p$)}
Let \(C\) be a smooth proper connected curve of genus \(g\) over an algebraically closed field of characteristic \(p>0\), and let \(\varphi\) be the Frobenius-linear endomorphism of \(H^1(C,\mathcal O_C)\) induced by \(a\mapsto a^p\) on \(\mathcal O_C\).  Then

\[
H^0(C,\mathbf F_p)=\mathbf F_p,\qquad
H^1(C,\mathbf F_p)=
\ker(\varphi-1:H^1(C,\mathcal O_C)\to H^1(C,\mathcal O_C)),
\]

and \(H^q(C,\mathbf F_p)=0\) for \(q>1\).  The group \(H^1(C,\mathbf F_p)\) is finite, of \(\mathbf F_p\)-dimension at most \(g\), and its pullback map is an isomorphism under algebraically closed extension.
\end{bpstatementit}

% proof: proof/proposition-p-torsion-smooth-proper-curve
\begin{proof}[{Proof}]\phantomsection\label{bp190}%
The Artin–Schreier sequence turns the question into one about \(\varphi-1\) on \(H^0(C,\mathcal O_C)=k\) and on \(H^1(C,\mathcal O_C)\). The key fact is that on a finite-dimensional space over an algebraically closed field, a Frobenius-linear map minus the identity is surjective with finite kernel.

\emph{Step 1: the Artin–Schreier sequence.} Write also \(\varphi\) for the endomorphism \(a\mapsto a^p\) of the sheaf \(\mathbf G_a\) on \(C_{\acute{e}t}\).  The Artin–Schreier sequence

\[
0\longrightarrow\mathbf F_p\longrightarrow\mathbf G_a
\xrightarrow{\varphi-1}\mathbf G_a\longrightarrow0
\]

is exact, because surjectivity is a local question and \(T^p-T-a\) has derivative \(-1\). The étale cohomology of \(\mathbf G_a\) is the coherent cohomology of \(\mathcal O_C\). Indeed, quasi-coherent sheaves have the same cohomology on the étale and on the Zariski site, by faithfully flat descent on affine opens and the Čech resolution. This coherent cohomology vanishes in degrees above one, since \(C\) is a curve. Since \(C\) is proper, \(H^1(C,\mathcal O_C)\) is a finite-dimensional \(k\)-vector space.

\emph{Step 2: linear algebra.} Let \(\varphi\) be a Frobenius-linear endomorphism of a finite-dimensional \(k\)-vector space \(V\); we show that \(\varphi-1\) is surjective and that its kernel is the set of \(k\)-points of a finite étale scheme of order \(p^f\) with \(f\leq\dim_k V\).  To prove this, we reduce along a finite filtration to the part where \(\varphi\) is nilpotent and the part where \(\varphi\) is bijective.  On the nilpotent part, \(\varphi-1\) is an automorphism given by an additive polynomial, with inverse \(-(1+\varphi+\varphi^2+\cdots)\), a finite sum; in particular its kernel there is zero.  On a bijective part of dimension \(f\), by contrast, we choose a basis, invert the matrix of \(\varphi\), and write the equations of the fibres as

\[
x_i^p-\sum_j b_{ij}x_j=c_i\qquad(1\leq i\leq f),
\]

where \((b_{ij})\) is invertible.  The coordinate algebra is free of dimension \(p^f\), with basis the monomials having each exponent less than \(p\).  Moreover, the Jacobian matrix of these equations is \(-(b_{ij})\), because \(p=0\) in \(k\), and it is invertible; so each fibre is finite étale of degree \(p^f\) and has \(p^f\) points over the algebraically closed field \(k\) [Stacks, \href{https://stacks.math.columbia.edu/tag/0A3L}{0A3L}]. In particular every fibre is nonempty, so \(\varphi-1\) is surjective, and the kernel has \(p^f\) elements.

\emph{Step 3: conclusion.} We apply Step 2 to \(H^0(C,\mathcal O_C)=k\) and to \(H^1(C,\mathcal O_C)\) in the long exact sequence associated to the Artin–Schreier sequence. Surjectivity of \(\varphi-1\) on \(H^0\) gives \(H^1(C,\mathbf F_p)=\ker(\varphi-1)\) on \(H^1(C,\mathcal O_C)\), with the bound on its dimension, and surjectivity on \(H^1\) together with \(H^2(C,\mathcal O_C)=0\) gives \(H^q(C,\mathbf F_p)=0\) for \(q>1\). Finally, the pullback map under an extension \(k'/k\) is an isomorphism [Stacks, \href{https://stacks.math.columbia.edu/tag/0HAN}{0HAN}, \href{https://stacks.math.columbia.edu/tag/0A3P}{0A3P}]: flat base change for coherent cohomology identifies the vector space and its Frobenius-linear map with their scalar extensions, and the kernel is a finite étale scheme, whose \(k\)-points and \(k'\)-points agree. As always, Theorem \bpout{2.35} carries the statement over to \(C_{\mathrm{pro}\acute{e}t}\).
\end{proof}

% text: prose-D-07
\phantomsection\label{bp191}%
On an affine curve the same sequence is even simpler to use, because coherent cohomology vanishes in positive degrees. On the other hand, \(H^1\) is now the cokernel of \(\varphi-1\) on an infinite-dimensional space.

% object: lemma-p-torsion-affine-curve-characteristic (blueprint/chap-Cohomology-of-torsion-sheaves-on-curves/sec-The-Artin-Schreier-sequence/lemma-p-torsion-affine-curve-characteristic/_section.md)
\phantomsection\label{bp192}%
\begin{bpstatementit}{Lemma~D.6}{ ($p$-torsion on an affine curve in characteristic $p$)}
If \(U\) is affine over an algebraically closed field \(k\) of characteristic \(p\), then, for the classical constant sheaf \(\mathbf F_p\) and \(\varphi(a)=a^p\) on \(\Gamma(U,\mathcal O_U)\), we have

\[
H^1(U,\mathbf F_p)=
\Gamma(U,\mathcal O_U)/(\varphi-1)\Gamma(U,\mathcal O_U),\qquad
H^q(U,\mathbf F_p)=0\quad(q>1).
\]
\end{bpstatementit}

% proof: proof/lemma-p-torsion-affine-curve-characteristic
\begin{proof}[{Proof}]\phantomsection\label{bp193}%
Quasi-coherent cohomology on an affine scheme vanishes in positive degrees, and, as in the proof of the proposition, étale and Zariski cohomology of \(\mathbf G_a\) agree. In the long exact sequence associated to the Artin–Schreier sequence on \(U_{\acute{e}t}\), the terms \(H^q(U,\mathbf G_a)\) for \(q\geq1\) are therefore zero. What remains is the exact sequence \(0\to H^0(U,\mathbf F_p)\to\Gamma(U,\mathcal O_U)\xrightarrow{\varphi-1}\Gamma(U,\mathcal O_U)\to H^1(U,\mathbf F_p)\to0\), together with \(H^q(U,\mathbf F_p)=0\) for \(q>1\).  By Theorem \bpout{2.35} these statements carry over to \(U_{\mathrm{pro}\acute{e}t}\).
\end{proof}

% text: prose-D-08
\phantomsection\label{bp194}%
For the affine line we can compute the cokernel explicitly, and we see that it is infinite and grows with the field, unlike in the proper case.

% object: example-affine-line-characteristic-p (blueprint/chap-Cohomology-of-torsion-sheaves-on-curves/sec-The-Artin-Schreier-sequence/example-affine-line-characteristic-p/_section.md)
\phantomsection\label{bp195}%
\begin{bpstatement}{Example~D.7}{ (The affine line in characteristic $p$)}
For \(U=\mathbf A^1_k\) in characteristic \(p\) we get

\[
H^1(U,\mathbf F_p)=k[t]/\{h^p-h:h\in k[t]\}.
\]

Every class has a unique representative with zero constant term and all positive exponents prime to \(p\).  Indeed, we can remove exponents divisible by \(p\) using \(p\)th roots in the perfect field, and we can remove the constant term since \(a\mapsto a^p-a\) is surjective on \(k\). For uniqueness, note that a nonconstant polynomial \(h^p-h\) has degree divisible by \(p\).  In particular the quotient is infinite, because the classes of \(t^m\) for positive \(m\) prime to \(p\) are \(\mathbf F_p\)-linearly independent.  For an algebraically closed extension \(k\subsetneq k'\) and \(a\in k'\setminus k\), the class of \(at\) is not in the image of the quotient over \(k\). Indeed, the reduction of a polynomial in \(k[t]\) uses only \(p\)th roots and roots of \(x^p-x=c\) in \(k\), so its reduced representative lies in \(k[t]\), whereas \(at\) is already reduced. Thus finiteness and invariance under field extension in degree one fail for \(p\)-torsion coefficients on an affine curve.
\end{bpstatement}

% text: prose-D-09
\phantomsection\label{bp196}%
For the constant sheaf \(\mathbf F_p\) on smooth curves, we obtain \((C_1)\), \((C_3)\), \((C_5)\) and \((C_7)\) from the proposition on proper curves, and \((C_1)\), \((C_2)\) and \((C_3)\) from the lemma on affine curves, while \((C_8)\) holds because \(H^2\) vanishes in both cases.


% object: sec-Normalization (blueprint/chap-Cohomology-of-torsion-sheaves-on-curves/sec-Normalization/_section.md), content hidden
\subsection*{\texorpdfstring{D.4\quad Normalization}{D.4 Normalization}}
\phantomsection\label{bp197}%
\addcontentsline{toc}{subsection}{\texorpdfstring{D.4\quad Normalization}{D.4 Normalization}}
% text: prose-D-10
\phantomsection\label{bp198}%
So far we have computed with constant coefficients on smooth curves, whereas Theorem \bpout{4.20} concerns arbitrary curves of dimension at most one and arbitrary torsion sheaves. Two reductions close the gap: normalization and removing finitely many points. Both rest on Lemma \hyperref[bp202]{D.9}, which says that a sheaf supported on finitely many points has no higher cohomology. Hence passing between a curve, its normalization and a dense open subset changes cohomology only in degrees zero and one, and only by finite groups that do not change under field extension.

For example, let \(C:y^2=x^3+x^2\) be the nodal cubic. Its normalization \(\mathbf A^1\to C\) has two points over the node. For \(M=\mathbf Z/n\) the normalization sequence of Lemma \hyperref[bp205]{D.10} has cokernel \(M^2/\Delta M\simeq M\) at the node. If \(n\) is invertible in \(k\), then \(H^1(\mathbf A^1,\mathbf Z/n)=0\), and the long exact sequence gives \(H^1(C,\mathbf Z/n)\simeq\mathbf Z/n\). This is the analogue of the loop created by the node.

% object: lemma-finite-morphisms-cohomology (blueprint/chap-Cohomology-of-torsion-sheaves-on-curves/sec-Normalization/lemma-finite-morphisms-cohomology/_section.md)
\phantomsection\label{bp199}%
\begin{bpstatementit}{Lemma~D.8}{ (Finite morphisms and cohomology)}
Let \(h:Y\to C\) be a finite morphism of schemes. The direct image \(h_*\) of classical sheaves is exact, its formation commutes with arbitrary base change, and \(H^q(C,h_*G)=H^q(Y,G)\) for every classical abelian sheaf \(G\) on \(Y_{\mathrm{pro}\acute{e}t}\) and every \(q\).
\end{bpstatementit}

% proof: proof/lemma-finite-morphisms-cohomology
\begin{proof}[{Proof}]\phantomsection\label{bp200}%
On the étale sites, Lemma \bpout{4.27} shows that \(h_*\) is exact, that its stalks are finite products of stalks of \(G\), and that its formation commutes with arbitrary base change.  Since \(h_*\) is exact, the Leray spectral sequence gives \(H^q(C,h_*G)=H^q(Y,G)\).  The comparison of classical direct images and Theorem \bpout{2.35} transport these statements to the pro-étale site; see [Stacks, \href{https://stacks.math.columbia.edu/tag/03QP}{03QP}, \href{https://stacks.math.columbia.edu/tag/0959}{0959}].
\end{proof}

% text: prose-D-11
\phantomsection\label{bp201}%
Since a sheaf supported on finitely many points is a direct image from those points, it has no higher cohomology. All the reductions below use this vanishing.

% object: lemma-sheaves-finite-support-finite-morphisms (blueprint/chap-Cohomology-of-torsion-sheaves-on-curves/sec-Normalization/lemma-sheaves-finite-support-finite-morphisms/_section.md)
\phantomsection\label{bp202}%
\begin{bpstatementit}{Lemma~D.9}{ (Sheaves with finite support)}
Let \(F\) be a classical abelian sheaf supported on finitely many closed points of a scheme \(C\) of finite type over an algebraically closed field \(k\). Then \(H^q(C,F)=0\) for \(q>0\), and \(H^0(C,F)\) is the direct sum of the stalks of \(F\) at these points; the pullback map on \(H^0\) is an isomorphism under algebraically closed extension of \(k\).
\end{bpstatementit}

% proof: proof/lemma-sheaves-finite-support-finite-morphisms
\begin{proof}[{Proof}]\phantomsection\label{bp203}%
The sheaf \(F\) is \(h_*\) of a sheaf on a finite disjoint union of copies of \(\operatorname{Spec}k\), where \(h\) is the inclusion of the points. On these points taking sections is exact because \(k\) is algebraically closed, and the sections are the product, here a finite direct sum, of the stalks, which do not change under extension of \(k\). Finally, Lemma \hyperref[bp199]{D.8} carries these assertions over to \(C\).
\end{proof}

% text: prose-D-12
\phantomsection\label{bp204}%
The normalization of a curve is finite and is an isomorphism away from finitely many points, so the constant sheaf on \(C\) differs from the direct image of the constant sheaf on \(C^{\mathrm{nor}}\) only by a sheaf with finite support.

% object: lemma-normalization-constant-sheaf (blueprint/chap-Cohomology-of-torsion-sheaves-on-curves/sec-Normalization/lemma-normalization-constant-sheaf/_section.md)
\phantomsection\label{bp205}%
\begin{bpstatementit}{Lemma~D.10}{ (Normalization of a constant sheaf)}
Let \(C\) be a reduced curve of finite type over an algebraically closed field, and let \(\pi:C^{\mathrm{nor}}\to C\) be its normalization.  For a finite abelian group \(M\), the natural map

\[
\underline M_C\longrightarrow \pi_*\underline M_{C^{\mathrm{nor}}}
\]

is injective with cokernel supported on finitely many points.  At a point with \(b\) geometric branches, the cokernel is \(M^b/\Delta M\), where \(\Delta\) is the diagonal. These cokernels and the resulting exact sequences commute with algebraically closed field extension.
\end{bpstatementit}

% proof: proof/lemma-normalization-constant-sheaf
\begin{proof}[{Proof}]\phantomsection\label{bp206}%
The map is computed on stalks by the formula for stalks of finite direct images, and the cokernel has finite support because \(\pi\) is an isomorphism away from finitely many points. Compatibility with field extension holds because normalization commutes with algebraically closed extension.

Normalization is finite for reduced schemes of finite type over a field.  On each one-dimensional component, its normal model is regular, hence smooth over the perfect field, while the isolated components of dimension zero are reduced points, that is, copies of the spectrum of the ground field.  Since the normalization is an isomorphism over a dense open of each component, the locus where it is not an isomorphism is finite.

By the calculation of stalks in Lemma \bpout{4.27}, the displayed map at a geometric point is the diagonal \(M\to M^b\), with \(b\geq1\), so its kernel is zero and its cokernel has the asserted form.  Where \(\pi\) is an isomorphism we have \(b=1\) and the cokernel vanishes; hence the cokernel is supported on the finite locus found above.

After algebraically closed extension, the normalization remains smooth and is a finite birational normal model of the reduced base change, so it is the normalization of the base change.  Base change for finite morphisms identifies the fibres and their sets of branches, and therefore the diagonal maps and their cokernels are obtained by extension of scalars.
\end{proof}

% text: prose-D-13
\phantomsection\label{bp207}%
The normalization sequence reduces constant coefficients on a reduced curve to the smooth case. By Lemma \hyperref[bp202]{D.9}, its cokernel term has finite sections, unchanged under field extension, and no positive cohomology. Moreover, by Lemma \hyperref[bp199]{D.8}, \(H^q(C,\pi_*\underline M)=H^q(C^{\mathrm{nor}},\underline M)\). The long exact sequence therefore yields isomorphisms \(H^q(C,\underline M)\simeq H^q(C^{\mathrm{nor}},\underline M)\) for \(q\geq2\), while in degrees zero and one it gives a natural exact sequence relating \(C\) and \(C^{\mathrm{nor}}\), whose remaining term is a finite group. From these sequences, vanishing and finiteness pass from \(C^{\mathrm{nor}}\) to \(C\), and by the five lemma so does invariance under field extension.

The normalization of an affine or proper curve is again affine or proper. Hence all assertions of Theorem \bpout{4.20} for constant \(\mathbf F_\ell\) follow from the calculations on smooth curves. For \((C_8)\), we apply this comparison to \(C\) and to the inverse image of an open subset. The commutative square of restriction maps then reduces surjectivity to the smooth components.


% object: sec-Extension-by-zero (blueprint/chap-Cohomology-of-torsion-sheaves-on-curves/sec-Extension-by-zero/_section.md), content hidden
\subsection*{\texorpdfstring{D.5\quad Extension by zero}{D.5 Extension by zero}}
\phantomsection\label{bp208}%
\addcontentsline{toc}{subsection}{\texorpdfstring{D.5\quad Extension by zero}{D.5 Extension by zero}}
% text: prose-D-14
\phantomsection\label{bp209}%
A constructible sheaf on a curve is locally constant on a dense open subset. Removing the finitely many remaining points replaces the sheaf by its extension by zero from this open, and changes cohomology only in degrees zero and one. Let \(j:U\hookrightarrow C\) be a dense open in a reduced curve, and let \(i:Z\hookrightarrow C\) be its finite complementary closed subscheme. For a classical sheaf \(F\) we have the exact sequence

\[
0\longrightarrow j_!j^*F\longrightarrow F
\longrightarrow i_*i^*F\longrightarrow0
\]

of an open and its closed complement. We check its exactness on geometric stalks on \(C_{\acute{e}t}\), and then pull it back.

% object: lemma-removing-finitely-many-points (blueprint/chap-Cohomology-of-torsion-sheaves-on-curves/sec-Extension-by-zero/lemma-removing-finitely-many-points/_section.md)
\phantomsection\label{bp210}%
\begin{bpstatementit}{Lemma~D.11}{ (Removing finitely many points)}
Let \(C\) be a reduced separated curve of finite type over an algebraically closed field \(k\), let \(j:U\hookrightarrow C\) be a dense open with finite complement \(i:Z\hookrightarrow C\), let \(\ell\) be a prime, and let \(F\) be a classical constructible sheaf with \(\ell F=0\). For each of the assertions \((C_1)\)–\((C_8)\) of Theorem \bpout{4.20}, under its hypotheses on \(C\) and on the prime \(\ell\), the assertion holds for \(F\) if and only if it holds for \(j_!j^*F\).
\end{bpstatementit}

% proof: proof/lemma-removing-finitely-many-points
\begin{proof}[{Proof}]\phantomsection\label{bp211}%
The third term of the exact sequence displayed above is supported on the finite set \(Z\), so by Lemma \hyperref[bp202]{D.9} it has finite sections and no higher cohomology. Comparing the long exact sequences of \(F\) and \(j_!j^*F\) then shows that each assertion passes in both directions.

The sheaf \(j_!j^*F\) is again killed by \(\ell\), so the hypotheses on the coefficient prime are the same for both sheaves. Since the last term of the exact sequence has no positive cohomology, the long exact sequence yields

\[
\begin{split}
0\to H^0(C,j_!j^*F)\to H^0(C,F)\to H^0(Z,i^*F)\to{}\qquad\qquad\\
H^1(C,j_!j^*F)\to H^1(C,F)\to0,
\end{split}
\]

and isomorphisms in degrees at least two. Since the sequence is natural under field extension and its term supported on \(Z\) is unchanged, the two long exact sequences show that the pullback maps are isomorphisms for \(F\) in all degrees if and only if they are isomorphisms for \(j_!j^*F\) in all degrees. For \((C_8)\), we apply the degree-two identifications after restriction to every open subset of \(C\). See also [Stacks, \href{https://stacks.math.columbia.edu/tag/0GJA}{0GJA}].
\end{proof}

% text: prose-D-15
\phantomsection\label{bp212}%
When \(j\) is dense, Lemma \hyperref[bp210]{D.11}, together with the calculation for constant \(\mathbf F_\ell\) on \(C\) via Lemma \hyperref[bp205]{D.10}, shows that \(j_!\underline{\mathbf F_\ell}\) satisfies the assertions of Theorem \bpout{4.20}. An arbitrary open is dense in the reduced union of the components it meets, a closed subscheme of \(C\), so we can apply the dense case there and pass back to \(C\): direct image along a finite morphism is exact [Stacks, \href{https://stacks.math.columbia.edu/tag/0A3Q}{0A3Q}] and, by Lemma \hyperref[bp199]{D.8}, does not change cohomology.


% object: sec-Reduction-to-constant-coefficients (blueprint/chap-Cohomology-of-torsion-sheaves-on-curves/sec-Reduction-to-constant-coefficients/_section.md), content hidden
\subsection*{\texorpdfstring{D.6\quad Reduction to constant coefficients}{D.6 Reduction to constant coefficients}}
\phantomsection\label{bp213}%
\addcontentsline{toc}{subsection}{\texorpdfstring{D.6\quad Reduction to constant coefficients}{D.6 Reduction to constant coefficients}}
% text: prose-D-16
\phantomsection\label{bp214}%
We reduce a local system to constant sheaves by pulling it back along a cover of degree prime to \(\ell\), on which it becomes unipotent. To find such a cover, note that a local system of finite \(\mathbf F_\ell\)-vector spaces has finite monodromy: the fundamental group acts on a fibre through a finite quotient. The local system then becomes a successive extension of constant sheaves on the cover defined by an \(\ell\)-Sylow subgroup of the monodromy group.

On the full cover of its monodromy the local system even becomes constant. But that cover can have degree divisible by \(\ell\), and then the local system need not be a direct summand of the pushforward of its pullback. For example, let \(v:y^2=x\) be the double cover of \(\mathbf G_m\), in characteristic different from two, and let \(\ell=2\). The composite of the unit \(\mathbf F_2\to v_*\mathbf F_2\) and the trace is multiplication by \(2\), which is zero. In fact \(\mathbf F_2\) is not a direct summand of \(v_*\mathbf F_2\) at all. The monodromy representation of \(v_*\mathbf F_2\) is the regular representation \(\mathbf F_2[\mathbf Z/2]\), which is indecomposable.

For a cover of degree prime to \(\ell\), on the other hand, Lemma \hyperref[bp215]{D.12} shows that the trace divided by the degree splits the unit. It therefore suffices to prove the vanishing, finiteness and invariance statements for the pullback. The pullback is unipotent, that is, built from constant sheaves by extensions.

% object: lemma-reduction-unipotent-local-system (blueprint/chap-Cohomology-of-torsion-sheaves-on-curves/sec-Reduction-to-constant-coefficients/lemma-reduction-unipotent-local-system/_section.md)
\phantomsection\label{bp215}%
\begin{bpstatementit}{Lemma~D.12}{ (Reduction to a unipotent local system)}
Let \(\ell\) be a prime, let \(U\) be a connected scheme, and let \(G\) be a finite classical local system of \(\mathbf F_\ell\)-vector spaces.  There is a finite étale cover \(v:V\to U\) of degree \(d\) prime to \(\ell\) such that \(v^*G\) has a finite filtration with successive quotients the constant sheaf \(\mathbf F_\ell\).  Moreover, \(G\) is a direct summand of \(v_*v^*G\).
\end{bpstatementit}

% proof: proof/lemma-reduction-unipotent-local-system
\begin{proof}[{Proof}]\phantomsection\label{bp216}%
The filtration comes from the representation theory of \(\ell\)-groups, whose only simple module over \(\mathbf F_\ell\) is the trivial one, and the splitting from the trace, whose composite with the unit is multiplication by the degree.

By Proposition \bpout{2.42}, \(G\) is the pullback of a finite étale local system, which corresponds to a finite continuous representation of the fundamental group.  Let \(H\) be the finite image of this representation, choose an \(\ell\)-Sylow subgroup \(P\), and take the cover associated to its inverse image in the fundamental group.  Its degree is \([H:P]\), which is prime to \(\ell\), and the restricted action factors through \(P\).  Since the augmentation ideal of \(\mathbf F_\ell[P]\) is nilpotent, the simple modules of the finite \(\ell\)-group \(P\) are trivial of dimension one, and a composition series yields the required filtration [Stacks, \href{https://stacks.math.columbia.edu/tag/0A3R}{0A3R}].

On a fibre, the unit \(G\to v_*v^*G\) is the diagonal and the trace map \(\operatorname{Tr}:v_*v^*G\to G\) adds the \(d\) entries. Their composite is therefore multiplication by \(d\), which is invertible on an \(\mathbf F_\ell\)-sheaf, so \(d^{-1}\operatorname{Tr}\) is a retraction.
\end{proof}

% text: prose-D-17
\phantomsection\label{bp217}%
To apply this to \(j_!G\) on a reduced curve \(C\), we take for \(U\) a dense smooth open on which \(G\) is locally constant.  Normalizing \(C\) in the finite extensions of function fields defined by \(V\), we get a finite morphism \(\bar v:Y\to C\) with \(V\) the inverse image of \(U\), and an open immersion \(j':V\hookrightarrow Y\).  From the cartesian square formed by these maps and the formula for stalks of direct images along finite morphisms, we find that

\[
j_!v_*v^*G=\bar v_*j'_!v^*G,
\]

which is also the exchange formula of [Stacks, \href{https://stacks.math.columbia.edu/tag/0GJ1}{0GJ1}].  If \(C\) is proper, then \(Y\) is proper, and if \(C\) is affine, then \(Y\) is affine.

The assertions travel from constant sheaves to \(j'_!v^*G\) on \(Y\), then to \(j_!v_*v^*G\) on \(C\), and from there to the direct summand \(j_!G\). Since \(j'_!\) is exact, the filtration of \(v^*G\) gives a filtration of \(j'_!v^*G\) with graded pieces \(j'_!\mathbf F_\ell\). The assertions of Theorem \bpout{4.20} hold for these pieces by the calculation for constant sheaves extended by zero. As we show at the start of the next section, they pass along extensions, so they hold for \(j'_!v^*G\). Since \(\bar v_*\) is exact, they hold for \(j_!v_*v^*G\). By the splitting given by the trace, they hold for its direct summand \(j_!G\). For invariance under field extension we use that the unit and the trace commute with extension of the ground field. Then the pullback map for \(j_!G\) is a retract of the pullback map for \(j_!v_*v^*G\), and a retract of an isomorphism is an isomorphism [Stacks, \href{https://stacks.math.columbia.edu/tag/03SD}{03SD}].

If \(C\) is not equidimensional, we choose a dense open that is smooth and is the disjoint union of the regular loci of the one-dimensional components and the isolated points.  On the isolated points taking sections is exact, each connected one-dimensional piece is treated as above, and finite direct sums preserve all the assertions.


% object: sec-Devissage (blueprint/chap-Cohomology-of-torsion-sheaves-on-curves/sec-Devissage/_section.md), content hidden
\subsection*{\texorpdfstring{D.7\quad Dévissage}{D.7 Dévissage}}
\phantomsection\label{bp218}%
\addcontentsline{toc}{subsection}{\texorpdfstring{D.7\quad Dévissage}{D.7 Dévissage}}
% text: prose-D-18
\phantomsection\label{bp219}%
The assertions of Theorem \bpout{4.20} pass along short exact sequences.  To see this, let \(0\to F_1\to F\to F_2\to0\) be exact, and suppose that the assertions are known for \(F_1\) and \(F_2\).  Vanishing and finiteness follow from the long exact sequence, and invariance under field extension follows from the map between the two long exact sequences and the five lemma. For \((C_8)\) an extra argument is needed. Let \(b\in H^2(U,F)\), and lift its image in \(H^2(U,F_2)\) to \(c_2\in H^2(C,F_2)\).  The obstruction to lifting \(c_2\) to \(H^2(C,F)\) lies in \(H^3(C,F_1)\), which vanishes by \((C_1)\), so such a lift \(c\) exists. The difference \(b-c|_U\) comes from \(H^2(U,F_1)\), and by \((C_8)\) for \(F_1\) it lifts to \(H^2(C,F_1)\); adding that lift to \(c\) gives the required global class [Stacks, \href{https://stacks.math.columbia.edu/tag/0A5C}{0A5C}].

% proof: proof/theorem-cohomology-curves/2
\begin{proof}[{Proof of Theorem 4.20}]\phantomsection\label{bp220}%
For constructible sheaves we reduce to the explicit calculations, and arbitrary torsion sheaves are filtered colimits of these. The reduction has four steps: dévissage to constructible \(\mathbf F_\ell\)-sheaves, extension by zero from an open on which the sheaf is locally constant, a cover of degree prime to \(\ell\) that makes it unipotent, and normalization, which leaves constant coefficients on smooth curves.

By Theorem \bpout{2.35}, we may work on \(C_{\acute{e}t}\), and since a universal homeomorphism does not change the étale topos, we may replace \(C\) by \(C_{\mathrm{red}}\).  A constructible torsion sheaf has a finite exponent, because it has finitely many isomorphism types of finite stalks.  It is therefore the direct sum of its finitely many primary parts, and a finite filtration of an \(\ell\)-primary part has constructible \(\mathbf F_\ell\)-sheaves as graded pieces.  By the stability under extensions just proved, it suffices to treat such sheaves.

Choose a dense smooth open \(j:U\hookrightarrow C\) on which the sheaf is finite locally constant.  By Lemma \hyperref[bp210]{D.11} we may replace the sheaf by its extension by zero.  On the connected components of \(U\), Lemma \hyperref[bp215]{D.12} and finite direct images reduce us to constant \(\mathbf F_\ell\) extended by zero on curves finite over \(C\); this is allowed because the assertions pass to direct summands and because, by Lemma \hyperref[bp199]{D.8}, finite direct images are exact and do not change cohomology.  Finally, normalization and Lemma \hyperref[bp202]{D.9} reduce these sheaves to the calculations with constant coefficients on smooth curves, namely the Kummer sequence for \(\ell\ne p\) and the Artin–Schreier sequence for \(\ell=p\).  Since all reductions preserve the hypothesis that \(C\) is proper or affine, the theorem holds for constructible torsion sheaves.

An arbitrary torsion sheaf on the qcqs scheme \(C\) is a filtered colimit of constructible torsion sheaves, and cohomology commutes with filtered colimits [Stacks, \href{https://stacks.math.columbia.edu/tag/03SA}{03SA}, \href{https://stacks.math.columbia.edu/tag/03Q5}{03Q5}].  For torsion prime to the characteristic or \(p\)-primary torsion, we may take the constructible sheaves with the same kind of torsion, by splitting off primary parts before passing to the colimit. Pullback and restriction also commute with filtered colimits. Hence vanishing and the pullback isomorphisms pass to the colimit. So does surjectivity in \((C_8)\), because a representative of a class already occurs at some stage. Finiteness is asserted only for constructible sheaves. Finally, Theorem \bpout{2.35} carries these conclusions and their maps over to \(C_{\mathrm{pro}\acute{e}t}\).
\end{proof}



% object: chap-Specialization-and-contracting-correspondences (blueprint/chap-Specialization-and-contracting-correspondences/_section.md), content hidden
\section*{\texorpdfstring{E\quad Specialization and contracting correspondences}{E Specialization and contracting correspondences}}
\phantomsection\label{bp221}%
\addcontentsline{toc}{section}{\texorpdfstring{E\quad Specialization and contracting correspondences}{E Specialization and contracting correspondences}}
% text: prose-E-01
\phantomsection\label{bp222}%
The trace formula needs the local term at an isolated fixed point to be the trace on the stalk. This is not a formal fact. A local term is computed from \(\cHom(c_2^*K,c_1^!K)\) near the fixed point, and the stalk does not determine this complex (Remark \bpout{12.4}). Instead, Chapter \bpout{12} proves the equality by a computation for Frobenius on a smooth curve (Proposition \bpout{12.16}) and reduces everything else to that case. In particular, its trace formula (Theorem \bpout{12.24}) does not depend on this appendix.

The computation leaves open why Frobenius behaves so well. Topology suggests that the reason is contraction. If a map \(\varphi\) of a manifold contracts a neighbourhood of a fixed point \(y\) into itself, then the local index at \(y\) is \(1\). With coefficients in a sheaf, the local index is the trace of \(\varphi\) on the stalk at \(y\). Frobenius contracts in the strongest possible sense: its differential is zero. In this appendix we make this explanation into a theorem, in any dimension and for a whole class of correspondences at once.

In algebraic geometry we measure contraction by powers of ideals. Let \(c=(c_1,c_2):C\to X\times X\) be a correspondence, \(Z\subset X\) a closed subscheme with ideal \(I\), and put \(I_i=c_i^{-1}(I)\mathcal O_C\). We say, as in [\hyperref[bp347]{Var05}], that \(c\) contracts towards \(Z\) if \(I_2\subset I_1\) and \(I_2^n\subset I_1^{n+1}\) for some \(n\geq1\). Frobenius is the model case. Indeed, take \(c=(\mathrm{id},\Phi_X^m)\) and \(Z\) a point \(y\in X_0(\mathbf F_{q^m})\), and choose coordinates \(x_i\) of an affine neighbourhood of \(y\) defined over \(\mathbf F_q\). Then \(y\) has coordinates \(a_i\in\mathbf F_{q^m}\), and \((\Phi_X^m)^*(x_i-a_i)=x_i^{q^m}-a_i=(x_i-a_i)^{q^m}\). Hence \(I_2\subset I_1^{q^m}\), and \(n=1\) works.

The main theorem, Theorem \hyperref[bp311]{E.30}, says that near a component of its fixed locus that meets \(Z\), the trace class of a contracting correspondence can be computed on \(Z\) alone, from the correspondence restricted to \(Z\). For \(Z\) a point, Corollary \hyperref[bp315]{E.31} says that the local term at an isolated fixed point is the trace on the stalk, even when \(X\) is singular. The theorem and its proof are due to Varshavsky [\hyperref[bp347]{Var05}]. They generalize Fujiwara's theorem, conjectured by Deligne, that the local terms of a correspondence composed with a sufficiently high power of Frobenius are the naive ones.

The proof makes the contraction literal by deforming \(X\) to the normal cone of \(Z\). Over the trait \(\operatorname{Spec}k[t]_{(t)}\), the deformation to the normal cone \(\mathcal D_ZX\) has generic fibre \(X\) and special fibre the normal cone \(N_ZX\). On the special fibre, the contraction hypothesis says that the second leg maps everything into the zero section (Lemma \hyperref[bp302]{E.26}). There the part of \(K\) away from \(Z\) contributes nothing. What remains is the specialization of \(K\) restricted to the zero section, and by Verdier's theorem (Theorem \hyperref[bp244]{E.8}) this is \(K|_Z\).

It remains to transport the answer back to \(X\). We identify the trace class on the special fibre with the specialization of the trace class of \(c\) (Proposition \hyperref[bp289]{E.22}). The specialization map is an isomorphism, because the fixed locus does not move in the family up to nilpotents (Lemma \hyperref[bp306]{E.28}). Before any of this, we must be able to specialize a trace class at all. A trace class is built from \(K\boxtimes\mathbb DK\), so we need nearby cycles to commute with external products and with Verdier duality. Propositions \hyperref[bp248]{E.9} and \hyperref[bp276]{Proposition E.18} prove both, by a dimension argument of Illusie [\hyperref[bp351]{Ill94}].

We use the six operations and the notation for correspondences and trace classes of Chapter \bpout{11}, and classical results on nearby cycles over a trait, which we recall where they are first needed. Throughout this appendix, \(k\) is algebraically closed, \(\ell\) is invertible in \(k\), and \(A\) is a finite commutative ring annihilated by a power of \(\ell\). The only exception to these conventions is the base of nearby cycles in Sections \hyperref[bp246]{E.2}–\hyperref[bp287]{Section E.5}, which may be any excellent strictly henselian trait, with residue field only separably closed. As in the rest of the notes, statements about constructible complexes on \(k\)-schemes are made on the pro-étale site, but we construct nearby cycles on the classical étale site and transport them by the comparison for constructible complexes (Theorem \bpout{2.43} and Corollary \bpout{2.44}).

% object: sec-Specialization-to-the-normal-cone (blueprint/chap-Specialization-and-contracting-correspondences/sec-Specialization-to-the-normal-cone/_section.md), content hidden
\subsection*{\texorpdfstring{E.1\quad Specialization to the normal cone}{E.1 Specialization to the normal cone}}
\phantomsection\label{bp223}%
\addcontentsline{toc}{subsection}{\texorpdfstring{E.1\quad Specialization to the normal cone}{E.1 Specialization to the normal cone}}
% text: prose-E-02
\phantomsection\label{bp224}%
Let \(S\) be the strict henselization of \(\operatorname{Spec}k[t]_{(t)}\), with generic point \(\eta\). Nearby cycles move a complex from the generic fibre of a family over \(S\) to its special fibre. Since the special fibre does not lie in the generic fibre, we cannot simply restrict. The obvious attempt is to extend by \(Rj_*\) along the open generic fibre and restrict by \(i^*\) to the special fibre. This fails already for \(S\) itself, because the fraction field of \(S\) is not separably closed. Its absolute Galois group is the inertia group, whose tame quotient is \(\prod_{\ell'\neq p}\mathbf Z_{\ell'}(1)\). For \(n\) invertible in \(k\), the stalk of \(R^1j_*(\mathbf Z/n)\) at the closed point is therefore \(H^1(\eta,\mathbf Z/n)=\mathbf Z/n(-1)\neq0\). This class comes from the arithmetic of \(\eta\), not from the geometry of the family. Passing first to the geometric generic fibre removes such Galois classes. After this change, the constant sheaf specializes to itself, in degree zero, by the first calculation below.

Both specialization to the normal cone and the comparison of trace classes in Proposition \hyperref[bp289]{E.22} are built from nearby cycles on the deformation to the normal cone. The key input is Verdier's theorem, Theorem \hyperref[bp244]{E.8} (see [\hyperref[bp347]{Var05}, §§1.3–1.4 and 3]).

The basic computation is that nearby cycles of a constant family \(Y\times S\) give back the complex we started with (Lemma \hyperref[bp227]{E.2}). We use it twice. It makes the map \(\tau_{X,Z}\) below land in \(i^*L\), and it makes the specialization of trace classes the identity when the fixed scheme does not move (Corollary \hyperref[bp298]{E.25}). Its proof uses only the definition of nearby cycles, their base change maps and two calculations, which we now recall.

Choose a geometric generic point \(\bar\eta\) of \(S\), and let \(\mathcal Y\to S\) be separated of finite type. To define nearby cycles, we pull \(\mathcal Y\) back to the integral closure of \(S\) in the chosen separable closure of its fraction field, and let \(\bar i\) and \(\bar j\) be the inclusions of the special fibre and of the geometric generic fibre of this pullback.  Then, for \(L\) on \(\mathcal Y_\eta\) with pullback \(L_{\bar\eta}\) to the geometric generic fibre, the geometric nearby cycles are

\[
R\Psi_{\mathcal Y}(L)=\bar i^{*}R\bar j_*L_{\bar\eta}.
\]

This defines a functor \(R\Psi_{\mathcal Y}:D^b_c(\mathcal Y_{\eta,\acute et},A)\to
D(\mathcal Y_{s,\acute et},A)\). We suppress the inertia action, except where invariants under the wild inertia group of a finite extension of traits are taken below.

% object: example-nearby-cycles-node (blueprint/chap-Specialization-and-contracting-correspondences/sec-Specialization-to-the-normal-cone/example-nearby-cycles-node/_section.md)
\phantomsection\label{bp225}%
\begin{bpstatement}{Example~E.1}{ (A node)}
Nearby cycles see more than the special fibre when the family degenerates. For example, for \(\mathcal Y=V(xy-t)\subset\mathbf A^2_S\) the special fibre is the pair of axes \(xy=0\), while the generic fibre is smooth. To see what to expect, note that over \(\mathbf C\) the part of the fibre \(xy=\varepsilon\) near the origin is an annulus \(\{(x,\varepsilon/x)\}\), which retracts onto a circle that vanishes as \(\varepsilon\to0\). Correspondingly, in our setting the stalk of \(R\Psi_{\mathcal Y}(A)\) at the node is \(A\) in degree zero and \(A(-1)\) in degree one; the twist is that of \(H^1(\mathbf G_m)\) (see [\hyperref[bp352]{SGA7II}, Exposé XV]). At every other point of the special fibre, by contrast, the family is smooth over \(S\), and there the nearby cycles are \(A\) in degree zero, by smooth base change and the first calculation below.
\end{bpstatement}

% text: prose-E-03
\phantomsection\label{bp226}%
For a morphism \(g:\mathcal Y'\to\mathcal Y\) of separated \(S\)-schemes of finite type, there is a base change map for pullback

\[
g_s^*R\Psi_{\mathcal Y}
\longrightarrow R\Psi_{\mathcal Y'}g_\eta^*,
\]

which is invertible for smooth \(g\). For proper \(g\) there is a base change map for proper direct image

\[
Rg_{s*}R\Psi_{\mathcal Y'}
\xrightarrow{\sim}R\Psi_{\mathcal Y}Rg_{\eta*}.
\]

These base change maps commute with composition and with the specialization map defined below. Their invertibility follows from two base change theorems for torsion coefficients. First, proper base change [Stacks, \href{https://stacks.math.columbia.edu/tag/095T}{095T}] says that restriction to a fibre commutes with proper direct image. Second, smooth base change [Stacks, \href{https://stacks.math.columbia.edu/tag/0F09}{0F09}] (for sheaves [Stacks, \href{https://stacks.math.columbia.edu/tag/0EYU}{0EYU}]) says that pullback along a smooth map commutes with direct image along a qcqs morphism, when the coefficients are torsion of order invertible on the base. Suppose first that \(g\) is proper. The two direct images on the generic fibres agree because direct image is compatible with composition, and \(\bar i^*\) commutes with the proper direct image by proper base change. Suppose next that \(g\) is smooth. Since \(R\bar j_*\) is direct image along a qcqs morphism, we may apply smooth base change to it and then restrict to the special fibre. Smooth base change applies because the coefficients are annihilated by a power of \(\ell\), which is invertible on \(S\). The same base change maps appear in [\hyperref[bp352]{SGA7II}], Exposé XIII, (2.1.7.1)–(2.1.7.3).

We also use two calculations for a finite \(A\)-module \(M\). The first is \(R\Psi_S(M_{\eta})=M\). It holds because the integral closure \(\bar S\) of \(S\) in the separable closure is strictly henselian, so the stalk of \(R\bar j_*M\) at its closed point is \(R\Gamma(\bar\eta,M)=M\). The second is the usual calculation for the affine line with coefficients prime to the characteristic: \(H^r(\mathbf A^1_k,M)=0\) for \(r>0\) and \(H^0(\mathbf A^1_k,M)=M\). It follows from [Stacks, \href{https://stacks.math.columbia.edu/tag/03RR}{03RR}] applied to the compactification \(\mathbf A^1\subset\mathbf P^1\), which has genus zero and one point at infinity. To apply that result, note that every finite \(M\) is a finite direct sum of cyclic modules, and over the algebraically closed field \(k\) the cyclic constant coefficients are isomorphic to the corresponding roots of unity. By truncation the two calculations extend to a bounded complex \(M^\bullet\) of finite \(A\)-modules, viewed as a constant complex:

\[
R\Psi_S(M^\bullet_\eta)\simeq M^\bullet,
\qquad R\Gamma(\mathbf A^1_k,M^\bullet_{\mathbf A^1_k})\simeq M^\bullet.
\]

Finally, all complexes stay bounded and constructible. By the finiteness theorem for nearby cycles with torsion coefficients ([\hyperref[bp353]{SGA4½}], \emph{Théorèmes de finitude}, §3; [illusie94, §4.1, pp. 44–45]), \(R\Psi_{\mathcal Y}\) takes bounded constructible complexes of finite Tor dimension to complexes of the same kind. Finite Tor dimension is what we need when we take traces. The same holds for an arbitrary bounded constructible complex, without the Tor condition. To see this, we may forget the \(A\)-action, since the underlying torsion sheaves of \(A\)-linear nearby cycles are the usual nearby cycles. Bounded constructible complexes are stable under extensions and truncation, so the filtration by powers of \(\ell\) reduces us to constructible \(\mathbf F_\ell\)-sheaves. These have finite Tor dimension over \(\mathbf F_\ell\), and the finiteness theorem applies.

For a separated \(k\)-scheme \(Y\) of finite type, the constant family \(Y\times S\) has special fibre \(Y\). For \(L\) on \(Y\), let \(L_\eta\) be its pullback to the generic fibre.  Restricting the adjunction unit of \(\bar j\) on the pullback of \(L\) to the special fibre, we get the specialization map \(\sigma_Y(L)\) below.

% object: lemma-nearby-cycles-constant-family (blueprint/chap-Specialization-and-contracting-correspondences/sec-Specialization-to-the-normal-cone/lemma-nearby-cycles-constant-family/_section.md)
\phantomsection\label{bp227}%
\begin{bpstatementit}{Lemma~E.2}{ (Nearby cycles of a constant family)}
For a separated \(k\)-scheme \(Y\) of finite type and \(L\in\Dbc(Y,A)\), the specialization map

\[
\sigma_Y(L):L\longrightarrow
R\Psi_{Y\times S}(L_\eta)
\]

is an isomorphism. It is natural in \(L\), and it is compatible with the base change maps above in the following sense.

\begin{enumerate}
\item If \(f:Y'\to Y\) is proper and \(L'\in\Dbc(Y',A)\), then the composite of \(Rf_*\sigma_{Y'}(L')\) with the base change map for proper direct image is \(\sigma_Y(Rf_*L')\), where \((Rf_*L')_\eta\) is identified with \(Rf_{\eta*}(L'_\eta)\) by proper base change.
\item If \(f:Y'\to Y\) is smooth, then the composite of \(f^*\sigma_Y(L)\) with the base change map for pullback is \(\sigma_{Y'}(f^*L)\).
\end{enumerate}
\end{bpstatementit}

% proof: proof/lemma-nearby-cycles-constant-family
\begin{proof}[{Proof}]\phantomsection\label{bp228}%
The compatibilities (1) and (2) are formal. For the isomorphism we show that the cone \(C_Y(L)\) of \(\sigma_Y(L)\) vanishes. To do this we use that, like nearby cycles, the cone commutes with smooth pullback and proper direct image. Smooth pullback lets us check vanishing locally, and it also shows that the cone vanishes generically, by reduction to constant coefficients and the calculation \(R\Psi_S(M_\eta)=M\). Proper direct image along a compactified linear projection then shows that the remaining support is empty, by induction on the dimension of an affine space containing \(Y\).

By the comparison for constructible complexes, we may work on the small étale site. There \(\sigma\) is \(\bar i^*\) of the unit of \(\bar j^*\dashv R\bar j_*\) and the base change maps are \(\bar i^*\) of mates, so (1), (2) and naturality follow from the triangle identities.

We take cofibres in the stable \(\infty\)-category underlying the derived category, where they are functorial, unlike cones in a triangulated category, and set \(C_Y(L)=\operatorname{cofib}(\sigma_Y(L))\).  This is an exact functor with bounded constructible values.  By (1) and (2), the base change maps yield

\[
C_Y(Rf_*L')\simeq Rf_*C_{Y'}(L')\quad(f\text{ proper}),\qquad
C_{Y'}(f^*L)\simeq f^*C_Y(L)\quad(f\text{ smooth}).
\]

By the base change isomorphism for étale maps, the question is local on \(Y\).  For a closed immersion \(Y\hookrightarrow\mathbf A^n\), the base change isomorphism for proper maps identifies the direct image of \(C_Y(L)\) with \(C_{\mathbf A^n}\) of the direct image of \(L\), and since direct image along a closed immersion is conservative, it suffices to treat affine spaces.

For \(Y=\mathbf A^n\), the complex \(C_Y(L)\) vanishes on a dense open subset.  Indeed, by the truncation triangles we may assume that \(L\) is a sheaf. After shrinking, it is locally constant, and after pulling back to a finite étale cover, it is constant. The pullback does no harm: \(C\) commutes with étale pullback, so vanishing can be checked after a surjective étale pullback.  For a constant sheaf \(M\) on a smooth \(k\)-scheme, the base change isomorphism for its smooth structure map identifies \(C\) with the pullback of \(C_{\operatorname{Spec}k}(M)\), which is zero because \(R\Psi_S(M_\eta)=M\).

We argue by induction on \(n\), the case of a point being the calculation above.  By the preceding paragraph, the closure \(T\) of the support of \(C_{\mathbf A^n}(L)\) is a proper closed subset.  By the linear form of Noether normalization there is a linear projection \(q:\mathbf A^n\to\mathbf A^{n-1}\) that is finite on \(T\): we choose its kernel line so that it avoids the intersection of the projective closure of \(T\) with the hyperplane at infinity.  Compactifying that line, we get

\[
\bar q:\bar Y=\mathbf P^1\times\mathbf A^{n-1}
\longrightarrow\mathbf A^{n-1},
\qquad v:\mathbf A^n\hookrightarrow\bar Y,
\qquad \bar L=v_!L.
\]

We compactified the line so that the support becomes finite over \(\mathbf A^{n-1}\) and can be detected by direct image there. As \(C\) commutes with restriction to the open subscheme \(\mathbf A^n\) of \(\bar Y\), the complex \(C_{\bar Y}(\bar L)\) is supported on \(T\cup(\bar Y\setminus\mathbf A^n)\).  This union is closed in \(\bar Y\) and finite over \(\mathbf A^{n-1}\). Indeed, \(T\) is finite, hence proper, over \(\mathbf A^{n-1}\), and therefore closed in \(\bar Y\); the boundary \(\bar Y\setminus\mathbf A^n\) is a section.  By the base change isomorphism for proper maps and the induction hypothesis,

\[
R\bar q_*C_{\bar Y}(\bar L)
\simeq C_{\mathbf A^{n-1}}(R\bar q_*\bar L)=0.
\]

Direct image is conservative on complexes supported on a closed subscheme finite over the base, since on a geometric stalk the direct image of each cohomology sheaf is the finite product of its stalks on the finite fibre.  Therefore \(C_{\bar Y}(\bar L)=0\), and restricting to \(\mathbf A^n\) gives \(C_{\mathbf A^n}(L)=0\).
\end{proof}

% text: prose-E-04
\phantomsection\label{bp229}%
The contracting argument needs a complex on the normal cone \(N_ZX\) that still remembers \(K\) near \(Z\). Pulling \(K|_Z\) back along the projection \(N_ZX\to Z\) is not enough, since it forgets the directions leaving \(Z\). For example, take \(X=\mathbf A^1\), \(Z\) the origin, and \(K=j_!A\) for \(j:\mathbf A^1\setminus\{0\}\hookrightarrow\mathbf A^1\). Then \(K|_Z=0\), so the pullback is zero. By contrast, the deformation to the normal cone is \(\operatorname{Spec}k[x,t,x/t]\) over the trait, that is, the constant family \(\mathbf A^1\times S\) in the coordinate \(x'=x/t\), and on its generic fibre \(K\) is \(j_!A\) for \(x'\neq0\). Hence, by Lemma \hyperref[bp227]{E.2}, the nearby cycles are \(j_!A\) on \(N_ZX=\mathbf A^1\), which is not zero. This construction is the specialization functor \(\operatorname{Sp}_{X,Z}\) defined below, and by Verdier's theorem (Theorem \hyperref[bp244]{E.8}) it still recovers \(K|_Z\) on the zero section. Theorem \hyperref[bp311]{E.30} uses both properties.

Let \(X\) be a separated \(k\)-scheme of finite type, and let \(i:Z\hookrightarrow X\) be a closed subscheme with ideal \(I\). The deformation to its normal cone, over the chosen trait, is

\[
\mathcal D_ZX=
\operatorname{Spec}_{X\times S}
\bigl(\mathcal O_{X\times S}[I/t]\bigr).
\]

Here \(\mathcal O_{X\times S}[I/t]\) is the base change to \(S\) of the subalgebra of \(\mathcal O_X[t,t^{-1}]\) generated by \(\mathcal O_X[t]\) and the fractions \(a/t\) with \(a\in I\).  Its generic fibre is \(X_\eta\) and its special fibre is

\[
N_ZX=\operatorname{Spec}_Z
\left(\bigoplus_{r\geq0}I^r/I^{r+1}\right).
\]

The quotient by \(I\) and all \(a/t\) defines a closed immersion \(\widetilde z:Z\times S\hookrightarrow\mathcal D_ZX\), whose special fibre \(z:Z\hookrightarrow N_ZX\) is the zero section.  For \(L\) on \(X\), let \(L_\eta\) be its pullback to the generic fibre \(X_\eta\) of \(\mathcal D_ZX\), and define the specialization functor \(\operatorname{Sp}_{X,Z}:D^b_c(X_{\acute et},A)\to
D^b_c((N_ZX)_{\acute et},A)\) by

\[
\operatorname{Sp}_{X,Z}(L)=R\Psi_{\mathcal D_ZX}(L_\eta).
\]

The base change map for pullback along \(\widetilde z\), followed by the inverse of the isomorphism of Lemma \hyperref[bp227]{E.2}, defines a natural transformation

\[
\tau_{X,Z}(L):z^*\operatorname{Sp}_{X,Z}(L)
\longrightarrow R\Psi_{Z\times S}((i^*L)_\eta)
\xrightarrow{\sim}i^*L.
\tag{E.1}
\]

We use this map, in this direction, to specialize trace classes.

% object: example-normal-cone-point (blueprint/chap-Specialization-and-contracting-correspondences/sec-Specialization-to-the-normal-cone/example-normal-cone-point/_section.md)
\phantomsection\label{bp230}%
\begin{bpstatement}{Example~E.3}{ (Normal cones of a point)}
If \(X\) is smooth at a point \(Z=\{y\}\), then \(\bigoplus_rI^r/I^{r+1}\) is the symmetric algebra of \(I/I^2\) and \(N_ZX\) is the tangent space at \(y\). At a singular point, however, the normal cone is the tangent cone, and it need not be reduced. For example, for the cusp \(X=V(y^2-x^3)\subset\mathbf A^2\) and \(Z\) the origin, with \(I=(x,y)\), the graded algebra is \(k[\xi,\eta]/(\eta^2)\), where \(\xi,\eta\) are the classes of \(x,y\) in \(I/I^2\) and \(\eta^2\) is the lowest-degree part of \(y^2-x^3\); so \(N_ZX\) is a double line. Similarly, for a nonreduced point \(Z=V(x^m)\subset\mathbf A^1\), each \(I^r/I^{r+1}=(x^{mr})/(x^{m(r+1)})\) is free of rank one over \(\mathcal O_Z=k[x]/(x^m)\), and \(N_ZX=Z\times\mathbf A^1\) is a line thickened along \(Z\); this is the case treated in Lemma \hyperref[bp238]{E.6}.
\end{bpstatement}

% text: prose-E-05
\phantomsection\label{bp231}%
We will apply the base change maps for nearby cycles to the deformations of different pairs, and for this we need to know how the deformation depends on the pair.

% object: lemma-functoriality-deformation (blueprint/chap-Specialization-and-contracting-correspondences/sec-Specialization-to-the-normal-cone/lemma-functoriality-deformation/_section.md)
\phantomsection\label{bp232}%
\begin{bpstatementit}{Lemma~E.4}{ (Functoriality of the deformation)}
Let \(f:X'\to X\) be a morphism of separated \(k\)-schemes of finite type, let \(Z'\subset X'\) be a closed subscheme with ideal \(I'\), and suppose that \(f^{-1}(I)\mathcal O_{X'}\subset I'\).  Then \(f\) induces \(\widetilde f:\mathcal D_{Z'}X'\to\mathcal D_ZX\). If \(Z'=f^{-1}Z\) schematically, then

\[
\mathcal D_{Z'}X'\longrightarrow X'\times_X\mathcal D_ZX
\]

is a closed immersion, and it is an isomorphism if \(f\) is flat. Moreover, \(\widetilde f_s^{-1}(Z)=Z'\) schematically. In this case \(\widetilde f\) is proper when \(f\) is proper and smooth when \(f\) is smooth.
\end{bpstatementit}

% proof: proof/lemma-functoriality-deformation
\begin{proof}[{Proof}]\phantomsection\label{bp233}%
All assertions are statements about the algebras \(\mathcal O[I/t]\), and we check them on generators. The first morphism is defined by \(a/t\mapsto f^\#(a)/t\), which makes sense because \(f^\#(a)\in I'\) for \(a\in I\). If \(I'=I\mathcal O_{X'}\), the natural map from the base change of the Rees algebra onto \(\mathcal O_{X'\times S}[I'/t]\) is surjective, and it is injective if \(f\) is flat.  On the associated graded algebras, the generators of \(\operatorname{gr}_I\) of positive degree generate the ideal of elements of positive degree of \(\operatorname{gr}_{I'}\).  The quotient by them is \(\mathcal O_{Z'}\), and this proves the assertion about the zero section.  If \(f\) is proper, then \(\widetilde f\) is the closed immersion followed by the base change of \(f\), which is proper. If \(f\) is smooth, it is flat, so \(\widetilde f\) is the base change \(X'\times_X\mathcal D_ZX\to\mathcal D_ZX\) of \(f\), which is smooth.
\end{proof}

% text: prose-E-06
\phantomsection\label{bp234}%
Theorem \hyperref[bp244]{E.8} asserts that \(\tau_{X,Z}\) is an isomorphism, that is, that the cofibre

\[
C_{X,Z}(L)=\operatorname{cofib}\bigl(
\tau_{X,Z}(L)\bigr)
\]

vanishes. As for \(C_Y\), we compute the cofibre in the stable \(\infty\)-category, so that \(C_{X,Z}\) is an exact functor. The proof has three steps.  First, the base change maps of the next lemma reduce us to the case \(i^*L=0\) and, after a blow-up, to the case of an effective Cartier divisor.  Then, for a principal divisor in a normal scheme, an explicit contraction of the affine line shows that \(C_{X,Z}(L)\) vanishes at the generic points of \(Z\).  Finally, after an embedding into affine space, projection along a line and induction on the dimension show that the remaining support is empty.

% object: lemma-base-change-cone-x-z (blueprint/chap-Specialization-and-contracting-correspondences/sec-Specialization-to-the-normal-cone/lemma-base-change-cone-x-z/_section.md)
\phantomsection\label{bp235}%
\begin{bpstatementit}{Lemma~E.5}{ (Base change for the cone of $\tau_{X,Z}$)}
Let \(f:X'\to X\) be a morphism of separated \(k\)-schemes of finite type, let \(Z'=f^{-1}Z\) schematically, and let \(f_Z:Z'\to Z\) be the restriction of \(f\).  There are isomorphisms

\[
\begin{aligned}
C_{X,Z}(Rf_*L')&\simeq Rf_{Z*}C_{X',Z'}(L')
&&(f\text{ proper}),\\
C_{X',Z'}(f^*L)&\simeq f_Z^*C_{X,Z}(L)
&&(f\text{ smooth}).
\end{aligned}
\]

Moreover, \(C_{X,Z}(i_*M)=0\) for every \(M\in D^b_c(Z_{\acute et},A)\). Consequently, writing \(j:X\setminus Z\hookrightarrow X\), we have

\[
C_{X,Z}(j_!j^*L)\xrightarrow{\sim}C_{X,Z}(L).
\]
\end{bpstatementit}

% proof: proof/lemma-base-change-cone-x-z
\begin{proof}[{Proof}]\phantomsection\label{bp236}%
Both the source and the target of \(\tau\) are compatible with proper direct image and smooth pullback, so the cone is as well. To prove the vanishing on \(i_*M\), we reduce to the pair \((Z,Z)\), where the deformation is a constant family.

By Lemma \hyperref[bp232]{E.4}, \(f\) induces a morphism \(\widetilde f\) of deformations, which is proper if \(f\) is proper and smooth if \(f\) is smooth. To this morphism we apply the base change maps of nearby cycles, together with proper or smooth base change for the restriction to the zero section.  The resulting squares for \(\tau\) then commute because the base change map for pullback is compatible with composition, and their cofibres give the two isomorphisms. In the special case \(f=i\), the source pair is \((Z,Z)\) and \(\mathcal D_ZZ=Z\times S\), so \(\tau_{Z,Z}\) is the inverse of the isomorphism of Lemma \hyperref[bp227]{E.2}. Hence \(C_{Z,Z}=0\), and since \(i\) is proper, the first isomorphism gives \(C_{X,Z}(i_*M)=0\). The last assertion then follows by applying the exact functor \(C_{X,Z}\) to \(j_!j^*L\to L\to i_*i^*L\).
\end{proof}

% text: prose-E-07
\phantomsection\label{bp237}%
To prove the vanishing at the generic points of \(Z\), we first treat a nonreduced point on the affine line, where an explicit contraction is available.

% object: lemma-case-z-v-x-m (blueprint/chap-Specialization-and-contracting-correspondences/sec-Specialization-to-the-normal-cone/lemma-case-z-v-x-m/_section.md)
\phantomsection\label{bp238}%
\begin{bpstatementit}{Lemma~E.6}{ (The case $Z=V(x^m)\subset\mathbf A^1$)}
For \(X=\mathbf A^1_k\), \(Z=V(x^m)\) with \(m\geq1\), and every finite \(A\)-module \(M\), one has \(C_{X,Z}(M_X)=0\).
\end{bpstatementit}

% proof: proof/lemma-case-z-v-x-m
\begin{proof}[{Proof}]\phantomsection\label{bp239}%
We give a homotopy argument. The contraction \(x\mapsto ax\) joins the identity (\(a=1\)) to a map that factors through \(Z\) (\(a=0\)), and the endomorphisms of \(z^*\operatorname{Sp}_{X,Z}(M_X)\) that these maps induce form a family over \(\mathbf A^1\). Since \(\mathbf A^1\) has the cohomology of a point, the family is constant, and comparing its two ends shows that \(\tau_{X,Z}(M_X)\) has an inverse.

Consider the contraction \(h:X\times\mathbf A^1\to X\), \(h(x,a)=ax\), and write \(h_a=h(-,a):X\to X\) for \(a\in k\).  It maps \(Z\times\mathbf A^1\) schematically into \(Z\), since \(h^*(x^m)=a^m x^m\in(x^m)\), and by Lemma \hyperref[bp232]{E.4}, whose hypothesis is this inclusion of ideals, all \(h_a\) lift to the deformations.  Moreover, \(h_1=\operatorname{id}\), and \(h_0\) factors through \(Z\) and restricts to the identity on \(Z_{\mathrm{red}}\).  Set \(N=z^*\operatorname{Sp}_{X,Z}(M_X)\). The base change maps for pullback along these lifts induce endomorphisms \(\phi_a:N\to N\), with \(\phi_1=\operatorname{id}_N\).

There is one endomorphism \(\phi\) of the pullback of \(N\) to \(Z\times\mathbf A^1\) whose fibre at \(a\) is \(\phi_a\).  To construct it, we apply the base change map for pullback along \(h:X\times\mathbf A^1\to X\) and then invert the base change map for the projection to \(X\), which is smooth; by flat base change, the deformation of \(X\times\mathbf A^1\) is \(\mathcal D_ZX\times\mathbf A^1\).  The restrictions of \(h\) and of the projection to \(Z\times\mathbf A^1\) induce the same map to \(Z\) on reduced schemes, so invariance under nilpotent thickenings identifies \(h_Z^*N\) with \(\operatorname{pr}_Z^*N\), naturally in \(N\), and we may view the family as an endomorphism of a single pullback.  By compatibility with composition, its fibres are the endomorphisms \(\phi_a\) constructed above.

The family \(\phi\) is constant; this is the algebraic form of homotopy invariance. Since \(Z_{\mathrm{red}}\) is a point, invariance under nilpotent thickenings identifies \(D^b_c(Z_{\acute et},A)\) with the bounded derived category of finite \(A\)-modules.  Pullback along the projection \(p:Z\times\mathbf A^1\to Z\) is fully faithful on this category. Indeed, the calculation for the affine line and the hypercohomology spectral sequence show that \(R\Gamma(\mathbf A^1,p^*B)=B\) for every such bounded complex \(B\), and full faithfulness follows by adjunction. Thus \(\phi\) is the pullback of a single endomorphism of \(N\). Its two restrictions \(\phi_0\) and \(\phi_1\) therefore coincide, and \(\phi_0=\operatorname{id}_N\).

Because \(h_0\) factors through \(Z\), we can also express \(\phi_0\) in terms of \(\tau_{X,Z}(M_X)\).  Let \(\alpha:M_Z\to N\) be the restriction of the unit of nearby cycles for the constant sheaf on \(\mathcal D_ZX\).  By naturality, and since \(h_0\) factors as \(X\to Z\hookrightarrow X\), we get

\[
\phi_0=\alpha\tau_{X,Z}(M_X),\qquad
\tau_{X,Z}(M_X)\alpha=\operatorname{id}_{M_Z}.
\]

The second equality holds because the composite is the specialization map for the constant family \(Z\times S\), followed by its fixed inverse.  Since \(\phi_0=\phi_1=\operatorname{id}_N\), the maps \(\alpha\) and \(\tau_{X,Z}(M_X)\) are inverse isomorphisms.
\end{proof}

% text: prose-E-08
\phantomsection\label{bp240}%
Near a generic point of a principal divisor in a normal scheme, the pair \((X,Z)\) is a smooth pullback of \((\mathbf A^1,V(x^m))\), so by Lemma \hyperref[bp238]{E.6} \(C_{X,Z}\) vanishes there on constant coefficients. The next lemma extends this to all sheaves that vanish on \(Z\).

% object: lemma-vanishing-generic-points-principal-divisor (blueprint/chap-Specialization-and-contracting-correspondences/sec-Specialization-to-the-normal-cone/lemma-vanishing-generic-points-principal-divisor/_section.md)
\phantomsection\label{bp241}%
\begin{bpstatementit}{Lemma~E.7}{ (Vanishing at the generic points of a principal divisor)}
Suppose that \(X\) is normal, that \(Z\subset X\) is an effective principal Cartier divisor, and that \(i^*L=0\).  Then \(C_{X,Z}(L)\) vanishes at every generic point of \(Z\).
\end{bpstatementit}

% proof: proof/lemma-vanishing-generic-points-principal-divisor
\begin{proof}[{Proof}]\phantomsection\label{bp242}%
For a constant sheaf extended by zero, a smooth map \(X\to\mathbf A^1\) near the generic point reduces the claim to Lemma \hyperref[bp238]{E.6}. Near the generic point, a general sheaf is \(j_!G\) with \(G\) locally constant. After a finite étale cover, \(G\) becomes constant, while the cokernel of the unit of the cover is again of the form \(j_!G'\). Passing to this cokernel shifts the cohomological degree by one, so an induction on the degree proves the claim for all \(G\) at once.

Since the question is local at a generic point \(\zeta\) of \(Z\), we may discard the other irreducible components of \(X\) and \(Z\).  Consider first \(L=j_!M_U\), where \(U=X\setminus Z\) and \(M\) is a finite \(A\)-module. Since \(X\) is normal, it is regular in codimension one.  After shrinking around \(\zeta\), both \(X\) and \(Z_{\mathrm{red}}\) are smooth over \(k\), and a local equation \(h\) of \(Z_{\mathrm{red}}\) defines a smooth map \(h:X\to\mathbf A^1\).  The local ring \(\mathcal O_{X,\zeta}\) is a discrete valuation ring with uniformizer \(h\), so a local equation \(g\) of \(Z\) can be written \(g=wh^m\) with \(m\geq1\) and \(w\) a unit of \(\mathcal O_{X,\zeta}\).  The function \(w\) is regular and invertible on a neighbourhood of \(\zeta\), and there the ideal of \(Z\) is \((h^m)\), that is, \(Z=h^{-1}(V(x^m))\) schematically.  By the base change isomorphism for smooth maps and Lemma \hyperref[bp238]{E.6}, we see that \(C_{X,Z}(M_X)=0\) there, and the last assertion of Lemma \hyperref[bp235]{E.5} yields the same vanishing for \(j_!M_U\).

For general \(L\), truncation triangles reduce us to a constructible sheaf \(F\) with \(F|_Z=0\).  Shrinking around \(\zeta\), we may assume that \(F=j_!G\) for a finite locally constant sheaf \(G\) on \(U\), since the closure in \(X\) of the locus in \(U\) where \(F\) is not locally constant cannot contain a generic point of the divisor \(Z\). We prove by induction on \(r\), simultaneously for all such normal pairs and all \(G\), that \(H^r(C_{X,Z}(j_!G))\) vanishes generically. As nearby cycles are left \(t\)-exact and \(i^*j_!G=0\), we have

\[
C_{X,Z}(j_!G)=
z^*\operatorname{Sp}_{X,Z}(j_!G)[1]\in D^{\geq-1}.
\]

Therefore the statement holds for every \(r<-1\), which starts the induction.

For the induction step, we compare \(F=j_!G\) with a sheaf to which the first part of the proof applies. To construct it, choose a finite étale surjection \(f_U:U'\to U\) trivializing \(G\), and let \(f:X'\to X\) be the normalization of \(X\) in \(U'\); it is finite since \(X\) is excellent, and \(X'\) is normal.  With \(Z'=f^{-1}Z\) and \(j':U'\hookrightarrow X'\), the sheaf \(f^*F=j'_!f_U^*G\) is \(j'_!\) of a sheaf that is constant on each component of \(U'\), and \(Z'\) is a principal divisor in the normal scheme \(X'\). Since \(f\) is finite, the points of \(Z'\) over a generic point of \(Z\) are generic points of \(Z'\), so the first part of the proof, applied to \((X',Z')\), and the base change isomorphism for proper maps show that \(C_{X,Z}(f_*f^*F)\) vanishes generically on \(Z\).

On a geometric stalk, the unit of the finite étale cover is the diagonal inclusion into a nonempty finite product, so it is injective and there is a short exact sequence

\[
0\longrightarrow F\longrightarrow f_*f^*F
\longrightarrow j_!G'\longrightarrow0,
\]

where \(G'\) is again finite locally constant on \(U\).  Applying \(C_{X,Z}\) and taking the long exact sequence at a generic point where the middle term vanishes, we find

\[
H^r(C_{X,Z}(F))\simeq
H^{r-1}(C_{X,Z}(j_!G')).
\]

By the induction hypothesis, the right side vanishes.
\end{proof}

% text: prose-E-09
\phantomsection\label{bp243}%
On \(D^b_{\mathrm{ctf}}(X,A)\) we define the functor \(\operatorname{Sp}_{X,Z}\) and the map (E.1) by transport of the étale constructions above along the comparison for constructible complexes with finite coefficients (Theorem \bpout{2.43}).

% object: theorem-verdier-s-theorem-specialization-zero (blueprint/chap-Specialization-and-contracting-correspondences/sec-Specialization-to-the-normal-cone/theorem-verdier-s-theorem-specialization-zero/_section.md)
\phantomsection\label{bp244}%
\begin{bpstatementit}{Theorem~E.8}{ (Verdier's theorem on specialization to the zero section)}
For every separated \(k\)-scheme \(X\) of finite type, every closed subscheme \(Z\subset X\), and \(K\in D^b_{\mathrm{ctf}}(X,A)\), the map (E.1)

\[
\tau_{X,Z}(K):z^*\operatorname{Sp}_{X,Z}(K)\longrightarrow i^*K
\]

is an isomorphism.  In particular, \(i^*K=0\) implies \(z^*\operatorname{Sp}_{X,Z}(K)=0\).
\end{bpstatementit}

% proof: proof/theorem-verdier-s-theorem-specialization-zero
\begin{proof}[{Proof}]\phantomsection\label{bp245}%
The main step is that \(C_{X,Z}(L)=0\) for every \(L\in
D^b_c(X_{\acute et},A)\), which the last paragraph transports to \(D^b_{\mathrm{ctf}}\).  We need the vanishing for all bounded constructible complexes, because the cokernels \(j_!G'\) in the dimension shifting of Lemma \hyperref[bp241]{E.7} may have infinite Tor dimension.

By Lemma \hyperref[bp235]{E.5}, we may replace \(L\) by \(j_!j^*L\) and so assume that \(i^*L=0\).  We may then replace \(X\) by its blow-up \(\pi:X'=\operatorname{Bl}_ZX\to X\) along \(Z\). The reason is that a complex that vanishes on the centre is recovered from its pullback to the blow-up, that is, the unit \(L\to R\pi_*\pi^*L\) is an isomorphism. Indeed, away from \(Z\) this holds because \(\pi\) is an isomorphism there, and on \(Z\) it holds because proper base change identifies the target with the cohomology of a fibre with zero coefficients.  Using this, the base change isomorphism for proper maps gives \(C_{X,Z}(L)\simeq C_{X,Z}(R\pi_*\pi^*L)\simeq R\pi_{Z*}C_{X',\pi^{-1}Z}(\pi^*L)\), and it suffices to treat the pair \((X',\pi^{-1}Z)\).  Since the pullback of the ideal is invertible on the blow-up, we are reduced to the case of a Cartier divisor.

To have linear projections available, we move the pair into affine space. By the base change isomorphism for open immersions the assertion is local, so we may assume that \(X\) is affine and \(Z=V(h)\) for a function \(h\).  Choosing algebra generators for \(\mathcal O(X)\) and appending \(h\) as the last coordinate gives a closed immersion

\[
\iota:X\hookrightarrow\mathbf A^n,\qquad
Z=\iota^{-1}(\mathbf A^{n-1}\times\{0\})
\]

schematically.  The base change isomorphism for proper maps identifies \(R\iota_{Z*}C_{X,Z}(L)\) with \(C_{\mathbf A^n,\mathbf A^{n-1}\times\{0\}}(R\iota_*L)\), and since direct image along a closed immersion is conservative, it suffices to treat the pair \((\mathbf A^n,\mathbf A^{n-1}\times\{0\})\). This is a principal divisor in a normal scheme, so Lemma \hyperref[bp241]{E.7} applies to it.

We argue by induction on \(n\geq1\), again assuming that \(L|_Z=0\). By Lemma \hyperref[bp241]{E.7}, the closure \(Y\) of the support of \(C_{X,Z}(L)\) is a proper closed subset of the hyperplane \(Z\); for \(n=1\) it is empty, and the assertion follows.  For \(n>1\), choose a line \(\lambda\) in the vector space \(Z\) such that the linear projection

\[
q:X\longrightarrow X'=X/\lambda\simeq\mathbf A^{n-1}
\]

is finite on \(Y\).  Such a line exists: if the projective point \(\mathbf P(\lambda)\) lies outside the part at infinity of the projective closure of \(Y\) in \(\mathbf P(Z\oplus k)\), the usual proof of linear Noether normalization shows that the projection is finite on \(Y\). Set \(Z'=Z/\lambda\), a coordinate hyperplane in \(X'\), and compactify the \(\lambda\)-direction:

\[
\bar q:\bar X=\mathbf P^1\times X'\longrightarrow X',\qquad
\bar Z=\bar q^{-1}(Z'),\qquad
v:X\hookrightarrow\bar X,\qquad \bar L=v_!L.
\]

Let \(\bar q_Z:\bar Z\to Z'\) be the restriction of \(\bar q\). By the base change isomorphism for proper maps and the induction hypothesis,

\[
R\bar q_{Z*}C_{\bar X,\bar Z}(\bar L)
\simeq C_{X',Z'}(R\bar q_*\bar L)=0.
\]

Since the cone commutes with restriction to the open subscheme \(X\) of \(\bar X\), the complex \(C_{\bar X,\bar Z}(\bar L)\) is supported on \(Y\cup(\bar Z\setminus Z)\).  This is a closed subset finite over \(Z'\). Indeed, \(Y\) is finite over \(Z'\) and therefore remains closed in \(\bar Z\), and \(\bar Z\setminus Z\) is the section at infinity of \(\bar Z\to Z'\). As in Lemma \hyperref[bp227]{E.2}, direct image along a finite map is conservative on complexes with this support. Hence \(C_{\bar X,\bar Z}(\bar L)=0\), and restriction to \(X\) gives \(C_{X,Z}(L)=0\) on the étale site.

By the comparison for finite coefficients, \(K\) and its maps descend to the classical étale constructible category, where \(C_{X,Z}\) vanishes.  The comparison is exact and fully faithful, and it identifies pullback, proper direct image and their adjunction units, so the étale map (E.1) corresponds to the map in the statement, which is therefore an isomorphism.
\end{proof}


% object: sec-Kunneth-formula-for-nearby-cycles (blueprint/chap-Specialization-and-contracting-correspondences/sec-Kunneth-formula-for-nearby-cycles/_section.md), content hidden
\subsection*{\texorpdfstring{E.2\quad Künneth formula for nearby cycles}{E.2 Künneth formula for nearby cycles}}
\phantomsection\label{bp246}%
\addcontentsline{toc}{subsection}{\texorpdfstring{E.2\quad Künneth formula for nearby cycles}{E.2 Künneth formula for nearby cycles}}
% text: prose-E-10
\phantomsection\label{bp247}%
To specialize a trace class we have to specialize the external product \(K\boxtimes \mathbb DK\), so we need the comparison map from the external product of nearby cycles to the nearby cycles of the external product to be an isomorphism. We prove this on the classical étale site by the dimension argument of [\hyperref[bp351]{Ill94}, Theorem 4.7, pp. 51–53].

In this section and the next three, \(S\) is an excellent strictly henselian trait, as announced at the start of this appendix, with separably closed residue field \(k\), generic point \(\eta\), and a chosen geometric generic point \(\bar\eta\).  Let \(p\) be the characteristic exponent of \(k\).  The coefficients lie in a finite commutative \(\ell\)-primary ring \(A\), with \(\ell\) invertible on \(S\).  On the geometric fibres we use the ordinary external Künneth formula, Proposition \bpout{8.18}. It applies over the separably closed field \(k\), since passing to an algebraic closure is a purely inseparable extension and does not change the étale topos.

We use two further facts about traits, both from [\hyperref[bp353]{SGA4½}], \emph{Théorèmes de finitude}, §§3.4 and 3.7. First, geometric nearby cycles are invariant under extension of traits, with the corresponding restriction of the inertia action (loc. cit., 3.7). Second, the calculation in loc. cit., 3.4, describes them by restriction to the generic point of a coordinate line, as follows. Let \(s'\) be a geometric generic point of the special fibre \(\mathbf A^1_k\) of \(\mathbf A^1_S\), and let \(S'=(\mathbf A^1_S)_{(s')}\) be the strict localization of \(\mathbf A^1_S\) at \(s'\).  This is a strictly henselian trait with the same uniformizer as \(S\).  With compatible geometric generic points, its inertia group fits into an exact sequence

\[
1\longrightarrow P\longrightarrow I'\longrightarrow I
\longrightarrow1,
\]

where \(P\) is pro-\(p\) (and is trivial when \(p=1\)). For \(a:W\to\mathbf A^1_S\), put \(W'=W\times_{\mathbf A^1_S}S'\), and let \(N'\) be the pullback of \(N\in D^b_{\mathrm{ctf}}(W_\eta,A)\).  There is a canonical identification

\[
(R\Psi_{W/S}N)|_{W_s\times_{\mathbf A^1_k}s'}
\simeq (R\Psi_{W'/S'}N')^P.
\tag{E.2}
\]

Here \((-)^P=R\Gamma(P,-)\) is exact on the discrete torsion coefficients prime to \(p\) under consideration.  The formula comes from factoring the map from the geometric generic fibre through the geometric generic point of \(S'\), which lies over \(\bar\eta\). The Galois group of the second step is \(P\). It is pro-\(p\) because the common uniformizer identifies the tame inertia quotients. The formula is compatible with the base change maps for pullback and with the product maps obtained from \(\bar i^*R\bar j_*\). To see this, we factor the maps of generic fibres in the same way and use that the adjunctions for direct image are compatible with composition.

The functor of \(P\)-invariants commutes with tensoring by an object on which \(P\) acts trivially: for a complex \(Q\) with continuous \(P\)-action and a complex \(B\) with trivial \(P\)-action, the projection map is an isomorphism

\[
(Q\mathbin{\otimes_A^{\mathbf L}}B)^P
\simeq Q^P\mathbin{\otimes_A^{\mathbf L}}B.
\tag{E.3}
\]

Indeed, \(P\)-invariants are exact on discrete torsion \(A\)-modules.  For a finite quotient \(G\) of \(P\), the invariants are the image of the averaging idempotent \(|G|^{-1}\sum_{g\in G}g\), since the \(p\)-power \(|G|\) is invertible in \(A\), and this idempotent commutes with tensoring by \(B\).  By continuity we pass from the finite quotients to \(P\), and resolving \(B\) by flat \(A\)-modules gives the derived identity.

For \(X,Y\) separated of finite type over \(S\), \(L\) on \(X_\eta\) and \(M\) on \(Y_\eta\), the comparison map \(\Pi_{L,M}\) below is the restriction to the special fibre of the map

\[
R\bar j_{X*}L_{\bar\eta}\boxtimes
R\bar j_{Y*}M_{\bar\eta}
\longrightarrow
R\bar j_{X\times Y,*}(L_{\bar\eta}\boxtimes M_{\bar\eta})
\]

defined by the base change maps for pullback and the lax monoidal structure of direct image.

% object: proposition-external-kunneth-formula-nearby-cycles (blueprint/chap-Specialization-and-contracting-correspondences/sec-Kunneth-formula-for-nearby-cycles/proposition-external-kunneth-formula-nearby-cycles/_section.md)
\phantomsection\label{bp248}%
\begin{bpstatementit}{Proposition~E.9}{ (External Künneth formula for nearby cycles)}
Let \(X,Y\) be separated of finite type over \(S\), and let \(L\in D^b_{\mathrm{ctf}}(X_\eta,A)\) and \(M\in D^b_{\mathrm{ctf}}(Y_\eta,A)\).  The comparison map

\[
\Pi_{L,M}:R\Psi_{X/S}L\boxtimes R\Psi_{Y/S}M
\longrightarrow
R\Psi_{(X\times_SY)/S}(L\boxtimes M)
\tag{E.4}
\]

is an isomorphism.
\end{bpstatementit}

% text: prose-E-11
\phantomsection\label{bp249}%
The proof has the same shape as that of Lemma \hyperref[bp227]{E.2}: the cone of \(\Pi_{L,M}\) has no global sections and finite support, so it vanishes. When \(X\) and \(Y\) are projective, the ordinary Künneth formula on the geometric generic fibres gives zero global sections, once we know how \(\Pi\) behaves under proper direct image. Restricting to the generic point of a coordinate line by (E.2) lowers the dimension of the generic fibre, and induction on the dimension shows that the support is finite. The next four lemmas supply these ingredients.

% object: lemma-kunneth-map-proper-direct-image (blueprint/chap-Specialization-and-contracting-correspondences/sec-Kunneth-formula-for-nearby-cycles/lemma-kunneth-map-proper-direct-image/_section.md)
\phantomsection\label{bp250}%
\begin{bpstatementit}{Lemma~E.10}{ (The Künneth map and proper direct image)}
Let \(f:X\to X_1\) and \(g:Y\to Y_1\) be proper morphisms of separated \(S\)-schemes of finite type, and let \(L\in D^b_{\mathrm{ctf}}(X_\eta,A)\) and \(M\in D^b_{\mathrm{ctf}}(Y_\eta,A)\). Identify direct images of external products by the ordinary Künneth formula, and direct images of nearby cycles by the base change isomorphism for proper maps. Then \(R(f_s\times g_s)_*\Pi_{L,M}\) becomes \(\Pi_{Rf_{\eta*}L,Rg_{\eta*}M}\).
\end{bpstatementit}

% proof: proof/lemma-kunneth-map-proper-direct-image
\begin{proof}[{Proof}]\phantomsection\label{bp251}%
Write the two product maps in terms of the tensor maps of \(R\bar j_*\), commute proper direct image past these maps on the generic fibres, and take the adjoints of the ordinary Künneth arrows.  Both adjoints are \(f_\eta^*Rf_{\eta*}L\boxtimes
g_\eta^*Rg_{\eta*}M\to L\boxtimes M\), the tensor product of the two counits. By the triangle identities of the adjunctions and their compatibility with composition, the maps themselves agree.
\end{proof}

% text: prose-E-12
\phantomsection\label{bp252}%
A complex supported on finitely many closed points is determined by its stalks at these points, and we can read off these stalks from global sections.

% object: lemma-finite-support-global-sections (blueprint/chap-Specialization-and-contracting-correspondences/sec-Kunneth-formula-for-nearby-cycles/lemma-finite-support-global-sections/_section.md)
\phantomsection\label{bp253}%
\begin{bpstatementit}{Lemma~E.11}{ (Complexes with finite support)}
Let \(Y\) be a scheme of finite type over a separably closed field, and let \(Q\in D^b_c(Y_{\acute et},A)\) be supported on a finite set of closed points. If \(R\Gamma(Y,Q)=0\), then \(Q=0\).
\end{bpstatementit}

% proof: proof/lemma-finite-support-global-sections
\begin{proof}[{Proof}]\phantomsection\label{bp254}%
Let \(i_T:T\hookrightarrow Y\) be the finite closed reduced subscheme with underlying set the support of \(Q\), so that \(Q\simeq i_{T*}i_T^*Q\). Since the residue fields of the points of \(T\) are finite and purely inseparable over the separably closed base field, \(T\) is étale locally a finite disjoint union of geometric points, and \(R\Gamma(Y,Q)=\bigoplus_{t\in T}Q_{\bar t}\). Every stalk of \(Q\) therefore vanishes, and \(Q=0\).
\end{proof}

% text: prose-E-13
\phantomsection\label{bp255}%
To restrict to coordinate lines we need coordinates in which the special fibre is quasi-finite over an affine space of the right dimension.

% object: lemma-quasi-finite-coordinates (blueprint/chap-Specialization-and-contracting-correspondences/sec-Kunneth-formula-for-nearby-cycles/lemma-quasi-finite-coordinates/_section.md)
\phantomsection\label{bp256}%
\begin{bpstatementit}{Lemma~E.12}{ (Quasi-finite coordinates)}
Let \(X\) be flat, separated and of finite type over \(S\), and let \(x\) be a closed point of \(X_s\). Put \(q=\dim_xX_s\). Then \(q\leq\dim X_{\bar\eta}\), and there are an open neighbourhood \(V\) of \(x\) in \(X\) and a quasi-finite morphism \(b:V\to\mathbf A^q_S\). In particular \(\dim V_{\bar\eta}\leq q\).
\end{bpstatementit}

% proof: proof/lemma-quasi-finite-coordinates
\begin{proof}[{Proof}]\phantomsection\label{bp257}%
The inequality \(q\leq\dim X_{\bar\eta}\) holds by flatness. Choose a system of parameters \(f_1,\ldots,f_q\) of the local ring \(\mathcal O_{X_s,x}\) of the special fibre, defined on an affine neighbourhood, and lift them to functions on the corresponding affine open of \(X\). The resulting map \(b:X\to\mathbf A^q_S\) is quasi-finite at \(x\), because its fibre through \(x\) is the fibre of \(b_s=(f_1,\ldots,f_q)\), in which \(x\) is isolated. Let \(V\) be the open locus where \(b\) is quasi-finite; as \(b_\eta\) is quasi-finite on \(V_\eta\), the generic fibre of \(V\) has dimension at most \(q\).
\end{proof}

% text: prose-E-14
\phantomsection\label{bp258}%
In the slicing step we compare \(\Pi\) with the Künneth map over the trait \(S'\) of (E.2). Let \(X,Y\) be separated of finite type over \(S\), let \(b:X\to\mathbf A^1_S\) be a morphism, and put

\[
X'=X\times_{\mathbf A^1_S}S',\qquad Y'=Y\times_S S',\qquad
L'=L|_{X'_{\eta'}},\quad M'=M|_{Y'_{\eta'}}.
\]

% object: lemma-kunneth-map-slice (blueprint/chap-Specialization-and-contracting-correspondences/sec-Kunneth-formula-for-nearby-cycles/lemma-kunneth-map-slice/_section.md)
\phantomsection\label{bp259}%
\begin{bpstatementit}{Lemma~E.13}{ (The Künneth map on a slice)}
In the situation just described, let \(L\in D^b_{\mathrm{ctf}}(X_\eta,A)\) and \(M\in D^b_{\mathrm{ctf}}(Y_\eta,A)\). Identify the restrictions over \(s'\) of the nearby cycles of \(L\), \(M\) and \(L\boxtimes M\) by

\[
\begin{aligned}
R\Psi_{X/S}L|_{s'}&\simeq(R\Psi_{X'/S'}L')^P,\\
R\Psi_{Y/S}M|_{s'}&\simeq R\Psi_{Y'/S'}M',\\
R\Psi_{(X\times Y)/S}(L\boxtimes M)|_{s'}
&\simeq
(R\Psi_{(X'\times Y')/S'}(L'\boxtimes M'))^P,
\end{aligned}
\]

given by (E.2) for \(b\), by invariance of nearby cycles under change of trait, and by (E.2) for \(b\circ\operatorname{pr}_1\). Then the restriction of \(\Pi_{L,M}\) over \(s'\) is \((\Pi_{L',M'})^P\). In particular, if \(\Pi_{L',M'}\) is an isomorphism, then so is the restriction of \(\Pi_{L,M}\) over \(s'\).
\end{bpstatementit}

% proof: proof/lemma-kunneth-map-slice
\begin{proof}[{Proof}]\phantomsection\label{bp260}%
In the middle line \(P\) acts trivially, because the action on this complex of nearby cycles is the restriction of an action of \(I\) and \(P\) is the kernel of \(I'\to I\). By (E.3), therefore, \((-)^P\) of the source of \(\Pi_{L',M'}\) is \((R\Psi_{X'/S'}L')^P\boxtimes R\Psi_{Y'/S'}M'\), the restriction of the source of \(\Pi_{L,M}\). By the compatibility of (E.2) with the product maps obtained from \(\bar i^*R\bar j_*\), the map \((\Pi_{L',M'})^P\) is the restriction of \(\Pi_{L,M}\).
\end{proof}

% proof: proof/proposition-external-kunneth-formula-nearby-cycles
\begin{proof}[{Proof of Proposition E.9}]\phantomsection\label{bp261}%
We argue by induction on \(n=\dim(X_{\bar\eta}\times_{\bar\eta}Y_{\bar\eta})\), simultaneously for all such traits.  Before slicing we need two reductions: global sections require proper families, and bounding fibre dimensions after slicing requires flatness, so we arrange both.  For properness, note first that the assertion is local on \(X_s\times_kY_s\).  We may therefore pass to affine neighbourhoods of the two factors, take projective closures over \(S\) and extend the coefficients by zero on their generic fibres. By the base change isomorphism for nearby cycles along open immersions, the restriction of the new map is the original map, so we may assume that \(X\) and \(Y\) are projective.  For flatness, the base change isomorphism for the closed immersions of the closures of their generic fibres, together with invariance under nilpotent thickenings, lets us also assume that they are reduced with dense generic fibres, hence flat over the trait.  Moreover, none of these reductions increases \(n\).

Let \(Q\) be the cone of (E.4). For proper families, global sections of nearby cycles on the special fibre are cohomology of the geometric generic fibre, so by the base change isomorphism for the proper structure maps and Lemma \hyperref[bp250]{E.10}, the map of global complexes induced by \(\Pi_{L,M}\) is the ordinary Künneth isomorphism

\[
R\Gamma(X_{\bar\eta},L)\otimes_A^{\mathbf L}
R\Gamma(Y_{\bar\eta},M)
\xrightarrow{\sim}
R\Gamma(X_{\bar\eta}\times Y_{\bar\eta},L\boxtimes M),
\]

and therefore

\[
R\Gamma(X_s\times_kY_s,Q)=0.
\tag{E.5}
\]

By Lemma \hyperref[bp253]{E.11}, \(Q=0\) as soon as \(Q\) has finite support.  When \(n=0\), the special fibres are zero-dimensional by flatness, so the support is finite. Let \(n>0\). By Lemma \hyperref[bp256]{E.12}, \(X_s\times_kY_s\) is covered by finitely many products \(V_s\times W_s\) of open subschemes with quasi-finite maps \(b:V\to\mathbf A^q_S\) and \(c:W\to\mathbf A^r_S\), where \(q+r\leq n\) and \(\dim V_{\bar\eta}\leq q\), \(\dim W_{\bar\eta}\leq r\). The restriction of \(Q\) to \(V_s\times W_s\) is the cone of \(\Pi_{L|_V,M|_W}\). If \(q+r<n\), it vanishes by the induction hypothesis. By symmetry we may therefore assume \(q+r=n\) and \(q>0\), and compose \(b\) with one of the coordinate projections \(\mathbf A^q_S\to\mathbf A^1_S\). In the notation of Lemma \hyperref[bp259]{E.13} for \(V\), \(W\) and this map, the map \(V'\to\mathbf A^q_S\times_{\mathbf A^1_S}S'=\mathbf
A^{q-1}_{S'}\) is a base change of \(b\), hence quasi-finite, and therefore

\[
\dim V'_{\bar\eta'}\leq q-1,\qquad
\dim W'_{\bar\eta'}\leq r.
\]

By the induction hypothesis \(\Pi_{L',M'}\) is an isomorphism. By Lemma \hyperref[bp259]{E.13}, the restriction of \(Q\) over \(s'\) is therefore zero. In other words, the support of \(Q\) has no point over the generic point of this coordinate line in the special fibre. Its image in that line is constructible and does not contain the generic point, so it is a finite set of closed points.

We repeat this argument for each of the \(q\) coordinates of \(b\). After interchanging the factors, we do the same for each of the \(r\) coordinates of \(c\); the analogue of Lemma \hyperref[bp259]{E.13} with the roles of \(X\) and \(Y\) exchanged holds by symmetry of \(\boxtimes\). The support of \(Q\) on \(V_s\times W_s\) is then contained in the inverse image under \(b_s\times c_s\) of a finite set. This inverse image is finite, because \(b_s\times c_s\) is quasi-finite. With finite support, (E.5) and Lemma \hyperref[bp253]{E.11} show that \(Q=0\).
\end{proof}


% object: sec-Base-change-for-nearby-cycles (blueprint/chap-Specialization-and-contracting-correspondences/sec-Base-change-for-nearby-cycles/_section.md), content hidden
\subsection*{\texorpdfstring{E.3\quad Base change for nearby cycles}{E.3 Base change for nearby cycles}}
\phantomsection\label{bp262}%
\addcontentsline{toc}{subsection}{\texorpdfstring{E.3\quad Base change for nearby cycles}{E.3 Base change for nearby cycles}}
% text: prose-E-15
\phantomsection\label{bp263}%
Let \(c\) be a correspondence over the trait \(S\) and \(u\) a cohomological correspondence on its generic fibre. Nearby cycles turn \(u\) into a map between \(\Psi\) of \(c_2^*K\) and \(\Psi\) of \(c_1^!K\). To obtain a cohomological correspondence \(u_s\) on the special fibre, we must move \(\Psi\) past \(c_2^*\) and \(c_1^!\). This needs base change maps between nearby cycles and the operations \(f^*\), \(Rf_*\), \(Rf_!\) and \(f^!\). We make these maps explicit here. We also prove the identities between them that Proposition \hyperref[bp289]{E.22} uses to show that the trace of \(u_s\) is the specialization of the trace of \(u\).

We write \(\Psi_X=R\Psi_{X/S}\) for a separated \(S\)-scheme \(X\) of finite type. For a morphism \(f:V\to W\) of separated \(S\)-schemes of finite type, put

\[
\begin{array}{rcl}
\alpha_f:f_s^*\Psi_W&\longrightarrow&\Psi_Vf_\eta^*,\\
\chi_f:\Psi_WRf_{\eta*}&\longrightarrow&Rf_{s*}\Psi_V,\\
\nu_f:Rf_{s!}\Psi_V&\longrightarrow&\Psi_WRf_{\eta!},\\
\beta_f:\Psi_Vf_\eta^!&\longrightarrow&f_s^!\Psi_W.
\end{array}
\]

The map \(\alpha_f\) is the base change map for pullback of \(\bar i^*R\bar j_*\), and \(\chi_f\) is its adjoint mate for \(f_s^*\dashv Rf_{s*}\) and \(f_\eta^*\dashv Rf_{\eta*}\).

To define \(\nu_f\), we choose a relative Nagata compactification \(f=ph\), with \(h:V\hookrightarrow\overline V\) open and \(p\) proper. We define \(\nu_h:h_{s!}\Psi_VM\to \Psi_{\overline V}h_{\eta!}M\) using the adjunction for the open immersion, by requiring that its adjoint be

\[
\Psi_VM=\Psi_Vh_\eta^*h_{\eta!}M
\xrightarrow{\alpha_h^{-1}}h_s^*\Psi_{\overline V}h_{\eta!}M.
\]

Since \(p\) is proper, \(\chi_p\) is invertible, and we put \(\nu_f=\chi_p^{-1}\,Rp_{s*}(\nu_h)\), using \(Rf_{s!}=Rp_{s*}h_{s!}\) and \(Rf_{\eta!}=Rp_{\eta*}h_{\eta!}\). The map \(\nu_f\) does not depend on the compactification. Indeed, if \(q\) is a common proper refinement of two compactifications, proper base change identifies \(Rq_*h'_!\) with \(h_!\) on both fibres. After taking adjoints for the open immersions, compatibility with this identification reduces to the composition rule for \(\alpha\) and the triangle identity of the ordinary adjunction. Common refinements also show that \(\nu\) is compatible with composition.

We can also describe \(\nu_f\) on the total families. Write \(i_V\) and \(j_V\) for the inclusions \(\bar i\) and \(\bar j\) of the special and geometric generic fibres of \(V\), so that \(\Psi_V=i_V^*Rj_{V*}\). By proper base change, \(i_W^*Rf_!\simeq Rf_{s!}i_V^*\) and \(j_W^*Rf_!\simeq Rf_{\eta!}j_V^*\). Consider the map \(Rf_!Rj_{V*}\to Rj_{W*}Rf_{\eta!}\) whose \(j_W^*\)-adjoint is

\[
j_W^*Rf_!Rj_{V*}\simeq Rf_{\eta!}j_V^*Rj_{V*}
\xrightarrow{Rf_{\eta!}\varepsilon_{j_V}}Rf_{\eta!}
\]

Its restriction by \(i_W^*\) is a map \(Rf_{s!}\Psi_V\to \Psi_WRf_{\eta!}\), and we claim that this map is \(\nu_f\). As both constructions are compatible with composition, it suffices to check this for \(p\) and for \(h\). For proper \(p\) the restricted map is \(\chi_p^{-1}\), by proper base change for \(i_W^*\). For the open immersion \(h\) we have \(h_\eta^*h_{\eta!}=1\), and the \(h_s^*\)-adjoint of the restricted map is \(\alpha_h^{-1}\), as in the definition of \(\nu_h\).

The map \(\beta_f(M)\) is the adjoint, for \(Rf_{s!}\dashv f_s^!\), of

\[
Rf_{s!}\Psi_Vf_\eta^!M\xrightarrow{\nu_f}
\Psi_WRf_{\eta!}f_\eta^!M
\xrightarrow{\Psi_W\varepsilon_{f_\eta}}\Psi_WM.
\tag{E.6}
\]

For proper \(f\), the maps \(\chi_f\) and \(\nu_f\) are inverse isomorphisms after the identification \(Rf_*=Rf_!\).

The lax monoidal map \(\mu_Y:\Psi_YM\otimes^{\mathbf L} \Psi_YN\to \Psi_Y(M\otimes^{\mathbf L} N)\) of \(R\bar j_*\) is adjoint to the tensor product of the two counits \(j^*Rj_*\to1\). It is associative, symmetric and unital, with unit \(\sigma_Y:A_{Y_s}\to \Psi_YA_{Y_\eta}\), and \(\alpha_f\) is monoidal. For \(f:V\to W\) and \(g:W\to Y\), we have

\[
\alpha_{gf}=\alpha_f(g_\eta^*)f_s^*\alpha_g,
\qquad \sigma_V=\alpha_f(A)f_s^*\sigma_W,
\qquad\mu(\sigma\otimes1)=1.
\tag{E.7}
\]

For the structure morphism \(a_Y:Y\to S\), set \(\theta_Y=\beta_{a_Y}(A):\Psi_Y\omega_{Y_\eta}\to\omega_{Y_s}\), using \(\Psi_SA=A\).

The proof of Proposition \hyperref[bp289]{E.22} moves the base change maps past the projection formula, composition, external duality and the cartesian square defining the fixed point scheme. The following four lemmas give the identities it needs. Each is proved by taking adjoints until both sides become the same composite of base change maps and counits, and then using the composition laws for the base change maps and for the counits. For a morphism \(f\) on either fibre, write \(t_f:f^!M\otimes^{\mathbf L} f^*N\to f^!(M\otimes^{\mathbf L} N)\) for the projection-formula mate of the exceptional counit.

% object: lemma-compatibilities-base-change-maps (blueprint/chap-Specialization-and-contracting-correspondences/sec-Base-change-for-nearby-cycles/lemma-compatibilities-base-change-maps/_section.md)
\phantomsection\label{bp264}%
\begin{bpstatementit}{Lemma~E.14}{ (Base change maps and the projection formula)}
Let \(f:V\to W\) be a morphism of separated \(S\)-schemes of finite type, and let \(M,N\in D^b_c(W_\eta,A)\). Then the two maps \(\Psi_Vf_\eta^!M\otimes^{\mathbf L} f_s^*\Psi_WN\to f_s^!\Psi_W(M\otimes^{\mathbf L} N)\)

\[
\beta_f \Psi_V(t_{f_\eta})\mu_V(1\otimes\alpha_f)
=f_s^!\mu_W\ t_{f_s}(\beta_f\otimes1)
\tag{E.8}
\]

are equal.
\end{bpstatementit}

% proof: proof/lemma-compatibilities-base-change-maps
\begin{proof}[{Proof}]\phantomsection\label{bp265}%
The identity follows, after taking \(Rf_{s!}\)-adjoints, from the compatibility of \(\nu_f\) with the projection formula: for \(L\) on \(V_\eta\) and \(N\) on \(W_\eta\), the composite

\[
Rf_{s!}(\Psi_VL\otimes^{\mathbf L} f_s^*\Psi_WN)\xrightarrow{\mu_V(1\otimes\alpha_f)}
Rf_{s!}\Psi_V(L\otimes^{\mathbf L} f_\eta^*N)\xrightarrow{\nu_f}
\Psi_WRf_{\eta!}(L\otimes^{\mathbf L} f_\eta^*N)\simeq \Psi_W(Rf_{\eta!}L\otimes^{\mathbf L} N)
\]

equals the projection formula on the special fibre followed by \(\mu_W(\nu_f\otimes1)\). To check this, we use the description of \(\nu_f\) on the total families and compare \(j_W^*\)-adjoints. On both sides the adjoint is the counit \(j_V^*Rj_{V*}\to1\) tensored with the identity, followed by the projection formula on the generic fibre.

Taking \(Rf_{s!}\)-adjoints in (E.8) and applying the projection formula, the left side becomes the following composite by this compatibility, and the right side becomes it by naturality of \(\mu_W\):

\[
\begin{split}
Rf_{s!}\Psi_Vf_\eta^!M\otimes^{\mathbf L} \Psi_WN
&\xrightarrow{\nu_f\otimes1}
\Psi_W(Rf_{\eta!}f_\eta^!M)\otimes^{\mathbf L} \Psi_WN\\
&\xrightarrow{\mu_W}\Psi_W(Rf_{\eta!}f_\eta^!M\otimes^{\mathbf L} N)
\xrightarrow{\Psi_W(\varepsilon_{f_\eta}\otimes1)}\Psi_W(M\otimes^{\mathbf L} N).
\end{split}
\]

Since the adjoints agree, so do the maps.
\end{proof}

% object: lemma-composition-exceptional-base-change-maps (blueprint/chap-Specialization-and-contracting-correspondences/sec-Base-change-for-nearby-cycles/lemma-composition-exceptional-base-change-maps/_section.md)
\phantomsection\label{bp266}%
\begin{bpstatementit}{Lemma~E.15}{ (Composition of the exceptional base change maps)}
Let \(f:V\to W\) and \(g:W\to Y\) be morphisms of separated \(S\)-schemes of finite type. Then

\[
\beta_{gf}=f_s^!\beta_g\ \beta_f(g_\eta^!),
\qquad\theta_V=f_s^!\theta_W\ \beta_f(\omega_{W_\eta}).
\tag{E.9}
\]
\end{bpstatementit}

% proof: proof/lemma-composition-exceptional-base-change-maps
\begin{proof}[{Proof}]\phantomsection\label{bp267}%
The first identity follows by taking adjoints for the adjunctions \(Rh_{s!}\dashv h_s^!\) (\(h=f,g,gf\)) and using the compatibility of \(\nu\) with composition together with the composition of counits \(\varepsilon_{gf}=\varepsilon_g\,Rg_!\varepsilon_f\). The second is the case \(g=a_W\), applied to \(A\), of the first, since \(a_V=a_Wf\).
\end{proof}

% text: prose-E-16
\phantomsection\label{bp268}%
Let \(p_1,p_2\) be the projections of \(X\times_SX\). On either fibre, let \(\xi:p_2^*\omega_X\to p_1^!A_X\) be the isomorphism of Lemma \bpout{11.7} for \(K=A_X\), and for a complex \(M\) put \(\kappa_M=t_{p_1}(\xi\otimes1):p_2^*\omega_X\otimes^{\mathbf L} p_1^*M\to p_1^!M\), where \(t_{p_1}:p_1^!A_X\otimes^{\mathbf L} p_1^*M\to p_1^!M\) is the projection-formula mate of the exceptional counit.

The next identity says that, under the base change maps, the map \(\kappa_M\) of external duality on the generic fibre corresponds to \(\kappa_{\Psi_XM}\) on the special fibre, up to \(\theta_X\).

% object: lemma-base-change-maps-external-duality (blueprint/chap-Specialization-and-contracting-correspondences/sec-Base-change-for-nearby-cycles/lemma-base-change-maps-external-duality/_section.md)
\phantomsection\label{bp269}%
\begin{bpstatementit}{Lemma~E.16}{ (Base change maps and external duality)}
Let \(X\) be a separated \(S\)-scheme of finite type with projections \(p_1,p_2\) of \(X\times_SX\), and let \(M\in D^b_c(X_\eta,A)\). Then the two maps \(p_{2s}^*\Psi_X\omega_{X_\eta}\otimes^{\mathbf L} p_{1s}^*\Psi_XM\to p_{1s}^!\Psi_XM\)

\[
\beta_{p_1}\Psi(\kappa_M)\mu
(\alpha_{p_2}\otimes\alpha_{p_1})
=\kappa_{\Psi_XM}(p_{2s}^*\theta_X\otimes1)
\tag{E.10}
\]

are equal.
\end{bpstatementit}

% proof: proof/lemma-base-change-maps-external-duality
\begin{proof}[{Proof}]\phantomsection\label{bp270}%
Since \(\kappa_M=t_{p_1}(\xi\otimes1)\), the main point is the case of the constant complex,

\[
\beta_{p_1}\Psi(\xi)\alpha_{p_2}
=p_{1s}^!\sigma_X\,\xi_s\,p_{2s}^*\theta_X,
\tag{C}
\]

and the general case follows from (C) by Lemma \hyperref[bp264]{E.14}. Since the proof uses proper base change, we prove (C) for a proper family and then extend it by compactification.

Suppose first that \(a=a_X:X\to S\) is proper. On the fibre over \(\eta\), let

\[
Q_\eta(L):a_\eta^*Ra_{\eta*}L\xrightarrow{\sim}Rp_{1\eta*}p_{2\eta}^*L
\]

be the proper base change isomorphism, and define \(Q_s\) in the same way. On each fibre, the \(Rp_{1!}\)-adjoint of \(\xi\) is \(Q(\omega_X)^{-1}\) followed by \(a^*\varepsilon_a\). We compare these adjoints using the identity

\[
\nu_{p_1}Rp_{1s*}(\alpha_{p_2})Q_s(\Psi_XL)
=\Psi_X(Q_\eta(L))\alpha_a(Ra_{\eta*}L)\,a_s^*\nu_a(L).
\]

To prove it, we take adjoints for \(p_{1s}^*\dashv Rp_{1s*}\). By the composition rule for \(\alpha\) and \(ap_1=ap_2\), both adjoints are the composite of \(a_s^*\nu_a\), the base change map \(\alpha_{ap_1}=\alpha_{ap_2}\), and the pullback by \(p_{2\eta}\) of the counit \(a_\eta^*Ra_{\eta*}L\to L\). Substituting this identity into the \(Rp_{1s!}\)-adjoint of the left side of (C), and using naturality of \(\alpha_a\) with respect to \(\varepsilon_{a_\eta}\) and \(\alpha_a(A)=\sigma_X\), we find that this adjoint is the composite

\[
a_s^*Ra_{s!}\Psi_X\omega_{X_\eta}
\xrightarrow{a_s^*\nu_a}
a_s^*Ra_{\eta!}\omega_{X_\eta}
\xrightarrow{a_s^*\varepsilon_{a_\eta}}A_{X_s}
\xrightarrow{\sigma_X}\Psi_XA_{X_\eta}.
\]

By (E.6) for \(a\), this composite equals \(\sigma_X a_s^*(\varepsilon_{a_s}Ra_{s!}\theta_X)\), which is the \(Rp_{1s!}\)-adjoint of the right side. This proves (C) for proper \(X\).

If \(X\) is not proper over \(S\), we choose an open immersion \(j:X\hookrightarrow\overline X\) into a proper \(S\)-scheme and restrict both factors of the product from \(\overline X\) to \(X\). Both sides of (C) are compatible with this restriction. For the maps \(\xi\) this follows from the composition rule for exceptional inverse images. For the base change maps it holds because \(j^!=j^*\), the map \(\alpha_j\) is invertible by smooth base change, and \(\beta_j=\alpha_j^{-1}\). So (C) on \(\overline X\) implies it on \(X\). Applying Lemma \hyperref[bp264]{E.14} to \(p_1\) and using \(\mu(\sigma\otimes1)=1\), we deduce (E.10).
\end{proof}

% text: prose-E-17
\phantomsection\label{bp271}%
For the last identity, let \(c:C\to X\times_SX\) be a correspondence of separated \(S\)-schemes of finite type, with fixed scheme \(F\), closed immersion \(i:F\hookrightarrow C\), map \(d:F\to X\), and closed diagonal \(\Delta:X\hookrightarrow X\times_SX\), so that \(F=C\times_{X\times X}X\). Let \(\zeta:c^!R\Delta_*\xrightarrow\sim Ri_*d^!\) be the base change isomorphism for exceptional inverse image and proper direct image in this cartesian square, on either fibre.

% object: lemma-base-change-maps-fixed-point-scheme (blueprint/chap-Specialization-and-contracting-correspondences/sec-Base-change-for-nearby-cycles/lemma-base-change-maps-fixed-point-scheme/_section.md)
\phantomsection\label{bp272}%
\begin{bpstatementit}{Lemma~E.17}{ (Base change maps and the fixed point scheme)}
In the situation just described,

\[
Ri_{s*}\beta_d\ \chi_i\ \Psi_C(\zeta_\eta)
=\zeta_s\ c_s^!(\chi_\Delta)\ \beta_c(R\Delta_{\eta*}).
\tag{E.11}
\]
\end{bpstatementit}

% proof: proof/lemma-base-change-maps-fixed-point-scheme
\begin{proof}[{Proof}]\phantomsection\label{bp273}%
The inverse of the base change map, \(Ri_*d^!\to c^!R\Delta_*\), is adjoint to the explicit map

\[
Rc_!Ri_*d^!=R\Delta_*Rd_!d^!
\xrightarrow{R\Delta_*\varepsilon_d}R\Delta_*.
\]

This identification uses that \(i\) and \(\Delta\) are proper. Substituting it in both sides, taking \(Rc_{s!}\)-adjoints, and replacing the comparison maps \(\chi_i,\chi_\Delta\) for proper maps by the inverses of \(\nu_i,\nu_\Delta\), we find by the composition rules for the base change maps and for the counits that both sides are the base change map for the composite \(ci=\Delta d\), followed by \(\Psi_X\varepsilon_d\).
\end{proof}


% object: sec-Duality-for-nearby-cycles (blueprint/chap-Specialization-and-contracting-correspondences/sec-Duality-for-nearby-cycles/_section.md), content hidden
\subsection*{\texorpdfstring{E.4\quad Duality for nearby cycles}{E.4 Duality for nearby cycles}}
\phantomsection\label{bp274}%
\addcontentsline{toc}{subsection}{\texorpdfstring{E.4\quad Duality for nearby cycles}{E.4 Duality for nearby cycles}}
% text: prose-E-18
\phantomsection\label{bp275}%
Since the second factor of \(K\boxtimes \mathbb DK\) is a Verdier dual, specializing a trace class also requires nearby cycles to commute with Verdier duality, by an explicit map \(\delta_{X,M}\) defined below [\hyperref[bp351]{Ill94}, Theorem 4.2].

The trait \(S\) and the coefficient ring \(A\) are those of Proposition \hyperref[bp248]{E.9}, and \(\mathbb D_{X_{\bar\eta}}\) and \(\mathbb D_{X_s}\) are the Verdier duals relative to the geometric generic point \(\bar\eta\) and the closed point \(s\) of \(S\). We work on the classical étale site and transport the result to the pro-étale site at the end.

We use the base change maps \(\chi_f\), \(\nu_f\), \(\beta_f\), the lax monoidal map \(\mu_X\) and the map \(\theta_X:\Psi_X\omega_{X_\eta}\to\omega_{X_s}\) of the preceding section.  Recall that \(\beta_f\) is the adjoint, for \(Rf_{s!}\dashv f_s^!\), of (E.6) and that \(\nu_f=\chi_f^{-1}\) for proper \(f\). For \(M\in D^b_{\mathrm{ctf}}(X_\eta,A)\), define the morphism \(\delta_{X,M}:\Psi_X\mathbb D_{X_\eta}M\to\mathbb D_{X_s}(\Psi_XM)\) as the tensor–Hom adjoint of the pairing

\[
\Psi_XM\otimes \Psi_X\mathbb D_{X_\eta}M
\xrightarrow{\mu_X}\Psi_X(M\otimes \mathbb D_{X_\eta}M)
\xrightarrow{\Psi_X\operatorname{ev}}\Psi_X\omega_{X_\eta}
\xrightarrow{\theta_X}\omega_{X_s}.
\tag{E.12}
\]

% object: proposition-nearby-cycles-commute-verdier-duality (blueprint/chap-Specialization-and-contracting-correspondences/sec-Duality-for-nearby-cycles/proposition-nearby-cycles-commute-verdier-duality/_section.md)
\phantomsection\label{bp276}%
\begin{bpstatementit}{Proposition~E.18}{ (Nearby cycles commute with Verdier duality)}
For every separated \(S\)-scheme \(X\) of finite type and \(M\in D^b_{\mathrm{ctf}}(X_\eta,A)\), the map

\[
\delta_{X,M}:\Psi_X\mathbb D_{X_\eta}M
\longrightarrow\mathbb D_{X_s}(\Psi_XM)
\]

defined by (E.12) is an isomorphism.
\end{bpstatementit}

% text: prose-E-19
\phantomsection\label{bp277}%
The argument is the one used for Proposition \hyperref[bp248]{E.9}. As there, on a projective \(X\) global duality shows that the cone of \(\delta\) has zero cohomology on \(X_s\). Restriction to the generic point of a coordinate line (E.2) and induction on dimension show that the cone has finite support. For this we need two compatibilities of \(\delta\): with proper direct image, and with restriction to the generic point of a coordinate line. In the second we keep track of the twist \(A(1)[2]\) on both fibres.

% object: lemma-nearby-duality-proper-direct-image (blueprint/chap-Specialization-and-contracting-correspondences/sec-Duality-for-nearby-cycles/lemma-nearby-duality-proper-direct-image/_section.md)
\phantomsection\label{bp278}%
\begin{bpstatementit}{Lemma~E.19}{ (Duality for nearby cycles and proper direct image)}
Let \(f:X\to Y\) be a proper morphism of separated \(S\)-schemes of finite type, and let \(M\in D^b_{\mathrm{ctf}}(X_\eta,A)\). Under the isomorphism \(\chi_f\) and the Verdier duality isomorphisms \(Rf_*\mathbb D\simeq\mathbb DRf_*\) of Chapter \bpout{10} on both fibres,

\[
Rf_{s*}\delta_{X,M}\quad\text{is identified with}\quad
\delta_{Y,Rf_{\eta*}M}.
\tag{E.13}
\]

For an open immersion \(j:U\to X\), the restriction \(j_s^*\delta_{X,M}\) is \(\delta_{U,j_\eta^*M}\) under the base change isomorphism \(\alpha_j\). If \(X\) is proper over \(S\), then under global Verdier duality on both fibres

\[
R\Gamma(X_s,\delta_{X,M})
\quad\text{is identified with the identity of}\quad
R\operatorname{Hom}_A(R\Gamma(X_{\bar\eta},M),A).
\tag{E.14}
\]
\end{bpstatementit}

% proof: proof/lemma-nearby-duality-proper-direct-image
\begin{proof}[{Proof}]\phantomsection\label{bp279}%
Under Verdier duality for the proper map \(f_\eta\), evaluation on \(Rf_{\eta*}M\) is the composite of the lax monoidal map of \(Rf_{\eta*}\), the map \(Rf_{\eta*}\operatorname{ev}_M\), and the trace \(Rf_{\eta*}\omega_{X_\eta}\to\omega_{Y_\eta}\). Hence, after taking tensor–Hom adjoints, both sides of (E.13) are the pairing

\[
\begin{split}
Rf_{s*}\Psi_XM\otimes Rf_{s*}\Psi_X\mathbb D M
&\xrightarrow{\chi_f^{-1}\otimes\chi_f^{-1}}
\Psi_YRf_{\eta*}M\otimes \Psi_YRf_{\eta*}\mathbb DM\xrightarrow{\mu_Y}
\Psi_Y(Rf_{\eta*}M\otimes Rf_{\eta*}\mathbb DM)\\
&\longrightarrow \Psi_YRf_{\eta*}(M\otimes \mathbb DM)
\longrightarrow \Psi_YRf_{\eta*}\omega_{X_\eta}
\longrightarrow \Psi_Y\omega_{Y_\eta}
\xrightarrow{\theta_Y}\omega_{Y_s}.
\end{split}
\]

To identify the pairing obtained from \(Rf_{s*}\delta_{X,M}\) with this one, we move the lax monoidal maps past the base change maps for proper morphisms, and we use the identity \(\theta_X=f_s^!\theta_Y\circ\beta_f(\omega_{Y_\eta})\) of (E.9). The lax monoidal maps are compatible with the base change maps because, after taking ordinary adjoints, both become the tensor product of the two counits for the open immersions of the generic fibres. This proves (E.13).

For an open immersion \(j\), the map \(\alpha_j\) is an isomorphism, \(j^!=j^*\) and \(\beta_j=\alpha_j^{-1}\). The argument above therefore applies with \(\alpha_j\) in place of \(\chi_f^{-1}\).

For (E.14) we apply (E.13) to the structure morphism \(a\) of \(X\). On the trait, \(\Psi_S\) takes the stalk at the geometric generic point, and the pairing (E.12) is ordinary evaluation. So \(\delta_{S,Ra_{\eta*}M}\) is the identity of \(R\operatorname{Hom}_A(R\Gamma(X_{\bar\eta},M),A)\). Global Verdier duality on both fibres now gives (E.14).
\end{proof}

% text: prose-E-20
\phantomsection\label{bp280}%
For the slice we need that taking invariants under the pro-\(p\) group \(P\) of (E.2) commutes with duality.

% object: lemma-duality-invariants-pro-p (blueprint/chap-Specialization-and-contracting-correspondences/sec-Duality-for-nearby-cycles/lemma-duality-invariants-pro-p/_section.md)
\phantomsection\label{bp281}%
\begin{bpstatementit}{Lemma~E.20}{ (Duality and invariants under $P$)}
Let \(Y\) be a scheme of finite type over a separably closed field, and let \(L\in D^b_{\mathrm{ctf}}(Y_{\acute et},A)\) carry a continuous action of \(P\). There is a natural isomorphism

\[
(\mathbb DL)^P\xrightarrow{\sim}\mathbb D(L^P),
\tag{E.15}
\]

under which the evaluation pairing of \(L^P\) is the restriction of the evaluation pairing of \(L\) to the direct summands \(L^P\) and \((\mathbb DL)^P\).
\end{bpstatementit}

% proof: proof/lemma-duality-invariants-pro-p
\begin{proof}[{Proof}]\phantomsection\label{bp282}%
Let \(G\) be a finite quotient of \(P\). Since \(G\) is a \(p\)-group, the order \(|G|\) is invertible in \(A\), and the invariants are the image of the idempotent \(e=|G|^{-1}\sum_{g\in G}g\).  Verdier duality sends this idempotent to the same averaging idempotent for the contragredient action, since inversion permutes the summands. Hence the image of \(\mathbb D(e)\) is the dual of \(\operatorname{im}e\), and the evaluation pairing restricts as stated.  As invariants are exact, continuity gives the assertion for the pro-\(p\) group \(P\).
\end{proof}

% text: prose-E-21
\phantomsection\label{bp283}%
We fix the notation for the slice. Let \(b:X\to\mathbf A^1_S\) be a coordinate map and let \(S'=(\mathbf A^1_S)_{(s')}\) be the trait at the generic point of the special fibre used in (E.2).  Set \(X'=X\times_{\mathbf A^1_S}S'\), \(M'=M|_{X'_{\bar\eta'}}\), and write \(g_s,g_\eta\) for the two pullbacks to the geometric generic point of the coordinate line.  Restriction to the geometric generic point of the coordinate line is compatible with internal Hom, and Lemma \bpout{10.6} and Lemma \bpout{10.7}, applied on each geometric fibre, give, locally on \(X\), isomorphisms

\[
g_\eta^*\mathbb D_{X_{\bar\eta}}M\simeq \mathbb D_{X'_{\bar\eta'}}M'\otimes A(1)[2],
\qquad
g_s^*\mathbb D_{X_s}L\simeq \mathbb D_{X'_{s'}}(g_s^*L)\otimes A(1)[2]
\tag{E.16}
\]

for \(L\in D^b_{\mathrm{ctf}}(X_s,A)\). The twist comes from the dualizing complexes. For a closed embedding over \(\mathbf A^1\) into \(\mathbf A^r_{\mathbf A^1}\), smooth purity gives \(A(r+1)[2r+2]\) before generic restriction and \(A(r)[2r]\) after it. Moreover, the comparison of the complexes with closed support for the embedding preserves evaluation. Thus on \(\mathbf A^r_{\mathbf A^1_S}\) the identifications of the dualizing complexes before and after slicing differ by \(A(1)[2]\) on both fibres. Finally, the group \(P\) acts trivially on \(A(1)[2]\), because roots of unity of order prime to the residue characteristic lift over a strictly henselian trait. So the twist commutes with taking \(P\)-invariants.

% object: lemma-nearby-duality-slice (blueprint/chap-Specialization-and-contracting-correspondences/sec-Duality-for-nearby-cycles/lemma-nearby-duality-slice/_section.md)
\phantomsection\label{bp284}%
\begin{bpstatementit}{Lemma~E.21}{ (Duality for nearby cycles on a slice)}
In the situation just described, suppose that \(X\) is affine, and let \(i:X\hookrightarrow\mathbf A^r_{\mathbf A^1_S}\) be a closed immersion over \(\mathbf A^1_S\); we use the identifications (E.16) defined by \(i\). Under the identifications (E.2), (E.16) and (E.15), we have

\[
g_s^*\delta_{X,M}
\ =\ (\delta_{X',M'})^P\otimes A(1)[2].
\tag{E.17}
\]
\end{bpstatementit}

% proof: proof/lemma-nearby-duality-slice
\begin{proof}[{Proof}]\phantomsection\label{bp285}%
The map \(\delta\) is the adjoint of the pairing (E.12). By Lemma \hyperref[bp281]{E.20} it therefore suffices to compare the three arrows of this pairing with the corresponding arrows on the slice, restricted to \(P\)-invariants.  Since the slicing formula comes from factoring the inclusion of the geometric generic point and composing the adjunctions for the direct images, the lax monoidal map and the evaluation correspond under it to the maps on the slice, restricted to \(P\)-invariants.

For the last arrow \(\theta_X\), we claim that under (E.2) and (E.16), the map \(g_s^*\theta_X\) corresponds to \((\theta_{X'})^P\otimes A(1)[2]\). We check this on a smooth ambient space, where \(\theta\) is the identity by smooth purity, and descend along the closed immersion \(i:X\hookrightarrow\mathbf A^r_{\mathbf A^1_S}\) of the statement, which we abbreviate to \(i:X\hookrightarrow\mathbf A^r\). First, on \(\mathbf A^r\) and on its slice, smooth purity and the base change isomorphism of nearby cycles for smooth pullback identify \(\theta\) with the identity of the dualizing complex of a smooth scheme, and these identifications differ by \(A(1)[2]\). To descend to \(X\), we use the identity (E.9) for \(i\), which gives \(\theta_X=i_s^!\theta_{\mathbf A^r}\circ\beta_i(\omega_{\mathbf A^r_\eta})\). By (E.6), \(\theta_X\) is then the adjoint, for \(Ri_{s*}\dashv i_s^!\), of \(\nu_i\), followed by \(\Psi\) of the exceptional counit \(Ri_{\eta*}i_\eta^!\to1\) and by \(\theta_{\mathbf A^r}\). Now the base change map \(\nu_i\) for the proper map \(i\) and the factorization used for slicing commute with this counit, because their ordinary adjoints are the same composite of the counit for the open immersion of the generic fibre and the counit for \(i\). Since generic restriction is compatible with \(i^!\) and with evaluation, the claim for \(\mathbf A^r\) gives it for \(X\) before taking invariants.

It remains to pass to \(P\)-invariants. For a finite quotient \(G\) of \(P\), the additive functor \(i_{s'}^!\) commutes with the action of \(G\), and so it preserves the image of the averaging idempotent. Thus \(i_{s'}^!\) commutes with \((-)^P\), and the same descent works after taking \(P\)-invariants. Since the shift \([2]\) is even, no additional Koszul sign appears.
\end{proof}

% proof: proof/proposition-nearby-cycles-commute-verdier-duality
\begin{proof}[{Proof of Proposition E.18}]\phantomsection\label{bp286}%
We argue by induction on \(n=\dim X_{\bar\eta}\), simultaneously for all traits \(S\) as above, and show that the cone of \(\delta_{X,M}\) has zero global sections and finite support.

\emph{Reductions.} We may assume that \(X\) is projective, reduced, and flat over \(S\), by reductions none of which increases \(n\). The assertion is local and, by Lemma \hyperref[bp278]{E.19}, compatible with restriction to opens. So we may pass to a projective compactification and extend \(M\) by zero on its generic fibre. Next we use the compatibility with proper direct image (E.13) for the closure of that generic fibre, together with invariance under nilpotent thickenings. This lets us assume that \(X\) is reduced and flat over \(S\).

Put \(Q=\Cone(\delta_{X,M})\). It is bounded and constructible by the finiteness theorem for nearby cycles and ordinary Verdier duality. By (E.14), \(R\Gamma(X_s,Q)=0\).

\emph{Finite support.} By Lemma \hyperref[bp253]{E.11}, it remains to show that the support of \(Q\) is finite. By Lemma \hyperref[bp256]{E.12}, \(X_s\) is covered by finitely many opens \(V_s\) with quasi-finite coordinate maps \(V\to\mathbf A^q_S\), \(q\leq n\), and \(\dim V_{\bar\eta}\leq q\). As the vanishing of \(Q\) is local, we may replace each \(V\) by finitely many affine opens covering it, with the restricted coordinate maps. Then for each coordinate map \(V\to\mathbf A^1_S\) there is a closed immersion of \(V\) into some \(\mathbf A^r_{\mathbf A^1_S}\) over \(\mathbf A^1_S\), as required in Lemma \hyperref[bp284]{E.21}.

We claim that \(Q\) vanishes over the generic point of each coordinate line in the special fibre. The slice \(V'\) has generic fibre of dimension at most \(q-1<n\). Hence the claim follows from the induction hypothesis and Lemma \hyperref[bp284]{E.21}. The support of \(Q\) therefore lies over finitely many closed points of each coordinate line. Since the coordinate map is quasi-finite, the support is finite. This argument also covers \(n=0\), where \(q=0\) and \(V_s\) is finite.

\emph{Pro-étale formulation.} The comparison for constructible complexes with finite coefficients transports both the pairing and this isomorphism to the small pro-étale site.
\end{proof}


% object: sec-Specialization-of-the-trace-class (blueprint/chap-Specialization-and-contracting-correspondences/sec-Specialization-of-the-trace-class/_section.md), content hidden
\subsection*{\texorpdfstring{E.5\quad Specialization of the trace class}{E.5 Specialization of the trace class}}
\phantomsection\label{bp287}%
\addcontentsline{toc}{subsection}{\texorpdfstring{E.5\quad Specialization of the trace class}{E.5 Specialization of the trace class}}
% text: prose-E-22
\phantomsection\label{bp288}%
In a continuous family of maps of a compact manifold, the Lefschetz number does not change. The algebraic counterpart is the following proposition: the trace class of the specialized correspondence \(u_s\) is the specialization \(\operatorname{sp}_F\) of the trace class of \(u\). Here \(\operatorname{sp}_F\) is a map of sections of dualizing complexes on the fixed scheme. A trace class is built from three pieces: the cohomological correspondence \(u\), the external product \(K\boxtimes\mathbb DK\) in which \(u\) becomes a section, and the evaluation map along the diagonal to the dualizing complex. Nearby cycles commute with each piece. They commute with external products by Proposition \hyperref[bp248]{E.9}, with Verdier duality by Proposition \hyperref[bp276]{E.18}, and with the four operations through the base change maps of Section \hyperref[bp262]{E.3}. We prove the proposition by putting these compatibilities together. Theorem \hyperref[bp311]{E.30} uses it to transfer the computation to the normal cone. For this it also needs Corollary \hyperref[bp298]{E.25}: \(\operatorname{sp}_F\) is the identity when the fixed scheme is, up to nilpotents, a constant family \(Y\times S\).

% object: proposition-specialization-trace-class (blueprint/chap-Specialization-and-contracting-correspondences/sec-Specialization-of-the-trace-class/proposition-specialization-trace-class/_section.md)
\phantomsection\label{bp289}%
\begin{bpstatementit}{Proposition~E.22}{ (Specialization of the trace class)}
Let \(c=(c_1,c_2):C\to X\times_SX\) be a correspondence of separated schemes of finite type over the trait \(S\), and let \(F=\Fix(c)\) be its fixed scheme, with \(i=i_c:F\hookrightarrow C\) and \(d:F\to X\). For \(K\in D^b_{\mathrm{ctf}}(X_\eta,A)\) and a morphism \(u:c_{2\eta}^*K\to c_{1\eta}^!K\), put \(K'=\Psi_XK\) and let \(u_s\) be the composite

\[
u_s:c_{2s}^*K'\xrightarrow{\alpha_{c_2}}\Psi_Cc_{2\eta}^*K
\xrightarrow{\Psi_Cu}\Psi_Cc_{1\eta}^!K
\xrightarrow{\beta_{c_1}}c_{1s}^!K'.
\]

For a section \(z:A_{F_\eta}\to\omega_{F_\eta}\) of the dualizing complex, its specialization is

\[
\operatorname{sp}_F(z)=\theta_F \Psi_F(z)\sigma_F:
A_{F_s}\longrightarrow\omega_{F_s}.
\]

Then

\[
\operatorname{Tr}_{c_s}(u_s)
=\operatorname{sp}_F\bigl(\operatorname{Tr}_{c_\eta}(u)\bigr).
\]
\end{bpstatementit}

% text: prose-E-23
\phantomsection\label{bp290}%
Neither properness of the legs \(c_1,c_2\) nor smoothness of \(F\) is assumed. This matters, because we apply the proposition to the deformation of a correspondence to the normal cone. Its generic fibre is the original correspondence, but its fixed scheme, which we describe in Lemma \hyperref[bp306]{E.28} of Section \hyperref[bp300]{E.6}, may be singular and nonreduced.

A trace class has two ingredients besides \(u\): the external duality isomorphism, which turns \(u\) into a section of \(c^!(K\boxtimes\mathbb DK)\), and the evaluation along the diagonal. In the proof we compare both once and for all on \(X\times_SX\), for the universal kernel \(K\boxtimes\mathbb DK\), and then pull the comparison back along \(c\). After taking adjoints, each comparison reduces to the pairing (E.12) that defines \(\delta_K\).

We fix the maps. In the setting of Proposition \hyperref[bp289]{E.22}, let \(P_X=K\boxtimes\mathbb DK\) be the external product of Lemma \bpout{11.7}.  The comparison map for external products is

\[
\Pi_{M,N}=\mu(\alpha_{p_1}\otimes\alpha_{p_2}):
\Psi_XM\boxtimes \Psi_XN\longrightarrow \Psi_{X\times_SX}(M\boxtimes N),
\]

and it is invertible by Proposition \hyperref[bp248]{E.9}. Using it and the isomorphism \(\delta_K=\delta_{X,K}\) of Proposition \hyperref[bp276]{E.18}, we define the isomorphism

\[
\rho:\Psi_{X\times_SX}P_X\xrightarrow{\Pi_{K,\mathbb DK}^{-1}}
K'\boxtimes \Psi_X\mathbb DK\xrightarrow{1\boxtimes\delta_K}K'\boxtimes\mathbb DK'.
\]

Write \(\lambda_K:P_X\xrightarrow\sim
\cHom_A(p_2^*K,p_1^!K)\) for the external duality isomorphism of Lemma \bpout{11.7}, and \(e_K:P_X\to R\Delta_*\omega_X\) for the evaluation along the diagonal (Definition \bpout{11.8}). Let

\[
\mathfrak h:\Psi_{X\times_SX}\cHom_A(p_2^*K,p_1^!K)\longrightarrow
\cHom_A(p_{2s}^*K',p_{1s}^!K')
\]

be the comparison obtained from the tensor–Hom mate of \(\Psi(\operatorname{ev})\mu\) by precomposing with \(\alpha_{p_2}\) and postcomposing with \(\beta_{p_1}\). For the correspondence \(c\), let \(\lambda_{c,K}:c^!P_X\xrightarrow\sim\cHom_A(c_2^*K,c_1^!K)\) be the external duality isomorphism on \(C\), and define \(\mathfrak h_c\) as \(\mathfrak h\), with \(\alpha_{c_2}\) and \(\beta_{c_1}\) in place of \(\alpha_{p_2}\) and \(\beta_{p_1}\).

% object: lemma-nearby-cycles-universal-kernel (blueprint/chap-Specialization-and-contracting-correspondences/sec-Specialization-of-the-trace-class/lemma-nearby-cycles-universal-kernel/_section.md)
\phantomsection\label{bp291}%
\begin{bpstatementit}{Lemma~E.23}{ (Nearby cycles of the universal kernel)}
In the setting above,

\[
\mathfrak h\,\Psi(\lambda_K)=\lambda_{K'}\rho,
\qquad
e_{K'}\rho=R\Delta_{s*}\theta_X\,\chi_\Delta\,\Psi(e_K).
\tag{E.18}
\]
\end{bpstatementit}

% proof: proof/lemma-nearby-cycles-universal-kernel
\begin{proof}[{Proof}]\phantomsection\label{bp292}%
For the first equality, we precompose with \(\Pi_{K,\mathbb DK}\) and take tensor–Hom adjoints. Both maps then have source

\[
(p_{1s}^*K'\otimes^{\mathbf L} p_{2s}^*\Psi_X\mathbb DK)\otimes^{\mathbf L} p_{2s}^*K'
\]

and target \(p_{1s}^!K'\). The adjoint defining \(\lambda_K\) is

\[
(p_1^*K\otimes^{\mathbf L} p_2^*\mathbb DK)\otimes^{\mathbf L} p_2^*K
\xrightarrow{1\otimes\widetilde{\operatorname{ev}}_K}
p_1^*K\otimes^{\mathbf L} p_2^*\omega_X
\xrightarrow{\sim}p_2^*\omega_X\otimes^{\mathbf L} p_1^*K
\xrightarrow{\kappa_K}p_1^!K,
\]

where \(\widetilde{\operatorname{ev}}_K\) is \(\mathbb DK\otimes^{\mathbf L} K\xrightarrow{\sim}K\otimes^{\mathbf L} \mathbb DK
\xrightarrow{\operatorname{ev}_K}\omega_X\). In both displays the arrows marked \(\sim\) are the symmetry isomorphisms of the tensor product, and all Koszul signs come from these symmetries. Since \(\alpha\) is monoidal, we can move the ordinary comparison maps past \(\mu\). By (E.10) we know what \(\beta_{p_1}\) does to the last arrow \(\kappa_K\). What remains on the factors \(K'\) and \(\Psi_X\mathbb DK\) is the composite \(\theta_X\Psi_X(\operatorname{ev}_K)\mu_X\). This is the pairing (E.12), that is, evaluation against \(\delta_K\). The result is the adjoint of \(\lambda_{K'}(1\boxtimes\delta_K)\). Since \(\Pi\) is invertible, the first equality follows.

For the second equality, we again precompose with \(\Pi_{K,\mathbb DK}\) and take adjoints for \(\Delta_s^*\dashv R\Delta_{s*}\). Since \(\alpha\) is monoidal and compatible with composition, the map \(\alpha_\Delta\Delta_s^*\Pi_{K,\mathbb DK}\) is \(\mu_X:K'\otimes^{\mathbf L} \Psi_X\mathbb DK\to \Psi_X(K\otimes^{\mathbf L} \mathbb DK)\). Hence both adjoints are the pairing (E.12) defining \(\delta_K\). Since \(\Delta\) is proper, the comparison map for direct image is \(\chi_\Delta\).
\end{proof}

% text: prose-E-24
\phantomsection\label{bp293}%
By pulling the first identity back along \(c\), we compute the section of \(c_s^!P_{X_s}\) that corresponds to \(u_s\).

% object: lemma-section-specialized-correspondence (blueprint/chap-Specialization-and-contracting-correspondences/sec-Specialization-of-the-trace-class/lemma-section-specialized-correspondence/_section.md)
\phantomsection\label{bp294}%
\begin{bpstatementit}{Lemma~E.24}{ (The section defined by the specialized correspondence)}
In the setting above,

\[
\mathfrak h_c\Psi_C(\lambda_{c_\eta,K})
=\lambda_{c_s,K'}\,c_s^!\rho\,\beta_c(P_X).
\tag{E.19}
\]

Consequently, if \(k_u:A_{C_\eta}\to c_\eta^!P_X\) is the section corresponding to \(u\), then the section corresponding to \(u_s\) is

\[
k_{u_s}=c_s^!\rho\,\beta_c\,\Psi_C(k_u)\,\sigma_C.
\]
\end{bpstatementit}

% proof: proof/lemma-section-specialized-correspondence
\begin{proof}[{Proof}]\phantomsection\label{bp295}%
After taking tensor–Hom adjoints, we factor the ordinary comparison map through \(c_2=p_2c\) and the exceptional comparison map through \(c_1=p_1c\), and we apply (E.8) for \(f=c\).  What remains is the first equality of (E.18) pulled back by \(c_s^!\). The formula for \(k_{u_s}\) follows from (E.19), the definition of \(u_s\) and \(\mu(\sigma\otimes1)=1\).
\end{proof}

% proof: proof/proposition-specialization-trace-class
\begin{proof}[{Proof of Proposition E.22}]\phantomsection\label{bp296}%
We compare the trace morphisms on the two fibres and then evaluate them on the sections of Lemma \hyperref[bp294]{E.24}. By Definition \bpout{11.8}, the trace morphism is \(t_{c,K}=\zeta\,c^!e_K\). The second equality of (E.18), together with (E.11) and with (E.9) for \(d\), gives

\[
t_{c_s,K'}c_s^!\rho\,\beta_c
=Ri_{s*}\theta_F\,\chi_i\,\Psi_C(t_{c_\eta,K}).
\]

If we apply this to the section \(k_{u_s}\) and take adjoints for \(i_s^*\dashv Ri_{s*}\), the equality of trace classes reduces to the identity for the units

\[
\chi_i \Psi_C(\eta_{i_\eta})\sigma_C
=Ri_{s*}(\sigma_F)\eta_{i_s}.
\]

Its ordinary adjoint is the identity \(\sigma_F=\alpha_i i_s^*\sigma_C\) of (E.7), followed by the ordinary triangle identity.
\end{proof}

% text: prose-E-25
\phantomsection\label{bp297}%
The proposition is useful only if \(\operatorname{sp}_F\) loses no information. In the situation of the contracting argument the fixed scheme does not move in the family, and then specialization is the identity.

% object: corollary-specialization-when-fixed-point-scheme (blueprint/chap-Specialization-and-contracting-correspondences/sec-Specialization-of-the-trace-class/corollary-specialization-when-fixed-point-scheme/_section.md)
\phantomsection\label{bp298}%
\begin{bpstatementit}{Corollary~E.25}{ (Specialization when the fixed point scheme is a constant family)}
If \(F_{\mathrm{red}}=Y\times S\), with \(Y\) separated and of finite type over \(k\), then the specialization \(\operatorname{sp}_F\) of Proposition \hyperref[bp289]{E.22} is the identity under the identifications for constant families and invariance under nilpotent thickenings. In particular, it is injective.
\end{bpstatementit}

% proof: proof/corollary-specialization-when-fixed-point-scheme
\begin{proof}[{Proof}]\phantomsection\label{bp299}%
By naturality it suffices to show that \(\theta_Y\sigma_Y(\omega_Y)\) is the identity of \(\omega_Y\). For proper \(Y\) this follows from the adjunction defining \(\theta_Y\). We deduce the case of general \(Y\) by a compactification, and the case of nonreduced \(F\) by invariance under nilpotent thickenings.

\emph{Proper \(Y\) and \(F=Y\times S\).} Here \(\theta_Y\sigma_Y(\omega_Y)=\operatorname{id}_{\omega_Y}\). The unit \(\sigma_Y(\omega_Y):\omega_Y\to \Psi_Y\omega_{Y_\eta}\) is an isomorphism by Lemma \hyperref[bp227]{E.2}. For the proper structure map \(a\) of \(Y\), the base change isomorphism for proper maps and its compatibility with \(\sigma\) give

\[
\nu_a\,Ra_!\sigma_Y(\omega_Y)
=\sigma_{\operatorname{Spec}k}(Ra_!\omega_Y)
=\operatorname{id}_{Ra_!\omega_Y}.
\]

By the defining adjunction (E.6), \(\theta_Y\sigma_Y(\omega_Y)\) is therefore the adjoint of the counit \(Ra_!\omega_Y\to A\), which is \(\operatorname{id}_{\omega_Y}\). Naturality of \(\sigma\) then gives \(\theta_Y\Psi_Y(z)\sigma_Y(A)=z\) for every section \(z\) of the dualizing complex.

\emph{General \(Y\).} Choose a proper compactification \(j:Y\hookrightarrow\overline Y\), so that \(\omega_Y=j^*\omega_{\overline Y}\). We compare the two sides of the identity under restriction along \(j\). The units for the constant families correspond by base change along the open immersion, and the map \(\theta_Y\) corresponds to \(j^*\theta_{\overline Y}\) by the composition rule (E.9) together with \(\beta_j=\alpha_j^{-1}\) for the open immersion. Restricting the identity \(\theta_{\overline Y}\sigma_{\overline Y}(\omega_{\overline Y})=1\) therefore gives \(\theta_Y\sigma_Y(\omega_Y)=1\), and by naturality the assertion follows again for all sections of the dualizing complex.

\emph{Nonreduced \(F\).} If only \(F_{\mathrm{red}}=Y\times S\), we replace \(F\) by its reduction. This does not change \(\operatorname{sp}_F\), because invariance under nilpotent thickenings of the total family and of its fibres shows that the six operations, their units and counits, and nearby cycles agree with those of the reduction.
\end{proof}


% object: sec-Contracting-correspondences (blueprint/chap-Specialization-and-contracting-correspondences/sec-Contracting-correspondences/_section.md), content hidden
\subsection*{\texorpdfstring{E.6\quad Contracting correspondences}{E.6 Contracting correspondences}}
\phantomsection\label{bp300}%
\addcontentsline{toc}{subsection}{\texorpdfstring{E.6\quad Contracting correspondences}{E.6 Contracting correspondences}}
% text: prose-E-26
\phantomsection\label{bp301}%
Let \(c=(c_1,c_2):C\to X\times X\) be a correspondence, and let \(Z\hookrightarrow X\) be a closed subscheme. Suppose that the leg \(c_2\), along which we take ordinary pullback, contracts towards \(Z\) faster than \(c_1\), as measured by powers of ideals. Then the trace class near a connected component of the fixed point scheme meeting \(Z\) can be computed by restricting the correspondence to \(Z\). With our convention \(u:c_2^*K\to c_1^!K\) for the legs, the leg that contracts is \(c_2\); in [\hyperref[bp347]{Var05}] it is the first leg.

The contraction becomes geometric on the deformation to the normal cone (Lemma \hyperref[bp302]{E.26} below). Let \(I\) be the ideal of \(Z\) and \(I_i=c_i^{-1}(I)\mathcal O_C\). Under the hypothesis \(I_2^n\subset I_1^{n+1}\), the special fibre of \(c_2\) maps into the zero section, while the inverse image of the zero section under \(c_1\) is the zero section. On the special fibre the trace class therefore only sees \(Z\).

% object: lemma-normal-cones-contracting-ideals (blueprint/chap-Specialization-and-contracting-correspondences/sec-Contracting-correspondences/lemma-normal-cones-contracting-ideals/_section.md)
\phantomsection\label{bp302}%
\begin{bpstatementit}{Lemma~E.26}{ (Normal cones and contracting ideals)}
Let \(c=(c_1,c_2):C\to X\times X\) be a correspondence, let \(Z\hookrightarrow X\) have ideal \(I\), and put

\[
I_1=c_1^{-1}(I)\mathcal O_C,\qquad
I_2=c_2^{-1}(I)\mathcal O_C,\qquad Z'=V(I_1).
\]

Assume that \(I_2\subset I_1\) and \(I_2^n\subset I_1^{n+1}\) for some \(n\geq1\). Over the trait \(T=\operatorname{Spec}k[t]_{(t)}\), whose strict henselization is the trait \(S\) used for nearby cycles, the two legs lift to the deformations to the normal cones,

\[
\widetilde c_i:\mathcal D_{Z'}C\longrightarrow\mathcal D_ZX,
\qquad
\mathcal D_ZX=
\operatorname{Spec}_{X_T}\mathcal O_{X_T}[I/t],
\]

where \(\mathcal D_ZX\) is written over \(T\) and \(\mathcal D_{Z'}C\) is defined in the same way from \(I_1\). Their special fibres \(\gamma_i:N_{Z'}C\to N_ZX\) satisfy:

\begin{enumerate}
\item \(\gamma_1^{-1}(Z)=Z'\) scheme-theoretically, where \(Z\) and \(Z'\) are the zero sections of the two normal cones;
\item the image of \(\gamma_2\) is set-theoretically contained in the zero section \(Z\subset N_ZX\).
\end{enumerate}
\end{bpstatementit}

% proof: proof/lemma-normal-cones-contracting-ideals
\begin{proof}[{Proof}]\phantomsection\label{bp303}%
We read everything off from the graded algebras. The first leg sends \(I\) onto generators of \(I_1\), so it identifies the ideals of the zero sections. The second leg, on the other hand, sends \(I\) into \(I_2\), and by the hypothesis \(I_2^n\subset I_1^{n+1}\) the images of degree-one classes are nilpotent in \(\operatorname{gr}_{I_1}\mathcal O_C\).

The algebra of the normal cone is \(\operatorname{gr}_I\mathcal O_X=\bigoplus_{r\geq0}I^r/I^{r+1}\), and the algebra of the deformation is the subalgebra of \(\mathcal O_{X_T}[t^{-1}]\) generated by \(\mathcal O_{X_T}\) and the fractions \(a/t\), \(a\in I\). Since \(c_i^{-1}(I)\mathcal O_C\subset I_1\) for both legs (for \(i=2\) this is the hypothesis \(I_2\subset I_1\)), the maps on generic fibres preserve these subalgebras. This gives the lifts.

The zero section of \(N_ZX\) is defined by the ideal of elements of positive degree of \(\operatorname{gr}_I\mathcal O_X\). Since \(I_1\) is the ideal generated by \(c_1^{-1}(I)\), the images of elements of degree one under the first leg generate \(I_1/I_1^2\) as an \(\mathcal O_{Z'}\)-module, and hence they generate the ideal of elements of positive degree of \(\operatorname{gr}_{I_1}\mathcal O_C\). Therefore the quotient by the pullback of the ideal of the zero section is \(\mathcal O_C/I_1=\mathcal O_{Z'}\), which proves (1).

Under the second leg, a product of \(n\) images of elements of degree one is the class of an element of \(I_2^n\) in \(I_1^n/I_1^{n+1}\), and this class is zero since \(I_2^n\subset I_1^{n+1}\). The pullback of the ideal of the zero section is therefore nilpotent. Its vanishing set is all of \(N_{Z'}C\), which means that \((N_{Z'}C)_{\mathrm{red}}\to N_ZX\) factors through the zero section. This proves (2).
\end{proof}

% object: example-contracting-squaring (blueprint/chap-Specialization-and-contracting-correspondences/sec-Contracting-correspondences/example-contracting-squaring/_section.md)
\phantomsection\label{bp304}%
\begin{bpstatement}{Example~E.27}{ (Squaring on the line)}
Let \(X=C=\mathbf A^1=\operatorname{Spec}k[x]\), let \(Z\) be the origin, with \(I=(x)\), and let \(c_1=\mathrm{id}\) and \(c_2(x)=x^2\), so that \(c_2^*x=x^2\). Then \(I_1=(x)\) and \(I_2=(x^2)=I_1^2\), so \(c\) contracts towards \(Z\) with \(n=1\). The fixed scheme is \(V(x^2-x)\). It consists of the two reduced points \(0\) and \(1\), and only the first meets \(Z\).

The deformation to the normal cone is \(\operatorname{Spec}k[t]_{(t)}[x,x/t]\). This is the constant family \(\mathbf A^1_T\) in the coordinate \(x'=x/t\). In this coordinate the first leg is the identity and the second is \(x'\mapsto x^2/t=tx'^2\). Hence, on the special fibre \(t=0\), the second leg is the constant map to \(0\), as (2) of the lemma predicts, and the deformed fixed scheme there is the point \(x'=0\). On the generic fibre, by contrast, there are two fixed points: one is \(x'=0\), the constant family of the origin, and the other is \(x'=1/t\), the point \(x=1\). This second point leaves every bounded region of the \(x'\)-line as \(t\to0\), and so it does not reach the special fibre.

Corollary \hyperref[bp315]{E.31} below therefore gives \(\operatorname{LT}_0(u)=\operatorname{tr}_A(u_{\bar 0}\mid K_{\bar
0})\) for every \(K\) and \(u\). For instance, the local term is zero if \(K=j_!A\) with \(j:\mathbf G_m\hookrightarrow\mathbf A^1\), as \(K_{\bar 0}=0\). The condition fails for the identity correspondence. There \(I_2=I_1\), the inclusion \(I_1^n\subset I_1^{n+1}\) is false for every \(n\), and the fixed scheme is all of \(X\).
\end{bpstatement}

% text: prose-E-27
\phantomsection\label{bp305}%
As in Proposition \hyperref[bp289]{E.22}, write \(F=\Fix(c)\) for the fixed scheme, and let \(d=c_1|_F=c_2|_F:F\to X\) be the common map. We will show that a connected component of \(F\) that meets \(Z\) lies in \(Z\) set-theoretically, and that the deformed fixed scheme is constant up to nilpotents. Both follow from the hypothesis \(I_2^n\subset I_1^{n+1}\), because the two legs agree on \(F\).

The deformation of \(F\) to the normal cone of \(d^{-1}(Z)\) embeds in the fixed scheme of the deformed correspondence. Put \(I_F=d^{-1}(I)\mathcal O_F\), so that \(d^{-1}(Z)=V(I_F)\). Since \(d=c_1|_F\), the ideal \(I_F^r\) is the image of \(I_1^r\) in \(\mathcal O_F\) for every \(r\). Hence the closed immersion \(F\hookrightarrow C\) induces a surjection of deformation algebras

\[
\mathcal O_{C_T}[I_1/t]\twoheadrightarrow
\mathcal O_{F_T}[I_F/t],
\]

that is, a closed immersion \(\mathcal D_{d^{-1}(Z)}F\hookrightarrow\mathcal D_{Z'}C\). Its two composites with \(\widetilde c_1\) and \(\widetilde c_2\) are the maps induced by \(c_1|_F\) and \(c_2|_F\), which are equal. So it factors through the fixed scheme \(\widetilde F=\Fix(\widetilde c)\) of \(\widetilde c=(\widetilde c_1,\widetilde c_2)\). We write

\[
\iota_F:\mathcal D_{d^{-1}(Z)}F\longrightarrow\widetilde F
\tag{E.20}
\]

for the resulting closed immersion.

% object: lemma-fixed-point-scheme-deformation (blueprint/chap-Specialization-and-contracting-correspondences/sec-Contracting-correspondences/lemma-fixed-point-scheme-deformation/_section.md)
\phantomsection\label{bp306}%
\begin{bpstatementit}{Lemma~E.28}{ (The fixed point scheme of the deformation)}
In the situation of Lemma \hyperref[bp302]{E.26}, suppose that \(F\) is connected and that \(d(F)\) meets \(Z\). Then \(d^{-1}(Z)_{\mathrm{red}}=F_{\mathrm{red}}\), the reduction of \(\mathcal D_{d^{-1}(Z)}F\) is \((F_T)_{\mathrm{red}}\), and the map \((\iota_F)_{\mathrm{red}}:(F_T)_{\mathrm{red}}\to\widetilde F_{\mathrm{red}}\) induced by (E.20) is an isomorphism.
\end{bpstatementit}

% proof: proof/lemma-fixed-point-scheme-deformation
\begin{proof}[{Proof}]\phantomsection\label{bp307}%
Restricting \(I_2^n\subset I_1^{n+1}\) to \(F\), where the two legs agree, gives \(I_F^n\subset I_F^{n+1}\). The reverse inclusion always holds, so the coherent ideal \(M=I_F^n\) satisfies \(M=I_FM\). Hence, at a point of \(V(I_F)\), Nakayama's lemma shows that \(M=0\) in the local ring, and hence on a neighbourhood of the point, whereas at a point outside \(V(I_F)\) we have \(M=\mathcal O_F\). The closed set \(V(I_F)\) is therefore also open. It is nonempty by assumption, and \(F\) is connected, so \(V(I_F)=F\) and \(I_F^n=0\). This proves the first assertion. As \(I_F\) is nilpotent, each generator \(h/t\) of \(\mathcal O_{F_T}[I_F/t]\) over \(\mathcal O_{F_T}\) is nilpotent. Hence the reduction of \(\mathcal D_{d^{-1}(Z)}F\) is \((F_T)_{\mathrm{red}}\).

It remains to see that \(\widetilde F\) has no further points. On the generic fibre the deformed correspondence is \(c\) itself, so the generic fibre of \(\widetilde F\) is that of \(F_T\). Let \(x\) be a point of the special fibre of \(\widetilde F\). Then \(\gamma_2(x)\) lies in the zero section by Lemma lemma-normal-cones-contracting-ideals, and so does \(\gamma_1(x)=\gamma_2(x)\). By Lemma lemma-normal-cones-contracting-ideals, the point \(x\) lies on the zero section \(Z'\subset N_{Z'}C\). On \(Z'\) the two maps are the restrictions of \(c_1\) and \(c_2\). Their equalizer is \(F\cap Z'\), whose underlying space is \(F\) by the first assertion. Thus every fixed point in the special fibre lies in the closed subscheme above, and the closed immersion of reduced schemes is surjective. Its ideal has zero radical and is therefore zero, so the immersion is an isomorphism.
\end{proof}

% text: prose-E-28
\phantomsection\label{bp308}%
Up to nilpotents, then, the fixed scheme of the deformed correspondence is the constant family \(F_T\). This is why the specialization of trace classes can be inverted in the proof of Theorem \hyperref[bp311]{E.30}. The hypothesis that \(F\) is connected and meets \(Z\) is needed. For instance, in Example \hyperref[bp304]{E.27}, the component \(F=\{0\}\) gives the constant family \(x'=0\), but the fixed point \(x=1\) does not meet \(Z\), and there the deformation gives only the generic-fibre point \(x'=1/t\), not the constant family \(\{1\}_T\).

% object: remark-reduction-connected-fixed-point-scheme (blueprint/chap-Specialization-and-contracting-correspondences/sec-Contracting-correspondences/remark-reduction-connected-fixed-point-scheme/_section.md)
\phantomsection\label{bp309}%
\begin{bpstatement}{Remark~E.29}{ (Reduction to a connected fixed point scheme)}
The hypothesis in the preceding lemma that \(F\) is connected is no loss of generality. Indeed, let \(F_0\subset \Fix(c)\) be a connected component that is open and closed. If we replace \(C\) by an open neighbourhood of \(F_0\) on which the ideal inclusions hold, and then remove \(\Fix(c)\setminus F_0\), we obtain an open subscheme of \(C\), with the same \(X\) and \(K\), whose fixed scheme is \(F_0\). Restriction of trace classes along this open immersion identifies the trace class of \(u\) on \(F_0\) with the full trace class of the restricted correspondence.
\end{bpstatement}

% text: prose-E-29
\phantomsection\label{bp310}%
The proof of the main theorem combines the two lemmas on the deformation with Verdier's theorem and with the specialization of trace classes. The statement of the main theorem involves the restriction of a correspondence to \(Z\), which we construct first.

Let \(c=(c_1,c_2):C\to X\times X\) be a correspondence, let \(Z\hookrightarrow X\) have ideal \(I\), and suppose that \(c_2^{-1}(I)\mathcal O_C\subset c_1^{-1}(I)\mathcal O_C\). Then \(c_1^{-1}(Z)\subset c_2^{-1}(Z)\). As before Lemma \bpout{11.17}, we put

\[
C_Z=(c_1^{-1}Z)_{\mathrm{red}},\qquad
c_Z=(c_{Z,1},c_{Z,2}):C_Z\longrightarrow Z_{\mathrm{red}}\times Z_{\mathrm{red}},
\]

with \(i:Z_{\mathrm{red}}\hookrightarrow X\) and \(i':C_Z\hookrightarrow C\). For \(K\in D^b_{\mathrm{ctf}}(X,A)\) put \(K_Z=i^*K\). The square formed by \(i'\), \(i\), \(c_1\) and \(c_{Z,1}\) is cartesian up to reduction. Hence invariance under nilpotent thickenings and base change for direct image with compact support give \(i^*Rc_{1!}\simeq Rc_{Z,1!}i'^*\). Since the second leg maps \(C_Z\) into \(Z_{\mathrm{red}}\), we also have \(i'^*c_2^*K\simeq c_{Z,2}^*K_Z\). For a cohomological correspondence \(u:c_2^*K\to c_1^!K\) with adjoint \(\widehat u:Rc_{1!}c_2^*K\to K\) under \(Rc_{1!}\dashv c_1^!\), the \emph{restriction} \(u_Z:c_{Z,2}^*K_Z\to c_{Z,1}^!K_Z\) of \(u\) to \(Z\) is the adjoint under \(Rc_{Z,1!}\dashv c_{Z,1}^!\) of

\[
Rc_{Z,1!}c_{Z,2}^*K_Z
\simeq i^*Rc_{1!}c_2^*K
\xrightarrow{i^*\widehat u}K_Z.
\tag{E.21}
\]

Equivalently, \(u_Z\) is \(i'^*u\) followed by the base change map \(i'^*c_1^!K\to c_{Z,1}^!K_Z\), which in general is not an isomorphism. This is the restriction to \(Z\) that occurs in Lemma \bpout{11.18}. If \(F_0\) is an open and closed subscheme of \(\Fix(c)\), put \(F_{0,Z}=\Fix(c_Z)\times_{\Fix(c)}F_0\), and let \(h_{F_0}:F_{0,Z}\to F_0\) be the projection.

% object: theorem-contracting-correspondences-local-terms (blueprint/chap-Specialization-and-contracting-correspondences/sec-Contracting-correspondences/theorem-contracting-correspondences-local-terms/_section.md)
\phantomsection\label{bp311}%
\begin{bpstatementit}{Theorem~E.30}{ (Contracting correspondences and local terms)}
Let \(c=(c_1,c_2):C\to X\times X\) be a correspondence, and let \(Z\hookrightarrow X\) have ideal \(I\). Let \(F_0\) be a connected component of \(\Fix(c)\) that is open and closed and whose image meets \(Z\). Choose an open neighbourhood \(W\subset C\) of \(F_0\) on which

\[
I_2=c_2^{-1}(I)\mathcal O_W\subset
I_1=c_1^{-1}(I)\mathcal O_W,
\qquad I_2^n\subset I_1^{n+1}\quad\text{for some }n\geq1.
\]

Let \(K\in D^b_{\mathrm{ctf}}(X,A)\), let \(u:c_2^*K\to c_1^!K\) be a cohomological correspondence, and let \(u_Z\) be the restriction (E.21) to \(Z\) of \(u|_W\), for the correspondence \(c|_W\), with \(c_Z:C_Z=(W\cap c_1^{-1}Z)_{\mathrm{red}}\to
Z_{\mathrm{red}}\times Z_{\mathrm{red}}\). Then \(h_{F_0}\) is an isomorphism after reduction, and

\[
\operatorname{Tr}_c(u)|_{F_0}
=(h_{F_0})_*\bigl(\operatorname{Tr}_{c_Z}(u_Z)|_{F_{0,Z}}\bigr),
\]

where \((h_{F_0})_*:H^0(F_{0,Z},\omega_{F_{0,Z}})\to H^0(F_0,\omega_{F_0})\) is the induced isomorphism. In particular, if \(F_0\) is proper, then \(\operatorname{LT}_{F_0}(u)=\operatorname{LT}_{F_{0,Z}}(u_Z)\).
\end{bpstatementit}

% text: prose-E-30
\phantomsection\label{bp312}%
For \(Z\) a point, the theorem specializes to Corollary \hyperref[bp315]{E.31}. Under its hypotheses the local term at \(F_0\) is the categorical trace of \(u_Z\). For \(c=(\mathrm{id}_X,h)\) this is \(\operatorname{tr}(u_{\bar y}\mid K_{\bar y})\).

% proof: proof/theorem-contracting-correspondences-local-terms
\begin{proof}[{Proof of Theorem E.30}]\phantomsection\label{bp313}%
We split \(K\) by the triangle \(j_!j^*K\to K\to Ri_*i^*K\) and use additivity of trace classes. The piece supported on \(Z\) gives the trace of \(u_Z\). The piece \(j_!j^*K\), on the other hand, has zero trace, which we show by deforming to the normal cone.

\emph{Reductions.} We may assume that \(W=C\), that \(\Fix(c)=F_0\) and that the ideal inclusions hold on \(C\). Indeed, by Remark \hyperref[bp309]{E.29}, restricting \(C\) to \(W\) and removing the other connected components of \(\Fix(c)\) changes neither the trace class on \(F_0\) nor \(u_Z\) near \(F_{0,Z}\). By Lemma \hyperref[bp306]{E.28}, the common map \(\Fix(c)\to X\) maps set-theoretically into \(Z\). The same lemma shows that after base change from \(T\) to \(S\) the reduction of the fixed scheme of the deformed correspondence is \((\Fix(c)\times S)_{\mathrm{red}}\), as Corollary \hyperref[bp298]{E.25} requires.

\emph{The case \(i^*K=0\).} We show that \(\operatorname{Tr}_c(u)=0\). First, the specialized complex vanishes on the zero section,

\[
\operatorname{Sp}_{X,Z}(K)|_Z=0,
\]

because Verdier's theorem on the zero section (Theorem \hyperref[bp244]{E.8}) identifies the restriction of \(\operatorname{Sp}_{X,Z}(K)\) to the zero section with the restriction of \(K\) to \(Z\), and invariance under the nilpotent thickening \(Z_{\mathrm{red}}\to Z\) passes from \(Z\) to \(Z_{\mathrm{red}}\).

Next, consider the leg \(\gamma_2\) of the correspondence on the special fibre, along which we take ordinary pullback. It maps set-theoretically into the zero section by Lemma lemma-normal-cones-contracting-ideals. Invariance under nilpotent thickenings therefore gives \(\gamma_2^*\operatorname{Sp}_{X,Z}(K)=0\). So the specialized morphism has zero source, and its trace is zero.

Finally, we relate this trace to \(\operatorname{Tr}_c(u)\). The generic fibre of the deformed correspondence is \(c\) over \(\eta\). So Proposition \hyperref[bp289]{E.22} identifies the zero trace with the specialization of the pullback of \(\operatorname{Tr}_c(u)\) to the generic fibre. By Corollary \hyperref[bp298]{E.25}, specialization is the identity under the identifications for constant families, and these carry that pullback to \(\operatorname{Tr}_c(u)\). Hence \(\operatorname{Tr}_c(u)=0\).

\emph{General \(K\).} We apply Proposition proposition-additivity-trace-classes to the triangle

\[
j_!j^*K\longrightarrow K\longrightarrow Ri_*i^*K,
\qquad j:X\setminus Z_{\mathrm{red}}\hookrightarrow X,
\]

whose hypothesis \(c_1^{-1}(Z_{\mathrm{red}})\subset c_2^{-1}(Z_{\mathrm{red}})\) follows from \(I_2\subset I_1\). The open piece has zero trace by the case just treated, while the closed graded piece is the proper pushforward of \(u_Z\) along \(i\) and \(i'\), by Lemma \bpout{11.18}, and Theorem \bpout{11.13} identifies its trace with the pushforward of the trace class of \(u_Z\). To compare with \(F_0\), note that the map \(h_{F_0}\) is an isomorphism after reduction, because \(\Fix(c)\) lies in \(c_1^{-1}(Z)\) set-theoretically by Lemma \hyperref[bp306]{E.28}. This proves the displayed formula. If \(F_0\) is proper, then so is \(F_{0,Z}\), and pushing both sides forward to \(\operatorname{Spec}k\) and composing with the trace map of Definition \bpout{11.9} gives the equality of local terms.
\end{proof}

% text: prose-E-31
\phantomsection\label{bp314}%
If \(Z\) is a point, the restricted correspondence is a correspondence of the point, and the theorem computes the local term as a trace on a stalk.

% object: corollary-contracting-local-terms-point (blueprint/chap-Specialization-and-contracting-correspondences/sec-Contracting-correspondences/corollary-contracting-local-terms-point/_section.md)
\phantomsection\label{bp315}%
\begin{bpstatementit}{Corollary~E.31}{ (Contracting local terms at a point)}
Let \(y\in X(k)\), and suppose that \(c_1^{-1}(y)\) is a single reduced point \(z\), with \(c_2(z)=y\). Suppose that \(F_0=\{z\}\) is an isolated reduced component of \(\Fix(c)\), open and closed, and that on a neighbourhood of \(F_0\) the ideal \(I_y\) satisfies

\[
c_2^{-1}(I_y)\mathcal O_C\subset c_1^{-1}(I_y)\mathcal O_C,
\qquad
\bigl(c_2^{-1}(I_y)\mathcal O_C\bigr)^n
\subset\bigl(c_1^{-1}(I_y)\mathcal O_C\bigr)^{n+1}
\quad\text{for some }n\geq1.
\]

Then for every \(K\in D^b_{\mathrm{ctf}}(X,A)\) and every morphism \(u:c_2^*K\to c_1^!K\), the local term at \(F_0\) is the categorical trace of the restriction \(u_Z\) to \(Z=\{y\}\) of \(u\) restricted to that neighbourhood, defined in (E.21). In particular, for \(c=(\operatorname{id}_X,h)\),

\[
\operatorname{LT}_y(u)
=\operatorname{tr}_A(u_{\bar y}\mid K_{\bar y}).
\]
\end{bpstatementit}

% text: prose-E-32
\phantomsection\label{bp316}%
Here \(X\) may be singular.

% proof: proof/corollary-contracting-local-terms-point
\begin{proof}[{Proof of Corollary E.31}]\phantomsection\label{bp317}%
We apply Theorem \hyperref[bp311]{E.30} with \(Z=\{y\}\) and \(W\) the neighbourhood of the statement; \(F_0\) meets \(Z\) because \(d(z)=c_1(z)=y\). Here \(C_Z\) is the reduced point \(z\) and \(F_{0,Z}=F_0\), so that \(c_Z\) is the identity correspondence of the point, with coefficient \(K_{\bar y}\), and by Definition \bpout{11.9} its local term is the categorical trace. If the first leg is the identity, the base change map in the construction of \(u_Z\) is the identity, and hence the restricted morphism is the stalk map \(u_{\bar y}\).
\end{proof}



% object: chap-Algebraic-groups-in-characteristic-zero (blueprint/chap-Algebraic-groups-in-characteristic-zero/_section.md), content hidden
\section*{\texorpdfstring{F\quad Algebraic groups in characteristic zero}{F Algebraic groups in characteristic zero}}
\phantomsection\label{bp318}%
\addcontentsline{toc}{section}{\texorpdfstring{F\quad Algebraic groups in characteristic zero}{F Algebraic groups in characteristic zero}}
% text: prose-F-01
\phantomsection\label{bp319}%
Rankin's argument (Chapter \bpout{14}) bounds the eigenvalues of a real sheaf by the largest determinant weight of its constituents, that is, the weight of the determinant divided by the rank. For this it needs the determinant weights of the factors of tensor and exterior powers of a lisse sheaf (Proposition \bpout{14.19}), before anyone knows how these powers decompose. The tool is an element of the arithmetic monodromy group, of positive degree, that acts by a scalar on each simple factor of every tensor construction. This scalar then gives the determinant weight.

In a simple example Frobenius itself does this. Take \(V=\overline{\mathbf Q}_\ell\oplus\overline{\mathbf Q}_\ell(-1)\) on a curve over \(\mathbf F_q\), on which a Frobenius element \(s\) of degree one acts by \(\operatorname{diag}(1,q)\). Here the geometric monodromy is trivial, so the arithmetic monodromy group \(G\) is the torus \(\{\operatorname{diag}(1,t)\}\), and \(\rho(s)\) is central in it. Hence on every tensor construction \(\rho(s)\) acts by a scalar on each simple factor. For instance, it acts by \(1\) and \(q\) on the factors of \(V\), of weights \(0\) and \(2\), and by \(q\) on \(\bigwedge^2V\).

In general \(\rho(s)\) is not central, and we must correct a power of it by an element of the geometric monodromy group \(H\). An element of \(H\) alone would not do, since its determinants on these factors are roots of unity. What we want instead is a central element of the form \(h\rho(s)^N\) with \(h\in H\) and \(N>0\).

Such an element exists by Lemma \bpout{14.15}. For a semisimple Weil representation \(V\), the identity component \(G^\circ\) is generated by \(Z(G^\circ)^\circ\) and the identity component \(H^\circ\). So a power of Frobenius lying in \(G^\circ\) is an element of \(Z(G^\circ)^\circ\) times an element of \(H^\circ\). Averaging over the components of \(G\) then makes it central in \(G\). This is the algebraic form of a familiar fact about compact Lie groups: a connected compact group is the product of the identity component of its center and its derived group.

For algebraic groups over \(\overline{\mathbf Q}_\ell\) we argue with the Lie algebra \(\mathfrak g\) of \(G^\circ\). Since \(\mathfrak g\) acts faithfully and semisimply on \(V\), it is reductive, that is,

\[
\mathfrak g=Z(\mathfrak g)+[\mathfrak g,\mathfrak g]
\]

(Proposition \hyperref[bp345]{F.8}). The derived algebra lies in the Lie algebra of \(H\), because commutators of \(G\) lie in \(H\). It remains to use that a closed connected subgroup of \(G^\circ\) with the full Lie algebra is all of \(G^\circ\).

The tools behind this argument need characteristic zero. Cartier's theorem that algebraic groups are smooth (Proposition \hyperref[bp322]{F.1}) lets us recognize a closed subgroup by its Lie algebra; this fails for \(\mu_p\subset\mathbf G_m\) in characteristic \(p\). Similarly, the decomposition of \(\mathfrak g\) rests on Cartan's criterion, which fails there too. Since the monodromy groups are defined over \(\overline{\mathbf Q}_\ell\), this does not restrict us. For the theorems of Lie, Engel and Cartan we refer to Humphreys [\hyperref[bp354]{Hum72}].

Throughout, \(\Omega\) is an algebraically closed field of characteristic zero, and all Lie algebras and representations are finite-dimensional. By an algebraic group we mean a group scheme of finite type over \(\Omega\). In particular, stabilizers and centralizers are taken as subgroup schemes, and by Cartier's theorem below they are smooth. The Lie algebra \(\operatorname{Lie}(G)\) of an algebraic group \(G\) is its tangent space at the identity; for a closed subgroup of \(\GL(V)\) we describe it by dual numbers in Section \hyperref[bp327]{F.2}.

% object: sec-Smoothness-and-components (blueprint/chap-Algebraic-groups-in-characteristic-zero/sec-Smoothness-and-components/_section.md), content hidden
\subsection*{\texorpdfstring{F.1\quad Smoothness and components}{F.1 Smoothness and components}}
\phantomsection\label{bp320}%
\addcontentsline{toc}{subsection}{\texorpdfstring{F.1\quad Smoothness and components}{F.1 Smoothness and components}}
% text: prose-F-02
\phantomsection\label{bp321}%
To pass from Lie algebras back to groups we use the following principle: a closed subgroup of a connected group \(G\) with the same Lie algebra is all of \(G\) (Lemma \hyperref[bp325]{F.2}). For reduced groups this is a dimension count: the dimension of a smooth group equals that of its tangent space at \(1\), and a closed subset of the irreducible variety \(G\) of the same dimension is all of \(G\). The difficulty is that stabilizers and centralizers are defined by equations and need not be reduced. For example, in characteristic \(p\), the subgroup \(\mu_p\subset\mathbf G_m\) defined by \(x^p=1\) is a single nonreduced point with tangent space \(\operatorname{Lie}(\mathbf G_m)\), and the principle fails. However, Cartier's theorem (Proposition \hyperref[bp322]{F.1}) shows that this cannot happen in characteristic zero: every algebraic group is smooth. In the proof we apply the comultiplication to a nilpotent function at the identity and divide by an integer, and this division is where we use characteristic zero. The facts about connected components and images of homomorphisms in Lemma \hyperref[bp325]{F.2} are standard consequences of smoothness and of Chevalley's theorem.

% object: proposition-cartier-s-theorem (blueprint/chap-Algebraic-groups-in-characteristic-zero/sec-Smoothness-and-components/proposition-cartier-s-theorem/_section.md)
\phantomsection\label{bp322}%
\begin{bpstatementit}{Proposition~F.1}{ (Cartier's theorem)}
Every algebraic group over \(\Omega\) is smooth and reduced. In particular

\[
\dim\operatorname{Lie}(G)=\dim G.
\]
\end{bpstatementit}

% proof: proof/proposition-cartier-s-theorem
\begin{proof}[{Proof}]\phantomsection\label{bp323}%
We first show, using translations, that the reduced subscheme \(G_{\mathrm{red}}\) is a smooth group scheme. It then suffices to show that \(G\) has the same tangent space at \(1\) as \(G_{\mathrm{red}}\), that is, that nilpotent functions at \(1\) lie in \(\mathfrak m^2\), and this is where we divide by an integer.

\emph{Step 1: \(G_{\mathrm{red}}\) is a smooth subgroup scheme of \(G\).} Multiplication and inversion restrict to \(G_{\mathrm{red}}\). This is because products of reduced schemes of finite type over the perfect field \(\Omega\) are reduced [Stacks, \href{https://stacks.math.columbia.edu/tag/035Z}{035Z}], and a morphism from a reduced scheme to \(G\) factors through \(G_{\mathrm{red}}\). The reduced scheme \(G_{\mathrm{red}}\) contains a dense smooth open subscheme [Stacks, \href{https://stacks.math.columbia.edu/tag/056V}{056V}]. It is smooth everywhere, since a translation moves a closed point of this open to any other closed point of \(G_{\mathrm{red}}\).

\emph{Step 2: nilpotent functions at \(1\) lie in \(\mathfrak m^2\).} Let \(R=\mathcal O_{G,1}\), with maximal ideal \(\mathfrak m\) and nilradical \(I\), and let \(\widehat R\) be its completion. At the point \((1,1)\) the multiplication induces a continuous comultiplication

\[
\Delta:\widehat R\longrightarrow
\widehat R\,\widehat\otimes_\Omega\widehat R.
\]

The counit is evaluation at \(1\), and composing \(\Delta\) with evaluation at \(1\) in either factor gives the identity (the counit identities), because \(1\) is a two-sided unit for the group law. We claim that \(I\subseteq\mathfrak m^2\). As completion is faithfully flat, it suffices to prove that every nilpotent \(b\in\widehat R\) lies in \(\widehat{\mathfrak m}^{\,2}\). Suppose that \(b\notin\widehat{\mathfrak m}^{\,2}\), and let \(n\ge2\) be the least integer with \(b^n=0\). By the counit identities,

\[
\Delta(b)=b\otimes1+1\otimes b+w,
\qquad
w\in\widehat{\mathfrak m}\,\widehat\otimes_\Omega
\widehat{\mathfrak m}.
\]

Now reduce the second factor modulo \(\widehat{\mathfrak m}^{\,2}\). Every term in the binomial expansion of \(\Delta(b)^n\) with at least two factors from \(1\otimes b+w\) vanishes, because \(1\otimes b+w\) has its second factor in \(\widehat{\mathfrak m}\), and as \(b^n=0\) we are left with

\[
n b^{n-1}\otimes b+n(b^{n-1}\otimes1)w=0.
\]

If we now also reduce the first factor modulo \(b^{n-1}\widehat{\mathfrak m}\), the second term vanishes, and we get

\[
n\,\overline{b^{n-1}}\otimes\overline b=0
\quad\text{in}\quad
(\widehat R/b^{n-1}\widehat{\mathfrak m})
\otimes_\Omega(\widehat R/\widehat{\mathfrak m}^{\,2}).
\]

This identity is impossible, which proves the claim. Indeed, the class \(\overline{b^{n-1}}\) is nonzero by Nakayama's lemma, because the module \((b^{n-1})\) is nonzero by the choice of \(n\). The class \(\overline b\) is nonzero by assumption. A tensor product of nonzero vectors over a field is nonzero, and \(n\ne0\) in characteristic zero.

\emph{Step 3: \(G\) is smooth.} The schemes \(G\) and \(G_{\mathrm{red}}\) have the same cotangent space at \(1\). Indeed, by Step 2 the kernel of \(\mathfrak m/\mathfrak m^2\to(\mathfrak m/I)/(\mathfrak m/I)^2\) is \((I+\mathfrak m^2)/\mathfrak m^2=0\). They also have the same local dimension at \(1\), because a nilpotent ideal does not change the underlying space. As the local ring \(R/I\) of \(G_{\mathrm{red}}\) at \(1\) is regular by Step 1, we get

\[
\dim_\Omega\mathfrak m/\mathfrak m^2=\dim R,
\]

so \(R\) is regular, and in particular a domain [Stacks, \href{https://stacks.math.columbia.edu/tag/00NP}{00NP}]. By translation the local rings of \(G\) at all closed points are regular. Since \(\Omega\) is perfect, \(G\) is smooth at every closed point. The nonsmooth locus is closed and contains no closed point, and hence it is empty. Being smooth, \(G\) is reduced, and its tangent space at \(1\) has dimension \(\dim G\). Finally, the open subgroup \(G^\circ\) has the same tangent space at \(1\).
\end{proof}

% text: prose-F-03
\phantomsection\label{bp324}%
Since algebraic groups over \(\Omega\) are smooth, their topology is as simple as that of a Lie group: irreducible and connected components coincide, and the dimension of a closed subgroup can be computed from its Lie algebra.

% object: lemma-connected-components-images-homomorphisms (blueprint/chap-Algebraic-groups-in-characteristic-zero/sec-Smoothness-and-components/lemma-connected-components-images-homomorphisms/_section.md)
\phantomsection\label{bp325}%
\begin{bpstatementit}{Lemma~F.2}{ (Connected components and images of homomorphisms)}
Let \(G\) be an algebraic group.

\begin{enumerate}
\item \(G\) has finitely many connected components. Its identity component \(G^\circ\) is an open, closed, irreducible normal subgroup, the other components are its cosets, and \(\operatorname{Lie}(G^\circ)=\operatorname{Lie}(G)\).
\item The image of a homomorphism \(G\to G'\) of algebraic groups is a closed subgroup of \(G'\).
\item If \(G\) is connected and \(K\subseteq G\) is a closed subgroup scheme with \(\operatorname{Lie}(K)=\operatorname{Lie}(G)\), then \(K=G\).
\end{enumerate}
\end{bpstatementit}

% proof: proof/lemma-connected-components-images-homomorphisms
\begin{proof}[{Proof}]\phantomsection\label{bp326}%
For (1) we use smoothness to see that irreducible and connected components coincide; for (2), Chevalley's theorem gives an open subset of the closure inside the image, and translations spread it over the closure; (3) is a dimension count.

\emph{Components.} The irreducible components of \(G\) are finite in number, as \(G\) is Noetherian, and they are exactly the connected components. Indeed, \(G\) is smooth by Proposition \hyperref[bp322]{F.1}, and a regular local ring is a domain. So distinct irreducible components of \(G\) do not meet, and they are open and closed. Let \(G^\circ\) be the component containing \(1\). It is a subgroup. To see this, note that \(G^\circ\times G^\circ\) is irreducible, \(\Omega\) being algebraically closed. Hence the image of the morphism \(G^\circ\times G^\circ\to G\), \((x,y)\mapsto xy^{-1}\), is irreducible. It contains \(1\), so it lies in \(G^\circ\). For the same reason the image of \(G^\circ\) under conjugation by a point of \(G\) is irreducible and contains \(1\), so \(G^\circ\) is normal. Translating \(G^\circ\) by points of \(G\), we see that the other components are its cosets. Since \(G^\circ\) is open in \(G\) and contains \(1\), both have the same local ring at \(1\), so \(\operatorname{Lie}(G^\circ)=\operatorname{Lie}(G)\).

For example, the homomorphism \(\mathbf G_m\to\mathbf G_m\), \(x\mapsto x^2\), has image all of \(\mathbf G_m\) and kernel \(\mu_2\), a closed subgroup with two components.

\emph{Images.} Let \(f:G\to G'\) be a homomorphism. We show that its image equals the closure \(J\) of the image, which is a subgroup by continuity of multiplication and inversion. First, the subgroup \(f(G(\Omega))\) of \(J(\Omega)\) contains the \(\Omega\)-points of a nonempty open \(U\subseteq J\). Indeed, \(f(G)\) is constructible by Chevalley's theorem [Stacks, \href{https://stacks.math.columbia.edu/tag/054K}{054K}], a dense constructible subset contains a nonempty open subset of its closure, and \(\Omega\) is algebraically closed. Translating \(U\) by the inverse of one of its points in the image, we obtain an open neighbourhood of \(1\) contained in the image. Translating further by points of the image, we see that the image is open in \(J\). It is then also closed, since its complement is a union of cosets of this open subgroup. As it is dense in \(J\), it is equal to \(J\).

\emph{Lie algebras.} Let \(K\) be as in (3). Cartier's theorem gives \(\dim K=\dim\operatorname{Lie}(K)=\dim G\), so \(K\) and \(G\) have the same underlying closed subset, as a proper closed subset of the irreducible variety \(G\) has smaller dimension. They are equal as schemes, since both are reduced by Proposition \hyperref[bp322]{F.1}.
\end{proof}


% object: sec-Lie-algebras (blueprint/chap-Algebraic-groups-in-characteristic-zero/sec-Lie-algebras/_section.md), content hidden
\subsection*{\texorpdfstring{F.2\quad Lie algebras}{F.2 Lie algebras}}
\phantomsection\label{bp327}%
\addcontentsline{toc}{subsection}{\texorpdfstring{F.2\quad Lie algebras}{F.2 Lie algebras}}
% text: prose-F-04
\phantomsection\label{bp328}%
Let \(G\) be a connected algebraic group with Lie algebra \(\mathfrak g\). We need two facts (Lemma \hyperref[bp329]{F.3}). First, in a representation, \(G\) and \(\mathfrak g\) have the same stable subspaces. Second, the Lie algebra of the identity component of the center of \(G\) is the center of \(\mathfrak g\). In the proof of Lemma \bpout{14.15} we compare the monodromy groups through their Lie algebras, and both facts are used there. By the first, a representation that is semisimple for \(G^\circ\) is semisimple for its Lie algebra. This is what lets Proposition \hyperref[bp345]{F.8} apply to the Weil representation. With the second we translate the decomposition of the Lie algebra into the identity \(G^\circ=Z(G^\circ)^\circ H^\circ\).

We prove both facts by the principle of the previous section, applied to a stabilizer and to the kernel of the adjoint representation. Their Lie algebras are computed with dual numbers. The first fact needs \(G\) to be connected. For example, the finite group \(\{\pm1\}\subset\GL_1\) has zero Lie algebra, so every subspace of a representation is stable under its Lie algebra. But if \(-1\) acts on \(M=\Omega^2\) by \(\operatorname{diag}(1,-1)\), then it does not stabilize the line spanned by \((1,1)\).

For a closed subgroup \(G\subseteq\GL(V)\), the tangent space \(\operatorname{Lie}(G)\subseteq\operatorname{End}(V)\) is described by dual-number points at \(1\):

\[
X\in\operatorname{Lie}(G)
\quad\Longleftrightarrow\quad
1+\varepsilon X\in G(\Omega[\varepsilon]/\varepsilon^2).
\]

The bracket is \([X,Y]=XY-YX\). To see that \(\operatorname{Lie}(G)\) is closed under the bracket, note that over \(\Omega[\varepsilon,\eta]/(\varepsilon^2,\eta^2)\) the commutator of \(1+\varepsilon X\) and \(1+\eta Y\) is \(1+\varepsilon\eta[X,Y]\). Its entries lie in the subring \(\Omega[\varepsilon\eta]\), which is a ring of dual numbers. The equations of \(G\) hold for this matrix in the larger ring, hence in this subring. Therefore \([X,Y]\in\operatorname{Lie}(G)\). The differential of a representation \(r:G\to\GL(M)\) is characterized by

\[
r(1+\varepsilon X)=1+\varepsilon\,dr(X).
\]

Applying \(r\) to the commutator of \(1+\varepsilon X\) and \(1+\eta Y\), we see that \(dr[X,Y]=[dr(X),dr(Y)]\), so \(dr\) is a Lie algebra homomorphism.

% object: lemma-stable-subspaces-center (blueprint/chap-Algebraic-groups-in-characteristic-zero/sec-Lie-algebras/lemma-stable-subspaces-center/_section.md)
\phantomsection\label{bp329}%
\begin{bpstatementit}{Lemma~F.3}{ (Stable subspaces and the center)}
Let \(G\subseteq\GL(V)\) be a connected closed subgroup and \(r:G\to\GL(M)\) an algebraic representation. A subspace of \(M\) is \(G\)-stable if and only if it is \(\operatorname{Lie}(G)\)-stable. Moreover

\[
\operatorname{Lie}\bigl(Z(G)^\circ\bigr)
=Z\bigl(\operatorname{Lie}(G)\bigr).
\]
\end{bpstatementit}

% proof: proof/lemma-stable-subspaces-center
\begin{proof}[{Proof}]\phantomsection\label{bp330}%
In both assertions we compare a closed subgroup with \(G\), namely the stabilizer of a subspace in the first and the kernel of the adjoint representation in the second. We compute the Lie algebra of this subgroup with dual numbers and conclude with Lemma lemma-connected-components-images-homomorphisms.

Let \(U\subseteq M\) be a subspace. In a basis of \(M\) extending a basis of \(U\), the scheme-theoretic stabilizer \(S\subseteq G\) of \(U\) is the closed subgroup defined by the vanishing of the lower-left block of \(r(g)\). Since the lower-left block of \(r(1+\varepsilon X)=1+\varepsilon\,dr(X)\) is \(\varepsilon\) times that of \(dr(X)\), we find

\[
\operatorname{Lie}(S)
=\{X\in\operatorname{Lie}(G):dr(X)U\subseteq U\}.
\]

If \(U\) is \(\operatorname{Lie}(G)\)-stable, then \(\operatorname{Lie}(S)=\operatorname{Lie}(G)\), and Lemma lemma-connected-components-images-homomorphisms shows that \(S=G\). Conversely, if \(U\) is \(G\)-stable, then \(S=G\), and the displayed formula shows that \(U\) is \(\operatorname{Lie}(G)\)-stable.

Put \(\mathfrak g=\operatorname{Lie}(G)\). We claim that the kernel of the adjoint representation \(G\to\GL(\mathfrak g)\) is the center. Indeed, if \(g\) acts trivially on \(\mathfrak g\), consider the centralizer \(C_G(g)\). It is a closed subgroup scheme whose Lie algebra is the fixed space of \(\operatorname{Ad}(g)\), that is, all of \(\mathfrak g\). So \(C_G(g)=G\) by Lemma lemma-connected-components-images-homomorphisms. Conversely, a central element acts trivially. Since the kernel and the center are reduced by Cartier's theorem, their equality on \(\Omega\)-points is an equality of schemes.

The Lie algebra of the center is now the kernel of a differential. The formula

\[
\operatorname{Ad}(1+\varepsilon X)(Y)=Y+\varepsilon[X,Y]
\]

shows that the differential of \(\operatorname{Ad}\) is \(\operatorname{ad}\). The center is a closed subgroup of \(\GL(V)\), so the dual-number description applies to it. By this description, the Lie algebra of the kernel of a representation is the kernel of its differential. Hence \(\operatorname{Lie}(Z(G))=\ker(\operatorname{ad})=Z(\mathfrak g)\). By Proposition \hyperref[bp322]{F.1}, \(Z(G)^\circ\) has the same Lie algebra as \(Z(G)\).
\end{proof}


% object: sec-Cartans-criterion (blueprint/chap-Algebraic-groups-in-characteristic-zero/sec-Cartans-criterion/_section.md), content hidden
\subsection*{\texorpdfstring{F.3\quad Cartan's criterion}{F.3 Cartan's criterion}}
\phantomsection\label{bp331}%
\addcontentsline{toc}{subsection}{\texorpdfstring{F.3\quad Cartan's criterion}{F.3 Cartan's criterion}}
% text: prose-F-05
\phantomsection\label{bp332}%
The decomposition \(\mathfrak l=Z(\mathfrak l)+[\mathfrak l,\mathfrak l]\) fails for Lie algebras in general. Consider the algebra \(\mathfrak b\) of upper triangular \(2\times2\) matrices. Its center consists of the scalars, and \([\mathfrak b,\mathfrak b]\) consists of the strictly upper triangular matrices. So their sum has dimension two and misses \(\operatorname{diag}(1,0)\). Here \(\mathfrak b\) is solvable and is its own radical, much larger than its center. Its standard representation on \(\Omega^2\) is faithful but not semisimple, since the line \(\Omega e_1\) has no stable complement. In general the obstruction is the radical, that is, the largest solvable ideal. In Proposition \hyperref[bp345]{F.8} we show that a faithful semisimple representation forces the radical to be central. This gives the decomposition.

The derived series of a Lie algebra is \(D^0\mathfrak l=\mathfrak l\), \(D^{n+1}\mathfrak l=[D^n\mathfrak l,D^n\mathfrak l]\), and \(\mathfrak l\) is solvable if this series reaches zero. The radical \(\operatorname{rad}\mathfrak l\) is the sum of all solvable ideals of \(\mathfrak l\). To see that it is itself solvable, note that an extension of solvable Lie algebras is solvable. Indeed, if the quotient of \(\mathfrak l\) by a solvable ideal is solvable, the derived series of \(\mathfrak l\) lies in that ideal after finitely many steps, and then reaches zero. Hence the sum of two solvable ideals is solvable, and \(\operatorname{rad}\mathfrak l\) is the largest solvable ideal of \(\mathfrak l\) [\hyperref[bp354]{Hum72}, §3.1]. From the same book we use three classical theorems, and we only sketch their proofs.

% object: lemma-lie-s-theorem (blueprint/chap-Algebraic-groups-in-characteristic-zero/sec-Cartans-criterion/lemma-lie-s-theorem/_section.md)
\phantomsection\label{bp333}%
\begin{bpstatementit}{Lemma~F.4}{ (Lie's theorem)}
A solvable Lie algebra acting on a nonzero vector space \(M\) has a common eigenvector. Equivalently, there are \(v\ne0\) and a linear functional \(\lambda\) such that \(av=\lambda(a)v\) for every \(a\).
\end{bpstatementit}

% text: prose-F-06
\phantomsection\label{bp334}%
See [\hyperref[bp354]{Hum72}, §4.1]. By induction one finds a common eigenvector for an ideal of codimension one, and then shows by a trace computation, which we redo in the proof of Proposition \hyperref[bp345]{F.8}, that the corresponding eigenspace is stable under the remaining direction.

% object: lemma-engel-s-theorem (blueprint/chap-Algebraic-groups-in-characteristic-zero/sec-Cartans-criterion/lemma-engel-s-theorem/_section.md)
\phantomsection\label{bp335}%
\begin{bpstatementit}{Lemma~F.5}{ (Engel's theorem)}
Let \(\mathfrak n\subseteq\operatorname{End}(M)\) be a Lie algebra consisting of nilpotent endomorphisms. If \(M\ne0\), then the common kernel of \(\mathfrak n\) is nonzero. There is a basis of \(M\) in which all elements of \(\mathfrak n\) are strictly upper triangular. In particular, \(\mathfrak n\) is nilpotent and solvable.
\end{bpstatementit}

% text: prose-F-07
\phantomsection\label{bp336}%
This is Engel's theorem in its matrix form [\hyperref[bp354]{Hum72}, §3.3]. The first assertion is proved by induction on \(\dim\mathfrak n\), using that a maximal proper subalgebra is an ideal of codimension one. The triangular basis then follows by induction on \(\dim M\).

% object: lemma-cartan-s-criterion (blueprint/chap-Algebraic-groups-in-characteristic-zero/sec-Cartans-criterion/lemma-cartan-s-criterion/_section.md)
\phantomsection\label{bp337}%
\begin{bpstatementit}{Lemma~F.6}{ (Cartan's criterion)}
Let \(M\) be a finite-dimensional \(\Omega\)-vector space, and let \(\mathfrak a\subseteq\operatorname{End}(M)\) be a Lie subalgebra satisfying

\[
\operatorname{tr}(XY)=0
\qquad(X\in[\mathfrak a,\mathfrak a],\ Y\in\mathfrak a).
\]

Then \(\mathfrak a\) is solvable.
\end{bpstatementit}

% text: prose-F-08
\phantomsection\label{bp338}%
See [\hyperref[bp354]{Hum72}, §4.3]. It suffices to show that every element \(X\) of \([\mathfrak a,\mathfrak a]\) is nilpotent. Then \([\mathfrak a,\mathfrak a]\) is solvable by Lemma \hyperref[bp335]{F.5}, and hence so is \(\mathfrak a\). To see that \(X\) is nilpotent, one pairs \(X\) with operators that are polynomials in its semisimple part. Let \(\lambda_i\) be the eigenvalues of the semisimple part, and let \(f\) be a \(\mathbf Q\)-linear functional on their \(\mathbf Q\)-span with values in \(\mathbf Q\). The hypothesis gives \(\sum_i\lambda_i f(\lambda_i)=0\). Applying \(f\) to this identity gives \(\sum_i f(\lambda_i)^2=0\). This is a vanishing sum of squares of rational numbers, so \(f\) vanishes on all \(\lambda_i\), and \(f=0\). As \(f\) was arbitrary, all \(\lambda_i\) are zero. This step is where we use characteristic zero.


% object: sec-Reductive-Lie-algebras (blueprint/chap-Algebraic-groups-in-characteristic-zero/sec-Reductive-Lie-algebras/_section.md), content hidden
\subsection*{\texorpdfstring{F.4\quad Reductive Lie algebras}{F.4 Reductive Lie algebras}}
\phantomsection\label{bp339}%
\addcontentsline{toc}{subsection}{\texorpdfstring{F.4\quad Reductive Lie algebras}{F.4 Reductive Lie algebras}}
% text: prose-F-09
\phantomsection\label{bp340}%
For \(\mathfrak{gl}_n\) the decomposition into scalars and trace-zero matrices, \(\mathfrak{gl}_n=\Omega\cdot1\oplus\mathfrak{sl}_n\), has the form \(\mathfrak l=Z(\mathfrak l)+[\mathfrak l,\mathfrak l]\). We prove this decomposition for every Lie algebra with a faithful semisimple representation (Proposition \hyperref[bp345]{F.8}), as Lemma \bpout{14.15} requires. The proof has two steps. First, the radical acts by scalars on each simple summand. To see this, we take a common eigenvector given by Lie's theorem. A trace computation, in which we divide by the dimension of a subspace, shows that its eigenspace is stable under all of \(\mathfrak l\). Since the representation is faithful, the radical is then central. Second, the quotient by the center has zero radical. For such a Lie algebra the Killing form is nondegenerate by Cartan's criterion. This forces the algebra to equal its derived algebra (Lemma \hyperref[bp341]{F.7}). The upper triangular algebra \(\mathfrak b\) of the previous section shows that the hypothesis of semisimplicity cannot be dropped.

% object: lemma-killing (blueprint/chap-Algebraic-groups-in-characteristic-zero/sec-Reductive-Lie-algebras/lemma-killing/_section.md)
\phantomsection\label{bp341}%
\begin{bpstatementit}{Lemma~F.7}{ (Killing form)}
If \(\operatorname{rad}\mathfrak l=0\), the Killing form \(B(x,y)=\operatorname{tr}(\operatorname{ad}x\operatorname{ad}y)\) is nondegenerate and \([\mathfrak l,\mathfrak l]=\mathfrak l\).
\end{bpstatementit}

% text: prose-F-10
\phantomsection\label{bp342}%
This is the semisimplicity criterion of [\hyperref[bp354]{Hum72}, §§5.1–5.2]; we include the short deduction from Cartan's criterion.

% proof: proof/lemma-killing
\begin{proof}[{Proof of Lemma F.7}]\phantomsection\label{bp343}%
We first use Cartan's criterion to show that the kernel of \(B\) is a solvable ideal, hence zero; nondegeneracy then identifies the orthogonal of the derived algebra with the center. By invariance of \(B\), that is, \(B([x,y],z)=B(x,[y,z])\), the kernel \(\mathfrak k\) of \(B\) is an ideal. For \(x\in[\mathfrak k,\mathfrak k]\) and \(y\in\mathfrak k\), the operators \(\operatorname{ad}x\) and \(\operatorname{ad}y\) map \(\mathfrak l\) into \(\mathfrak k\), so the trace of \(\operatorname{ad}x\operatorname{ad}y\) may be computed on \(\mathfrak k\). Since \(x\in\mathfrak k\), we obtain

\[
\operatorname{tr}_{\mathfrak k}
(\operatorname{ad}(x)|_{\mathfrak k}
\operatorname{ad}(y)|_{\mathfrak k})=B(x,y)=0.
\]

By Cartan's criterion (Lemma \hyperref[bp337]{F.6}), the image of \(\mathfrak k\) in \(\operatorname{End}(\mathfrak k)\) is solvable. The kernel of \(\mathfrak k\to\operatorname{End}(\mathfrak k)\) is the center of \(\mathfrak k\), which is abelian. So \(\mathfrak k\) is a solvable ideal of \(\mathfrak l\). Since \(\operatorname{rad}\mathfrak l=0\), we get \(\mathfrak k=0\), that is, \(B\) is nondegenerate. Invariance and nondegeneracy then give

\[
[\mathfrak l,\mathfrak l]^\perp=Z(\mathfrak l).
\]

The center is a solvable ideal, hence zero, and a subspace with zero orthogonal complement is all of \(\mathfrak l\).
\end{proof}

% text: prose-F-11
\phantomsection\label{bp344}%
For \(\mathfrak{sl}_2\), with basis \(e,h,f\), the brackets \([h,e]=2e\), \([h,f]=-2f\) and \([e,f]=h\) already show that \([\mathfrak{sl}_2,\mathfrak{sl}_2]=\mathfrak{sl}_2\). The lemma gives the second half of the next proposition, once we know that the quotient by the center has zero radical.

% object: proposition-lie-algebras-faithful-semisimple-representation (blueprint/chap-Algebraic-groups-in-characteristic-zero/sec-Reductive-Lie-algebras/proposition-lie-algebras-faithful-semisimple-representation/_section.md)
\phantomsection\label{bp345}%
\begin{bpstatementit}{Proposition~F.8}{ (Lie algebras with a faithful semisimple representation)}
If \(\mathfrak l\) has a faithful semisimple representation, then

\[
\operatorname{rad}\mathfrak l=Z(\mathfrak l),
\qquad
\mathfrak l=Z(\mathfrak l)+[\mathfrak l,\mathfrak l].
\]
\end{bpstatementit}

% proof: proof/proposition-lie-algebras-faithful-semisimple-representation
\begin{proof}[{Proof}]\phantomsection\label{bp346}%
We follow the two halves described at the start of this section: first the radical is central, and then the quotient by the center equals its derived algebra.

Put \(\mathfrak r=\operatorname{rad}\mathfrak l\), and write the faithful module as a direct sum of simple modules. Let \(M\) be one of them. By Lemma \hyperref[bp333]{F.4}, there is a nonzero simultaneous \(\lambda\)-eigenspace \(M_\lambda\) for \(\mathfrak r\). We claim that \(M_\lambda\) is \(\mathfrak l\)-stable. To prove this, fix \(x\in\mathfrak l\) and \(0\ne v\in M_\lambda\), and consider \(U=\operatorname{span}\{v,xv,x^2v,\ldots\}\). For every \(a\in\mathfrak r\) we have \([a,x]\in\mathfrak r\), and the identity

\[
ax^j=x^ja+\sum_{i=0}^{j-1}x^i[a,x]x^{j-1-i}
\]

shows, by induction on \(j\), that \(a\) preserves \(U\) and acts on it by a triangular matrix with constant diagonal \(\lambda(a)\) in the basis \(v,xv,x^2v,\ldots\) of \(U\). Indeed, the identity gives \(ax^jv=\lambda(a)x^jv\) plus terms \(x^i[a,x]x^{j-1-i}v\), and by induction applied to \([a,x]\in\mathfrak r\) these lie in the span of \(v,\ldots,x^{j-1}v\). With this, the trace of \([x,a]\) on \(U\) can be computed in two ways. Since \(x\) and \(a\) preserve \(U\), the restriction of \([x,a]\) to \(U\) is a commutator, whose trace is zero. On the other hand, \([x,a]\in\mathfrak r\) acts on \(U\) with constant diagonal \(\lambda([x,a])\). Comparing,

\[
0=\operatorname{tr}_U[x,a]=\dim(U)\lambda([x,a]),
\]

and \(\lambda([x,a])=0\) because \(\Omega\) has characteristic zero. It follows that \(a(xv)=\lambda(a)xv+\lambda([a,x])v=\lambda(a)xv\), which proves the claim.

Since \(M\) is simple, the claim gives \(M_\lambda=M\). Thus \(\mathfrak r\) acts on \(M\) by scalars, and in particular \([\mathfrak l,\mathfrak r]\) acts as zero on \(M\). This holds for every simple summand, and the module is faithful, so \([\mathfrak l,\mathfrak r]=0\), that is, \(\mathfrak r\subseteq Z(\mathfrak l)\). The reverse inclusion holds because the center is an abelian ideal.

For the second equality, we show that the quotient \(\overline{\mathfrak l}=\mathfrak l/Z(\mathfrak l)\) has zero radical and apply Lemma \hyperref[bp341]{F.7}. The inverse image of a solvable ideal of the quotient is a central extension of a solvable Lie algebra, hence solvable. So it is contained in \(\mathfrak r=Z(\mathfrak l)\). The lemma now gives \([\overline{\mathfrak l},\overline{\mathfrak l}]
 =\overline{\mathfrak l}\). In other words, every element of \(\mathfrak l\) is congruent modulo \(Z(\mathfrak l)\) to a sum of brackets, and \(\mathfrak l=Z(\mathfrak l)+[\mathfrak l,\mathfrak l]\).
\end{proof}




\clearpage\phantomsection\hypertarget{bp-outside}{}%
\section*{Outside this export}
\addcontentsline{toc}{section}{Outside this export}
References marked \ensuremath{{}^\ast} are to parts of the blueprint this document leaves out:
\begin{description}
\item[Definition 2.30] Classical pro-étale sheaves, in Ⅰ The six operations (\bpslug{definition-{}\allowbreak{}classical-{}\allowbreak{}pro-{}\allowbreak{}etale-{}\allowbreak{}sheaves})
\item[Theorem 2.32] Classical sheaves as the essential image, in Ⅰ The six operations (\bpslug{theorem-{}\allowbreak{}classical-{}\allowbreak{}sheaves-{}\allowbreak{}essential-{}\allowbreak{}image})
\item[Theorem 2.35] Bounded-below cohomology comparison, in Ⅰ The six operations (\bpslug{theorem-{}\allowbreak{}bounded-{}\allowbreak{}below-{}\allowbreak{}cohomology-{}\allowbreak{}comparison})
\item[Proposition 2.42] Locally constant sheaves of finite presentation, in Ⅰ The six operations (\bpslug{proposition-{}\allowbreak{}locally-{}\allowbreak{}constant-{}\allowbreak{}sheaves-{}\allowbreak{}finite-{}\allowbreak{}presentation})
\item[Theorem 2.43] Comparison for constructible sheaves, in Ⅰ The six operations (\bpslug{theorem-{}\allowbreak{}comparison-{}\allowbreak{}constructible-{}\allowbreak{}sheaves})
\item[Corollary 2.44] Finite Tor amplitude under comparison, in Ⅰ The six operations (\bpslug{corollary-{}\allowbreak{}finite-{}\allowbreak{}tor-{}\allowbreak{}amplitude-{}\allowbreak{}under-{}\allowbreak{}comparison})
\item[Proposition 2.58] Extension by zero and restriction to a closed subscheme, in Ⅰ The six operations (\bpslug{proposition-{}\allowbreak{}extension-{}\allowbreak{}zero-{}\allowbreak{}restriction-{}\allowbreak{}closed-{}\allowbreak{}subscheme})
\item[Section 4.4] The projective line, in Ⅰ The six operations (\bpslug{sec-{}\allowbreak{}The-{}\allowbreak{}projective-{}\allowbreak{}line})
\item[Theorem 4.20] Cohomology of curves, in Ⅰ The six operations (\bpslug{theorem-{}\allowbreak{}cohomology-{}\allowbreak{}curves})
\item[Lemma 4.21] Separably closed field extension, in Ⅰ The six operations (\bpslug{lemma-{}\allowbreak{}extension-{}\allowbreak{}separably-{}\allowbreak{}closed-{}\allowbreak{}ground-{}\allowbreak{}field})
\item[Section 4.5] Proper morphisms, in Ⅰ The six operations (\bpslug{sec-{}\allowbreak{}Proper-{}\allowbreak{}morphisms})
\item[Lemma 4.27] Finite morphisms, in Ⅰ The six operations (\bpslug{lemma-{}\allowbreak{}finite-{}\allowbreak{}morphisms})
\item[Lemma 4.29] Composition and cancellation, in Ⅰ The six operations (\bpslug{lemma-{}\allowbreak{}composition-{}\allowbreak{}cancellation})
\item[Proposition 5.26] Construction and independence of $Rf_!$, in Ⅰ The six operations (\bpslug{proposition-{}\allowbreak{}construction-{}\allowbreak{}independence-{}\allowbreak{}rf-{}\allowbreak{}shriek})
\item[Proposition 8.18] Ordinary external Künneth, in Ⅰ The six operations (\bpslug{proposition-{}\allowbreak{}ordinary-{}\allowbreak{}external-{}\allowbreak{}kunneth})
\item[Chapter 10] Verdier and Poincaré duality, in Ⅱ Duality and the trace formula (\bpslug{chap-{}\allowbreak{}Verdier-{}\allowbreak{}and-{}\allowbreak{}Poincare-{}\allowbreak{}duality})
\item[Lemma 10.6] Restriction to a geometric generic fibre, in Ⅱ Duality and the trace formula (\bpslug{lemma-{}\allowbreak{}restriction-{}\allowbreak{}geometric-{}\allowbreak{}generic-{}\allowbreak{}fibre})
\item[Lemma 10.7] Double duals on the geometric generic fibre, in Ⅱ Duality and the trace formula (\bpslug{lemma-{}\allowbreak{}double-{}\allowbreak{}dual-{}\allowbreak{}geometric-{}\allowbreak{}generic-{}\allowbreak{}fibre})
\item[Chapter 11] The Lefschetz–Verdier formula, in Ⅱ Duality and the trace formula (\bpslug{chap-{}\allowbreak{}The-{}\allowbreak{}Lefschetz-{}\allowbreak{}Verdier-{}\allowbreak{}formula})
\item[Lemma 11.7] External duality, in Ⅱ Duality and the trace formula (\bpslug{lemma-{}\allowbreak{}external-{}\allowbreak{}duality-{}\allowbreak{}trace-{}\allowbreak{}morphism})
\item[Definition 11.8] Trace morphism, in Ⅱ Duality and the trace formula (\bpslug{definition-{}\allowbreak{}trace-{}\allowbreak{}morphism})
\item[Definition 11.9] Trace class and local terms, in Ⅱ Duality and the trace formula (\bpslug{definition-{}\allowbreak{}trace-{}\allowbreak{}class-{}\allowbreak{}local-{}\allowbreak{}terms})
\item[Theorem 11.13] Proper pushforward and Lefschetz–Verdier, in Ⅱ Duality and the trace formula (\bpslug{theorem-{}\allowbreak{}proper-{}\allowbreak{}pushforward-{}\allowbreak{}lefschetz-{}\allowbreak{}verdier})
\item[Lemma 11.17] Lifting a correspondence to an invariant decomposition, in Ⅱ Duality and the trace formula (\bpslug{lemma-{}\allowbreak{}lift-{}\allowbreak{}correspondence-{}\allowbreak{}invariant-{}\allowbreak{}decomposition})
\item[Lemma 11.18] The closed piece as a pushforward, in Ⅱ Duality and the trace formula (\bpslug{lemma-{}\allowbreak{}closed-{}\allowbreak{}piece-{}\allowbreak{}pushforward})
\item[Chapter 12] The trace formula, in Ⅱ Duality and the trace formula (\bpslug{chap-{}\allowbreak{}The-{}\allowbreak{}trace-{}\allowbreak{}formula})
\item[Remark 12.4] Local terms are not computed stalk by stalk, in Ⅱ Duality and the trace formula (\bpslug{remark-{}\allowbreak{}locality-{}\allowbreak{}warning})
\item[Proposition 12.16] The Frobenius local term on a curve, in Ⅱ Duality and the trace formula (\bpslug{proposition-{}\allowbreak{}frobenius-{}\allowbreak{}local-{}\allowbreak{}term-{}\allowbreak{}curve})
\item[Theorem 12.24] Grothendieck–Lefschetz trace formula, in Ⅱ Duality and the trace formula (\bpslug{theorem-{}\allowbreak{}grothendieck-{}\allowbreak{}lefschetz-{}\allowbreak{}trace-{}\allowbreak{}formula})
\item[Chapter 14] Purity of real sheaves, in Ⅲ Weights and the Riemann hypothesis (\bpslug{chap-{}\allowbreak{}Purity-{}\allowbreak{}of-{}\allowbreak{}real-{}\allowbreak{}sheaves})
\item[Lemma 14.1] A proper smooth model, in Ⅲ Weights and the Riemann hypothesis (\bpslug{lemma-{}\allowbreak{}proper-{}\allowbreak{}smooth-{}\allowbreak{}model})
\item[Proposition 14.3] No fixed class in degree one, in Ⅲ Weights and the Riemann hypothesis (\bpslug{proposition-{}\allowbreak{}no-{}\allowbreak{}fixed-{}\allowbreak{}class-{}\allowbreak{}degree-{}\allowbreak{}one})
\item[Lemma 14.15] The identity component, in Ⅲ Weights and the Riemann hypothesis (\bpslug{lemma-{}\allowbreak{}identity-{}\allowbreak{}component})
\item[Proposition 14.19] Determinant weights of tensor operations, in Ⅲ Weights and the Riemann hypothesis (\bpslug{proposition-{}\allowbreak{}determinant-{}\allowbreak{}weights-{}\allowbreak{}tensor-{}\allowbreak{}operations})
\end{description}
\clearpage\phantomsection%
\section*{References}
\addcontentsline{toc}{section}{References}
\bpbibitem{BBD82}{bp355}{A. Beilinson, J. Bernstein, and P. Deligne. \emph{Faisceaux pervers}, volume 100 of \emph{Astérisque}. 1982.}

\bpbibitem{BS15}{bp356}{B. Bhatt and P. Scholze. The pro-étale topology for schemes. \emph{Astérisque}, 369:99–201, 2015. \href{https://people.mpim-bonn.mpg.de/scholze/proetale.pdf}{Author's manuscript}.}

\bpbibitem{Con07}{bp348}{B. Conrad. Deligne's notes on Nagata compactifications. \emph{Journal of the Ramanujan Mathematical Society}, 22:205–257, 2007. \href{https://math.stanford.edu/\~conrad/papers/nagatafinal.pdf}{Author's manuscript}.}

\bpbibitem{Del74}{bp357}{P. Deligne. La conjecture de Weil. I. \emph{Publ. Math. IHÉS}, 43:273–307, 1974.}

\bpbibitem{Del80}{bp358}{P. Deligne. La conjecture de Weil. II. \emph{Publ. Math. IHÉS}, 52:137–252, 1980.}

\bpbibitem{Den22}{bp359}{C. Deninger. There is no “Weil-”cohomology theory with real coefficients for arithmetic curves. \emph{Ann. Sc. Norm. Super. Pisa Cl. Sci.}, 2022. To appear. arXiv:\href{https://arxiv.org/abs/2204.02714}{2204.02714}.}

\bpbibitem{Deu41}{bp360}{M. Deuring. Die Typen der Multiplikatorenringe elliptischer Funktionenkörper. \emph{Abh. Math. Sem. Univ. Hamburg}, 14:197–272, 1941.}

\bpbibitem{Eke90}{bp361}{T. Ekedahl. On the adic formalism. In \emph{The Grothendieck Festschrift, Vol. II}, volume 87 of \emph{Progress in Mathematics}, pages 197–218. Birkhäuser, 1990.}

\bpbibitem{Has36}{bp362}{H. Hasse. Zur Theorie der abstrakten elliptischen Funktionenkörper. I, II, III. \emph{J. Reine Angew. Math.}, 175:55–62, 69–88, 193–208, 1936.}

\bpbibitem{Hum72}{bp354}{J. E. Humphreys. \emph{Introduction to Lie Algebras and Representation Theory}, volume 9 of \emph{Graduate Texts in Mathematics}. Springer, 1972.}

\bpbibitem{Ill94}{bp351}{L. Illusie. Autour du théorème de monodromie locale. \emph{Astérisque}, 223:9–57, 1994. §4. \href{https://www.numdam.org/item/AST_1994__223__9_0/}{Original article}.}

\bpbibitem{Kat88}{bp363}{N. M. Katz. Travaux de Laumon. \emph{Astérisque}, 161–162:105–132, 1988. Séminaire Bourbaki, exp. 691.}

\bpbibitem{KW01}{bp364}{R. Kiehl and R. Weissauer. \emph{Weil conjectures, perverse sheaves and \(\ell\)-adic Fourier transform}, volume 42 of \emph{Ergebnisse der Mathematik}. Springer, 2001.}

\bpbibitem{Lau87}{bp365}{G. Laumon. Transformation de Fourier, constantes d'équations fonctionnelles et conjecture de Weil. \emph{Publ. Math. IHÉS}, 65:131–210, 1987.}

\bpbibitem{Mil86a}{bp349}{J. S. Milne. \href{https://www.jmilne.org/math/xnotes/AVs.pdf}{Abelian Varieties}. In \emph{Arithmetic Geometry}. 1986. Corrected online version.}

\bpbibitem{Mil86b}{bp350}{J. S. Milne. \href{https://www.jmilne.org/math/xnotes/JVs.pdf}{Jacobian Varieties}. In \emph{Arithmetic Geometry}. 1986. Corrected online version.}

\bpbibitem{Ran39}{bp366}{R. A. Rankin. Contributions to the theory of Ramanujan's function \(\tau(n)\) and similar arithmetical functions. II. \emph{Proc. Cambridge Philos. Soc.}, 35:357–372, 1939.}

\bpbibitem{SGA4}{bp367}{M. Artin, A. Grothendieck, and J.-L. Verdier, editors. \emph{Théorie des topos et cohomologie étale des schémas}, volumes 269, 270, 305 of \emph{Lecture Notes in Mathematics}. Springer, 1972–1973. SGA 4.}

\bpbibitem{SGA4½}{bp353}{P. Deligne. \emph{Cohomologie étale}, volume 569 of \emph{Lecture Notes in Mathematics}. Springer, 1977. SGA 4½; \emph{Théorèmes de finitude}, §3, and \emph{Rapport sur la formule des traces}, §3, and \emph{Sommes trigonométriques}, §1.}

\bpbibitem{SGA5}{bp368}{A. Grothendieck, editor. \emph{Cohomologie \(\ell\)-adique et fonctions \(L\)}, volume 589 of \emph{Lecture Notes in Mathematics}. Springer, 1977. SGA 5.}

\bpbibitem{SGA7II}{bp352}{P. Deligne and N. Katz, editors. \emph{Groupes de monodromie en géométrie algébrique}, volume 340 of \emph{Lecture Notes in Mathematics}. Springer, 1973. SGA 7 II; Exposé XIII, (2.1.7).}

\bpbibitem{Stacks}{bp369}{The Stacks Project Authors. \href{https://stacks.math.columbia.edu}{The Stacks Project}. Individual results are cited by their stable tags in the text.}

\bpbibitem{Var05}{bp347}{Y. Varshavsky. Lefschetz–Verdier trace formula and a generalization of a theorem of Fujiwara. 2005. Version 2. arXiv:\href{https://arxiv.org/abs/math/0505564}{math/0505564}. \href{https://arxiv.org/abs/math/0505564v2}{Author's preprint}.}

\bpbibitem{Wei48a}{bp370}{A. Weil. \emph{Sur les courbes algébriques et les variétés qui s'en déduisent}. Hermann, Paris, 1948.}

\bpbibitem{Wei48b}{bp371}{A. Weil. \emph{Variétés abéliennes et courbes algébriques}. Hermann, Paris, 1948.}

\bpbibitem{Wei49}{bp372}{A. Weil. Numbers of solutions of equations in finite fields. \emph{Bull. Amer. Math. Soc.}, 55:497–508, 1949.}

\end{document}
